\documentclass[11pt,a4paper]{report}
\usepackage{fontspec,amsmath,amssymb,amsthm,microtype,xcolor}
\setmainfont{lmroman10-regular.otf}[BoldFont=lmroman10-bold.otf,ItalicFont=lmroman10-italic.otf,BoldItalicFont=lmroman10-bolditalic.otf]
\setmonofont{lmmono10-regular.otf}[Scale=MatchLowercase]
\usepackage[margin=25mm]{geometry}
\usepackage[unicode,hidelinks]{hyperref}
\setlength{\emergencystretch}{3em}
\setlength{\parskip}{0.4em}
\setlength{\parindent}{0pt}
\setcounter{tocdepth}{1}
\makeatletter
\renewcommand*\l@section{\@dottedtocline{1}{1.5em}{3.8em}}
\renewcommand*\@pnumwidth{2.5em}
\renewcommand*\@tocrmarg{3.5em}
\makeatother
\setcounter{secnumdepth}{2}
\allowdisplaybreaks[2]
\raggedbottom
\title{\texorpdfstring{Interacting Particle Langevin under Strong Convexity}{Interacting Particle Langevin under Strong Convexity}\\[0.6em]\large Mathematical proofs}
\author{LangevinLean}
\date{}
\begin{document}
\pagenumbering{roman}
\maketitle
\chapter*{About this document}
\addcontentsline{toc}{chapter}{About this document}
This volume collects the theorem statements, mathematical arguments and supporting proofs for this development. Every locally authored proof module required by the displayed results, proof roadmap and linked extensions is included once, in dependency order. Foundational results from Mathlib are used as standard mathematical results.



The final strongly convex theorem derives its stationary reference, target moments, and equilibrium concentration from the joint-potential assumptions. Parameter and latent dimensions are independent. The entropy extension proves its bound from a lifted-target LSI and finite initial joint KL.



\textbf{Published source and unpublished extension}



Source: Theorem 1, equation (34), and Appendix B of arXiv:2409.07101v2 (27 May 2025), using the ULA estimate of Dalalyan and Karagulyan. The independent parameter- and latent-dimension extension and its proof presented here are unpublished; equal dimensions recover the paper’s displayed bound. The assumption audit explains the joint-potential smoothness condition.


\begin{itemize}

\item Alex Glyn-Davies, Connor Duffin, Ieva Kazlauskaite, Mark Girolami, and Ö. Deniz Akyildiz (2025). Statistical Finite Elements via Interacting Particle Langevin Dynamics. SIAM/\allowbreak{}ASA Journal on Uncertainty Quantification 13(3), 1200–1227. Audited version: arXiv:2409.07101v2
\item Arnak S. Dalalyan and Avetik G. Karagulyan (2019). User-friendly guarantees for the Langevin Monte Carlo with inaccurate gradient. Stochastic Processes and their Applications 129(12), 5278–5311. arXiv
\end{itemize}
\tableofcontents
\clearpage
\pagenumbering{arabic}
\chapter{Assumptions and main results}
\section{\texorpdfstring{Parameter error for the actual IPLA recursion}{Parameter error for the actual IPLA recursion}}


Fix an observation \(y\) and a differentiable joint potential \(U_y:\mathbb R^{d_\theta}\times\mathbb R^{d_x}\to\mathbb R\). Assume \(U_y\) is \(m\)-strongly convex with \(m>0\) and has \(L\)-Lipschitz gradient. Let \(N\ge1\); the independent dimensions \(d_\theta,d_x\ge0\) may be zero. Define


\[\ell_y(\theta)=\int e^{-U_y(\theta,x)}dx,\quad
\theta_\star\in\operatorname*{argmax}_\theta\ell_y(\theta).\]


In scaled coordinates \(z=(\theta,x^1/\sqrt N,\ldots,x^N/\sqrt N)\), set


\[V_N(z)=\sum_{i=1}^NU_y(\theta,\sqrt N z_i),\qquad
\Pi_N(dz)\propto e^{-V_N(z)}dz.\]


Use a square-integrable initial state and standard Gaussian innovations independent in every coordinate and time and independent of initialization. For \(0<\gamma\le2/(m+L)\), the actual IPLA updates are


\[\begin{aligned}
\theta_{k+1}&=\theta_k-\frac\gamma N\sum_i\nabla_\theta U_y(\theta_k,x_k^i)+\sqrt{2\gamma/N}\,\xi_k^0,\\
x_{k+1}^i&=x_k^i-\gamma\nabla_xU_y(\theta_k,x_k^i)+\sqrt{2\gamma}\,\xi_k^i.
\end{aligned}\]


Writing \(\mu_0\) for the initial scaled-state law, for every \(k\in\mathbb N\),


\[\begin{aligned}
\bigl(\mathbb E\|\theta_k-\theta_\star\|^2\bigr)^{1/2}
\le{}&\sqrt{\frac{2d_\theta}{Nm}}+(1-m\gamma)^kW_2(\mu_0,\Pi_N)\\
&+\frac{33}{20}\frac Lm\sqrt{d_x+\frac{d_\theta}{N}}\sqrt\gamma.
\end{aligned}\]


The target normalization, finite transport moments, stationary reference, marginal strong convexity, and equilibrium radius are derived. The selected center is a maximizer of the integrated joint density; including a prior in \(U_y\) gives the prior-inclusive interpretation.


{\small\color{gray}Supporting proof: Section~\ref{proof-module-IPLA-Theorem-lean}.\par}
\section{Proof organization}
The supporting proofs below establish the following steps.
\begin{enumerate}
\item \textbf{Lift and rescale particles}. Dimension \(d_\theta+Nd_x\); curvature \(Nm\) See Section~\ref{proof-module-IPLA-Geometry-lean}.
\item \textbf{Normalize the lifted target}. Derived probability and finite moments See Section~\ref{proof-module-IPLA-Foundations-lean}.
\item \textbf{Compute the parameter marginal}. Density proportional to \(\ell_y(\theta)^N\) See Section~\ref{proof-module-IPLA-MarginalConcentration-lean}.
\item \textbf{Identify IPLA with ULA}. Exact algorithm law; step \(\gamma/N\) See Section~\ref{proof-module-IPLA-ULAAdapter-lean}.
\item \textbf{Apply the shared ULA theorem}. Stationarity and Gaussian law proved See Section~\ref{proof-module-UnadjustedLangevin-GaussianULA-lean}.
\item \textbf{Prove marginal strong convexity}. Prékopa–Leindler and tensorization See Section~\ref{proof-module-IPLA-StrongLogConcavity-lean}.
\item \textbf{Simplify particle constants}. Exact \(N\), step, and dimension factors See Section~\ref{proof-module-IPLA-Scaling-lean}.
\item \textbf{Project the transport estimate}. Parameter projection: \(\|P_\theta\|\le1\) See Section~\ref{proof-module-IPLA-Transport-lean}.
\item \textbf{Derive the equilibrium radius}. RMS \(\le\sqrt{2d_\theta/(Nm)}\) See Section~\ref{proof-module-IPLA-Concentration-lean}.
\item \textbf{Final IPLA parameter bound}. General independent dimensions See Section~\ref{proof-module-IPLA-Theorem-lean}.
\end{enumerate}
\chapter{Unpublished entropy extension}


\textbf{Unpublished proof}



The IPLA LSI proof presented here is unpublished. It derives the joint entropy bound, per-particle normalization, and parameter-marginal consequence by applying the shared ULA entropy theorem to the actual lifted IPLA target. The papers below are background references for that construction.


\begin{itemize}

\item Santosh S. Vempala and Andre Wibisono (2019). Rapid Convergence of the Unadjusted Langevin Algorithm: Isoperimetry Suffices. Advances in Neural Information Processing Systems 32, 8094–8106. Full paper with appendices
\item Alex Glyn-Davies, Connor Duffin, Ieva Kazlauskaite, Mark Girolami, and Ö. Deniz Akyildiz (2025). Statistical Finite Elements via Interacting Particle Langevin Dynamics. SIAM/\allowbreak{}ASA Journal on Uncertainty Quantification 13(3), 1200–1227. Audited version: arXiv:2409.07101v2
\end{itemize}


The lifted-target entropy bound follows from the full actual Gaussian ULA proof and finite initial joint KL. The LSI hypothesis concerns the actual lifted target; all interpolation and moment estimates are derived.



The adjacent theorem assumes the lifted-target LSI and finite initial joint KL. No regular interpolation data or compatibility premise is supplied. Parameter and latent dimensions are independent.

\section{\texorpdfstring{Joint KL bound and the linear lifted-\allowbreak{}LSI specialization}{Joint KL bound and the linear lifted-LSI specialization}}


Let \(N\ge1\), \(D=d_\theta+Nd_x\), assume \(D\ge1\), and let the joint potential \(U(\theta,x)\) have an \(L\)-Lipschitz gradient with \(L>0\). In scaled coordinates set


\[z=(\theta,x^1/\sqrt N,\ldots,x^N/\sqrt N),\qquad
V_N(z)=\sum_{i=1}^NU(\theta,\sqrt N z_i),\quad
\Pi_N(dz)\propto e^{-V_N(z)}dz.\]


Assume an LSI with constant \(\beta_N>0\) for this actual lifted target. Compatibility gives \(\beta_N\le NL\) by the positive-dimensional argument on the shared ULA page; the final Lean theorem derives this compatibility internally. For finite initial joint entropy and


\[0<\gamma\le\frac{\beta_N}{4NL^2},\]


the IPLA laws \(\mu_k\), equivalently ULA for \(V_N\) at step \(\gamma/N\), satisfy


\[\frac{H_{\Pi_N}(\mu_k)}N\le
e^{-\beta_N\gamma k/N}\frac{H_{\Pi_N}(\mu_0)}N
+\frac{8\gamma DL^2}{\beta_N}.\]


In particular, if the lifted target has the admissible constant \(\beta_N=N\alpha\), then


\[\boxed{\frac{H_{\Pi_N}(\mu_k)}N\le
e^{-\alpha\gamma k}\frac{H_{\Pi_N}(\mu_0)}N
+\frac{8\gamma L^2}{\alpha}\left(d_x+\frac{d_\theta}N\right).}\]


The condition is now \(0<\gamma\le\alpha/(4L^2)\). An LSI for the original pair target is not substituted for the lifted-target hypothesis. The Lean result derives all interpolation estimates from the shared ULA theorem and finite initial joint KL. If \(D=0\), both laws are the unique point mass and their entropy is zero.


{\small\color{gray}Supporting proof: Section~\ref{proof-module-IPLA-LSI-UnconditionalTheorem-lean}.\par}
\section{Mathematical argument}
\subsection{\texorpdfstring{The geometry fixes the dimension, gradient constant, and step}{The geometry fixes the dimension, gradient constant, and step}}


Each physical pair \((\theta,x^i)\) is extracted from the scaled state by \((\theta,\sqrt N z_i)\). Differentiating the lifted sum gives the averaged parameter drift and ordinary latent drifts at step \(\gamma/N\). The corresponding Euclidean state has


\[D=d_\theta+Nd_x,\qquad L_N=NL,\qquad h_N=\gamma/N.\]


The parameter block of the lifted gradient is a sum over particles; each latent block contributes a factor \(\sqrt N\). Squaring and summing the blocks, then applying Cauchy–Schwarz to the common-parameter sum, proves the gradient Lipschitz constant \(NL\). The coordinate Gaussian laws and independence identify the actual noise-driven IPLA law with this ULA transition. No equality between \(d_\theta\) and \(d_x\) is needed.



\textbf{Scope.} The lifted geometry and identification with the actual Gaussian particle algorithm are checked.

\subsection{\texorpdfstring{Apply the ULA one-\allowbreak{}step estimate with the lifted target}{Apply the ULA one-step estimate with the lifted target}}


Apply the ULA entropy argument to \((V_N,\Pi_N)\), with constant \(\beta_N\), dimension \(D\), smoothness \(NL\), and step \(\gamma/N\). Its step restriction becomes


\[\frac\gamma N\le\frac{\beta_N}{4(NL)^2}
\quad\Longleftrightarrow\quad
\gamma\le\frac{\beta_N}{4NL^2}.\]


The one-step forcing simplifies without losing a particle factor:


\[6h_N^2D L_N^2=6(\gamma/N)^2D(NL)^2=6\gamma^2DL^2.\]


Consequently the unnormalized joint entropy obeys


\[H_{k+1}\le e^{-\beta_N\gamma/N}H_k+6\gamma^2DL^2.\]


The Lean proof constructs probability and finite-KL guards first, applies the shared one-step theorem, and then rewrites the rescaled constants. Every analytic input is supplied by the actual Gaussian ULA construction.



\textbf{Scope.} The particle substitution and scalar algebra use the shared ULA proof from smoothness, LSI, and finite initial KL.

\subsection{\texorpdfstring{Normalize by the particle number only after deriving the joint bound}{Normalize by the particle number only after deriving the joint bound}}


Iteration gives the unnormalized floor


\[\frac{8h_NDL_N^2}{\beta_N}
=\frac{8\gamma NDL^2}{\beta_N}.\]


Divide the entire inequality by positive \(N\), including the initial entropy and the floor. This proves


\[\frac{H_k}N\le e^{-\beta_N\gamma k/N}\frac{H_0}N
+\frac{8\gamma DL^2}{\beta_N}.\]


Now, and only under the separate lifted-target assumption \(\beta_N=N\alpha\), use \(D/N=d_x+d_\theta/N\). The rate is \(e^{-\alpha\gamma k}\), and the per-particle floor is \((8\gamma L^2/\alpha)(d_x+d_\theta/N)\). Shared parameter coordinates prevent a direct tensorization argument from a pair LSI; the target for which LSI is assumed remains \(\Pi_N\).



\textbf{Scope.} The general lifted-constant bound, division by N, and linear-LSI specialization are checked with their assumptions visible.

\subsection{\texorpdfstring{Project to the parameter marginal using data processing}{Project to the parameter marginal using data processing}}


Let \(P_\theta\) be the parameter projection, \(\mu_k^\theta=P_{\theta\#}\mu_k\), and \(\Pi_N^\theta=P_{\theta\#}\Pi_N\). Relative entropy decreases under a measurable map:


\[H_{\Pi_N^\theta}(\mu_k^\theta)\le H_{\Pi_N}(\mu_k).\]


This controls the unnormalized parameter KL. In particular, under \(\beta_N=N\alpha\),


\[H_{\Pi_N^\theta}(\mu_k^\theta)
\le e^{-\alpha\gamma k}H_{\Pi_N}(\mu_0)
+\frac{8\gamma NL^2}{\alpha}\left(d_x+\frac{d_\theta}N\right).\]


The initial entropy remains the full joint entropy. Data processing does not allow it to be replaced by the smaller initial parameter entropy on the right. Likewise, a normalized joint bound must be multiplied by \(N\) before it bounds the unscaled parameter KL. The Lean proof establishes finite joint KL, transfers finiteness to the marginal, and only then compares the real values.



\textbf{Scope.} The parameter data-processing argument and all normalization factors are checked. This conclusion concerns a probability marginal, not distance to a likelihood maximizer.

\chapter{Supporting proofs}
The following mathematical arguments are reproduced in full from the authored exposition of the local proof dependencies. Each module appears after its local imports.

\section{\texorpdfstring{Scaled particle geometry and exact coordinate updates}{Scaled particle geometry and exact coordinate updates}}
\label{proof-module-IPLA-Geometry-lean}
{\small\color{gray}Source module: \texttt{IPLA/\allowbreak{}Geometry.\allowbreak{}lean}.\par}


Fix a particle count \(N\), independent dimensions \(d_\theta,d_x\), and a joint potential \(U\). Write \(z=(\theta,z_1,\ldots,z_N)\), where the physical particles are \(x_i=\sqrt N z_i\). The lifted potential is \(V_N(z)=\sum_iU(\theta,\sqrt N z_i)\). This module derives its dimension, strong convexity, smoothness and exact Euler coordinates; none is an additional hypothesis.

\subsection{\texorpdfstring{The coordinate spaces and their dimensions}{The coordinate spaces and their dimensions}}


The scaled state is the Euclidean space indexed by \(\mathrm{Fin}(d_\theta)\sqcup(\mathrm{Fin}(N)\times\mathrm{Fin}(d_x))\), and a joint pair is indexed by \(\mathrm{Fin}(d_\theta)\sqcup\mathrm{Fin}(d_x)\). Counting these orthogonal coordinates proves

\[\dim(\mathrm{State})=d_\theta+Nd_x,\qquad \dim(\mathrm{Joint})=d_\theta+d_x.\]

The block map extracts \(z_i\), while \(A_i z=\operatorname{joint}_i(z)=(\theta,\sqrt N z_i)\) restores the physical particle. The two coordinate evaluation identities follow directly from this definition. The definition of \(V_N\) then sums the original potential over these pairs.




\textbf{Correspondence with the displayed declarations.}


\begin{itemize}

\item \texttt{State}: The scaled IPLA state: the common parameter followed by \texttt{N} particle blocks.
\item \texttt{Joint}: A parameter-particle pair, with parameter block first.
\item \texttt{finrank\_\allowbreak{}state}: IPLA has one common parameter and \texttt{N} latent vectors, in arbitrary, independent parameter and latent dimensions.
\item \texttt{finrank\_\allowbreak{}joint}: The original joint potential has one parameter block and one latent block.
\item \texttt{block}: Extract the scaled latent vector of particle \texttt{i}.
\item \texttt{joint}: Pair the common parameter with the unscaled \texttt{i}th particle.
\item \texttt{potential}: The potential used by ULA after scaling the particles by \texttt{1 /\allowbreak{} sqrt N}.
\end{itemize}
\subsection{\texorpdfstring{Replication preserves a precise quadratic identity}{Replication preserves a precise quadratic identity}}


Each \(A_i\) is linear: checking a parameter coordinate and a latent coordinate separately proves preservation of addition, scalar multiplication and subtraction. Every parameter coordinate occurs \(N\) times, while each latent coordinate is multiplied by \(\sqrt N\) once. Therefore

\[\sum_i\|A_i z\|^2=N\|\theta\|^2+N\sum_i\|z_i\|^2=N\|z\|^2.\]

The proof expands squared Euclidean norms into finite sums and uses \((\sqrt N)^2=N\), which is valid also when \(N=0\). This exact identity supplies the constants in both the convexity and smoothness arguments.




\textbf{Correspondence with the displayed declarations.}


\begin{itemize}

\item \texttt{joint\_\allowbreak{}norm\_\allowbreak{}sq\_\allowbreak{}sum}: Squared Euclidean norms of all replicated pairs sum to exactly \texttt{N} times the squared norm in scaled particle coordinates.
\end{itemize}
\subsection{\texorpdfstring{Joint strong convexity lifts to N times the constant}{Joint strong convexity lifts to N times the constant}}


If \(U\) is \(m\)-strongly convex, then for \(a,b\ge0\), \(a+b=1\), each pair satisfies

\[U(aA_ix+bA_iy)\le aU(A_ix)+bU(A_iy)-\frac{mab}{2}\|A_i(x-y)\|^2.\]

Sum over \(i\), use linearity to identify the arguments with \(A_i(ax+by)\), and use the quadratic replication identity on \(x-y\). The result is the strong-convexity inequality for \(V_N\) with constant \(Nm\). The continuous linear map \(A_i\) is obtained from its finite-dimensional linear map; its evaluation theorem is definitional. If \(U\) is differentiable, the finite sum of \(U\circ A_i\) is differentiable by the chain rule.




\textbf{Correspondence with the displayed declarations.}


\begin{itemize}

\item \texttt{potential\_\allowbreak{}strong\allowbreak{}Convex}: Strong convexity of the original pair potential gives \texttt{N*m} strong convexity of the scaled particle potential, with no extra hypothesis.
\item \texttt{joint\allowbreak{}CLM}: The linear joint-pair extraction map.
\item \texttt{potential\_\allowbreak{}differentiable}: Differentiability survives the finite sum and the coordinate scaling.
\end{itemize}
\subsection{\texorpdfstring{The gradient chain rule and the lifted Lipschitz constant}{The gradient chain rule and the lifted Lipschitz constant}}


Differentiating in a direction \(w\) gives

\[\langle\nabla V_N(z),w\rangle=\sum_i\langle\nabla U(A_i z),A_iw\rangle.\]

Cauchy–Schwarz in the particle index, followed by the quadratic replication identity, proves \(\sum_i\|A_iv\|\|A_iw\|\le N\|v\|\|w\|\). If \(\nabla U\) is \(L\)-Lipschitz, subtract the two gradient identities to obtain

\[\langle\nabla V_N(x)-\nabla V_N(y),w\rangle\le NL\|x-y\|\|w\|.\]

Choose \(w=\nabla V_N(x)-\nabla V_N(y)\). If its norm is zero the desired bound is immediate; otherwise cancel its positive norm. Thus \(\nabla V_N\) is \(NL\)-Lipschitz, with the same particle factor as the convexity constant.




To track the constants without a coordinatewise overestimate, test the score difference against an arbitrary \(w\). The chain rule gives a sum of pair-gradient differences paired with \(A_iw\). Smoothness bounds it by \(L\sum_i\|A_i(x-y)\|\|A_iw\|\), and Cauchy–Schwarz over particles makes this \(NL\|x-y\|\|w\|\). Taking \(w\) equal to the score difference and cancelling its norm, with the zero case separate, proves the exact \(NL\) Lipschitz constant. This is the content of \texttt{inner\_\allowbreak{}gradient\_\allowbreak{}potential}, \texttt{sum\_\allowbreak{}joint\_\allowbreak{}norm\_\allowbreak{}mul} and the final norm comparison.




\textbf{Correspondence with the displayed declarations.}


\begin{itemize}

\item \texttt{inner\_\allowbreak{}gradient\_\allowbreak{}potential}: Chain rule for the lifted potential, stated as an inner-product identity.
\item \texttt{sum\_\allowbreak{}joint\_\allowbreak{}norm\_\allowbreak{}mul}: Cauchy--Schwarz in the replicated pair coordinates.
\item \texttt{potential\_\allowbreak{}gradient\_\allowbreak{}lipschitz}: Smoothness scales by \texttt{N}, exactly as strong convexity does.
\end{itemize}
\subsection{\texorpdfstring{The two gradient blocks}{The two gradient blocks}}


For a unit parameter basis vector \(e_j\), all \(A_i e_j\) equal the matching joint parameter vector. For a latent basis vector belonging to particle \(i\), \(A_qe_{i,j}=0\) for \(q\ne i\), and \(A_i e_{i,j}=\sqrt N e_j\). Insert these basis vectors into the directional gradient formula. The Euclidean inner product reads off coordinates, yielding

\[\nabla_\theta V_N(z)=\sum_i\nabla_\theta U(\theta,\sqrt Nz_i),\qquad (\nabla V_N(z))_i=\sqrt N\nabla_xU(\theta,\sqrt Nz_i).\]

This calculation supplies the parameter average and ordinary physical-particle drift after the step size is divided by \(N\).




\textbf{Correspondence with the displayed declarations.}


\begin{itemize}

\item \texttt{gradient\_\allowbreak{}potential\_\allowbreak{}parameter}: The common parameter derivative is the sum of the pair derivatives.
\item \texttt{gradient\_\allowbreak{}potential\_\allowbreak{}particle}: The scaled particle derivative is \texttt{sqrt N} times the pair derivative.
\end{itemize}
\subsection{\texorpdfstring{The parameter projection is a contraction}{The parameter projection is a contraction}}


The continuous linear map \(p_\theta z=\theta\) simply selects the parameter coordinates. Expanding squared norms gives \(\|p_\theta z\|^2\le\|z\|^2\), since the omitted latent squared coordinates are nonnegative. Both norms are nonnegative, so taking square roots proves \(\|p_\theta z\|\le\|z\|\). The definition of operator norm then gives \(\|p_\theta\|\le1\). The coordinate evaluation theorem identifies each output coordinate directly. These inequalities permit quadratic transport to control the parameter RMS.




\textbf{Correspondence with the displayed declarations.}


\begin{itemize}

\item \texttt{parameter\allowbreak{}CLM}: Project the scaled state to its common parameter vector.
\item \texttt{norm\_\allowbreak{}parameter\allowbreak{}CLM\_\allowbreak{}apply\_\allowbreak{}le}: Forgetting the particle coordinates is a contraction.
\item \texttt{norm\_\allowbreak{}parameter\allowbreak{}CLM\_\allowbreak{}le\_\allowbreak{}one}: The parameter projection has operator norm at most one.
\end{itemize}
\subsection{\texorpdfstring{One ULA step is exactly one physical IPLA step}{One ULA step is exactly one physical IPLA step}}


Define the scaled update with physical step \(h\) by

\[z'=z-\frac hN\nabla V_N(z)+\sqrt{2h/N}\,\xi.\]

Its parameter coordinate is \(\theta'=\theta-(h/N)\sum_i\nabla_\theta U(\theta,x_i)+\sqrt{2h/N}\,\xi_\theta\). For \(N>0\), multiply the \(i\)-th scaled latent coordinate by \(\sqrt N\). The drift coefficient becomes \((\sqrt N)^2h/N=h\), and \(\sqrt N\sqrt{2h/N}=\sqrt{2h}\). Consequently

\[x_i'=x_i-h\nabla_xU(\theta,x_i)+\sqrt{2h}\,\xi_i.\]

The square-root identity follows by multiplying radicands and cancelling positive \(N\); Lean's real square-root convention makes the algebraic identity valid for the stated arbitrary real \(h\) as well.




The score-block lemmas are obtained by testing the derivative against single basis vectors: a parameter basis vector is replicated across all pairs, while a latent basis vector appears in one pair with factor \(\sqrt N\). In the particle update, multiplying the scaled coordinate by \(\sqrt N\) turns the drift coefficient \((\gamma/N)\sqrt N\sqrt N\) into \(\gamma\), and the noise coefficient into \(\sqrt{2\gamma}\). The parameter block keeps the averaged force and noise coefficient \(\sqrt{2\gamma/N}\). Positive \(N\) protects exactly these unscaling cancellations.




\textbf{Correspondence with the displayed declarations.}


\begin{itemize}

\item \texttt{scaled\allowbreak{}Update}: One IPLA step represented in the scaled joint coordinates.
\item \texttt{scaled\allowbreak{}Update\_\allowbreak{}parameter}: The common parameter update is exactly the average-gradient IPLA update.
\item \texttt{scaled\allowbreak{}Update\_\allowbreak{}particle}: Unscaling a particle coordinate gives its usual IPLA update, including the \texttt{sqrt (2*h)} noise coefficient.
\end{itemize}
\section{\texorpdfstring{The quadratic upper model}{The quadratic upper model}}
\label{proof-module-UnadjustedLangevin-Analysis-Smooth-lean}
{\small\color{gray}Source module: \texttt{Unadjusted\allowbreak{}Langevin/\allowbreak{}Analysis/\allowbreak{}Smooth.\allowbreak{}lean}.\par}


We next obtain the complementary upper bound from a Lipschitz gradient. This argument uses only first derivatives: restrict the potential to a segment, subtract its affine approximation and the proposed quadratic error, and show that the resulting scalar function decreases.



\texttt{smooth\_\allowbreak{}upper\_\allowbreak{}bound} concerns the gradient of the same differentiable potential occurring in its function values. The argument uses only first derivatives and a Lipschitz gradient; the scalar line restriction and its interval integral provide the quadratic remainder, with no Hessian assumption.

\subsection{\texorpdfstring{The descent lemma}{The descent lemma}}


Let \(H\) be a real Hilbert space and \(L\ge0\). Suppose \(f:H\to\mathbb R\) is differentiable and

\[\|\nabla f(u)-\nabla f(v)\|\le L\|u-v\|\quad(u,v\in H).\]

Then for every \(x,y\),

\[f(y)\le f(x)+\langle\nabla f(x),y-x\rangle+\frac L2\|y-x\|^2.\]

Set \(d=y-x\), \(a=\langle\nabla f(x),d\rangle\), \(b=\|d\|^2\), and introduce the corrected segment function

\[q(t)=f(x+td)-f(x)-ta-\frac L2t^2b.\]

It suffices to show \(q(1)\le q(0)=0\).



The segment has derivative \(d\). The chain rule therefore gives, for every real \(t\),

\[\frac{d}{dt}f(x+td)=\langle\nabla f(x+td),d\rangle.\]

Subtract the derivatives of the affine and quadratic terms:

\[q'(t)=\langle\nabla f(x+td),d\rangle-a-Ltb.\]

For \(t\) in the interior of \([0,1]\), we have \(t\ge0\). Lipschitz continuity gives

\[\|\nabla f(x+td)-\nabla f(x)\|\le L\|td\|=Lt\|d\|.\]

By Cauchy–Schwarz,

\[\langle\nabla f(x+td)-\nabla f(x),d\rangle\le\|\nabla f(x+td)-\nabla f(x)\|\,\|d\|\le Lt\|d\|^2.\]

Substituting \(a\) and \(b\) shows that \(q'(t)\le0\).



The derivative formula also makes \(q\) continuous on \([0,1]\) and differentiable in its interior. The derivative criterion for monotonicity now implies that \(q\) is nonincreasing on this interval.



Compare the two endpoints:

\[q(1)\le q(0),\qquad f(x+d)-f(x)-a-\frac L2b\le0.\]

Finally \(x+d=x+(y-x)=y\). Substitute the definitions of \(a,b,d\) and rearrange to obtain the quadratic upper bound.

\section{\texorpdfstring{Supporting planes and strong monotonicity}{Supporting planes and strong monotonicity}}
\label{proof-module-UnadjustedLangevin-StrongConvexity-lean}
{\small\color{gray}Source module: \texttt{Unadjusted\allowbreak{}Langevin/\allowbreak{}Strong\allowbreak{}Convexity.\allowbreak{}lean}.\par}


We begin with the deterministic geometry behind synchronous Langevin coupling. Throughout this part, \(\langle\cdot,\cdot\rangle\) and \(\|\cdot\|\) are the real inner product and its norm. Strong convexity uses the normalization \(m/2\): the function \(v\mapsto f(v)-\frac m2\|v\|^2\) is convex.



The formal normalization is \texttt{Strong\allowbreak{}Convex\allowbreak{}On Set.\allowbreak{}univ m f}, meaning convexity after subtracting \(m\|x\|^2/2\). \texttt{fderiv \ensuremath{\mathbb R} f x} is the continuous linear derivative; \texttt{gradient f x} is its Hilbert-space representative, so \(Df(x)[v]=\langle\nabla f(x),v\rangle\). The proofs explicitly move between these two representations.

\subsection{\texorpdfstring{The supporting-\allowbreak{}plane inequality}{The supporting-plane inequality}}


Let \(H\) be a real inner-product space and let \(f:H\to\mathbb R\) be differentiable and convex on all of \(H\). For any \(x,y\in H\),

\[f(x)+Df(x)[y-x]\le f(y).\]

Restrict \(f\) to the affine line \(\ell(t)=x+t(y-x)\). Since affine composition preserves convexity, \(\varphi=f\circ\ell\) is convex on \(\mathbb R\), with \(\varphi(0)=f(x)\) and \(\varphi(1)=f(y)\).



The derivative of the line is its constant direction:

\[\ell'(0)=y-x.\]

By the chain rule, the derivative of the restricted function is

\[\varphi'(0)=Df(x)[y-x].\]

For a differentiable convex function, the derivative at the left endpoint is bounded by the secant slope to a later point. Applying this at \(0<1\) gives

\[Df(x)[y-x]\le\frac{\varphi(1)-\varphi(0)}{1-0}.\]

Evaluate the two endpoints and rearrange:

\[Df(x)[y-x]\le f(y)-f(x),\qquad f(x)+Df(x)[y-x]\le f(y).\]\subsection{\texorpdfstring{Strong monotonicity of the gradient}{Strong monotonicity of the gradient}}


Let \(H\) be a real Hilbert space. If \(f:H\to\mathbb R\) is differentiable and \(m\)-strongly convex, where \(m\in\mathbb R\), then for every \(x,y\in H\),

\[m\|x-y\|^2\le\langle\nabla f(x)-\nabla f(y),x-y\rangle.\]

Remove the strong-convexity quadratic. Set

\[g(v)=f(v)-\frac m2\|v\|^2.\]

The definition of strong convexity says precisely that \(g\) is convex.



The squared norm has derivative \(D\|v\|^2[u]=2\langle v,u\rangle\). Consequently \(g\) is differentiable and

\[Dg(v)[u]=Df(v)[u]-m\langle v,u\rangle.\]

Apply the supporting-plane inequality in both directions, and substitute the derivative just computed:

\[g(x)+Dg(x)[y-x]\le g(y),\qquad g(y)+Dg(y)[x-y]\le g(x).\]

Express the derivative of \(f\) using its gradient. The two inequalities become

\[\begin{aligned}f(x)-\tfrac m2\|x\|^2+\langle\nabla f(x),y-x\rangle-m\langle x,y-x\rangle&\le f(y)-\tfrac m2\|y\|^2,\\ f(y)-\tfrac m2\|y\|^2+\langle\nabla f(y),x-y\rangle-m\langle y,x-y\rangle&\le f(x)-\tfrac m2\|x\|^2.\end{aligned}\]

Add these inequalities. The function values and quadratic values cancel. Using \(\langle x-y,x-y\rangle=\|x-y\|^2\), what remains is

\[0\le\langle\nabla f(x)-\nabla f(y),x-y\rangle-m\|x-y\|^2.\]

This is the required monotonicity inequality.

\subsection{\texorpdfstring{The quadratic lower model}{The quadratic lower model}}


Let \(H\) be a real Hilbert space, \(f:H\to\mathbb R\) differentiable, and \(m\in\mathbb R\). If \(f\) is \(m\)-strongly convex on all of \(H\), then for every \(x,y\in H\),

\[f(x)+\langle\nabla f(x),y-x\rangle+\frac m2\|y-x\|^2\le f(y).\]

Set \(g(v)=f(v)-\frac m2\|v\|^2\). Strong convexity makes \(g\) convex.



Differentiate this shifted potential:

\[Dg(v)[u]=Df(v)[u]-m\langle v,u\rangle.\]

In particular, \(g\) is differentiable on \(H\).



Its supporting-plane inequality at \(x\), together with \(Df(x)[u]=\langle\nabla f(x),u\rangle\), yields

\[f(x)-\frac m2\|x\|^2+\langle\nabla f(x),y-x\rangle-m\langle x,y-x\rangle\le f(y)-\frac m2\|y\|^2.\]

Collect the quadratic correction by the identity

\[\|y\|^2-\|x\|^2-2\langle x,y-x\rangle=\|y-x\|^2.\]

Moving terms gives the desired lower model.

\section{\texorpdfstring{A first Euler-\allowbreak{}step estimate}{A first Euler-step estimate}}
\label{proof-module-UnadjustedLangevin-GradientStep-lean}
{\small\color{gray}Source module: \texttt{Unadjusted\allowbreak{}Langevin/\allowbreak{}Gradient\allowbreak{}Step.\allowbreak{}lean}.\par}


Before using the full lower and upper models together, it is useful to record the estimate obtained directly from strong monotonicity and a Lipschitz bound. This estimate applies to arbitrary vectors \(g_x,g_y\); they need not arise as gradients.



The formal inputs \texttt{gx} and \texttt{gy} are arbitrary vectors satisfying the displayed monotonicity and Lipschitz inequalities. Thus \texttt{gradient\_\allowbreak{}step\_\allowbreak{}sq\_\allowbreak{}bound} is an algebraic Hilbert-space result; the potential's actual gradients are substituted later. \texttt{gradient\_\allowbreak{}step\_\allowbreak{}contract} takes a nonnegative square root and does not require a separate sign assumption on the scalar coefficient.

\subsection{\texorpdfstring{The squared Euler-\allowbreak{}step bound}{The squared Euler-step bound}}


Let \(x,y,g_x,g_y\) lie in a real inner-product space, let \(m\in\mathbb R\), \(L\ge0\), and \(h\ge0\). Assume

\[m\|x-y\|^2\le\langle g_x-g_y,x-y\rangle,\qquad\|g_x-g_y\|\le L\|x-y\|.\]

Then

\[\|(x-hg_x)-(y-hg_y)\|^2\le(1-2hm+h^2L^2)\|x-y\|^2.\]

Both sides of the Lipschitz estimate are nonnegative, so squaring preserves the inequality:

\[\|g_x-g_y\|^2\le L^2\|x-y\|^2.\]

Multiply strong monotonicity by \(2h\ge0\), and multiply the squared Lipschitz estimate by \(h^2\ge0\):

\[2hm\|x-y\|^2\le2h\langle g_x-g_y,x-y\rangle,\qquad h^2\|g_x-g_y\|^2\le h^2L^2\|x-y\|^2.\]

Write \(d=x-y\), \(v=g_x-g_y\). Expanding the squared norm and applying the preceding estimates gives

\[\begin{aligned}\|(x-hg_x)-(y-hg_y)\|^2&=\|d-hv\|^2\\&=\|d\|^2-2h\langle v,d\rangle+h^2\|v\|^2\\&\le(1-2hm+h^2L^2)\|d\|^2.\end{aligned}\]\subsection{\texorpdfstring{Taking the square root}{Taking the square root}}


Let \(H\) be a real inner-product space, \(x,y,g_x,g_y\in H\), \(m\in\mathbb R\), \(L\ge0\), and \(h\ge0\). Assume

\[m\|x-y\|^2\le\langle g_x-g_y,x-y\rangle,\qquad\|g_x-g_y\|\le L\|x-y\|.\]

Then

\[\|(x-hg_x)-(y-hg_y)\|\le\sqrt{1-2hm+h^2L^2}\,\|x-y\|.\]

No separate sign assumption on \(1-2hm+h^2L^2\) is required; Lean uses the real square root, which is zero on nonpositive inputs.



Apply the preceding squared estimate. With \(A=\|(x-hg_x)-(y-hg_y)\|\), \(r=\|x-y\|\), and \(c=1-2hm+h^2L^2\), it says \(A^2\le cr^2\).



Since \(A\ge0\), take square roots and use \(r^2\ge0\):

\[A\le\sqrt{cr^2}=\sqrt c\,\sqrt{r^2}=\sqrt c\,r.\]

The square-root product identity is valid here even without a sign hypothesis on \(c\), because the second factor \(r^2\) is nonnegative.

\section{\texorpdfstring{The Langevin update and cancellation of shared noise}{The Langevin update and cancellation of shared noise}}
\label{proof-module-UnadjustedLangevin-ULAStep-lean}
{\small\color{gray}Source module: \texttt{Unadjusted\allowbreak{}Langevin/\allowbreak{}ULAStep.\allowbreak{}lean}.\par}


We now work in finite-dimensional Euclidean space. A ULA step adds the same scaled noise vector to a gradient step. When two starting points use this same noise, it cancels exactly from their difference. The deterministic estimates can therefore be applied pointwise, before introducing any probability law.



\texttt{ula\allowbreak{}Update} is the pointwise Euler map. Its definition accepts arbitrary vector noise; Gaussian laws enter later in realization theorems. \texttt{ula\allowbreak{}Update\_\allowbreak{}sub} uses the same innovation on both sides, so that contribution cancels exactly. The monotonicity of the actual gradient is supplied by the strong-convexity theorem when applying the deterministic step bound.

\subsection{\texorpdfstring{The Euclidean state space}{The Euclidean state space}}


For a natural number \(d\), write \(E_d=\mathbb R^d\) with its Euclidean inner product and norm. Lean represents its vectors as real-valued functions on \(\operatorname{Fin}(d)\), the set of indices \(0,\ldots,d-1\); when \(d=0\), this index set is empty.



Introduce the abbreviation

\[E_d:=\operatorname{EuclideanSpace}_{\mathbb R}(\operatorname{Fin}(d)).\]

This equips the coordinate functions with the Euclidean norm \(\|x\|=(\sum_{i=0}^{d-1}x_i^2)^{1/2}\), so the subsequent gradient and distance estimates use the usual geometry of \(\mathbb R^d\).

\subsection{\texorpdfstring{One unadjusted Langevin step}{One unadjusted Langevin step}}


For a potential \(U:E_d\to\mathbb R\), a real step size \(h\), a state \(x\in E_d\), and a noise vector \(z\in E_d\), define

\[T_h^z(x)=x-h\nabla U(x)+\sqrt{2h}\,z.\]

The definition itself imposes no sign condition on \(h\); later contraction results assume \(h\ge0\).



The update consists of the current state, the drift increment \(-h\nabla U(x)\), and the noise increment \(\sqrt{2h}\,z\).

\subsection{\texorpdfstring{Shared noise cancels}{Shared noise cancels}}


For every potential \(U:E_d\to\mathbb R\), \(h\in\mathbb R\), and vectors \(x,y,z\in E_d\),

\[T_h^z(x)-T_h^z(y)=(x-h\nabla U(x))-(y-h\nabla U(y)).\]

Expand the two updates. The two copies of \(\sqrt{2h}\,z\) subtract to zero:

\[(x-h\nabla U(x)+\sqrt{2h}\,z)-(y-h\nabla U(y)+\sqrt{2h}\,z)=(x-h\nabla U(x))-(y-h\nabla U(y)).\]\subsection{\texorpdfstring{The initial shared-\allowbreak{}noise estimate}{The initial shared-noise estimate}}


Suppose \(U:E_d\to\mathbb R\) is differentiable and \(m\)-strongly convex, \(\nabla U\) is \(L\)-Lipschitz with \(L\ge0\), and \(h\ge0\). For every \(x,y,z\in E_d\),

\[\|T_h^z(x)-T_h^z(y)\|\le\sqrt{1-2hm+h^2L^2}\,\|x-y\|.\]

For \(h>0\) this factor is strictly below one whenever \(hL^2<2m\).



Cancel the common noise. The left-hand side becomes the distance between the two deterministic gradient steps.



Strong convexity gives

\[m\|x-y\|^2\le\langle\nabla U(x)-\nabla U(y),x-y\rangle,\]

and the Lipschitz assumption gives \(\|\nabla U(x)-\nabla U(y)\|\le L\|x-y\|\). Insert these two inequalities into the Euler-step square-root bound.

\section{\texorpdfstring{The sharp deterministic contraction}{The sharp deterministic contraction}}
\label{proof-module-UnadjustedLangevin-SharpContraction-lean}
{\small\color{gray}Source module: \texttt{Unadjusted\allowbreak{}Langevin/\allowbreak{}Sharp\allowbreak{}Contraction.\allowbreak{}lean}.\par}


The preceding square-root factor uses strong monotonicity and Lipschitz continuity separately. The quadratic lower and upper models contain more information. We combine them into the interpolation inequality, then interpret that inequality as a ball containing the gradient difference. This yields the two sharp factors \(1-mh\) and \(Lh-1\), with the change of regime at \(h=2/(m+L)\).



The proof separates the two regimes for the maximum contraction factor \(\max(|1-mh|,|1-Lh|)\). Each scalar bound follows after the strong monotonicity and cocoercivity estimates have supplied the needed inner-product control. The final norm theorem removes the square by nonnegativity, retaining the precise step-regime hypotheses.

\subsection{\texorpdfstring{Cocoercivity from two quadratic models}{Cocoercivity from two quadratic models}}


Let \(H\) be a real inner-product space, \(f:H\to\mathbb R\), \(g:H\to H\), and \(K\ge0\). Suppose that for all \(a,b\in H\),

\[f(a)+\langle g(a),b-a\rangle\le f(b)\le f(a)+\langle g(a),b-a\rangle+\frac K2\|b-a\|^2.\]

Then for all \(x,y\),

\[\|g(x)-g(y)\|^2\le K\langle g(x)-g(y),x-y\rangle.\]

The proof includes \(K=0\), and does not separately assume that \(g\) is a gradient.



Fix \(a,b\), \(t\ge0\), and set \(v=g(a)-g(b)\), \(z=a-tv\). Compare the lower model based at \(b\) with the upper model based at \(a\), both evaluated at \(z\):

\[f(b)+\langle g(b),z-b\rangle\le f(z)\le f(a)+\langle g(a),z-a\rangle+\frac K2\|z-a\|^2.\]

The displacements are \(z-a=-tv\) and \(z-b=(a-b)-tv\). Substitute them into the two models. Since \(t\ge0\), the quadratic term is \(\|z-a\|^2=t^2\|v\|^2\).



The difference of the linear terms satisfies

\[\langle g(a),v\rangle-\langle g(b),v\rangle=\|v\|^2.\]

Eliminating \(f(z)\) therefore gives

\[\left(t-\frac K2t^2\right)\|g(a)-g(b)\|^2\le f(a)-f(b)-\langle g(b),a-b\rangle.\]

Apply this estimate to \((a,b)=(x,y)\) and \((y,x)\), with the same \(t\ge0\). Adding will cancel the function values.



The two gradient norms are equal and the reversed displacement changes the sign of its inner product. Thus, writing \(A=\|g(x)-g(y)\|^2\) and \(B=\langle g(x)-g(y),x-y\rangle\), we have for every \(t\ge0\),

\[(2t-Kt^2)A\le B.\]

First consider \(K=0\). Taking \(t=1\) gives \(2A\le B\). As \(A\ge0\), this also shows \(B\ge0\).



If the desired conclusion \(A\le0\) were false, then \(A>0\). We could choose

\[t=\frac{B+1}{2A}\ge0.\]

The estimate \(2tA\le B\) would then force \(B+1\le B\).



To justify that substitution, \(A>0\) implies \(\|g(x)-g(y)\|\ne0\), so its square can be cancelled from the denominator. Hence indeed

\[2\frac{B+1}{2A}A=B+1.\]

The contradiction \(B+1\le B\) excludes \(A>0\). Thus \(A=0\), proving the claim when \(K=0\).



Now suppose \(K>0\). Choose \(t=K^{-1}\) in the common estimate:

\[\frac AK\le B.\]

Multiplying by \(K\) gives \(A\le KB\), as required.

\subsection{\texorpdfstring{The strong-\allowbreak{}smooth interpolation inequality}{The strong-smooth interpolation inequality}}


Let \(H\) be a real Hilbert space, \(f:H\to\mathbb R\) differentiable and \(m\)-strongly convex, and let \(\nabla f\) be \(L\)-Lipschitz, where \(L\ge0\) and \(m\le L\). Then for every \(x,y\),

\[mL\|x-y\|^2+\|\nabla f(x)-\nabla f(y)\|^2\le(m+L)\langle\nabla f(x)-\nabla f(y),x-y\rangle.\]

The inequality holds also at the boundary \(m=L\).



Subtract the strong-convexity quadratic and its gradient:

\[F(v)=f(v)-\frac m2\|v\|^2,\qquad G(v)=\nabla f(v)-mv.\]

The quadratic lower model for \(f\) becomes an ordinary supporting plane for \(F\):

\[F(a)+\langle G(a),b-a\rangle\le F(b).\]

This follows by expanding the squared distance and collecting the terms involving \(m\).



Apply the smooth upper model to \(f\), and substitute the same definitions of \(F\) and \(G\).



Expanding the quadratic terms shows that the upper curvature decreases from \(L\) to \(L-m\):

\[F(b)\le F(a)+\langle G(a),b-a\rangle+\frac{L-m}{2}\|b-a\|^2.\]

Since \(L-m\ge0\), cocoercivity for these two models gives

\[\|G(x)-G(y)\|^2\le(L-m)\langle G(x)-G(y),x-y\rangle.\]

Write \(d=x-y\) and \(v=\nabla f(x)-\nabla f(y)\). The shifted gradient difference is exactly

\[G(x)-G(y)=v-md.\]

Expand the cocoercivity inequality using this identity:

\[\|v\|^2-2m\langle v,d\rangle+m^2\|d\|^2\le(L-m)\bigl(\langle v,d\rangle-m\langle d,d\rangle\bigr).\]

Use \(\langle d,d\rangle=\|d\|^2\) and collect terms. This gives

\[\|v\|^2+mL\|d\|^2\le(m+L)\langle v,d\rangle.\]

When \(m=L\), the cocoercivity constant above is zero, forcing \(G(x)=G(y)\); consequently \(\nabla f(x)-\nabla f(y)=m(x-y)\). Thus no division by \(L-m\) or limiting argument is needed.

\subsection{\texorpdfstring{The gradient difference lies in a ball}{The gradient difference lies in a ball}}


Let \(H\) be a real inner-product space, \(x,y,g_x,g_y\in H\), and \(m,L\in\mathbb R\) with \(m\le L\). Suppose

\[mL\|x-y\|^2+\|g_x-g_y\|^2\le(m+L)\langle g_x-g_y,x-y\rangle.\]

Then

\[\left\|(g_x-g_y)-\frac{m+L}{2}(x-y)\right\|\le\frac{L-m}{2}\|x-y\|.\]

The proposed radius is nonnegative because \(m\le L\). Both sides of the desired norm bound are therefore nonnegative, so it is enough to prove the corresponding squared inequality.



Set \(c=(m+L)/2\), \(r=(L-m)/2\), \(d=x-y\), and \(v=g_x-g_y\). Completing the square in the interpolation inequality gives

\[\begin{aligned}\|v-cd\|^2&=\|v\|^2-(m+L)\langle v,d\rangle+c^2\|d\|^2\\&\le(c^2-mL)\|d\|^2=r^2\|d\|^2.\end{aligned}\]

Taking nonnegative square roots proves the claim. If \(m=L\), the radius is zero and \(v=md\).

\subsection{\texorpdfstring{The centered Euler-\allowbreak{}step bound}{The centered Euler-step bound}}


Let \(H\) be a real inner-product space, \(m,L,h\in\mathbb R\) with \(m\le L\) and \(h\ge0\), and let \(x,y,g_x,g_y\in H\) satisfy

\[mL\|x-y\|^2+\|g_x-g_y\|^2\le(m+L)\langle g_x-g_y,x-y\rangle.\]

Then

\[\|(x-hg_x)-(y-hg_y)\|\le\left(\left|1-h\frac{m+L}{2}\right|+h\frac{L-m}{2}\right)\|x-y\|.\]

The preceding ball bound gives \(\|v-cd\|\le r\|d\|\), where \(d=x-y\), \(v=g_x-g_y\), \(c=(m+L)/2\), and \(r=(L-m)/2\).



Decompose the Euler difference into its centered part and its deviation:

\[d-hv=(1-hc)d-h(v-cd).\]

The triangle inequality and homogeneity of the norm give

\[\|d-hv\|\le|1-hc|\,\|d\|+h\|v-cd\|,\]

where \(|-h|=h\) because \(h\ge0\).



Multiply the ball bound by \(h\ge0\) and substitute it:

\[\|d-hv\|\le\bigl(|1-hc|+hr\bigr)\|d\|.\]

For \(m=L\), this reduces to the exact scalar factor \(|1-mh|\).

\subsection{\texorpdfstring{The sharp small-\allowbreak{}step factor}{The sharp small-step factor}}


Let \(H\) be a real inner-product space, \(x,y,g_x,g_y\in H\), and \(m,L,h\in\mathbb R\). Suppose \(m\le L\), \(m+L>0\), \(0\le h\le2/(m+L)\), and

\[mL\|x-y\|^2+\|g_x-g_y\|^2\le(m+L)\langle g_x-g_y,x-y\rangle.\]

Then

\[\|(x-hg_x)-(y-hg_y)\|\le(1-mh)\|x-y\|.\]

Apply the centered bound. The step-size assumption and \(m+L>0\) imply \(h(m+L)\le2\), hence \(1-h(m+L)/2\ge0\).



Remove the absolute value with this sign and simplify its coefficient:

\[\left|1-h\frac{m+L}{2}\right|+h\frac{L-m}{2}=1-h\frac{m+L}{2}+h\frac{L-m}{2}=1-mh.\]\subsection{\texorpdfstring{The sharp large-\allowbreak{}step factor}{The sharp large-step factor}}


Let \(H\) be a real inner-product space, \(x,y,g_x,g_y\in H\), and \(m,L,h\in\mathbb R\). Suppose \(m\le L\), \(m+L>0\), \(h\ge0\), \(2/(m+L)\le h\), and

\[mL\|x-y\|^2+\|g_x-g_y\|^2\le(m+L)\langle g_x-g_y,x-y\rangle.\]

Then

\[\|(x-hg_x)-(y-hg_y)\|\le(Lh-1)\|x-y\|.\]

This is a distance bound; it is a strict contraction only when its coefficient is below one.



Apply the centered estimate. Now the step-size assumption gives \(2\le h(m+L)\), so \(1-h(m+L)/2\le0\).



Replace the absolute value by the negative of its argument. The coefficient becomes

\[\left|1-h\frac{m+L}{2}\right|+h\frac{L-m}{2}=h\frac{m+L}{2}-1+h\frac{L-m}{2}=Lh-1.\]\subsection{\texorpdfstring{Small-\allowbreak{}step contraction for a potential}{Small-step contraction for a potential}}


Let \(f:H\to\mathbb R\) be differentiable and \(m\)-strongly convex on a real Hilbert space, with an \(L\)-Lipschitz gradient, \(L\ge0\), \(m\le L\), and \(m+L>0\). If \(0\le h\le2/(m+L)\), then for every \(x,y\),

\[\|(x-h\nabla f(x))-(y-h\nabla f(y))\|\le(1-mh)\|x-y\|.\]

The quadratic models imply the strong-smooth interpolation inequality for \(g_x=\nabla f(x)\) and \(g_y=\nabla f(y)\). Apply the small-step Euler bound to these two vectors. Its hypotheses are exactly the stated curvature and step-size assumptions.

\subsection{\texorpdfstring{Large-\allowbreak{}step contraction for a potential}{Large-step contraction for a potential}}


Let \(H\) be a real Hilbert space and \(f:H\to\mathbb R\) differentiable and \(m\)-strongly convex, with an \(L\)-Lipschitz gradient. Suppose \(L\ge0\), \(m\le L\), \(m+L>0\), \(h\ge0\), and \(2/(m+L)\le h\). Then for every \(x,y\in H\),

\[\|(x-h\nabla f(x))-(y-h\nabla f(y))\|\le(Lh-1)\|x-y\|.\]

Use the potential's interpolation inequality in the large-step Euler bound. This gives the second sharp factor \(Lh-1\).

\subsection{\texorpdfstring{The endpoint-\allowbreak{}factor form}{The endpoint-factor form}}


Let \(f\) be differentiable and \(m\)-strongly convex on a real Hilbert space, with an \(L\)-Lipschitz gradient, \(L\ge0\), \(m\le L\), and \(m+L>0\). For every \(h\ge0\) and every \(x,y\),

\[\|(x-h\nabla f(x))-(y-h\nabla f(y))\|\le\max\{|1-mh|,|1-Lh|\}\,\|x-y\|.\]

If \(0<m\le L\), this factor is below one for \(0<h<2/L\). At \(m=L\), the two endpoint factors coincide.



Split according to whether \(h\le2/(m+L)\). In the first case the small-step theorem bounds the distance by \((1-mh)\|x-y\|\).



Since \(1-mh\le|1-mh|\le\max\{|1-mh|,|1-Lh|\}\), multiplying by \(\|x-y\|\ge0\) gives the desired bound in this case.



In the other case \(h>2/(m+L)\), so the large-step theorem bounds the distance by \((Lh-1)\|x-y\|\).



The elementary absolute-value inequality yields

\[Lh-1=-(1-Lh)\le|1-Lh|.\]

Finally \(|1-Lh|\le\max\{|1-mh|,|1-Lh|\}\). Multiply by the nonnegative distance to finish the second case.

\subsection{\texorpdfstring{Sharp shared-\allowbreak{}noise contraction: small steps}{Sharp shared-noise contraction: small steps}}


Let \(f:E_d\to\mathbb R\) be differentiable and \(m\)-strongly convex, with an \(L\)-Lipschitz gradient, \(L\ge0\), \(m\le L\), and \(m+L>0\). If \(0\le h\le2/(m+L)\), then for all starting states \(x,y\) and every common noise vector \(z\),

\[\|T_h^z(x)-T_h^z(y)\|\le(1-mh)\|x-y\|.\]

The common noise cancels exactly, leaving the distance between the gradient steps from \(x\) and \(y\).



Apply the potential-based small-step estimate. It gives the factor \(1-mh\) independently of the value or distribution of \(z\).

\subsection{\texorpdfstring{Sharp shared-\allowbreak{}noise contraction: large steps}{Sharp shared-noise contraction: large steps}}


Let \(f:E_d\to\mathbb R\) be differentiable and \(m\)-strongly convex, with an \(L\)-Lipschitz gradient, \(L\ge0\), \(m\le L\), and \(m+L>0\). If \(h\ge0\) and \(2/(m+L)\le h\), then for every \(x,y,z\in E_d\),

\[\|T_h^z(x)-T_h^z(y)\|\le(Lh-1)\|x-y\|.\]

Cancel the shared noise, reducing the two ULA updates to their deterministic gradient-step difference.



Apply the potential-based large-step estimate. This yields the factor \(Lh-1\) for every common noise vector.

\section{\texorpdfstring{The local Euler defect}{The local Euler defect}}
\label{proof-module-UnadjustedLangevin-LocalError-lean}
{\small\color{gray}Source module: \texttt{Unadjusted\allowbreak{}Langevin/\allowbreak{}Local\allowbreak{}Error.\allowbreak{}lean}.\par}


Let \(H\) be a real Banach space, with norm \(\|\cdot\|\). All Lipschitz constants \(L\) below are nonnegative. All time integrals are Bochner integrals. For the Langevin application we will take \(H=L^2(P;\mathbb R^d)\); the norms below then measure random variables in mean square. The scalar defect estimate permits a nonnegative real parameter \(d\); its application uses the natural-number dimension, viewed as a real number. Begin with a Lipschitz drift \(g\), a reference curve \(Y\), and its additive noise \(B\). We first derive the error of freezing the drift over one time step directly from the integral equation.



The analytic calculation is conditional on the displayed integral equation and norm bounds. It neither constructs a stationary Langevin diffusion nor proves its score moment. The Gaussian and Brownian noise moments used later are proved separately.



\texttt{local\allowbreak{}Defect} is a vector-valued Bochner interval integral. It exists as a totalized expression before integrability is proved; all norm and linearity steps below explicitly establish the continuous or integrable hypotheses needed for their ordinary mathematical meaning. The interval \(0..h\) is oriented, while the error estimates assume \(h\ge0\).

\subsection{\texorpdfstring{Definition of the local defect}{Definition of the local defect}}


For \(g:H\to H\), \(Y:\mathbb R\to H\), and \(h\in\mathbb R\), define the drift omitted when Euler freezes \(g\) at the initial point:

\[D_h(g,Y):=\int_0^h\bigl(g(Y_t)-g(Y_0)\bigr)\,dt.\]

Integrate the difference between the evolving drift and its frozen value. This defines the defect itself; the estimates below will be consequences of this definition.

\subsection{\texorpdfstring{Separate the exact and frozen drift}{Separate the exact and frozen drift}}


If \(t\mapsto g(Y_t)\) is integrable on the interval between \(0\) and \(h\), then

\[D_h(g,Y)=\int_0^h g(Y_t)\,dt-h\,g(Y_0).\]

Expand the definition. Both the actual drift and the constant function \(g(Y_0)\) are integrable, so linearity separates their difference into two integrals.



The integral of the constant vector is \((h-0)g(Y_0)=h\,g(Y_0)\), yielding the identity.

\subsection{\texorpdfstring{An exact step is Euler minus its defect}{An exact step is Euler minus its defect}}


Assume the preceding drift integrability and the endpoint integral equation

\[Y_h-Y_0=-\int_0^h g(Y_t)\,dt+B_h.\]

Then the exact reference step has the form

\[Y_h=Y_0-h\,g(Y_0)+B_h-D_h(g,Y).\]

Solve the assumed integral equation for \(Y_h\), then replace the defect by the difference between the exact drift integral and its frozen approximation.



Cancel the two frozen-drift terms. The resulting equality is precisely the Euler update with the signed remainder \(-D_h\).

\subsection{\texorpdfstring{Continuity supplies drift integrability}{Continuity supplies drift integrability}}


Let \(H\) be a normed additive group, \(g:H\to H\) be \(L\)-Lipschitz with \(L\ge0\), and \(Y:\mathbb R\to H\) be continuous. For every real \(h\), the function \(t\mapsto g(Y_t)\) is interval integrable from \(0\) to \(h\).



A Lipschitz map is continuous, so \(g\circ Y\) is continuous. A continuous function is integrable on every compact interval, proving the claim.

\subsection{\texorpdfstring{The square-\allowbreak{}root integral}{The square-root integral}}


For every \(h\ge0\),

\[\int_0^h\sqrt t\,dt=\frac23h\sqrt h.\]

Write \(\sqrt t=t^{1/2}\) and use the power-integral formula. The exponent satisfies \(1/2>-1\), so the formula applies at the zero endpoint.



The contribution at zero vanishes. Since \(h\ge0\), the other endpoint is \(h^{1/2+1}=h^{1/2}h=\sqrt h\,h\).



Divide by \(1/2+1=3/2\) and rearrange the factors to obtain \((2/3)h\sqrt h\).

\subsection{\texorpdfstring{Bound a reference increment}{Bound a reference increment}}


Let \(h\ge0\) and suppose that, for every \(0\le s\le h\),

\[\|g(Y_s)\|\le G,\qquad \|B_s\|\le N\sqrt s,\qquad Y_s-Y_0=-\int_0^s g(Y_r)\,dr+B_s.\]

Then, for each \(0\le t\le h\),

\[\|Y_t-Y_0\|\le tG+N\sqrt t.\]

First bound the drift integral on \([0,t]\) by its constant norm bound:

\[\left\|\int_0^t g(Y_s)\,ds\right\|\le G\,|t-0|=Gt,\]

where \(t\ge0\) removes the absolute value.



To justify the constant bound under that integral, note that \(s\in(0,t]\) implies \(0\le s\le t\le h\). Hence the assumed bound on \([0,h]\) applies to every such \(s\).



Insert the integral equation and apply the triangle inequality:

\[\|Y_t-Y_0\|\le\left\|\int_0^t g(Y_s)\,ds\right\|+\|B_t\|.\]

Negation leaves the norm of the drift integral unchanged.



Apply the two bounds to obtain \(Gt+N\sqrt t\), and commute the factors \(G\) and \(t\).

\subsection{\texorpdfstring{Integrate the Lipschitz increment bound}{Integrate the Lipschitz increment bound}}


Assume \(h\ge0\), \(g\) is \(L\)-Lipschitz with \(L\ge0\), and for \(0\le t\le h\) the same integral equation and bounds \(\|g(Y_t)\|\le G\), \(\|B_t\|\le N\sqrt t\) hold. Then

\[\|D_h(g,Y)\|\le L\left(\frac{h^2G}{2}+\frac23h\sqrt h\,N\right).\]

For \(0<t\le h\), Lipschitz continuity gives

\[\|g(Y_t)-g(Y_0)\|\le L\|Y_t-Y_0\|.\]

Multiply the increment estimate by \(L\ge0\) to obtain the pointwise majorant

\[\|g(Y_t)-g(Y_0)\|\le L(tG+N\sqrt t).\]

The scalar majorant \(t\mapsto L(tG+N\sqrt t)\) is continuous, hence integrable on \([0,h]\).



Expand \(D_h\). The norm estimate for a Bochner integral, together with the integrable majorant, gives

\[\|D_h\|\le\int_0^h L(tG+N\sqrt t)\,dt.\]

Pull out constants and split the two integrable summands. Evaluate

\[\int_0^h t\,dt=\frac{h^2}{2},\qquad \int_0^h\sqrt t\,dt=\frac23h\sqrt h.\]

This leaves the two drift and noise contributions in the claimed expression.



Rearrange the scalar factors to place the drift term first and the noise term second.

\subsection{\texorpdfstring{The exact rational constant}{The exact rational constant}}


The numerical factor in the corrected Dalalyan–Karagulyan estimate satisfies

\[\frac76\sqrt2\le\frac{33}{20}=1.65.\]

Use \((\sqrt2)^2=2\) and \(\sqrt2\ge0\). Comparing squares proves the bound, since \(49/18\le1089/400\).

\subsection{\texorpdfstring{The local error with stationary moment scales}{The local error with stationary moment scales}}


Let \(h,d\in\mathbb R\) with \(h,d\ge0\), let \(g\) be \(L\)-Lipschitz with \(L\ge0\), and suppose \(Lh\le2\). Assume that, for every \(0\le t\le h\),

\[Y_t-Y_0=-\int_0^t g(Y_s)\,ds+B_t,\quad \|g(Y_t)\|\le\sqrt{Ld},\quad \|B_t\|\le\sqrt{2d}\sqrt t.\]

Then

\[\boxed{\|D_h(g,Y)\|\le\frac{33}{20}Lh\sqrt{hd}.}\]

In \(L^2\), the drift bound corresponds to a score second moment at most \(Ld\); here that bound remains a hypothesis.



Compare the nonnegative quantities \(h\sqrt{Ld}\) and \(\sqrt2\sqrt{hd}\) by squaring. It suffices to show \(h^2Ld\le2hd\).



Replace each squared square root by its nonnegative argument. Multiply \(Lh\le2\) by \(hd\ge0\), obtaining \(h^2Ld\le2hd\), and thus

\[h\sqrt{Ld}\le\sqrt2\sqrt{hd}.\]

Split products under the square roots and rearrange:

\[h\sqrt h\sqrt{2d}=h\sqrt2\sqrt{hd}.\]

The nonnegativity of \(h\) and \(2\) justifies the square-root product identities.



Apply the two-term local-defect bound with \(G=\sqrt{Ld}\) and \(N=\sqrt{2d}\):

\[\|D_h\|\le L\left(\frac{h^2\sqrt{Ld}}2+\frac23h\sqrt h\sqrt{2d}\right).\]

Bound the first term using \(h\sqrt{Ld}\le\sqrt2\sqrt{hd}\), and rewrite the second with the preceding identity. Since \(1/2+2/3=7/6\),

\[\|D_h\|\le L\frac76\sqrt2\,h\sqrt{hd}.\]

Multiply \((7/6)\sqrt2\le33/20\) by the nonnegative factor \(Lh\sqrt{hd}\). This gives the claimed constant without numerical rounding.

\section{\texorpdfstring{The noise-\allowbreak{}driven iterates}{The noise-driven iterates}}
\label{proof-module-UnadjustedLangevin-Dynamics-lean}
{\small\color{gray}Source module: \texttt{Unadjusted\allowbreak{}Langevin/\allowbreak{}Dynamics.\allowbreak{}lean}.\par}


Write \(T_h(x)=x-h\nabla U(x)\), so one Euler update is \(T_h(x)+\sqrt{2h}\,z\). The noise vectors remain explicit. Common-noise algebra works for any supplied noise sequence; identifying the algorithm with Gaussian ULA also requires independent standard Gaussian noise hypotheses.



\texttt{ula\allowbreak{}Chain} recursively applies \texttt{ula\allowbreak{}Update} to a supplied sequence of vectors. Its zero and successor equations are definitional. \texttt{ula\allowbreak{}Chain\_\allowbreak{}measurable} proves measurable dependence on the sample by induction, combining the current iterate, the Lipschitz gradient's measurability, and the next supplied measurable innovation. No independence or Gaussian-law assumption is needed for this measurability assertion.

\subsection{\texorpdfstring{Definition of the chain}{Definition of the chain}}


Given \(U:\mathbb R^d\to\mathbb R\), \(h\in\mathbb R\), an initial vector \(x_0\), and noise \((\xi_k)\), define

\[X_0=x_0,\qquad X_{k+1}=X_k-h\nabla U(X_k)+\sqrt{2h}\,\xi_k.\]

Start the recursion at the supplied initial vector: \(X_0=x_0\).



Apply the Euler update to \(X_k\) using the next noise vector \(\xi_k\). This determines \(X_{k+1}\).

\subsection{\texorpdfstring{The initial iterate}{The initial iterate}}


For every potential, step, initial vector, and noise sequence, the recursively defined chain satisfies \(X_0=x_0\).



The equality is exactly the initial clause of the recursive definition.

\subsection{\texorpdfstring{The update identity}{The update identity}}


For every \(k\ge0\), the chain satisfies \(X_{k+1}=\operatorname{ulaUpdate}(U,h,X_k,\xi_k)\).



The successor clause of the recursion gives the asserted update formula.

\subsection{\texorpdfstring{Measurability of every iterate}{Measurability of every iterate}}


Let \(\Omega\) be a measurable space. If \(\nabla U\) is \(L\)-Lipschitz, the initial vector \(X_0:\Omega\to\mathbb R^d\) is measurable, and every \(\xi_k:\Omega\to\mathbb R^d\) is measurable, then \(X_n\) is measurable for every \(n\).



A Lipschitz gradient is continuous and therefore measurable.



Induct on \(n\). The initial iterate is measurable by hypothesis. For the successor step, assume \(X_n\) is measurable.



The composition \(\nabla U\circ X_n\) is measurable. Subtracting its scalar multiple from \(X_n\), then adding \(\sqrt{2h}\,\xi_n\), preserves measurability. Hence \(X_{n+1}\) is measurable.

\section{\texorpdfstring{The actual gradient on square-\allowbreak{}integrable random variables}{The actual gradient on square-integrable random variables}}
\label{proof-module-UnadjustedLangevin-LpDrift-lean}
{\small\color{gray}Source module: \texttt{Unadjusted\allowbreak{}Langevin/\allowbreak{}Lp\allowbreak{}Drift.\allowbreak{}lean}.\par}


Let \(P\) be a finite measure on \(\Omega\). The first five results need only a normed additive commutative group \(H\); all Lipschitz constants \(L\) are nonnegative. When \(H\) is a real Banach space, \(L^2(P;H)\) provides the Banach space needed for the local-defect estimate. For a Lipschitz vector field \(g:H\to H\), construct the operator \(\mathcal G[X]=[g\circ X]\), including the case \(g(0)\ne0\). Finiteness of \(P\) ensures the constant offset is square integrable.



The gradient results use a real Hilbert space \(H\). The two score-moment lemmas need only a real inner-product space, without completeness. Whenever \(g=\nabla U\), this is the potential’s actual gradient, and the resulting \(L^2\) recursion is proved to represent the original pointwise ULA chain.



\texttt{lp\allowbreak{}Drift} lifts a Lipschitz pointwise map to \(L^2\), using its linear growth and the finite probability measure to prove integrability. Its representative theorem identifies the lifted map almost everywhere with ordinary composition. All subsequent Hilbert-space estimates concern these quotient classes and are invariant under representative choices.

\subsection{\texorpdfstring{Center a Lipschitz vector field}{Center a Lipschitz vector field}}


In a normed additive group \(H\), if \(g:H\to H\) is \(L\)-Lipschitz, then so is

\[\widetilde g(x):=g(x)-g(0).\]

For arbitrary \(x,y\), the common offset cancels:

\[\|\widetilde g(x)-\widetilde g(y)\|=\|g(x)-g(y)\|\le L\|x-y\|.\]

This is precisely the required Lipschitz inequality.

\subsection{\texorpdfstring{Construct the lifted drift}{Construct the lifted drift}}


Let \(P\) be finite and \(g\) be \(L\)-Lipschitz. For \(X\in L^2(P;H)\), define

\[\mathcal G[X]:=[\widetilde g\circ X]+[g(0)].\]

The centered function satisfies \(\widetilde g(0)=0\) and \(\|\widetilde g(x)\|\le L\|x\|\), so composition produces an \(L^2\) class. Add the constant class \([g(0)]\), which belongs to \(L^2\) because \(P\) is finite.

\subsection{\texorpdfstring{The lift agrees with the original vector field}{The lift agrees with the original vector field}}


For finite \(P\), Lipschitz \(g\), and \(X\in L^2(P;H)\),

\[\mathcal G[X](\omega)=g(X(\omega))\quad\text{for }P\text{-almost every }\omega.\]

Expand \(\mathcal G\) into the centered composition class plus the constant class.



Intersect the full-measure sets on which the representative of the sum, the composition, and the constant class agree with their pointwise formulas.



On this set the value is \((g(X(\omega))-g(0))+g(0)\). Cancel the offset to recover \(g(X(\omega))\).

\subsection{\texorpdfstring{The actual pointwise drift is square integrable}{The actual pointwise drift is square integrable}}


For finite \(P\), Lipschitz \(g\), and \(X\in L^2(P;H)\),

\[g\circ X\in L^2(P;H).\]

The constructed class \(\mathcal G[X]\) has an \(L^2\) representative. Transfer its square integrability to \(g\circ X\) using their almost-everywhere equality.

\subsection{\texorpdfstring{The drift keeps its Lipschitz constant in L\ensuremath{{}^{2}}}{The drift keeps its Lipschitz constant in L²}}


For finite \(P\) and \(L\)-Lipschitz \(g\), the lifted drift is \(L\)-Lipschitz:

\[\|\mathcal G[X]-\mathcal G[Y]\|_{L^2}\le L\|X-Y\|_{L^2}.\]

Fix \(X,Y\). The two constant classes cancel in their drift difference. For the remaining centered compositions, the pointwise Lipschitz estimate passes to the \(L^2\) norm, giving the same constant \(L\).

\subsection{\texorpdfstring{Lift the small-\allowbreak{}step gradient contraction}{Lift the small-step gradient contraction}}


Let \(H\) be a real Hilbert space, \(P\) finite, and \(U:H\to\mathbb R\) differentiable and \(m\)-strongly convex on \(H\), with \(L\)-Lipschitz gradient. Assume \(m\le L\), \(m+L>0\), and \(0\le h\le2/(m+L)\). Write \(\mathcal G[X]=[\nabla U\circ X]\). Then, for all \(X,Y\in L^2(P;H)\),

\[\|(X-h\mathcal G[X])-(Y-h\mathcal G[Y])\|_{L^2}\le(1-mh)\|X-Y\|_{L^2}.\]

It suffices to prove the corresponding norm inequality almost everywhere between representatives: integrating the squared pointwise inequality gives the \(L^2\) bound.



Work on a common full-measure set where the outer difference and both Euler-step differences have their ordinary pointwise subtraction formulas.



Also require the scalar multiples, both lifted gradients, and \(X-Y\) to have their pointwise formulas. The two drift terms are then \(\nabla U(X(\omega))\) and \(\nabla U(Y(\omega))\).



Substitute these representative identities and apply the deterministic sharp contraction to \(x=X(\omega)\), \(y=Y(\omega)\):

\[\|(x-h\nabla U(x))-(y-h\nabla U(y))\|\le(1-mh)\|x-y\|.\]\subsection{\texorpdfstring{Lift the second step-\allowbreak{}size regime}{Lift the second step-size regime}}


Let \(H\) be a real Hilbert space, \(P\) finite, and \(U\) differentiable, \(m\)-strongly convex, with \(L\)-Lipschitz gradient. Assume \(m\le L\), \(m+L>0\), \(h\ge0\), and \(h\ge2/(m+L)\). Then

\[\|(X-h\mathcal G[X])-(Y-h\mathcal G[Y])\|_{L^2}\le(Lh-1)\|X-Y\|_{L^2}.\]

The inequality holds in this whole regime; it is a strict contraction when its factor is less than one.



Reduce the \(L^2\) inequality to its almost-everywhere pointwise counterpart with factor \(Lh-1\).



Choose a full-measure set on which the outer subtraction and each Euler-step subtraction agree with the pointwise vector formulas.



On the same set identify scalar multiplication, both lifted gradient representatives, and the representative of \(X-Y\).



After substitution, invoke the deterministic second-regime estimate at \(X(\omega),Y(\omega)\):

\[\|(x-h\nabla U(x))-(y-h\nabla U(y))\|\le(Lh-1)\|x-y\|.\]\subsection{\texorpdfstring{The local defect for the actual L\ensuremath{{}^{2}} gradient}{The local defect for the actual L² gradient}}


Let \(P\) be finite, \(H\) a real Hilbert space, and \(\nabla U\) be \(L\)-Lipschitz. Let \(Y,B:\mathbb R\to L^2(P;H)\), let \(h,d\in\mathbb R\), and assume \(h,d\ge0\), \(Lh\le2\). For every \(0\le t\le h\), suppose

\[\|\mathcal G[Y_t]\|_{L^2}\le\sqrt{Ld},\quad\|B_t\|_{L^2}\le\sqrt{2d}\sqrt t,\quad Y_t-Y_0=-\int_0^t\mathcal G[Y_s]\,ds+B_t.\]

Then

\[\|D_h(\mathcal G,Y)\|_{L^2}\le\frac{33}{20}Lh\sqrt{hd}.\]

Apply the Banach-space local-defect theorem to \(\mathcal G\). Its required Lipschitz property is the one just proved for the actual lifted gradient; the reference equation and moment-scale bounds are the displayed hypotheses.

\subsection{\texorpdfstring{Identify the drift norm with the score moment}{Identify the drift norm with the score moment}}


For finite \(P\), a real inner-product space \(H\), Lipschitz \(g:H\to H\), and \(X\in L^2(P;H)\),

\[\|\mathcal G[X]\|_{L^2}^2=\int_\Omega\|g(X(\omega))\|^2\,dP(\omega).\]

Write the squared \(L^2\) norm as a self-inner-product and expand the \(L^2\) inner product into the integral of the pointwise inner product.



Almost everywhere, \(\mathcal G[X](\omega)=g(X(\omega))\). Replace the representative under the integral and use \(\langle v,v\rangle=\|v\|^2\).

\subsection{\texorpdfstring{A second-\allowbreak{}moment bound gives the drift bound}{A second-moment bound gives the drift bound}}


Under the preceding finite-measure and inner-product assumptions, if

\[\int_\Omega\|g(X(\omega))\|^2\,dP(\omega)\le C,\]

then

\[\|\mathcal G[X]\|_{L^2}\le\sqrt C.\]

For \(g=\nabla U\) and \(C=Ld\), the moment inequality is an explicit input, not a consequence of this lemma.



The norm-square identity turns the assumed moment bound into \(\|\mathcal G[X]\|^2\le C\). Taking the nonnegative square root yields the conclusion.

\subsection{\texorpdfstring{Define one ULA update in L\ensuremath{{}^{2}}}{Define one ULA update in L²}}


For finite \(P\), a real Hilbert space \(H\), a Lipschitz gradient \(\nabla U\), step \(h\in\mathbb R\), and \(X,Z\in L^2(P;H)\), define

\[\mathsf T_h(X,Z):=X-h\mathcal G[X]+\sqrt{2h}\,Z.\]

Form the drift step and add the scaled noise inside \(L^2\). All terms are already equivalence classes of square-integrable variables, so the vector-space operations produce another such class.

\subsection{\texorpdfstring{The L\ensuremath{{}^{2}} update is the pointwise ULA formula}{The L² update is the pointwise ULA formula}}


For finite \(P\), a real Hilbert space \(H\), Lipschitz \(\nabla U\), and arbitrary \(X,Z\in L^2(P;H)\),

\[\mathsf T_h(X,Z)(\omega)=X(\omega)-h\nabla U(X(\omega))+\sqrt{2h}\,Z(\omega)\quad\text{a.e.}\]

This identity needs no Gaussian or independence assumption on \(Z\).



Expand the definition of \(\mathsf T_h\) into its subtraction, drift, and scaled-noise terms.



Intersect the full-measure sets on which addition, subtraction, both scalar multiples, and the lifted gradient agree with their pointwise formulas.



Substitute these identities. The representative is exactly \(X(\omega)-h\nabla U(X(\omega))+\sqrt{2h}\,Z(\omega)\).

\subsection{\texorpdfstring{Define the genuine L\ensuremath{{}^{2}} ULA chain}{Define the genuine L² ULA chain}}


For a finite measure \(P\), a Lipschitz gradient on \(E_d\), an initial class \(X_0\in L^2(P;E_d)\), and noise classes \(Z_n\in L^2(P;E_d)\), define recursively, with the supplied initial class at index zero,

\[X_{n+1}:=\mathsf T_h(X_n,Z_n).\]

The initial class is the supplied \(X_0\). At the successor index, apply the \(L^2\) ULA update to the preceding iterate and the next noise class \(Z_n\).

\subsection{\texorpdfstring{The initial iterate}{The initial iterate}}


The recursively defined chain satisfies \(X_0=\mathrm{initial}\).



This is the zero branch of the recursion, so both sides reduce to the same class.

\subsection{\texorpdfstring{The successor iterate}{The successor iterate}}


For every \(n\ge0\), the same finite-measure, Lipschitz-gradient, initial-class and noise-class data give

\[X_{n+1}=\mathsf T_h(X_n,Z_n).\]

This is the successor branch of the defining recursion; unfold that branch to obtain the equality.

\subsection{\texorpdfstring{Identify the chain with the original algorithm}{Identify the chain with the original algorithm}}


For finite \(P\), Lipschitz gradient \(\nabla U:E_d\to E_d\), square-integrable initial and noise classes, and every \(n\ge0\),

\[X_n(\omega)=\operatorname{ulaChain}\bigl(U,h,X_0(\omega),(Z_k(\omega))_{k\ge0},n\bigr)\quad\text{for }P\text{-a.e. }\omega.\]

Induct on \(n\). At \(n=0\), both recursions return the same initial representative, so equality holds for every outcome. For the successor case, assume the representatives agree at index \(n\).



On a common full-measure set, combine the one-step representative identity with the induction hypothesis. Replace \(X_n(\omega)\) by the pointwise chain’s preceding iterate and apply the two successor recursion equations. Both next iterates are the same ULA update.

\section{\texorpdfstring{Couplings and quadratic transport}{Couplings and quadratic transport}}
\label{proof-module-UnadjustedLangevin-Coupling-lean}
{\small\color{gray}Source module: \texttt{Unadjusted\allowbreak{}Langevin/\allowbreak{}Coupling.\allowbreak{}lean}.\par}


Let \(H\) be a normed additive group with its Borel measurable structure and a second countable topology. Let \(P\) be a measure on a measurable space \(\Omega\), and write \(\mathcal L_P(X)=X_\#P\). The symbol \(\|Z\|_{2,P}\) denotes the extended \(L^2\) seminorm, which may be infinite. In the probability setting the following definitions give the usual quadratic transport cost.



\texttt{Coupling} stores a joint measure with exact first and second pushforward marginals. \texttt{coupling\allowbreak{}Cost} is its extended \(L^2\) displacement norm, and \texttt{wasserstein2} is the infimum over all such couplings. Finiteness guards precede \texttt{to\allowbreak{}Real} conversions; no existence of an optimal coupling is assumed.

\subsection{\texorpdfstring{Definition of a coupling}{Definition of a coupling}}


A coupling of measures \(\mu,\nu\) on \(H\) is a joint measure \(\gamma\) on \(H\times H\) with those two marginals. Denote the family by \(\Gamma(\mu,\nu)\).



Choose the joint measure \(\gamma\) on pairs \((x,y)\).



Require \((\operatorname{pr}_1)_\#\gamma=\mu\) and \((\operatorname{pr}_2)_\#\gamma=\nu\).

\subsection{\texorpdfstring{Definition of quadratic coupling cost}{Definition of quadratic coupling cost}}


For a joint measure \(\gamma\), define

\[c(\gamma)=\left(\int\|x-y\|^2\,d\gamma(x,y)\right)^{1/2}\in[0,+\infty].\]

Take the extended \(L^2(\gamma)\) norm of the displacement map \((x,y)\mapsto x-y\).

\subsection{\texorpdfstring{Definition of the transport infimum}{Definition of the transport infimum}}


Define the extended quadratic transport distance by

\[W_2(\mu,\nu)=\inf_{\gamma\in\Gamma(\mu,\nu)}c(\gamma).\]

Take the infimum over all couplings without selecting a minimizer.

\subsection{\texorpdfstring{Real conversion of transport distance}{Real conversion of transport distance}}


Define \(W_2^{\mathbb R}(\mu,\nu)=\operatorname{toReal}(W_2(\mu,\nu))\). The comparison lemmas below establish finite upper bounds before using this conversion.



Apply real conversion to the extended distance; for a finite distance this returns its ordinary real value.

\subsection{\texorpdfstring{Every coupling supplies an upper bound}{Every coupling supplies an upper bound}}


For every \(\gamma\in\Gamma(\mu,\nu)\), \(W_2(\mu,\nu)\le c(\gamma)\).



An infimum is at most each member of the family. Apply this to the chosen coupling.

\subsection{\texorpdfstring{The joint law of two random vectors}{The joint law of two random vectors}}


If \(X,Y:\Omega\to H\) are measurable, then \((X,Y)_\#P\) is a coupling of \(X_\#P\) and \(Y_\#P\).



Set \(\gamma=(\omega\mapsto(X(\omega),Y(\omega)))_\#P\).



Compose the pushforward with the first projection. Since \(\operatorname{pr}_1(X(\omega),Y(\omega))=X(\omega)\), the first marginal is \(X_\#P\).



The same calculation with the second projection gives the second marginal \(Y_\#P\).

\subsection{\texorpdfstring{Joint laws from almost everywhere measurable vectors}{Joint laws from almost everywhere measurable vectors}}


If \(X,Y\) are measurable almost everywhere with respect to \(P\), their joint law is still a coupling of their two laws.



Again define \(\gamma=(X,Y)_\#P\).



Apply pushforward composition for almost everywhere measurable maps to identify the first marginal with \(X_\#P\).



Apply its second-projection version to identify the second marginal with \(Y_\#P\).

\subsection{\texorpdfstring{Cost of the joint-\allowbreak{}law coupling}{Cost of the joint-law coupling}}


For measurable \(X,Y:\Omega\to H\),

\[c((X,Y)_\#P)=\|X-Y\|_{2,P}.\]

Expand the two definitions. The cost is the \(L^2\) norm of \((x,y)\mapsto x-y\) under the joint law.



Change variables through the pushforward. The displacement composed with the paired random vector is \(\omega\mapsto X(\omega)-Y(\omega)\), giving the identity.

\subsection{\texorpdfstring{A coupled pair bounds transport}{A coupled pair bounds transport}}


For measurable \(X,Y\),

\[W_2(X_\#P,Y_\#P)\le\|X-Y\|_{2,P}.\]

Use the joint law as an admissible coupling, then replace its cost by the \(L^2\) distance just computed.

\subsection{\texorpdfstring{The almost everywhere measurable transport bound}{The almost everywhere measurable transport bound}}


The same transport comparison holds for almost everywhere measurable \(X,Y\):

\[W_2(X_\#P,Y_\#P)\le\|X-Y\|_{2,P}.\]

The almost everywhere measurable joint-law construction supplies an admissible coupling, whose cost bounds the infimum.



The pushforward change of variables identifies that cost with \(\|X-Y\|_{2,P}\).

\subsection{\texorpdfstring{Finite coupled cost gives a real bound}{Finite coupled cost gives a real bound}}


If \(X,Y\) are measurable and \(X-Y\in L^2(P;H)\), then

\[W_2^{\mathbb R}(X_\#P,Y_\#P)\le\operatorname{toReal}(\|X-Y\|_{2,P}).\]

The extended comparison has a finite upper bound because \(X-Y\in L^2\). Apply monotonicity of real conversion under a finite upper bound.

\subsection{\texorpdfstring{The transport bound for L\ensuremath{{}^{2}} classes}{The transport bound for L² classes}}


For \(X,Y\in L^2(P;H)\), their representatives define laws satisfying

\[W_2(X_\#P,Y_\#P)\le\operatorname{ofReal}(\|X-Y\|_{L^2}).\]

The representative of \(X-Y\) agrees almost everywhere with the difference of the representatives. Their extended \(L^2\) norms agree.



The L\ensuremath{{}^{2}} norm is the real value of the finite extended seminorm. Embedding it back into extended reals recovers that seminorm.



Apply the almost everywhere measurable transport bound to the two representatives and use the equality of displacement seminorms.

\subsection{\texorpdfstring{The real L\ensuremath{{}^{2}}-\allowbreak{}to-\allowbreak{}transport bridge}{The real L²-to-transport bridge}}


For \(X,Y\in L^2(P;H)\),

\[W_2^{\mathbb R}(X_\#P,Y_\#P)\le\|X-Y\|_{L^2}.\]

The preceding extended bound is dominated by the finite number \(\operatorname{ofReal}(\|X-Y\|_{L^2})\), so convert monotonically to real numbers.



Real conversion of the embedding of a nonnegative norm returns the norm itself.

\subsection{\texorpdfstring{Pointwise domination passes to L\ensuremath{{}^{2}}}{Pointwise domination passes to L²}}


Let \(q\ge0\), let \(A\) be almost everywhere strongly measurable, and suppose \(\|A(\omega)\|\le q\|D(\omega)\|\) almost everywhere. Then

\[\|A\|_{2,P}\le q\|D\|_{2,P}.\]

Apply monotonicity and positive homogeneity of the extended \(L^2\) seminorm to the almost everywhere domination. No independence is required.

\subsection{\texorpdfstring{One-\allowbreak{}step error from contraction and a defect}{One-step error from contraction and a defect}}


Suppose \(q\ge0\), \(A\) is almost everywhere strongly measurable, and almost everywhere

\[X'-Y'=A-V,\qquad\|A\|\le q\|X-Y\|.\]

Then

\[\|X'-Y'\|_{2,P}\le q\|X-Y\|_{2,P}+\|V\|_{2,P}.\]

Replace \(X'-Y'\) by \(A-V\) in the seminorm using the almost everywhere error identity.



Minkowski gives \(\|A-V\|_2\le\|A\|_2+\|V\|_2\). Pointwise contraction bounds the first term by \(q\|X-Y\|_2\).

\subsection{\texorpdfstring{One-\allowbreak{}step transport estimate}{One-step transport estimate}}


Under the preceding decomposition and contraction hypotheses, if \(X',Y'\) are measurable, then

\[W_2(X'_\#P,Y'_\#P)\le q\|X-Y\|_{2,P}+\|V\|_{2,P}.\]

The joint-law coupling bounds transport by \(\|X'-Y'\|_2\). Compose this with the synchronous \(L^2\) estimate.

\subsection{\texorpdfstring{The finite-\allowbreak{}moment one-\allowbreak{}step inequality}{The finite-moment one-step inequality}}


Assume \(q\ge0\), \(A\) is almost everywhere strongly measurable, \(X'-Y'=A-V\) and \(\|A\|\le q\|X-Y\|\) almost everywhere. If \(X-Y,V\in L^2(P;H)\), then

\[\operatorname{toReal}(\|X'-Y'\|_{2,P})\le q\operatorname{toReal}(\|X-Y\|_{2,P})+\operatorname{toReal}(\|V\|_{2,P}).\]

The product \(q\|X-Y\|_2\) is finite because both factors are finite.



Its sum with the finite defect norm is finite. Use this finite upper bound to convert the synchronous extended inequality monotonically to real numbers.



Distribute real conversion through the finite sum and the product, giving the stated real inequality.

\section{\texorpdfstring{Passing from a chosen initial coupling to an infimum}{Passing from a chosen initial coupling to an infimum}}
\label{proof-module-UnadjustedLangevin-InitialCoupling-lean}
{\small\color{gray}Source module: \texttt{Unadjusted\allowbreak{}Langevin/\allowbreak{}Initial\allowbreak{}Coupling.\allowbreak{}lean}.\par}


The analytic estimate first uses the cost of one initial coupled pair. The infimum calculation here replaces that cost by \(W_2(\mu,\nu)\) when all finite-cost couplings give the same output laws. \texttt{Product\allowbreak{}Reference.\allowbreak{}lean} proves this common-law property on concrete product extensions; \texttt{Canonical\allowbreak{}ULA.\allowbreak{}lean} supplies their stationary references.



The mathematical infimum here is the infimum in \texttt{wasserstein2}. The proof works with arbitrary finite-cost couplings and approximates that infimum when necessary. This gives the true initial transport term without postulating an optimal transport plan.

\subsection{\texorpdfstring{A uniform affine estimate passes to the infimum}{A uniform affine estimate passes to the infimum}}


Let \(a\ge0\) be finite, \(b\in[0,+\infty]\), and \(\Gamma(\mu,\nu)\ne\varnothing\). If the fixed output measures satisfy

\[W_2(\mu',\nu')\le ac(\gamma)+b\quad\text{for every }\gamma\in\Gamma(\mu,\nu),\]

then \(W_2(\mu',\nu')\le aW_2(\mu,\nu)+b\).



Use the nonempty coupling family and expand transport distances as infima.



Multiplication by finite \(a\ge0\) and addition of \(b\) commute with the nonempty extended-real infimum, including \(a=0\). Thus \(aW_2(\mu,\nu)+b=\inf_\gamma(ac(\gamma)+b)\).



The output distance lies below every member of this last family, hence below its infimum.

\subsection{\texorpdfstring{A finite infimum supplies a finite-\allowbreak{}cost coupling}{A finite infimum supplies a finite-cost coupling}}


If \(W_2(\mu,\nu)<+\infty\), then some \(\gamma\in\Gamma(\mu,\nu)\) has \(c(\gamma)<+\infty\). Attainment of the infimum is not asserted.



The coupling-cost infimum is strictly below \(+\infty\) by the finiteness hypothesis.



An infimum strictly below a threshold has a member strictly below that threshold. Apply this with threshold \(+\infty\).

\subsection{\texorpdfstring{The real uniform-\allowbreak{}coupling estimate}{The real uniform-coupling estimate}}


Let \(a,b\ge0\) be finite and assume \(W_2(\mu,\nu)<+\infty\). If \(W_2(\mu',\nu')\le ac(\gamma)+b\) for every initial coupling, then

\[W_2^{\mathbb R}(\mu',\nu')\le aW_2^{\mathbb R}(\mu,\nu)+b.\]

A finite initial distance supplies a coupling. Apply the extended uniform-coupling lemma to obtain the desired inequality in extended reals.



The product \(aW_2(\mu,\nu)\) is finite.



Its sum with finite \(b\) is finite, so the extended output inequality permits monotone real conversion.



Distribute real conversion through the finite sum and product to obtain the real inequality.

\subsection{\texorpdfstring{Finite-\allowbreak{}cost initial couplings suffice}{Finite-cost initial couplings suffice}}


Let \(a,b\ge0\) be finite. Assume both the initial and output transport distances are finite and that every finite-cost \(\gamma\in\Gamma(\mu,\nu)\) satisfies

\[W_2^{\mathbb R}(\mu',\nu')\le a\operatorname{toReal}(c(\gamma))+b.\]

Then \(W_2^{\mathbb R}(\mu',\nu')\le aW_2^{\mathbb R}(\mu,\nu)+b\).



For each finite-cost coupling, embed the real estimate in extended nonnegative reals. Finiteness of the output distance and coupling cost recovers their original extended values, giving \(W_2(\mu',\nu')\le ac(\gamma)+b\).



Apply the all-couplings theorem. For a finite-cost coupling, the needed extended estimate has just been proved.



For an infinite-cost coupling, split into the cases \(a=0\) and \(a\ne0\).



If \(a=0\), choose one finite-cost coupling supplied by the finite initial distance. Its estimate gives \(W_2(\mu',\nu')\le b\), which also supplies the required zero-coefficient estimate for the infinite-cost coupling.



If \(a\ne0\), then \(a>0\), so \(a\cdot(+\infty)+b=+\infty\) and the required inequality is automatic.

\section{\texorpdfstring{Iteration of the scalar error estimate}{Iteration of the scalar error estimate}}
\label{proof-module-UnadjustedLangevin-Recurrence-lean}
{\small\color{gray}Source module: \texttt{Unadjusted\allowbreak{}Langevin/\allowbreak{}Recurrence.\allowbreak{}lean}.\par}


Once coupling gives \(e_{k+1}\le qe_k+b\), the remaining calculation is deterministic. We retain the exact geometric sum before obtaining the simpler bias-floor bound. The upper envelope approaches \(b/(1-q)\); this does not assert that the actual error converges to that value.



The scalar sequence need not have an algorithmic interpretation for these lemmas. The proofs separate the exact geometric sum from its uniform bias-floor envelope and explicitly require the contraction factor in \([0,1)\). Later the coupling construction proves that the actual ULA error satisfies this recurrence.

\subsection{\texorpdfstring{Iteration gives a geometric sum}{Iteration gives a geometric sum}}


If \(q\ge0\) and \(e_{k+1}\le qe_k+b\) for every \(k\), then

\[e_n\le q^ne_0+b\sum_{j=0}^{n-1}q^j.\]

Induct on \(n\). At zero the sum is empty and both sides equal \(e_0\). Assume the bound at \(n\).



Multiply the induction inequality by \(q\ge0\), preserving its direction.



Write \(S_n=\sum_{j=0}^{n-1}q^j\). The geometric identity gives \(qS_n+1=S_n+q^n\); multiply this equality by \(b\).



Combine the one-step bound with the multiplied induction bound. Since \(S_{n+1}=S_n+q^n\), obtain \(e_{n+1}\le q^{n+1}e_0+bS_{n+1}\).

\subsection{\texorpdfstring{The exact finite-\allowbreak{}horizon envelope}{The exact finite-horizon envelope}}


If \(0\le q<1\) and \(e_{k+1}\le qe_k+b\), then

\[e_n\le q^ne_0+\frac{b(1-q^n)}{1-q}.\]

Apply the geometric-sum bound and evaluate the sum as \((q^n-1)/(q-1)\). Its denominator is nonzero because \(q<1\).



Change the signs of numerator and denominator simultaneously and put \(b\) in the numerator to obtain \(b(1-q^n)/(1-q)\).

\subsection{\texorpdfstring{A uniform bias-\allowbreak{}floor bound}{A uniform bias-floor bound}}


If \(0\le q<1\), \(b\ge0\), and \(e_{k+1}\le qe_k+b\), then

\[e_n\le q^ne_0+\frac b{1-q}.\]

Start from the exact finite-horizon bound. Its denominator \(1-q\) is positive.



Since \(q^n\ge0\) and \(b\ge0\), the numerator satisfies \(b(1-q^n)\le b\).



Divide by the positive denominator, add \(q^ne_0\), and compose with the exact bound.

\subsection{\texorpdfstring{Limit of the upper envelope}{Limit of the upper envelope}}


For \(0\le q<1\) and arbitrary real \(b,e_0\),

\[q^ne_0+\frac b{1-q}\longrightarrow\frac b{1-q}\quad(n\to\infty).\]

The geometric powers satisfy \(q^n\to0\).



Multiply this limit by \(e_0\) and add the constant \(b/(1-q)\).

\subsection{\texorpdfstring{Every accuracy above the floor is eventually attained}{Every accuracy above the floor is eventually attained}}


Assume \(0\le q<1\), \(b\ge0\), and \(e_{k+1}\le qe_k+b\). For every \(\varepsilon>b/(1-q)\), the errors satisfy \(e_n<\varepsilon\) for every sufficiently large \(n\).



The upper envelope tends to a value strictly below \(\varepsilon\), so it is eventually below \(\varepsilon\).



At each such index, place \(e_n\) below the envelope by the recurrence bound and then below \(\varepsilon\).

\section{\texorpdfstring{The shared Hilbert-\allowbreak{}space proof and both step regimes}{The shared Hilbert-space proof and both step regimes}}
\label{proof-module-UnadjustedLangevin-HilbertDalalyan-lean}
{\small\color{gray}Source module: \texttt{Unadjusted\allowbreak{}Langevin/\allowbreak{}Hilbert\allowbreak{}Dalalyan.\allowbreak{}lean}.\par}


Work on a probability space with \(\mathcal H=L^2(P;H)\), and write \(g=\widehat{\nabla U}\). This is the single analytic proof used by both Euclidean ULA and the lifted IPLA system. The reference record exposes the integral equation, stationary laws, and moment bounds used at this stage. \texttt{Canonical\allowbreak{}ULA.\allowbreak{}lean} constructs these inputs from strong convexity, a Lipschitz gradient, and independent Gaussian innovations.



\texttt{L2\allowbreak{}Reference} stores a stationary reference path in the same \(L^2\) space and with the same innovations as the numerical chain. The analytic theorem derives its recurrence and final transport estimate. The reference is a supplied argument at this layer; the canonical Brownian construction later proves its fields from the potential assumptions.

\subsection{\texorpdfstring{The common L\ensuremath{{}^{2}} algorithm}{The common L² algorithm}}


For an initial \(X_0\in L^2(P;H)\) and innovations \(\xi_k\in L^2(P;H)\), set \(X_{k+1}=X_k-h\widehat{\nabla U}(X_k)+\sqrt{2h}\,\xi_k\).



The initial iterate is the supplied random vector.



Each successor uses the actual-gradient Euler update.

\subsection{\texorpdfstring{Unfolding one step}{Unfolding one step}}


The next iterate is exactly one L\ensuremath{{}^{2}} ULA update of the current iterate.



Evaluate the successor clause of the definition.

\subsection{\texorpdfstring{The continuous reference assumptions}{The continuous reference assumptions}}


For an \(L\)-Lipschitz gradient, step \(h\), noise vectors \(\xi_k\in\mathcal H\), and a measure \(\pi\), a reference consists of segments \(Y_{k,t}\in\mathcal H\) and increments \(B_{k,t}\in\mathcal H\) satisfying the following conditions.



Supply the two families \((Y_{k,t})\) and \((B_{k,t})\) in the Hilbert space \(\mathcal H\).



Require continuity of \(t\mapsto Y_{k,t}\) in the \(L^2\) norm for every \(k\).



Require compatible endpoints, common endpoint noise, and stationary grid laws:

\[Y_{k,h}=Y_{k+1,0},\qquad B_{k,h}=\sqrt{2h}\,\xi_k,\qquad(Y_{k,0})_\#P=\pi.\]

For \(0\le t\le h\), require

\[Y_{k,t}-Y_{k,0}=-\int_0^t g(Y_{k,s})\,ds+B_{k,t}.\]

Require the score moment bound \(\|g(Y_{k,t})\|_2\le\sqrt{Ld}\) throughout each closed step interval.



Require the increment moment bound \(\|B_{k,t}\|_2\le\sqrt{2d}\sqrt t\) on the same interval.

\subsection{\texorpdfstring{Full convergence calculation from the integral reference}{Full convergence calculation from the integral reference}}


Suppose a reference with the preceding properties is supplied, \(h\ge0\), and \(Lh\le2\). Let \(0\le q<1\), and assume the lifted Euler map contracts:

\[\|(X-hg(X))-(Y-hg(Y))\|_2\le q\|X-Y\|_2.\]

For any initial vector \(X_0\in\mathcal H\), the \(L^2\) chain driven by the reference endpoint noise satisfies

\[W_2^{\mathbb R}((X_n)_\#P,\pi)\le q^n\|X_0-Y_{0,0}\|_2+\frac{(33/20)Lh\sqrt{hd}}{1-q}.\]

The statement is conditional on the supplied reference, and the initial term is the cost of the chosen coupling.



Set \(g=\widehat{\nabla U}\), name the actual iterates \(X_k\), and define

\[b=\frac{33}{20}Lh\sqrt{hd},\qquad D_k=\int_0^h\bigl(g(Y_{k,t})-g(Y_{k,0})\bigr)\,dt.\]

The bound \(b\) is nonnegative. Apply the local-defect theorem to the actual lifted gradient and the supplied equation and moment bounds, deriving \(\|D_k\|_2\le b\).



We claim the reference endpoints obey

\[Y_{k+1,0}=Y_{k,0}-hg(Y_{k,0})+\sqrt{2h}\,\xi_k-D_k.\]

The lifted gradient is Lipschitz and the reference segment is continuous, so their composition is continuous and Bochner integrable on \([0,h]\).



Apply the Euler-defect identity at \(t=h\), then substitute the segment join and endpoint noise identities. This proves the claimed reference update. Its defect has a minus sign.



Fix \(k\) and compare the actual update with the reference endpoint equation, both using the same \(\xi_k\).



Cancel their identical noise terms. Because the reference subtracts \(D_k\),

\[X_{k+1}-Y_{k+1,0}=\bigl(X_k-hg(X_k)\bigr)-\bigl(Y_{k,0}-hg(Y_{k,0})\bigr)+D_k.\]

Apply the triangle inequality, Euler contraction, and the defect bound. Setting \(e_k=\|X_k-Y_{k,0}\|_2\) gives \(e_{k+1}\le qe_k+b\).



The actual pair supplies a finite-cost coupling, so its transport distance is at most \(e_n\). The second marginal is \(\pi\) by the stationary grid-law assumption.



Apply the scalar recurrence bound to the derived inequality and compose it with the transport comparison.

\subsection{\texorpdfstring{The small-\allowbreak{}step bound}{The small-step bound}}


Suppose \(U\) is differentiable and \(m\)-strongly convex on \(H\), \(\nabla U\) is \(L\)-Lipschitz, and \(0<m\le L\). Let \(0<h\le2/(m+L)\), and supply a reference with the preceding properties. For any initial vector \(X_0\in\mathcal H\),

\[W_2^{\mathbb R}((X_n)_\#P,\pi)\le(1-mh)^n\|X_0-Y_{0,0}\|_2+\frac{33}{20}\frac Lm\sqrt{hd}.\]

This uses a chosen initial coupling and a supplied stationary reference. Replacing the initial cost by \(W_2(\mathcal L(X_0),\pi)\) needs the additional uniform family of reference couplings explained in the final section.



Since \(m+L>0\), clear the positive denominator in the step hypothesis to obtain \(h(m+L)\le2\).



Use \(m\le L\) to get \(2mh\le h(m+L)\le2\), hence \(mh\le1\).



The same threshold gives \(Lh\le2\). Thus \(q=1-mh\ge0\); since \(mh>0\), also \(q<1\).



The sharp small-step contraction of the lifted Euler map gives the contraction hypothesis with \(q=1-mh\). Apply the integral-reference theorem.



Because \(m,h>0\), simplify the bias term:

\[\frac{(33/20)Lh\sqrt{hd}}{1-(1-mh)}=\frac{(33/20)Lh\sqrt{hd}}{mh}=\frac{33}{20}\frac Lm\sqrt{hd}.\]

Substitute the bias identity into the integral-reference bound.

\subsection{\texorpdfstring{The large-\allowbreak{}step bound}{The large-step bound}}


Suppose \(U\) is differentiable and \(m\)-strongly convex on \(H\), \(\nabla U\) is \(L\)-Lipschitz, and \(0<m\le L\). Let \(h>0\), \(h\ge2/(m+L)\), and \(Lh<2\), and supply a reference with the stated properties. For any initial vector \(X_0\in\mathcal H\),

\[W_2^{\mathbb R}((X_n)_\#P,\pi)\le(Lh-1)^n\|X_0-Y_{0,0}\|_2+\frac{33}{20}\frac{Lh}{2-Lh}\sqrt{hd}.\]

The denominator is positive by strict stability. The reference remains an input, and the initial term remains the chosen coupling cost.



Since \(m+L>0\), clear the denominator in the lower step threshold to obtain \(2\le h(m+L)\).



Use \(m\le L\) and \(h\ge0\): \(2\le h(m+L)\le2Lh\). Thus \(Lh\ge1\) and \(q=Lh-1\ge0\).



Strict stability \(Lh<2\) gives \(q=Lh-1<1\).



Apply the sharp large-step Euler contraction with \(q=Lh-1\), then invoke the integral-reference theorem.



Use \(1-(Lh-1)=2-Lh\) and factor out \(\sqrt{hd}\):

\[\frac{(33/20)Lh\sqrt{hd}}{1-(Lh-1)}=\left(\frac{(33/20)Lh}{2-Lh}\right)\sqrt{hd}.\]

Substitute this identity into the integral-reference bound to obtain the large-step formula.

\section{\texorpdfstring{Compatibility exports for the shared Hilbert-\allowbreak{}space coupling proof}{Compatibility exports for the shared Hilbert-space coupling proof}}
\label{proof-module-IPLA-Dalalyan-lean}
{\small\color{gray}Source module: \texttt{IPLA/\allowbreak{}Dalalyan.\allowbreak{}lean}.\par}


This module introduces no new mathematical declarations. It exports the shared \texttt{ULA.\allowbreak{}Hilbert} recursion, integral reference interface and Dalalyan estimates under the historical IPLA namespace. The imported shared module contains their Lean proofs and parallel mathematical exposition.

\subsection{\texorpdfstring{The shared synchronous coupling calculation}{The shared synchronous coupling calculation}}


The common Euler and diffusion noises cancel in the coupled difference. If \(e_k=\|X_k-Y_k\|_{L^2}\), contraction of the deterministic gradient step and the integrated drift-defect estimate give

\[e_{k+1}\le q e_k+b,\qquad b=\frac{33}{20}Lh\sqrt{hd}.\]

Iteration gives \(e_k\le q^ke_0+b/(1-q)\). In the small-step regime \(0<h\le2/(m+L)\), \(q=1-mh\), yielding bias \((33/20)(L/m)\sqrt{hd}\). The reference structure explicitly requires continuous joined segments, their integral equation, stationary knot laws and score/\allowbreak{}noise moment bounds. This compatibility file proves none of those requirements anew; the final particle theorem obtains them from the shared Gaussian ULA construction.



The \texttt{export} commands expose the already proved declarations under the IPLA namespace. Their definitions and proofs are unchanged, so the mathematical argument is the one in the imported source modules; the dependency diagram below opens those full proofs. This file itself makes no additional analytic assumption and contains no replacement proof.

\section{\texorpdfstring{Statistical target, initialization and Gaussian innovations}{Statistical target, initialization and Gaussian innovations}}
\label{proof-module-IPLA-Model-lean}
{\small\color{gray}Source module: \texttt{IPLA/\allowbreak{}Model.\allowbreak{}lean}.\par}


This file names the original parameter–latent statistical objects and the law of the actual algorithm. Definitions identify a density and a candidate maximizer; normalization, existence and concentration are proved in later modules.

\subsection{\texorpdfstring{The marginal likelihood and its maximizer}{The marginal likelihood and its maximizer}}


The pairing map forms \(\operatorname{pair}(\theta,x)=(\theta,x)\) in the orthogonal joint space. The unnormalized parameter marginal is

\[M(\theta)=\int_{\mathbb R^{d_x}}e^{-U(\theta,x)}\,dx.\]

The definition uses the nonnegative extended integral, so it remains meaningful before finiteness has been established. A marginal maximizer is a point \(\theta_*\) such that \(M(\theta)\le M(\theta_*)\) for every \(\theta\). If \(U=-\log p_\theta(x,y)\), this is maximum marginal likelihood; a prior included in \(U\) gives the corresponding MAP interpretation. The deprecated name is an abbreviation for exactly the same predicate.




\textbf{Correspondence with the displayed declarations.}


\begin{itemize}

\item \texttt{pair}: The original parameter/\allowbreak{}latent pair \texttt{(theta,x)}, parameter first. For fixed observations \texttt{y}, the standard IPLA potential is \texttt{U y (pair theta x) = -\allowbreak{}log p\_\allowbreak{}theta(x,y)}. The geometry and dynamics work with the fixed slice \texttt{U y}.
\item \texttt{marginal\allowbreak{}Density}: Unnormalized marginal density of the parameter, integrating out the latent variable. For \texttt{U\_\allowbreak{}y(theta,x) = -\allowbreak{}log p\_\allowbreak{}theta(x,y)}, this is the marginal likelihood \texttt{p\_\allowbreak{}theta(y)}. A parameter prior may instead be included in \texttt{U\_\allowbreak{}y}.
\item \texttt{Is\allowbreak{}Marginal\allowbreak{}Maximizer}: The parameter maximizes the integrated joint density. This is maximum likelihood in standard IPLA notation, and MAP when the parameter prior is included in \texttt{U}. Its existence is not asserted by this definition.
\item \texttt{Is\allowbreak{}Marginal\allowbreak{}MAP}: Compatibility name for the prior-inclusive formulation.
\end{itemize}
\subsection{\texorpdfstring{The actual Gibbs target and all Gaussian assumptions}{The actual Gibbs target and all Gaussian assumptions}}


For \(V_N\) defined by the geometry module, the target is \(\Pi_N=Z_N^{-1}e^{-V_N}\,dz\), with \(Z_N=\int e^{-V_N}\,dz\). Both normalization and density are initially extended-valued; this definition alone asserts no probability property. The Gaussian innovation structure asks that every coordinate \(\xi_{k,a}\) have law \(\mathcal N(0,1)\), that all coordinates at all times be mutually independent, and that the entire noise path be independent of the initialization. These are explicit distributional hypotheses on the supplied random variables.




The Gaussian structure concerns the entire countable collection \((k,a)\mapsto\xi_{k,a}\), not only coordinate independence at one time. The separate initialization field states independence of \(Z_0\) from the whole innovation path. Together these conditions imply that every transition receives a fresh full state-space standard Gaussian. They are the probabilistic assumptions later used to equate a supplied implementation with the canonical Markov laws; moment membership in \(L^2\) alone would not establish either the Gaussian distribution or independence.




\textbf{Correspondence with the displayed declarations.}


\begin{itemize}

\item \texttt{gibbs\allowbreak{}Target}: The normalized density in (49), with \texttt{U\_\allowbreak{}N} the actual lifted potential. The integral and scalar are extended-valued. \texttt{gibbs\allowbreak{}Target\_\allowbreak{}is\allowbreak{}Probability} derives positive finite normalization from the original joint strong convexity, before any conversion to real numbers.
\item \texttt{Gaussian\allowbreak{}Innovations}: The actual update noise is independent standard Gaussian in every coordinate and time, and independent of initialization. These are distributional hypotheses, not a claim to have constructed a random sequence.
\end{itemize}
\subsection{\texorpdfstring{The law of the L\ensuremath{{}^{2}} recursion}{The law of the L² recursion}}


A square-integrable initial state and square-integrable noise vectors lie in \(L^2(P)\). With \(h=\gamma/N\), define the random sequence through the shared ULA recursion

\[Z_{k+1}=Z_k-h\nabla V_N(Z_k)+\sqrt{2h}\,\xi_k,\qquad \rho_k=(Z_k)_\#P.\]

The Lipschitz estimate for \(\nabla V_N\) ensures the drift update stays in \(L^2\). The definition of \texttt{laws} is precisely this pushforward. The dynamics module proves that its representatives agree almost everywhere with the ordinary pointwise IPLA recursion, so the eventual theorem concerns this actual algorithm.




\textbf{Correspondence with the displayed declarations.}


\begin{itemize}

\item \texttt{laws}: The laws of the actual noise-driven IPLA recursion, in scaled coordinates. The coordinate identities in \texttt{Geometry.\allowbreak{}lean} show that this is precisely the usual average-gradient parameter step and ordinary particle Langevin steps.
\end{itemize}
\section{\texorpdfstring{Pointwise recursion and the L\ensuremath{{}^{2}} algorithm law}{Pointwise recursion and the L² algorithm law}}
\label{proof-module-IPLA-Dynamics-lean}
{\small\color{gray}Source module: \texttt{IPLA/\allowbreak{}Dynamics.\allowbreak{}lean}.\par}


The formal convergence proof uses equivalence classes in \(L^2\). This module proves that their pushforward laws equal those of the displayed coordinate algorithm, including the exact square-root noise coefficients.

\subsection{\texorpdfstring{Recover the original pairs from scaled coordinates}{Recover the original pairs from scaled coordinates}}


For every particle \(i\), checking parameter and latent coordinates gives \(A_i z=\operatorname{pair}(p_\theta z,\sqrt N\,\operatorname{block}_i z)\). Substitution into the finite sum proves

\[V_N(z)=\sum_iU\bigl(p_\theta z,\sqrt N\,\operatorname{block}_iz\bigr).\]

This identifies the abstract lifted potential with the sum of the original pair potentials at the physical particles.




\textbf{Correspondence with the displayed declarations.}


\begin{itemize}

\item \texttt{joint\_\allowbreak{}eq\_\allowbreak{}pair}: A joint extraction is exactly the original parameter-particle pair after unscaling its particle.
\item \texttt{potential\_\allowbreak{}eq\_\allowbreak{}pair\_\allowbreak{}sum}: The lifted potential is the sum of the original pair potential, with physical particles equal to \texttt{sqrt N} times the scaled particle blocks.
\end{itemize}
\subsection{\texorpdfstring{The explicit recursion and its two block formulas}{The explicit recursion and its two block formulas}}


Set \(Z_0=z_0\) and recursively let \(Z_{k+1}=\operatorname{scaledUpdate}(U,\gamma,Z_k,\xi_k)\). The zero and successor identities are the two defining cases of this recursion. Applying the coordinate identities from the geometry module at \(Z_k\) proves

\[\theta_{k+1}=\theta_k-\frac\gamma N\sum_i\nabla_\theta U(\theta_k,x_k^i)+\sqrt{2\gamma/N}\,\xi_k^\theta,\]\[x_{k+1}^i=x_k^i-\gamma\nabla_xU(\theta_k,x_k^i)+\sqrt{2\gamma}\,\xi_k^i,\qquad x_k^i=\sqrt N(Z_k)_i.\]

The particle formula requires \(N>0\), exactly where unscaling cancels \(N\).




\textbf{Correspondence with the displayed declarations.}


\begin{itemize}

\item \texttt{scaled\allowbreak{}Chain}: IPLA as an ordinary recursion on scaled states and noise vectors.
\item \texttt{scaled\allowbreak{}Chain\_\allowbreak{}parameter}: The parameter recurrence in the pointwise IPLA chain.
\item \texttt{scaled\allowbreak{}Chain\_\allowbreak{}particle}: The physical-particle recurrence in the pointwise IPLA chain.
\end{itemize}
\subsection{\texorpdfstring{Induction identifies the L\ensuremath{{}^{2}} and pointwise algorithms}{Induction identifies the L² and pointwise algorithms}}


At time zero, the representative of the \(L^2\) chain is the initialization. Suppose its representative equals the pointwise chain almost surely at time \(k\). The shared update representative theorem gives the Euler formula almost surely; intersect that full-measure event with the induction event, substitute the previous-state equality, and rewrite \(2(\gamma/N)=2\gamma/N\). This proves almost-sure equality at \(k+1\). Finally, almost-everywhere equal measurable maps have equal pushforward measures, hence

\[\rho_k=P\circ\bigl[\omega\mapsto\operatorname{scaledChain}(U,\gamma,Z_0(\omega),\xi(\omega),k)\bigr]^{-1}.\]


The representative induction uses equality almost everywhere at each fixed finite time. In the successor case the intersection of the update's representative event and the previous induction event still has full measure, so substitution is legitimate. Once the random variables agree almost everywhere, \texttt{Measure.\allowbreak{}map\_\allowbreak{}congr} gives equality of their laws. This step is what permits results proved for the Hilbert-space recursion to be presented as results for the explicit coordinate algorithm displayed above.




\textbf{Correspondence with the displayed declarations.}


\begin{itemize}

\item \texttt{lp\allowbreak{}Chain\_\allowbreak{}coe}: The L\ensuremath{{}^{2}} chain represents the explicit IPLA recursion almost everywhere. This includes the exact identity between \texttt{sqrt (2 * (gamma/\allowbreak{}N))} and the noise coefficient \texttt{sqrt (2*gamma/\allowbreak{}N)} in the displayed IPLA update.
\item \texttt{laws\_\allowbreak{}eq\_\allowbreak{}scaled\allowbreak{}Chain}: The laws appearing in the final theorem are exactly the laws of the explicit, noise-driven, pointwise IPLA recursion.
\end{itemize}
\section{\texorpdfstring{Gibbs normalization and second moments}{Gibbs normalization and second moments}}
\label{proof-module-UnadjustedLangevin-Gibbs-lean}
{\small\color{gray}Source module: \texttt{Unadjusted\allowbreak{}Langevin/\allowbreak{}Gibbs.\allowbreak{}lean}.\par}


On a finite-dimensional real Hilbert space, differentiability and \(m\)-strong convexity with \(m>0\) imply that \(\pi(dx)=Z^{-1}e^{-U(x)}dx\) is a probability measure with finite second moment. This establishes existence of the target measure used throughout the coupling proof.

\subsection{\texorpdfstring{gibbs\allowbreak{}Weight}{gibbsWeight}}


The unnormalized density is \(w_U(x)=e^{-U(x)}\), embedded in \([0,\infty]\) with \texttt{ENNReal.\allowbreak{}of\allowbreak{}Real}. This choice lets the partition integral retain the value infinity; positivity of the exponential ensures that embedding does not discard any part of the density.

\subsection{\texorpdfstring{gibbs\allowbreak{}Partition}{gibbsPartition}}


The partition function is \(Z_U=\int e^{-U(x)}\,dx\), formed as a nonnegative extended integral. At this definition no finiteness assumption is made. Later normalization uses separate proofs of \(0<Z_U\) and \(Z_U<\infty\), which are needed to justify cancellation with its inverse.

\subsection{\texorpdfstring{gibbs\allowbreak{}Measure}{gibbsMeasure}}


Define \(\pi_U(dx)=Z_U^{-1}e^{-U(x)}dx\) using \texttt{volume.\allowbreak{}with\allowbreak{}Density}. The inverse is the extended nonnegative inverse. The definition therefore exists for every potential, while the later probability theorem proves that it has total mass one under the stated strong convexity assumptions.

\subsection{\texorpdfstring{integrable exp neg mul norm sq}{integrable exp neg mul norm sq}}


For \(a>0\), \(e^{-a\|x\|^2}\) is Lebesgue integrable in finite dimension. The Lean proof specializes \texttt{Gaussian\allowbreak{}Fourier.\allowbreak{}integrable\_\allowbreak{}cexp\_\allowbreak{}neg\_\allowbreak{}mul\_\allowbreak{}sq\_\allowbreak{}norm\_\allowbreak{}add} to real positive precision \(a\) and zero linear term, then takes the real part. The complex Gaussian integrand is the real embedding of the desired exponential, so the final simplification identifies the functions exactly. This also includes dimension zero.

\subsection{\texorpdfstring{strong\allowbreak{}Convex quadratic lower bound}{strongConvex quadratic lower bound}}


Apply the strong supporting-plane inequality at zero: \(U(x)\ge U(0)+\langle\nabla U(0),x\rangle+\frac m2\|x\|^2\). Cauchy–Schwarz and \(ab\le a^2/m+(m/4)b^2\), valid for \(m>0\), absorb the linear term. The result is \(U(x)\ge U(0)-\|\nabla U(0)\|^2/m+(m/4)\|x\|^2\). Lean isolates this scalar square inequality before its final arithmetic rearrangement.

\subsection{\texorpdfstring{integrable exp neg of strong\allowbreak{}Convex}{integrable exp neg of strongConvex}}


Exponentiating the quadratic lower bound gives \(e^{-U(x)}\le e^{-U(0)+\|\nabla U(0)\|^2/m}e^{-(m/4)\|x\|^2}\). The right side is an integrable Gaussian multiple. Differentiability supplies continuity and measurability of the left side, so \texttt{Integrable.\allowbreak{}mono'} transfers integrability from that explicit dominating function.

\subsection{\texorpdfstring{integrable norm sq mul exp neg mul norm sq}{integrable norm sq mul exp neg mul norm sq}}


For \(a>0\), split the Gaussian decay into two equal factors. The scalar function \(s\mapsto s e^{-(a/2)s}\) is bounded on \([0,\infty)\), so \(\|x\|^2e^{-a\|x\|^2}\le C_a e^{-(a/2)\|x\|^2}\). The remaining Gaussian is integrable; the proof checks measurability and nonnegativity before applying domination.

\subsection{\texorpdfstring{integrable norm sq mul exp neg of strong\allowbreak{}Convex}{integrable norm sq mul exp neg of strongConvex}}


Multiply the previous Gibbs Gaussian domination by \(\|x\|^2\). Its right side has the form \(C\|x\|^2e^{-(m/4)\|x\|^2}\), whose integrability was proved separately. The pointwise inequality uses \(\|x\|^2\ge0\), so multiplying it preserves order; this is the second-moment input for the normalized measure.

\subsection{\texorpdfstring{gibbs\allowbreak{}Partition pos}{gibbsPartition pos}}


The continuous function \(e^{-U}\) is strictly positive everywhere. Its extended integral cannot vanish: a zero integral of a nonnegative measurable function would force it to vanish almost everywhere, contradicting positivity on a nonzero-volume open set. This proves \(0<Z_U\) without requiring an upper bound on \(Z_U\).

\subsection{\texorpdfstring{gibbs\allowbreak{}Partition ne top}{gibbsPartition ne top}}


Strong convexity with \(m>0\) has already given ordinary integrability of \(e^{-U}\). For a nonnegative integrable function its extended integral is finite. Unfolding \texttt{gibbs\allowbreak{}Partition} identifies that integral exactly with \(Z_U\), giving the guard needed for inverse cancellation.

\subsection{\texorpdfstring{gibbs\allowbreak{}Partition eq of\allowbreak{}Real integral}{gibbsPartition eq ofReal integral}}


Nonnegativity and integrability identify the extended integral with the embedded real integral: \(Z_U=\operatorname{ofReal}(\int e^{-U}\,dx)\). The proof uses the genuine integral/\allowbreak{}lintegral comparison theorem, so it cannot identify an infinite partition with the default real value of an undefined integral.

\subsection{\texorpdfstring{integral exp neg pos}{integral exp neg pos}}


Combine \(0<Z_U\) with \(Z_U=\operatorname{ofReal}(\int e^{-U}\,dx)\). Positivity of \texttt{of\allowbreak{}Real} is equivalent to positivity of its real argument, proving that the ordinary normalizing constant is strictly positive and allowing real division by it later.

\subsection{\texorpdfstring{is\allowbreak{}Probability\allowbreak{}Measure gibbs\allowbreak{}Measure}{isProbabilityMeasure gibbsMeasure}}


The mass of a density measure is its extended density integral. Pull the constant \(Z_U^{-1}\) out to get \(\pi_U(H)=Z_U^{-1}Z_U\). The proofs that \(Z_U\ne0\) and \(Z_U\ne\infty\) justify this product being one. This theorem constructs the \texttt{Is\allowbreak{}Probability\allowbreak{}Measure} instance rather than assuming normalization.

\subsection{\texorpdfstring{integrable norm sq gibbs\allowbreak{}Measure}{integrable norm sq gibbsMeasure}}


Changing measure converts integrability under \(\pi_U\) to integrability of \(Z_U^{-1}e^{-U(x)}\|x\|^2\) under volume. The weighted second moment is already integrable, and the real normalizing factor is a finite scalar. The with-density theorem supplies the almost-everywhere measurability and finiteness guards for this conversion.

\subsection{\texorpdfstring{mem\allowbreak{}Lp id gibbs\allowbreak{}Measure}{memLp id gibbsMeasure}}


For the identity random variable, membership in \(L^2(\pi_U;H)\) is equivalent to integrability of \(\|x\|^2\). The identity is measurable, and the preceding weighted-moment theorem supplies the integral. Lean uses the \(p=2\) characterization of \texttt{Mem\allowbreak{}Lp}, retaining the actual function before passing to any quotient class.

\subsection{\texorpdfstring{mem\allowbreak{}Lp sub const gibbs\allowbreak{}Measure}{memLp sub const gibbsMeasure}}


Since \(\pi_U\) is a probability, the constant map \(x\mapsto z\) belongs to \(L^2\). Subtract it from the identity, already in \(L^2(\pi_U)\), to obtain \(x\mapsto x-z\in L^2(\pi_U)\). This is a moment assertion for every chosen center, without requiring that the center minimize the potential.

\section{\texorpdfstring{Subgaussian Gibbs score and its sharp second moment}{Subgaussian Gibbs score and its sharp second moment}}
\label{proof-module-UnadjustedLangevin-GibbsScore-lean}
{\small\color{gray}Source module: \texttt{Unadjusted\allowbreak{}Langevin/\allowbreak{}Gibbs\allowbreak{}Score.\allowbreak{}lean}.\par}


Differentiability, an \(L\)-Lipschitz gradient, and finite Gibbs normalization suffice for \(\int\|\nabla U\|^2d\pi\le Ld\). Strong convexity is used only by the final wrapper to derive normalization.

\subsection{\texorpdfstring{gibbs\allowbreak{}Weight mul exp score le}{gibbsWeight mul exp score le}}


The descent lemma gives \(U(x+v)\le U(x)+\langle\nabla U(x),v\rangle+(L/2)\|v\|^2\). Negate and exponentiate, then rearrange to \(e^{-U(x)}e^{-\langle\nabla U(x),v\rangle}\le e^{L\|v\|^2/2}e^{-U(x+v)}\). Lean first establishes the real inequality and then embeds it in the extended nonnegative numbers.

\subsection{\texorpdfstring{gibbs score lmgf le}{gibbs score lmgf le}}


Integrate the pointwise exponential-score bound. Translation invariance gives \(\int e^{-U(x+v)}dx=Z_U\); normalization cancels this factor, leaving \(\mathbb E_{\pi_U}e^{-\langle\nabla U,v\rangle}\le e^{L\|v\|^2/2}\). Finiteness of \(Z_U\) is explicit because extended inverse cancellation needs it.

\subsection{\texorpdfstring{gibbs score subgaussian}{gibbs score subgaussian}}


Apply the preceding bound with the vector \(-tv\), for every real \(t\). The exponent becomes \(t\langle\nabla U,v\rangle\) and the right side is \(e^{t^2L\|v\|^2/2}\). This is exactly \texttt{Has\allowbreak{}Subgaussian\allowbreak{}MGF} with variance proxy \(L\|v\|^2\); the proof also establishes exponential integrability required by that structure.

\subsection{\texorpdfstring{sq le exp add exp neg sub two}{sq le exp add exp neg sub two}}


The series \(\cosh a=\sum_{n\ge0}a^{2n}/(2n)!\) has nonnegative terms, so retaining its first two terms gives \(\cosh a\ge1+a^2/2\). Substitute \(2\cosh a=e^a+e^{-a}\) and rearrange to \(a^2\le e^a+e^{-a}-2\). Lean uses \texttt{Real.\allowbreak{}has\allowbreak{}Sum\_\allowbreak{}cosh} with \texttt{sum\_\allowbreak{}le\_\allowbreak{}has\allowbreak{}Sum}, explicitly proving nonnegativity of every discarded even-power term.

\subsection{\texorpdfstring{integral sq le of subgaussian}{integral sq le of subgaussian}}


Integrate \(t^2X^2\le e^{tX}+e^{-tX}-2\), giving \(t^2\mathbb EX^2\le2(e^{ct^2/2}-1)\). Subgaussianity provides the exponential integrability and hence the square integrability needed here. Divide by \(t^2\) for \(t\ne0\), and let \(t\to0\) to obtain \(\mathbb EX^2\le c\). Lean computes the scalar limit before applying closedness of order.

\subsection{\texorpdfstring{integral norm sq gradient gibbs le}{integral norm sq gradient gibbs le}}


Choose an orthonormal basis \((e_i)_{i=1}^d\) and use Parseval: \(\|\nabla U(x)\|^2=\sum_i\langle\nabla U(x),e_i\rangle^2\). Each coordinate is subgaussian with proxy \(L\), hence has second moment at most \(L\). Finite-sum integrability permits integration term by term, yielding \(\mathbb E_{\pi_U}\|\nabla U\|^2\le Ld\).

\subsection{\texorpdfstring{mem\allowbreak{}Lp gradient gibbs\allowbreak{}Measure}{memLp gradient gibbsMeasure}}


The coordinate subgaussian bounds prove integrability of every squared gradient coordinate. Summing over the finite orthonormal basis proves integrability of \(\|\nabla U\|^2\). Lipschitz continuity provides strong measurability, and the \texttt{Mem\allowbreak{}Lp} square criterion packages those facts as \(\nabla U\in L^2(\pi_U)\).

\subsection{\texorpdfstring{integral norm sq gradient gibbs le of strong\allowbreak{}Convex}{integral norm sq gradient gibbs le of strongConvex}}


Strong convexity with \(m>0\) supplies the finite-partition premise of the score theorem. Substitute that proof into the first-order smoothness result to obtain \(\mathbb E_{\pi_U}\|\nabla U\|^2\le Ld\). The strong convexity constant is used for normalization here; it does not appear in the sharp score bound.

\section{\texorpdfstring{Finite quadratic transport transfers second moments}{Finite quadratic transport transfers second moments}}
\label{proof-module-UnadjustedLangevin-TransportMoments-lean}
{\small\color{gray}Source module: \texttt{Unadjusted\allowbreak{}Langevin/\allowbreak{}Transport\allowbreak{}Moments.\allowbreak{}lean}.\par}


Finite \(W_2\) distance to a finite-second-moment law supplies the other second moment automatically. This removes a redundant integrability premise when passing an initial law to the L\ensuremath{{}^{2}} formulation.

\subsection{\texorpdfstring{mem\allowbreak{}Lp id of wasserstein2 ne top}{memLp id of wasserstein2 ne top}}


Finite Wasserstein distance supplies a finite-cost coupling. On that coupling space the second coordinate is in \(L^2\) by the reference moment, and the coordinate difference is in \(L^2\) by finite coupling cost. Their sum is the first coordinate, so its law \(\mu\) has a finite second moment. No optimal coupling is selected.

\subsection{\texorpdfstring{wasserstein2 ne top of mem\allowbreak{}Lp}{wasserstein2 ne top of memLp}}


Use the explicit independent coupling \(\mu\otimes\pi\). Each coordinate is in \(L^2\), hence their difference is in \(L^2\) and has finite cost. Wasserstein distance is at most this coupling cost, proving finiteness.

\subsection{\texorpdfstring{mem\allowbreak{}Lp of law of wasserstein2 ne top}{memLp of law of wasserstein2 ne top}}


The finite-transport theorem first proves \(\mathrm{id}\in L^2(\mu)\). Replace \(\mu\) by the supplied pushforward law of \(X\), then use the change-of-law \(L^2\) criterion to pull this moment back to \(X\) on its original sample space.

\subsection{\texorpdfstring{mem\allowbreak{}Lp of wasserstein2 map ne top}{memLp of wasserstein2 map ne top}}


Specialize the preceding theorem to \(\mu=P\circ X^{-1}\), whose law identity is reflexive. Finite transport of the actual pushforward to a finite-moment reference therefore gives \(X\in L^2(P)\).

\section{\texorpdfstring{Normalization, moments and score moments of the lifted target}{Normalization, moments and score moments of the lifted target}}
\label{proof-module-IPLA-Foundations-lean}
{\small\color{gray}Source module: \texttt{IPLA/\allowbreak{}Foundations.\allowbreak{}lean}.\par}


Apply the shared ULA Gibbs results to \(V_N\), whose differentiability and \(Nm\)-strong convexity were derived from \(U\). All normalization and moment facts are conclusions from the original joint assumptions.

\subsection{\texorpdfstring{The named target is the normalized Gibbs measure}{The named target is the normalized Gibbs measure}}


The definitions give \(\operatorname{gibbsTarget}(U)=\operatorname{gibbsMeasure}(V_N)\). If \(N>0\) and \(m>0\), then \(Nm>0\). The lifted strong-convexity theorem and chain-rule differentiability therefore meet the shared Gibbs theorem's hypotheses. Its quadratic lower bound gives integrable Gaussian tails and a positive partition, proving \(0<Z_N<\infty\) and thus \(\Pi_N(\mathbb R^D)=1\). The same tails imply \(\int\|z\|^2d\Pi_N<\infty\), expressed as membership of the identity in \(L^2(\Pi_N)\).




The geometry has already proved that \(V_N\) is differentiable and \(Nm\)-strongly convex. Its quadratic lower bound dominates \(e^{-V_N}\) and its norm-weighted powers by integrable Gaussian functions, giving a finite positive partition and finite second moment. The first declaration identifies the named target with the shared normalized Gibbs definition; the next two apply the normalization and moment theorems with modulus \(Nm>0\). Thus target mass one and finite moments are consequences of the original joint assumptions.




\textbf{Correspondence with the displayed declarations.}


\begin{itemize}

\item \texttt{gibbs\allowbreak{}Target\_\allowbreak{}eq\_\allowbreak{}gibbs\allowbreak{}Measure}: IPLA's target is exactly the general normalized Gibbs density.
\item \texttt{gibbs\allowbreak{}Target\_\allowbreak{}is\allowbreak{}Probability}: Normalization of the lifted target follows from strong convexity of the original joint potential; no separate Gibbs normalization hypothesis.
\item \texttt{gibbs\allowbreak{}Target\_\allowbreak{}mem\allowbreak{}Lp}: The lifted target has a finite second moment from the original joint strong-convexity assumption. Gradient Lipschitz continuity is not needed.
\end{itemize}
\subsection{\texorpdfstring{Finite transport supplies square integrability of initialization}{Finite transport supplies square integrability of initialization}}


Suppose a measurable initialization \(X\) has \(W_2(X_\#P,\Pi_N)<\infty\). The target already has a finite second moment. For any finite-cost coupling, \(\|x\|^2\le2\|x-y\|^2+2\|y\|^2\); integrating transfers the target moment and coupling cost to the initial law. The shared transport-moment theorem formalizes this argument and proves \(X\in L^2(P)\). Thus this alternative initialization interface introduces no independent moment assumption.




\textbf{Correspondence with the displayed declarations.}


\begin{itemize}

\item \texttt{initial\_\allowbreak{}mem\allowbreak{}Lp\_\allowbreak{}of\_\allowbreak{}finite\_\allowbreak{}transport}: An actual measurable initialization at finite W\ensuremath{{}_{2}} from the lifted Gibbs measure automatically belongs to L\ensuremath{{}^{2}}. This justifies the L\ensuremath{{}^{2}} representation used by the main analytic theorem without an independent moment assumption.
\end{itemize}
\subsection{\texorpdfstring{The lifted Gibbs score moment has the correct dimension}{The lifted Gibbs score moment has the correct dimension}}


The shared score estimate for a strongly convex potential with \(\ell\)-Lipschitz gradient is \(\int\|\nabla V\|^2d\pi\le\ell\dim H\). Substitute \(V=V_N\), \(\ell=NL\), and \(\dim H=d_\theta+Nd_x\) to obtain

\[\int\|\nabla V_N(z)\|^2d\Pi_N(z)\le NL(d_\theta+Nd_x).\]

The lifted smoothness and convexity hypotheses follow from the joint ones. The proof only rewrites the common Gibbs definition and exact dimension; it adds no score-moment or Hessian premise.




Finite initial Wasserstein distance to a target with finite second moment implies finite initial second moment via a finite-cost coupling and \(\|X\|^2\le2\|X-Y\|^2+2\|Y\|^2\). The next theorem applies the shared Gibbs-score estimate to the actual lifted score with constant \(NL\), obtaining \[\mathbb E_{\Pi_N}\|\nabla V_N\|^2\le NL(d_\theta+Nd_x).\] The dimension substitution uses \texttt{finrank\_\allowbreak{}state}; a separate Hessian or score-moment premise is unnecessary.




\textbf{Correspondence with the displayed declarations.}


\begin{itemize}

\item \texttt{gibbs\allowbreak{}Target\_\allowbreak{}score\_\allowbreak{}moment}: The lifted Gibbs score moment also follows from the original joint potential assumptions; no separate score hypothesis or Hessian is introduced.
\end{itemize}
\section{\texorpdfstring{An abstract family of compatible finite-\allowbreak{}coupling references}{An abstract family of compatible finite-coupling references}}
\label{proof-module-UnadjustedLangevin-ReferenceFamily-lean}
{\small\color{gray}Source module: \texttt{Unadjusted\allowbreak{}Langevin/\allowbreak{}Reference\allowbreak{}Family.\allowbreak{}lean}.\par}


This interface isolates the passage from a bound for each initial coupling to a theorem with the exact initial Wasserstein distance. A concrete product construction later supplies this interface without assuming an optimal plan.

\subsection{\texorpdfstring{Finite\allowbreak{}Coupling}{FiniteCoupling}}


Restrict the coupling type to those with \(\operatorname{couplingCost}\ne\infty\). This subtype records the integrability needed to realize initial class distances and convert them to real values while taking the Wasserstein infimum.

\subsection{\texorpdfstring{Reference\allowbreak{}Family}{ReferenceFamily}}


For each finite coupling, supply initial and noise \(L^2\) variables and an \texttt{L2\allowbreak{}Reference}. The fields identify initial law, bound initial class distance by that coupling's cost, and require one common output law across the family. These are explicit construction obligations, with no convergence estimate among them.

\subsection{\texorpdfstring{Reference\allowbreak{}Family.\allowbreak{}output finite}{ReferenceFamily.output finite}}


Choose one finite initial coupling using the finite initial Wasserstein premise. Its supplied numerical/\allowbreak{}reference \(L^2\) variables yield a finite-cost output coupling. The output-law and stationary-reference fields identify its marginals, proving \(W_2(\nu_n,\pi)<\infty\).

\subsection{\texorpdfstring{Pass the analytic bound through the infimum}{Pass the analytic bound through the infimum}}


Apply the analytic \(L^2\) estimate to every finite coupling realization in the family. Replace its initial error by the coupling cost and its numerical law by the common \(\nu_n\). Because \((1-mh)^n\ge0\), the affine finite-coupling infimum theorem gives the displayed bound with the true initial Wasserstein distance.

\section{\texorpdfstring{Compatibility exports for the initial-\allowbreak{}coupling infimum}{Compatibility exports for the initial-coupling infimum}}
\label{proof-module-IPLA-ReferenceFamily-lean}
{\small\color{gray}Source module: \texttt{IPLA/\allowbreak{}Reference\allowbreak{}Family.\allowbreak{}lean}.\par}


This module re-exports the shared finite-coupling and reference-family interface. It contains no separate definitions or proofs; the imported ULA module develops the passage from a coupled initial RMS to the genuine initial \(W_2\) term.

\subsection{\texorpdfstring{Take the infimum without assuming an optimal coupling}{Take the infimum without assuming an optimal coupling}}


A reference family supplies, for every finite-cost coupling \(c\) of \(\mu_0\) and \(\pi\), an \(L^2\) realization of the initial variables and common Euler noise. Its fields require that the initial \(L^2\) difference be at most \(\operatorname{cost}(c)\), and that every realization produce the same output law \(\mu_k\). The shared coupling theorem gives \(W_2(\mu_k,\pi)\le q^k\operatorname{cost}(c)+b\) for each such \(c\). Taking the cost infimum yields \(q^kW_2(\mu_0,\pi)+b\). A finite-cost coupling also proves output finiteness before real conversion. These fields are explicit construction obligations in the reusable interface; the public Gaussian theorem discharges them.



The \texttt{export} commands expose the already proved declarations under the IPLA namespace. Their definitions and proofs are unchanged, so the mathematical argument is the one in the imported source modules; the dependency diagram below opens those full proofs. This file itself makes no additional analytic assumption and contains no replacement proof.

\section{\texorpdfstring{Exact particle-\allowbreak{}number and time scaling in the transport bound}{Exact particle-number and time scaling in the transport bound}}
\label{proof-module-IPLA-Scaling-lean}
{\small\color{gray}Source module: \texttt{IPLA/\allowbreak{}Scaling.\allowbreak{}lean}.\par}


The lifted ULA problem has \(m_N=Nm\), \(L_N=NL\), step \(h=\gamma/N\), and dimension \(D=d_\theta+Nd_x\). This module performs the exact scalar substitutions into the shared Dalalyan theorem.

\subsection{\texorpdfstring{Step restriction and square-\allowbreak{}root bias}{Step restriction and square-root bias}}


Since \(N>0\), dividing both sides by \(N\) preserves order, and \(Nm+NL=N(m+L)\). Therefore \(\gamma/N\le2/(Nm+NL)\) if and only if \(\gamma\le2/(m+L)\). The ratio \(NL/(Nm)\) cancels to \(L/m\), and

\[\frac\gamma N(d_\theta+Nd_x)=\gamma\left(d_x+\frac{d_\theta}{N}\right).\]

The intensive dimension in parentheses is nonnegative. The real square-root multiplication rule therefore factors the square root, giving the exact identity

\[\frac{33}{20}\frac{NL}{Nm}\sqrt{\frac\gamma N D}=\frac{33}{20}\frac Lm\sqrt{d_x+d_\theta/N}\sqrt\gamma.\]


The numerical substitutions are \(m_N=Nm\), \(L_N=NL\), \(h=\gamma/N\), \(D=d_\theta+Nd_x\). Positive \(N\) cancels the common factor in the step restriction. In the discretization term, \(L_N/m_N=L/m\) and \[hD=\gamma(d_x+d_\theta/N).\] Since the intensive dimension is nonnegative, the square-root product rule gives \(\sqrt{hD}=\sqrt{d_x+d_\theta/N}\sqrt\gamma\). These are exact identities, so this rewrite loses no constant.




\textbf{Correspondence with the displayed declarations.}


\begin{itemize}

\item \texttt{scaled\_\allowbreak{}step\_\allowbreak{}iff}: The original and rescaled step-size restrictions are equivalent.
\item \texttt{scaled\_\allowbreak{}bias\_\allowbreak{}eq}: Exact equality of the Dalalyan discretization term with the term in (34).
\end{itemize}
\subsection{\texorpdfstring{Substitute in the Hilbert-\allowbreak{}space Dalalyan theorem}{Substitute in the Hilbert-space Dalalyan theorem}}


For an explicit finite-cost reference family and finite initial \(W_2\), the shared theorem gives

\[W_2(\nu_k,\pi)\le(1-m_Nh)^kW_2(\nu_0,\pi)+\frac{33}{20}\frac{L_N}{m_N}\sqrt{hD}.\]

The assumptions \(m\le L\), \(m>0\), \(\gamma>0\), and the base step restriction imply the corresponding lifted inequalities by multiplication or division by positive \(N\). Finally \(Nm(\gamma/N)=m\gamma\), and the preceding bias identity yields the displayed IPLA transport estimate. This reusable intermediate theorem accepts the reference family explicitly; the public Gaussian theorem constructs it through shared ULA results.




The shared Hilbert-space theorem first gives
\[W_2(\nu_k,\pi)\le(1-m_Nh)^kW_2(\nu_0,\pi)+(33/20)(L_N/m_N)\sqrt{hD}\]
. The local facts \texttt{h\allowbreak{}NL} and \texttt{hs} verify its compatibility and step premises from \(m\le L\) and the physical step restriction. Substituting \(m_Nh=m\gamma\) and the exact bias identity yields the displayed IPLA rate. This module expects an explicit stationary reference family; later modules construct that family before the public theorem invokes the estimate.




\textbf{Correspondence with the displayed declarations.}


\begin{itemize}

\item \texttt{scaled\_\allowbreak{}dalalyan}: Dalalyan in the rescaled IPLA coordinates. Geometry will discharge the lifted convexity, differentiability and smoothness hypotheses from the joint potential. The only probabilistic inputs are the explicit reference family.
\end{itemize}
\section{\texorpdfstring{Project transport convergence into parameter RMS error}{Project transport convergence into parameter RMS error}}
\label{proof-module-UnadjustedLangevin-Projection-lean}
{\small\color{gray}Source module: \texttt{Unadjusted\allowbreak{}Langevin/\allowbreak{}Projection.\allowbreak{}lean}.\par}


For a continuous linear contraction \(T:H\to K\), reference point \(b\), and equilibrium projected radius \(c\), the parameter error obeys \((\int\|Tx-b\|^2d\mu)^{1/2}\le W_2(\mu,\pi)+c\). This derives the estimate directly from couplings.

\subsection{\texorpdfstring{projected\allowbreak{}RMS}{projectedRMS}}


Define the extended projected error \(R_{T,b}(\mu)=\|x\mapsto Tx-b\|_{L^2(\mu)}\). Using an extended norm retains infinity for a divergent second moment and allows the following transport comparison before any finiteness assumption.

\subsection{\texorpdfstring{projected\allowbreak{}RMSReal}{projectedRMSReal}}


Take the real value of the extended projected norm. Later estimates first prove that the extended quantity is finite; that guard is essential because \texttt{ENNReal.\allowbreak{}to\allowbreak{}Real} alone does not distinguish infinity from its totalized real value.

\subsection{\texorpdfstring{projected\allowbreak{}Difference le}{projectedDifference le}}


If \(\|T\|\le1\), then \(\|T(x-y)\|\le\|x-y\|\) pointwise. Monotonicity of the extended \(L^2\) norm bounds the projected difference norm by the original coupling cost.

\subsection{\texorpdfstring{projected\allowbreak{}RMS le coupling}{projectedRMS le coupling}}


On a coupling write \(Tx-b=T(x-y)+(Ty-b)\). Minkowski bounds its \(L^2\) norm by the sum of the two norms; the contraction estimate controls the first by coupling cost. Marginal identities rewrite the outer norms as \(R_{T,b}(\mu)\) and \(R_{T,b}(\pi)\).

\subsection{\texorpdfstring{projected\allowbreak{}RMS le wasserstein2}{projectedRMS le wasserstein2}}


The preceding bound holds for every coupling, with the same equilibrium radius. Take the infimum of coupling costs to obtain \(R_{T,b}(\mu)\le W_2(\mu,\pi)+R_{T,b}(\pi)\). Nonemptiness supplies the indexed-infimum identity; no minimizing plan is assumed.

\subsection{\texorpdfstring{projected\allowbreak{}RMS ne top}{projectedRMS ne top}}


Finite transport guarantees a coupling exists and makes \(W_2(\mu,\pi)\) finite. The assumed equilibrium radius is bounded by a finite real embedding. The transport-radius inequality therefore bounds the current projected norm by a sum of two finite extended values, proving its finiteness.

\subsection{\texorpdfstring{projected\allowbreak{}RMSReal eq sqrt integral}{projectedRMSReal eq sqrt integral}}


The explicit finite extended norm and measurability give \(x\mapsto Tx-b\in L^2(\mu)\). The \(L^2\) norm formula then identifies its real norm with \(\sqrt{\int\|Tx-b\|^2d\mu}\). This derivation guards both the real integral and the \texttt{to\allowbreak{}Real} conversion against hidden divergence.

\subsection{\texorpdfstring{projected\allowbreak{}RMS le of second\allowbreak{}Moment}{projectedRMS le of secondMoment}}


The supplied \(L^2\) membership makes the squared-norm integral genuine. Its bound by \(c^2\), together with \(c\ge0\), implies the root-mean-square norm is at most \(c\). Re-embed this finite real inequality to obtain the extended radius bound.

\subsection{\texorpdfstring{Convert the finite bound to real quantities}{Convert the finite bound to real quantities}}


First prove current-radius finiteness using finite transport and the finite equilibrium bound. Then convert the extended transport-radius inequality to real values, using finiteness of every summand. The result is \(R_{T,b}^{\mathbb R}(\mu)\le W_2^{\mathbb R}(\mu,\pi)+c\).

\section{\texorpdfstring{Compatibility exports for projection of quadratic transport}{Compatibility exports for projection of quadratic transport}}
\label{proof-module-IPLA-Transport-lean}
{\small\color{gray}Source module: \texttt{IPLA/\allowbreak{}Transport.\allowbreak{}lean}.\par}


This module exports the shared projected RMS and transport inequalities. The proof lives in \texttt{Unadjusted\allowbreak{}Langevin/\allowbreak{}Projection.\allowbreak{}lean}; IPLA instantiates it with its norm-one parameter projection.

\subsection{\texorpdfstring{Minkowski on each coupling gives the projected RMS estimate}{Minkowski on each coupling gives the projected RMS estimate}}


For a linear contraction \(T\), any center \(b\), and a coupling \(c\) of \(\mu,\pi\), write \(Tx-b=T(x-y)+(Ty-b)\). Minkowski and \(\|T(x-y)\|\le\|x-y\|\) imply

\[\|Tx-b\|_{L^2(\mu)}\le\operatorname{cost}(c)+\|Ty-b\|_{L^2(\pi)}.\]

Taking the infimum over couplings proves the same inequality with \(W_2(\mu,\pi)\). Finite transport plus a finite equilibrium radius proves finiteness of the projected moment. Only then does the shared real conversion identify its value with \(\sqrt{\int\|Tx-b\|^2d\mu}\). A second-moment bound at equilibrium supplies the concentration radius. No optimal coupling or separate Wasserstein triangle inequality is needed.



The \texttt{export} commands expose the already proved declarations under the IPLA namespace. Their definitions and proofs are unchanged, so the mathematical argument is the one in the imported source modules; the dependency diagram below opens those full proofs. This file itself makes no additional analytic assumption and contains no replacement proof.

\section{\texorpdfstring{Marginal partitions, existence and the exact parameter target}{Marginal partitions, existence and the exact parameter target}}
\label{proof-module-IPLA-MarginalConcentration-lean}
{\small\color{gray}Source module: \texttt{IPLA/\allowbreak{}Marginal\allowbreak{}Concentration.\allowbreak{}lean}.\par}


For differentiable jointly \(m\)-strongly convex \(U\), \(m>0\), this module proves finite positive latent partitions, continuity and vanishing parameter tails, existence of a maximizer, and the exact parameter projection of \(\Pi_N\). It establishes these facts before any logarithm, real conversion or concentration estimate is used.

\subsection{\texorpdfstring{Pair coordinates split the squared norm}{Pair coordinates split the squared norm}}


The pairing map has parameter coordinates \(\theta_j\) and latent coordinates \(x_j\). Expanding the Euclidean squared norm over the disjoint index sets gives

\[\|\operatorname{pair}(\theta,x)\|^2=\|\theta\|^2+\|x\|^2.\]

Every output coordinate of pairing is a continuous product projection; the finite coordinate-to-Euclidean identification is continuous as well. Hence \((\theta,x)\mapsto\operatorname{pair}(\theta,x)\) is jointly continuous. These elementary identities justify the orthogonal Gaussian domination and later dominated convergence.




\textbf{Correspondence with the displayed declarations.}


\begin{itemize}

\item \texttt{norm\_\allowbreak{}pair\_\allowbreak{}sq}: The pair coordinates preserve the orthogonal decomposition into parameter and latent variables.
\item \texttt{continuous\_\allowbreak{}pair}: Joint continuity of the parameter/\allowbreak{}latent pairing.
\end{itemize}
\subsection{\texorpdfstring{Strong convexity supplies uniform Gaussian domination}{Strong convexity supplies uniform Gaussian domination}}


The shared strong-convexity lower bound, obtained by bounding the linear term at zero, is

\[U(v)\ge U(0)-\frac{\|\nabla U(0)\|^2}{m}+\frac m4\|v\|^2.\]

Set \(C=\exp(\|\nabla U(0)\|^2/m-U(0))\) and \(a=m/4>0\). Exponentiation and the norm splitting yield

\[e^{-U(\theta,x)}\le Ce^{-a\|\theta\|^2}e^{-a\|x\|^2}\le Ce^{-a\|x\|^2}.\]

The final majorant is independent of \(\theta\) and integrable over the latent Euclidean space. Continuity supplies measurability, so domination proves integrability of every latent slice. No bounded parameter neighborhood is needed.




Strong convexity and the affine score term at zero give a Gaussian majorant after completing a square. Split the joint squared norm as \(\|\operatorname{pair}(\theta,x)\|^2=\|\theta\|^2+\|x\|^2\). Dropping the nonnegative parameter contribution yields an integrable Gaussian bound in \(x\), uniform in \(\theta\) when needed for local parameter limits. This proves latent integrability before converting the extended marginal density to its real integral. Strict positivity follows because the exponential weight is positive everywhere.




\textbf{Correspondence with the displayed declarations.}


\begin{itemize}

\item \texttt{marginal\_\allowbreak{}weight\_\allowbreak{}gaussian\_\allowbreak{}bound}: A Gaussian domination uniform in both the parameter and latent variable.
\item \texttt{marginal\_\allowbreak{}weight\_\allowbreak{}uniform\_\allowbreak{}bound}: An integrable bound for every parameter slice, with no local parameter restriction.
\item \texttt{integrable\_\allowbreak{}marginal\_\allowbreak{}weight}: Every latent partition is integrable from joint assumptions alone.
\end{itemize}
\subsection{\texorpdfstring{The marginal is positive, finite and continuous}{The marginal is positive, finite and continuous}}


Let \(M(\theta)=\int e^{-U(\theta,x)}dx\). Slice integrability and nonnegativity identify the extended marginal with \(\operatorname{ofReal}(M(\theta))\). This excludes infinity. The integrand is strictly positive everywhere, and Lebesgue measure is nonzero, so \(M(\theta)>0\) and the extended marginal is strictly positive as well. If \(\theta_n\to\theta\), the integrands converge pointwise by continuity of \(U\), and the preceding common Gaussian majorant permits dominated convergence. Thus \(M\) is continuous; composing with the continuous \(\operatorname{ofReal}\) map proves continuity of the extended marginal.




\textbf{Correspondence with the displayed declarations.}


\begin{itemize}

\item \texttt{marginal\allowbreak{}Density\_\allowbreak{}eq\_\allowbreak{}of\allowbreak{}Real\_\allowbreak{}integral}: The finite extended integral agrees with the ordinary real marginal.
\item \texttt{marginal\allowbreak{}Density\_\allowbreak{}ne\_\allowbreak{}top}: The parameter marginal is finite, excluding infinite-likelihood maximizers.
\item \texttt{marginal\allowbreak{}Density\_\allowbreak{}pos}: Every parameter has strictly positive marginal density.
\item \texttt{continuous\_\allowbreak{}marginal\_\allowbreak{}integral}: Continuity of the real-valued marginal follows by dominated convergence.
\item \texttt{continuous\_\allowbreak{}marginal\allowbreak{}Density}: The actual extended-valued marginal is continuous as well.
\end{itemize}
\subsection{\texorpdfstring{Gaussian parameter tails force existence of a maximizer}{Gaussian parameter tails force existence of a maximizer}}


Integrating the stronger Gaussian domination gives

\[0<M(\theta)\le C\left(\int e^{-a\|x\|^2}dx\right)e^{-a\|\theta\|^2}.\]

The finite constant is independent of \(\theta\), and the right-hand side tends to zero as \(\theta\) leaves compact sets. The squeeze theorem proves \(M(\theta)\to0\) at infinity. Since \(M(0)>0\), outside a sufficiently large compact region the values are below \(M(0)\). Continuity then attains a global maximum within a compact region. Monotonicity of \(\operatorname{ofReal}\) transfers that maximum to the extended marginal predicate. Uniqueness is proved later from Prékopa strong convexity.




Integrating the stronger Gaussian bound while retaining the parameter term gives \(M(\theta)\le C e^{-c\|\theta\|^2}\), for finite \(C\) and \(c>0\). Dominated convergence proves continuity of \(M\), and this bound proves \(M(\theta)\to0\) as \(\|\theta\|\to\infty\). Since \(M\) is positive at any chosen point, outside a sufficiently large ball it is smaller than that value. Compact attainment on the ball proves existence of the actual global marginal maximizer.




\textbf{Correspondence with the displayed declarations.}


\begin{itemize}

\item \texttt{marginal\_\allowbreak{}integral\_\allowbreak{}gaussian\_\allowbreak{}bound}: Joint strong convexity gives a Gaussian parameter tail after integrating out the latent variable. The constant is finite and independent of \texttt{theta}.
\item \texttt{tendsto\_\allowbreak{}marginal\_\allowbreak{}integral\_\allowbreak{}cocompact}: The parameter marginal vanishes at infinity.
\item \texttt{exists\_\allowbreak{}marginal\allowbreak{}Maximizer}: Existence of the actual marginal maximizer is derived from the original joint assumptions. No maximizer or compactness premise is introduced.
\end{itemize}
\subsection{\texorpdfstring{Product integration and the particle-\allowbreak{}scaling Jacobian}{Product integration and the particle-scaling Jacobian}}


Tonelli induction over \(n\) proves the extended-integral identity

\[\int\prod_{i=1}^nf(x_i)\,d\mu^{\otimes n}=\left(\int f\,d\mu\right)^n.\]

The base case is the mass-one empty product; the successor case separates one coordinate and multiplies the remaining integral. This statement records divergence correctly because no real conversion is used. For a nonzero real dilation \(r\) in dimension \(d\), Haar measure scaling gives \(\int f(rx)dx=|r^d|^{-1}\int f(x)dx\). Applying these two identities at \(r=\sqrt N\), with \(N>0\), yields

\[\int\prod_i e^{-U(\theta,\sqrt N z_i)}\,dz_1\cdots dz_N=\left(N^{-d_x/2}M(\theta)\right)^N.\]

The formal statement retains the absolute-value inverse-power Jacobian, which is equal to the displayed factor.




For a fixed \(\theta\), the lifted weight factors as \(\prod_i e^{-U(\theta,\sqrt N z_i)}\). Tonelli factors the nonnegative product integral into a product of identical latent integrals. The substitution \(x_i=\sqrt N z_i\) contributes \(N^{-d_x/2}\) per particle, so the fiber integral is \(N^{-Nd_x/2}M(\theta)^N\). This Jacobian is finite and strictly positive when \(N>0\), and does not depend on \(\theta\).




\textbf{Correspondence with the displayed declarations.}


\begin{itemize}

\item \texttt{lintegral\_\allowbreak{}fin\_\allowbreak{}prod\_\allowbreak{}eq\_\allowbreak{}pow}: Tonelli factorization of identical nonnegative weights over finitely many independent latent coordinates. This extended-integral statement also records divergence correctly; no integrability premise is hidden.
\item \texttt{lintegral\_\allowbreak{}comp\_\allowbreak{}smul\_\allowbreak{}eq}: Change of variables for a nonnegative weight under scalar dilation.
\item \texttt{scaled\_\allowbreak{}particle\_\allowbreak{}partition}: Integrating the actual scaled particle potential at a fixed parameter gives the \texttt{N}th power of the original marginal, times the explicit scalar Jacobian. This is the factor that cancels in the normalized parameter law.
\end{itemize}
\subsection{\texorpdfstring{A coordinate equivalence groups the state into \ensuremath{\theta} and particle blocks}{A coordinate equivalence groups the state into θ and particle blocks}}


The configuration map \(C(\theta,(z_i)_i)\) inserts the parameter coordinates first and then all particle coordinates. Its coordinate evaluation theorems follow directly from insertion. The inverse is \(z\mapsto(p_\theta z,(\operatorname{block}_iz)_i)\); checking each coordinate proves both inverse identities and linearity, hence a finite-dimensional continuous linear equivalence. The evaluation theorem identifies this equivalence with the configuration definition, while \(p_\theta C(\theta,z)=\theta\) and \(A_iC(\theta,z)=\operatorname{pair}(\theta,\sqrt N z_i)\). These identities give the exact arguments of the fiber integrals.




\textbf{Correspondence with the displayed declarations.}


\begin{itemize}

\item \texttt{configuration}: Assemble the scaled state from its parameter and particle blocks.
\item \texttt{configuration\allowbreak{}CLE}: Grouping coordinates is a constructed continuous linear equivalence.
\end{itemize}
\subsection{\texorpdfstring{The weight factors and its parameter projection integrates each fiber}{The weight factors and its parameter projection integrates each fiber}}


In grouped coordinates, exponentiation converts the sum in \(V_N\) into a product:

\[e^{-V_N(C(\theta,z))}=\prod_i e^{-U(\theta,\sqrt N z_i)}.\]

Two measure identities justify using this pointwise formula. First, for a measurable equivalence \(e\), pushing \(\mu\) with weight \(w\circ e\) through \(e\) gives \((e_\#\mu)\) with weight \(w\); prove this by evaluating any measurable set and the pushforward integral formula. Second, projecting \(w(\theta,z)d\mu(\theta)d\nu(z)\) to its first coordinate gives density \(\int w(\theta,z)d\nu(z)\) relative to \(\mu\), by Tonelli on the indicator of an arbitrary measurable parameter set.




\textbf{Correspondence with the displayed declarations.}


\begin{itemize}

\item \texttt{weight\_\allowbreak{}configuration}: In grouped coordinates the actual Gibbs weight factors over particles.
\item \texttt{map\_\allowbreak{}with\allowbreak{}Density\_\allowbreak{}comp\_\allowbreak{}equiv}: A density commutes with a measurable coordinate equivalence.
\item \texttt{map\_\allowbreak{}fst\_\allowbreak{}with\allowbreak{}Density}: Tonelli's theorem identifies the first-coordinate density of a weighted product measure.
\end{itemize}
\subsection{\texorpdfstring{All finite positive volume constants cancel on normalization}{All finite positive volume constants cancel on normalization}}


A continuous linear equivalence carries Euclidean volume to a Haar measure, hence to \(c\) times the target volume for a finite positive \(c\). This proves the required coordinate-volume relation without assuming its scalar is one. For any measure \(\mu\) and \(0<c<\infty\), extended-real inverse multiplication gives

\[\frac{c\mu}{(c\mu)(E)}=\frac{\mu}{\mu(E)}.\]

Thus both the coordinate-grouping scalar and the particle-dilation Jacobian cancel in a normalized law. The formal normalization lemma proves this algebra for arbitrary \(\mu\), keeping the nonzero and finite conditions on \(c\) explicit.




\textbf{Correspondence with the displayed declarations.}


\begin{itemize}

\item \texttt{normalize\_\allowbreak{}smul\_\allowbreak{}eq}: Normalization is unchanged by multiplication by any finite positive constant. This is why grouping coordinates and the particle Jacobian cancel.
\item \texttt{configuration\_\allowbreak{}map\_\allowbreak{}volume}: The finite-dimensional coordinate grouping changes volume by a constructed finite positive constant. Its value is irrelevant after normalization.
\end{itemize}
\subsection{\texorpdfstring{The actual parameter marginal is the powered likelihood target}{The actual parameter marginal is the powered likelihood target}}


Define the normalized power target \(\pi_N^\theta(d\theta)=Z_\theta^{-1}M(\theta)^N d\theta\). To identify it, let \(\mu=e^{-V_N}dz\) and \(\rho=M(\theta)^Nd\theta\). Grouping coordinates changes volume by \(c\), and the fiber calculation contributes \(J^N\), \(J=N^{-d_x/2}\). The preceding pushforward identities prove \(c\,p_{\theta\#}\mu=J^N\rho\). Both scalars are finite and positive. Normalize this equality, cancel both scalars, and use preservation of total mass under measurable projection to conclude

\[p_{\theta\#}\Pi_N=\pi_N^\theta.\]

Every change of variables and normalization factor in this identity is proved; no formula for the marginal is assumed.




The configuration equivalence groups the orthogonal state coordinates into the parameter block and all latent blocks. The volume theorem tracks the finite-dimensional coordinate-change constant, while the density pushforward lemmas integrate each parameter fiber. The preceding factorization leaves a density proportional to \(M(\theta)^N\), with the scaling and coordinate-change factors independent of \(\theta\). Normalization cancels those positive finite factors exactly, yielding \[(p_\theta)_\#\Pi_N(d\theta)=\frac{M(\theta)^N\,d\theta}{\int M(u)^Ndu}.\] The final target probability and moment lemmas are consequences of this identity and the already proved lifted target facts.




\textbf{Correspondence with the displayed declarations.}


\begin{itemize}

\item \texttt{marginal\allowbreak{}Target}: The parameter law with density proportional to the \texttt{N}th power of the original latent marginal. This names the statistical target without assuming anything about its concentration.
\end{itemize}
\subsection{\texorpdfstring{Normalize the power target and guard its moments}{Normalize the power target and guard its moments}}


The established projection identity and the probability property of \(\Pi_N\) prove that \(\pi_N^\theta\) is a probability measure. Its unnormalized density \(M^N\) is everywhere positive, so its extended partition is strictly positive. If that partition were infinite, the definition of the normalized target would be the zero measure, contradicting mass one; thus it is also finite. Finally the lifted Gibbs law has a finite second moment and \(p_\theta\) is a bounded linear map. The pushforward integral identity therefore gives \(\theta\mapsto\theta\in L^2(\pi_N^\theta)\). This supplies finite marginal moments before the later sharp bound is applied.




\textbf{Correspondence with the displayed declarations.}


\begin{itemize}

\item \texttt{marginal\allowbreak{}Target\_\allowbreak{}is\allowbreak{}Probability}: The marginal power density normalizes to a probability measure as a consequence of its proved equality with the actual Gibbs projection.
\item \texttt{marginal\allowbreak{}Partition\_\allowbreak{}pos}: The normalizer of the powered marginal density is strictly positive.
\item \texttt{marginal\allowbreak{}Partition\_\allowbreak{}ne\_\allowbreak{}top}: The powered marginal partition is finite. Its proof uses the previously established probability identity, so divergence cannot disappear in a real conversion or an infinite normalizing constant.
\item \texttt{marginal\allowbreak{}Target\_\allowbreak{}mem\allowbreak{}Lp}: The parameter marginal has finite second moment, without a separate marginal moment or concentration assumption.
\end{itemize}
\section{\texorpdfstring{One-\allowbreak{}dimensional Prékopa–Leindler by constructed monotone transport}{One-dimensional Prékopa–Leindler by constructed monotone transport}}
\label{proof-module-IPLA-PrekopaLeindler-lean}
{\small\color{gray}Source module: \texttt{IPLA/\allowbreak{}Prekopa\allowbreak{}Leindler.\allowbreak{}lean}.\par}


The proof constructs transport directly from cumulative integrals of positive continuous densities. It derives the one-dimensional Prékopa–Leindler inequality without postulating a transport map or marginal log-concavity theorem.

\subsection{\texorpdfstring{Cumulative integrals are differentiable and strictly increasing}{Cumulative integrals are differentiable and strictly increasing}}


A density here is a continuous strictly positive real function \(f\), integrable on \(\mathbb R\), with \(\int f=1\). Define \(F(x)=\int_{(-\infty,x]}f\). The difference formula

\[F(y)-F(0)=\int_0^y f(t)dt\]

and the fundamental theorem of calculus give the strict derivative \(F'(x)=f(x)\). Differentiability implies continuity, and positivity of the derivative implies strict monotonicity. The density assumptions are precisely the inputs required for these assertions.




\textbf{Correspondence with the displayed declarations.}


\begin{itemize}

\item \texttt{Density}: A positive, continuous probability density on the real line.
\item \texttt{cumulative}: Cumulative integral of an ordinary real density.
\end{itemize}
\subsection{\texorpdfstring{The cumulative range is exactly the open unit interval}{The cumulative range is exactly the open unit interval}}


Integrability implies \(F(x)\to0\) as \(x\to-\infty\). Splitting the total integral at \(x\) gives \(F(x)=1-\int_{(x,\infty)}f\), so \(F(x)\to1\) as \(x\to+\infty\). Nonnegativity gives \(0\le F\le1\); strict monotonicity, comparing \(x\) with \(x-1\) and \(x+1\), improves this to \(0<F(x)<1\). For every \(u\in(0,1)\), the endpoint limits supply points on either side of \(u\), and the intermediate value theorem supplies \(F(x)=u\). Therefore \(F(\mathbb R)=(0,1)\).

\subsection{\texorpdfstring{Construct and differentiate the monotone transport}{Construct and differentiate the monotone transport}}


For two densities with cumulative maps \(F,G\), set \(T=G^{-1}\circ F\), using the inverse function on its range. Since both ranges are \((0,1)\), \(G(T(x))=F(x)\) everywhere. Strict monotonicity makes \(G\) injective. The one-dimensional inverse function theorem applies because \(G'(T(x))=g(T(x))>0\), and the chain rule yields

\[T'(x)=\frac{f(x)}{g(T(x))}>0.\]

Thus \(T\) is strictly increasing, and the density ratio is a proved derivative of an explicitly constructed function.




For positive continuous unit-mass densities, \(F(x)=\int_{-\infty}^xf\) and \(G(y)=\int_{-\infty}^yg\) are differentiable with strictly positive derivatives, hence strictly increasing, and their limits give range \((0,1)\). Define \(T=G^{-1}\circ F\), so \(G(T(x))=F(x)\). The inverse-function derivative gives \(T'(x)=f(x)/g(T(x))>0\). The cumulative identities and strict monotonicity declarations justify this change of variables on the whole real line.




\textbf{Correspondence with the displayed declarations.}


\begin{itemize}

\item \texttt{transport}: The monotone transport from density \texttt{f} to density \texttt{g}.
\end{itemize}
\subsection{\texorpdfstring{The normalized inequality is a change-\allowbreak{}of-\allowbreak{}variables calculation}{The normalized inequality is a change-of-variables calculation}}


Assume \(0\le t\le1\) and \(h((1-t)x+ty)\ge f(x)^{1-t}g(y)^t\). Put \(S(x)=(1-t)x+tT(x)\) and \(r=f(x)/g(T(x))>0\). Then \(S'=1-t+tr>0\), and convexity of the exponential gives \(r^t\le1-t+tr=S'\). At \(y=T(x)\), the pointwise premise therefore implies

\[f(x)=f(x)^{1-t}g(T(x))^t r^t\le h(S(x))S'(x).\]

Since \(S\) is strictly increasing, the one-dimensional Jacobian theorem changes the right-hand integral to the integral of \(h\) over \(S(\mathbb R)\). Nonnegativity then bounds it by the integral over all \(\mathbb R\), proving \(1\le\int h\). The formal conclusion uses an extended nonnegative integral, so \(h\) need not itself be integrable or continuous.




Set \(S(x)=(1-t)x+tT(x)\). The interpolation hypothesis bounds \(h(S(x))\) below by \(f(x)^{1-t}g(T(x))^t\). Weighted arithmetic–geometric mean gives \(S'(x)=(1-t)+tT'(x)\ge T'(x)^t\). Since \(T'=f/g(T)\), multiplication yields \(h(S(x))S'(x)\ge f(x)\). The strictly increasing change-of-variables formula then proves \(\int h\ge\int f=1\). For arbitrary positive finite masses, divide \(f,g\) by their integrals, apply the normalized statement, and multiply back; logarithmic mass terms recover \((\int f)^{1-t}(\int g)^t\).




\textbf{Correspondence with the displayed declarations.}


\begin{itemize}

\item \texttt{normalized\_\allowbreak{}prekopa\_\allowbreak{}leindler}: The normalized one-dimensional Prékopa--Leindler inequality, proved by the constructed monotone transport. The third weight need not be continuous or integrable; the conclusion is an extended nonnegative integral.
\end{itemize}
\subsection{\texorpdfstring{Remove the unit-\allowbreak{}integral normalization}{Remove the unit-integral normalization}}


An everywhere positive integrable function on a nonzero measure space has strictly positive integral, since its support has positive measure. For the Prékopa weights write \(A=\int f>0\), \(B=\int g>0\), and \(C=A^{1-t}B^t>0\). Apply the normalized inequality to \(f/A\), \(g/B\), and \(h/C\). Expanding \(\log(f/A)=\log f-\log A\) verifies its pointwise premise. If \(h\) is nonnegative and integrable, the extended conclusion converts to \(1\le(\int h)/C\). Multiplication by \(C\) gives

\[\exp\!\left((1-t)\log\int f+t\log\int g\right)\le\int h.\]

The positivity lemma also applies to the zero-dimensional product measure needed later.




\textbf{Correspondence with the displayed declarations.}


\begin{itemize}

\item \texttt{integral\_\allowbreak{}positive}: Strict positivity of the real integral of an everywhere positive integrable function, including the zero-dimensional product space used in tensorization.
\item \texttt{prekopa\_\allowbreak{}leindler\_\allowbreak{}real}: Prékopa--Leindler on the real line for arbitrary positive continuous integrable first and second weights. No normalization is assumed.
\end{itemize}
\section{\texorpdfstring{Finite-\allowbreak{}dimensional Prékopa–Leindler with Gaussian tails}{Finite-dimensional Prékopa–Leindler with Gaussian tails}}
\label{proof-module-IPLA-PrekopaTensorization-lean}
{\small\color{gray}Source module: \texttt{IPLA/\allowbreak{}Prekopa\allowbreak{}Tensorization.\allowbreak{}lean}.\par}


Positive continuous weights dominated by an integrable Gaussian form a class closed under integrating out one coordinate. This supplies all continuity and integrability facts needed to iterate the one-dimensional transport proof.

\subsection{\texorpdfstring{Gaussian domination proves integrability and positive integral}{Gaussian domination proves integrability and positive integral}}


Write \(\mathrm{Coordinates}(n)=\mathbb R^n\). A Gaussian weight satisfies continuity, strict positivity, and \(f(x)\le Ce^{-a\sum_i x_i^2}\), with \(C,a>0\). The Euclidean Gaussian integral is finite, and the coordinate identification preserves volume, so domination proves \(f\in L^1\). Positivity on the whole nonzero volume space proves \(\int f>0\). Multiplying by \(c>0\) preserves these properties and changes the envelope constant to \(cC\).




\textbf{Correspondence with the displayed declarations.}


\begin{itemize}

\item \texttt{Gaussian\allowbreak{}Weight}: The class of positive continuous weights with a Gaussian upper bound.
\end{itemize}
\subsection{\texorpdfstring{A coordinate slice and its marginal are integrable and continuous}{A coordinate slice and its marginal are integrable and continuous}}


The insertion map \((s,x)\mapsto(s,x_1,\ldots,x_n)\) is continuous coordinate by coordinate. Define \(f_1(x)=\int_\mathbb R f(s,x)ds\). Since \(\sum_i(s,x)_i^2=s^2+\sum_i x_i^2\ge s^2\), every slice is bounded by \(Ce^{-as^2}\), which is integrable independently of \(x\). Slice continuity supplies measurability and domination proves slice integrability. If \(x_j\to x\), joint continuity gives pointwise convergence in \(s\), and the same Gaussian bound allows dominated convergence. Hence \(f_1\) is continuous.




Write a point in \(\mathbb R^{n+1}\) as \(\operatorname{cons}(z,x)\), where \(z\in\mathbb R\) is the coordinate being integrated. The coordinate-insertion continuity theorem checks its first and remaining coordinates separately. A Gaussian weight has constants \(C,a>0\) satisfying

\[0<f(\operatorname{cons}(z,x))\le C\exp\!\left[-a\left(z^2+\sum_{i=1}^{n}x_i^2\right)\right]\le C e^{-az^2}.\]

The last inequality drops a nonnegative sum of squares, so the majorant is independent of \(x\). It is integrable in \(z\) by the one-dimensional Gaussian integral theorem. The declaration \texttt{slice\_\allowbreak{}integrable} uses positivity to identify the absolute value of the slice with the slice itself, continuity to obtain measurability, and this majorant to prove integrability.



Now let \(x_m\to x\). For every fixed \(z\), continuity of \(f\) and of coordinate insertion gives \(f(\operatorname{cons}(z,x_m))\to f(\operatorname{cons}(z,x))\). The same Gaussian majorant controls every term and the limit. Dominated convergence therefore gives convergence of their real integrals, proving continuity of the first marginal. In the Lean invocation of \texttt{continuous\_\allowbreak{}of\_\allowbreak{}dominated}, the four obligations are slice measurability, the uniform almost-everywhere bound, integrability of that bound, and pointwise continuity in the remaining coordinates. Those are all established before the induction later treats the marginal as another Gaussian weight.




\textbf{Correspondence with the displayed declarations.}


\begin{itemize}

\item \texttt{continuous\_\allowbreak{}cons}: Continuity of inserting the first real coordinate.
\item \texttt{first\allowbreak{}Marginal}: Integrate out the first coordinate of a weight.
\end{itemize}
\subsection{\texorpdfstring{Coordinate integration preserves Gaussian tails and total integral}{Coordinate integration preserves Gaussian tails and total integral}}


Retaining the full Gaussian envelope and factoring its exponent gives

\[f_1(x)\le C\left(\int_\mathbb R e^{-as^2}ds\right)e^{-a\sum_i x_i^2}.\]

The new constant is positive and finite, and integrating the positive slice makes \(f_1\) positive everywhere. Together with continuity, this proves that \(f_1\) is again a Gaussian weight. The coordinate splitting equivalence preserves product volume. Its integral formula and Fubini, justified by the Gaussian integrability already proved, then show \(\int f_1=\int f\).

\subsection{\texorpdfstring{Induct on the dimension}{Induct on the dimension}}


For \(n=0\), the coordinate space is a singleton with volume one; each integral equals its function value and the result is the pointwise premise. For dimension \(n+1\), fix remaining-coordinate vectors \(x,y\). The one-dimensional theorem applied to the slices \(s\mapsto f(s,x)\), \(r\mapsto g(r,y)\), and \(u\mapsto h(u,(1-t)x+ty)\) gives the pointwise Prékopa inequality for \(f_1,g_1,h_1\). These marginals are Gaussian weights, so the induction hypothesis applies in dimension \(n\). Replace their integrals by those of the original weights to obtain

\[\left(\int f\right)^{1-t}\left(\int g\right)^t\le\int h.\]

The formal statement writes positive powers using exponentials and logarithms, including \(t=0\) and \(t=1\).




The induction splits \(\mathbb R^{d+1}\) into its first coordinate and \(d\) remaining coordinates. For fixed remaining coordinates, apply the proved one-dimensional inequality to the first-coordinate slices. Their integrals define marginal functions that satisfy the same interpolation inequality in \(d\) dimensions. Gaussian domination proves slice integrability, continuity of those marginal functions by dominated convergence, and Gaussian tails for them. The induction hypothesis applies, and Fubini identifies the iterated integrals with the original total integrals. Dimension zero is the one-point base case.




\textbf{Correspondence with the displayed declarations.}


\begin{itemize}

\item \texttt{prekopa\_\allowbreak{}leindler}: The finite-dimensional Prékopa--Leindler inequality. The induction uses the proved real-line transport theorem and integrates one coordinate at a time.
\end{itemize}
\section{\texorpdfstring{Derived strong convexity of the actual marginal potential}{Derived strong convexity of the actual marginal potential}}
\label{proof-module-IPLA-StrongLogConcavity-lean}
{\small\color{gray}Source module: \texttt{IPLA/\allowbreak{}Strong\allowbreak{}Log\allowbreak{}Concavity.\allowbreak{}lean}.\par}


Let \(M(\theta)=\int e^{-U(\theta,x)}dx\) and \(\Phi(\theta)=-\log M(\theta)\). The proved finite-dimensional Prékopa–Leindler theorem shows that \(\Phi\) retains the original joint strong-convexity constant \(m\).

\subsection{\texorpdfstring{Define a legitimate logarithmic marginal potential}{Define a legitimate logarithmic marginal potential}}


The real marginal integral \(M\) is defined from the original slice; the earlier domination theorem makes it finite and strictly positive. Thus \(\Phi=-\log M\) is an ordinary real-valued continuous function: continuity follows by composing the continuous marginal with logarithm away from zero. If \(\theta_*\) maximizes the extended marginal, its finite real representation and monotonicity of logarithm imply \(\Phi(\theta_*)\le\Phi(\theta)\) for every \(\theta\). No logarithm of zero or infinite partition is used.




\textbf{Correspondence with the displayed declarations.}


\begin{itemize}

\item \texttt{marginal\allowbreak{}Integral}: The real latent partition. Positivity and finiteness were proved from the joint assumptions in \texttt{Marginal\allowbreak{}Concentration.\allowbreak{}lean}.
\item \texttt{marginal\allowbreak{}Potential}: The actual marginal negative log-likelihood, with the observations fixed.
\item \texttt{marginal\allowbreak{}Maximizer\_\allowbreak{}minimizes\_\allowbreak{}potential}: The statistical marginal maximizer minimizes the actual negative log-likelihood. All real conversions use proved positive finite partitions.
\end{itemize}
\subsection{\texorpdfstring{Prepare Gaussian weights in product coordinates}{Prepare Gaussian weights in product coordinates}}


Define \(f_\theta(x)=e^{-U(\theta,x)}\) in ordinary finite-product coordinates, using the volume-preserving Euclidean identification for the latent block. The joint Gaussian domination gives \(f_\theta(x)\le Ce^{-m\|x\|^2/4}\), and continuity plus strict positivity make it a Gaussian weight. Its integral equals \(M(\theta)\) by the volume-preserving change of coordinates. Finally pairing preserves subtraction and affine combinations coordinatewise:

\[\operatorname{pair}(\theta,x)-\operatorname{pair}(\theta',y)=\operatorname{pair}(\theta-\theta',x-y),\]\[\operatorname{pair}(a\theta+b\theta',ax+by)=a\operatorname{pair}(\theta,x)+b\operatorname{pair}(\theta',y).\]


\textbf{Correspondence with the displayed declarations.}


\begin{itemize}

\item \texttt{latent\allowbreak{}Weight}: A latent slice written in ordinary product coordinates for applying the finite-dimensional Prékopa--Leindler theorem.
\end{itemize}
\subsection{\texorpdfstring{Apply Prékopa with the parameter convexity penalty}{Apply Prékopa with the parameter convexity penalty}}


Let \(a,b\ge0\), \(a+b=1\), \(\theta_r=a\theta+b\theta'\), and \(c=(mab/2)\|\theta-\theta'\|^2\). Strong convexity of \(U\), the affine pair identity and norm splitting imply

\[aU(\theta,x)+bU(\theta',y)\ge U(\theta_r,ax+by)+c,\]

after discarding the nonnegative latent penalty \((mab/2)\|x-y\|^2\). Exponentiate and apply Prékopa–Leindler with third weight \(e^{-c}f_{\theta_r}\), which remains a Gaussian weight. This gives \(M(\theta)^aM(\theta')^b\le e^{-c}M(\theta_r)\). All factors are positive, so take logarithms and negate to conclude

\[\Phi(a\theta+b\theta')\le a\Phi(\theta)+b\Phi(\theta')-\frac{mab}{2}\|\theta-\theta'\|^2.\]

Thus marginal strong convexity is a theorem of the joint assumptions.




For \(\theta_r=a\theta+b\vartheta\), \(a+b=1\), joint strong convexity yields \[U(\theta_r,ax+by)\le aU(\theta,x)+bU(\vartheta,y)-\tfrac{mab}{2}\|\theta-\vartheta\|^2,\] after dropping the additional nonnegative latent deficit. Thus the latent weights satisfy the Prékopa hypothesis with the multiplier \(e^{-c}\), \(c=mab\|\theta-\vartheta\|^2/2\), on the interpolated weight. Integrating gives \(M(\theta)^aM(\vartheta)^b\le e^{-c}M(\theta_r)\).



Every marginal integral is already finite and positive, so logarithms preserve the inequality. Negating gives \(F(\theta_r)\le aF(\theta)+bF(\vartheta)-c\), exactly \(m\)-strong convexity of \(F=-\log M\). The code's \texttt{hpoint} is the weight inequality, \texttt{hpl} its Prékopa integral consequence, and \texttt{hlog} its logarithmic conversion. Strict convexity for \(m>0\) then proves uniqueness of the previously constructed maximizer.




\textbf{Correspondence with the displayed declarations.}


\begin{itemize}

\item \texttt{marginal\allowbreak{}Potential\_\allowbreak{}strong\allowbreak{}Convex}: Integrating out the latent variables preserves precisely the original strong-convexity constant. This is the derived Prékopa theorem for IPLA.
\end{itemize}
\subsection{\texorpdfstring{The marginal maximizer exists uniquely}{The marginal maximizer exists uniquely}}


Because \(m>0\), the proved \(m\)-strong convexity makes \(\Phi\) strictly convex. Two marginal maximizers are both global minimizers of \(\Phi\); if distinct, the midpoint strict-convexity inequality would lie strictly below their common minimum, a contradiction. Hence they coincide. The earlier continuity-and-tail argument already supplied existence of a maximizer. Combining existence with this uniqueness yields \(\exists!\theta_*\) satisfying the original statistical maximizer predicate.




\textbf{Correspondence with the displayed declarations.}


\begin{itemize}

\item \texttt{marginal\allowbreak{}Maximizer\_\allowbreak{}unique}: Joint strong convexity implies uniqueness of the marginal maximizer.
\item \texttt{exists\allowbreak{}Unique\_\allowbreak{}marginal\allowbreak{}Maximizer}: The marginal maximizer exists uniquely under exactly the original joint potential assumptions.
\end{itemize}
\section{\texorpdfstring{Concentration around a minimizer without differentiability}{Concentration around a minimizer without differentiability}}
\label{proof-module-UnadjustedLangevin-GibbsConcentration-lean}
{\small\color{gray}Source module: \texttt{Unadjusted\allowbreak{}Langevin/\allowbreak{}Gibbs\allowbreak{}Concentration.\allowbreak{}lean}.\par}


For a continuous \(q\)-strongly convex potential \(F\), \(q>0\), and an explicitly supplied minimizer \(z\), this module proves the equilibrium radius \(\int\|x-z\|^2d\pi_F\le2d/q\). The proof uses dilation and convexity, so it also applies to the nonsmooth potentials used by proximal algorithms.

\subsection{\texorpdfstring{strong\allowbreak{}Convex mode quadratic bound}{strongConvex mode quadratic bound}}


Use strong convexity at the midpoint of \(z\) and \(x\), then use minimality \(F(z)\le F((z+x)/2)\). Rearranging yields \(F(x)\ge F(z)+(q/4)\|x-z\|^2\). This argument uses values of \(F\), not a derivative or a first-order optimality assertion at the minimizer.

\subsection{\texorpdfstring{integrable exp neg of continuous strong\allowbreak{}Convex mode}{integrable exp neg of continuous strongConvex mode}}


For \(q>0\), the midpoint bound gives \(e^{-F(x)}\le e^{-F(z)}e^{-(q/4)\|x-z\|^2}\). Translation invariance reduces integrability of the majorant to the centered Gaussian integral. Continuity makes the density measurable, so domination proves integrability without differentiability of \(F\).

\subsection{\texorpdfstring{integrable sq dist mul exp neg of continuous strong\allowbreak{}Convex mode}{integrable sq dist mul exp neg of continuous strongConvex mode}}


Multiply the translated Gaussian bound by \(\|x-z\|^2\). Translation preserves Lebesgue measure, and a quadratic times a decaying Gaussian is integrable. Lean transports that integrability through \(x\mapsto x-z\) and then applies the nonnegative domination theorem to the desired weighted moment.

\subsection{\texorpdfstring{Contract toward the mode}{Contract toward the mode}}


For \(0<r<1\), compare \(F(z+r(x-z))\) with \(F(x)\) using strong convexity and \(F(z)\le F(x)\). Exponentiating and using \(e^u\ge1+u\) produces a linear correction in \(\|x-z\|^2\). The dilation \(x\mapsto z+r(x-z)\) multiplies volume by \(r^d\), hence \(Z+(q/2)r(1-r)M\le r^{-d}Z\), with \(M=\int\|x-z\|^2e^{-F(x)}dx\).

\subsection{\texorpdfstring{integral sq dist mul exp neg le of continuous strong\allowbreak{}Convex mode}{integral sq dist mul exp neg le of continuous strongConvex mode}}


Subtract \(Z\) from the dilation inequality and divide by \(1-r>0\). As \(r\uparrow1\), \((r^{-d}-1)/(1-r)\to d\), while \((q/2)rM\to(q/2)M\). The order-closed limit theorem yields \(M\le(2d/q)Z\). Lean proves the one-sided derivative limit of the scalar power before passing the inequality to the limit.

\subsection{\texorpdfstring{integrable sq dist gibbs of continuous strong\allowbreak{}Convex mode}{integrable sq dist gibbs of continuous strongConvex mode}}


The density-measure integrability criterion changes the desired \(\pi_F\)-integral into a constant multiple of \(\int\|x-z\|^2e^{-F(x)}dx\). Its finiteness follows from the continuous, strongly convex mode estimate. No score moment or smoothness assumption is introduced during normalization.

\subsection{\texorpdfstring{integral sq dist gibbs le of continuous strong\allowbreak{}Convex mode}{integral sq dist gibbs le of continuous strongConvex mode}}


First prove \(Z>0\), \(Z<\infty\), and integrability of the weighted distance. Then write \(\int\|x-z\|^2d\pi_F=M/Z\) and apply \(M\le(2d/q)Z\). Cancellation is justified by positivity of \(Z\); the result is \(\mathbb E_{\pi_F}\|X-z\|^2\le2d/q\), including \(d=0\).

\section{\texorpdfstring{The dimension-\allowbreak{}sharp equilibrium parameter concentration}{The dimension-sharp equilibrium parameter concentration}}
\label{proof-module-IPLA-Concentration-lean}
{\small\color{gray}Source module: \texttt{IPLA/\allowbreak{}Concentration.\allowbreak{}lean}.\par}


The exact parameter marginal has density proportional to \(M(\theta)^N=e^{-N\Phi(\theta)}\). Its potential is \(Nm\)-strongly convex in dimension \(d_\theta\), giving the equilibrium radius \(\sqrt{2d_\theta/(Nm)}\) without a separate concentration hypothesis.

\subsection{\texorpdfstring{Identify the scaled marginal potential and its Gibbs measure}{Identify the scaled marginal potential and its Gibbs measure}}


Multiply the strong-convexity inequality for \(\Phi\) by the nonnegative scalar \(N\). Every affine term and the quadratic penalty scale by \(N\), so \(N\Phi\) is \(Nm\)-strongly convex. Since \(M>0\),

\[e^{-N\Phi(\theta)}=e^{N\log M(\theta)}=M(\theta)^N.\]

The finite-real representation of the original extended marginal identifies this weight exactly with the one defining \texttt{marginal\allowbreak{}Target}. Its normalized law is therefore the shared Gibbs measure of \(N\Phi\).




\textbf{Correspondence with the displayed declarations.}


\begin{itemize}

\item \texttt{scaled\_\allowbreak{}marginal\allowbreak{}Potential\_\allowbreak{}strong\allowbreak{}Convex}: Multiplying the actual marginal potential by \texttt{N} multiplies its proved strong-convexity constant by \texttt{N}.
\item \texttt{marginal\allowbreak{}Target\_\allowbreak{}eq\_\allowbreak{}gibbs\allowbreak{}Measure}: The normalized powered marginal density is exactly the Gibbs law of the proved \texttt{N}-scaled negative log marginal.
\end{itemize}
\subsection{\texorpdfstring{A dilation estimate gives the sharp parameter moment}{A dilation estimate gives the sharp parameter moment}}


A marginal maximizer \(\theta_*\) minimizes \(\Phi\), hence also \(N\Phi\). The shared continuous strongly convex Gibbs theorem, which needs continuity and a mode rather than differentiability of \(\Phi\), bounds the squared distance by twice dimension over convexity:

\[\int\|\theta-\theta_*\|^2d\pi_N^\theta(\theta)\le\frac{2d_\theta}{Nm}.\]

Apply it to the continuous potential \(N\Phi\), use the proved \(Nm\)-strong convexity, and rewrite its Gibbs law by the preceding identity. Only the parameter dimension enters; latent dimension does not enlarge this equilibrium radius.




\textbf{Correspondence with the displayed declarations.}


\begin{itemize}

\item \texttt{marginal\allowbreak{}Target\_\allowbreak{}second\allowbreak{}Moment\_\allowbreak{}le}: Exact finite second-moment bound about any parameter maximizing the actual latent marginal. Its existence and uniqueness were proved separately.
\end{itemize}
\subsection{\texorpdfstring{Transfer the marginal moment to the lifted Gibbs target}{Transfer the marginal moment to the lifted Gibbs target}}


The exact projection identity \(p_{\theta\#}\Pi_N=\pi_N^\theta\) and the measurable pushforward integral formula yield

\[\int\|p_\theta z-\theta_*\|^2d\Pi_N(z)\le\frac{2d_\theta}{Nm}.\]

The projected Gibbs variable is already in \(L^2\), and subtracting the constant \(\theta_*\) preserves this. The shared projected-RMS theorem converts the real moment inequality to the extended RMS bound

\[\operatorname{RMS}_{\Pi_N}(p_\theta-\theta_*)\le\sqrt{2d_\theta/(Nm)}.\]

The radius is nonnegative and its square equals the moment bound, so the square-root step is legitimate. This term needs joint differentiability and strong convexity, with no gradient-Lipschitz premise.




Powering the marginal density by \(N\) multiplies its negative log potential by \(N\), so its modulus becomes \(Nm\). The actual marginal maximizer minimizes that potential. The dilation concentration theorem used in this module bounds its parameter second moment by \(2d_\theta/(Nm)\), and the exact parameter pushforward identity gives the same moment under the lifted Gibbs law. The final square-root estimate uses the proved finite second moment to interpret projected \(L^2\) error as an ordinary integral. The stronger \(d_\theta/(Nm)\) constant is separately proved in the proximal sharp-concentration extension.




\textbf{Correspondence with the displayed declarations.}


\begin{itemize}

\item \texttt{gibbs\allowbreak{}Target\_\allowbreak{}parameter\_\allowbreak{}second\allowbreak{}Moment\_\allowbreak{}le}: Transfer the proved marginal moment bound to the actual scaled IPLA Gibbs measure using the exact parameter-marginal identity.
\item \texttt{gibbs\allowbreak{}Target\_\allowbreak{}parameter\_\allowbreak{}concentration}: The equilibrium concentration input in IPLA Theorem 1 is a proved theorem of the original joint assumptions. No gradient-Lipschitz or marginal assumption is needed for this statistical term.
\end{itemize}
\section{\texorpdfstring{A noise-\allowbreak{}driven recursion has a law independent of its realization}{A noise-driven recursion has a law independent of its realization}}
\label{proof-module-UnadjustedLangevin-ChainLaw-lean}
{\small\color{gray}Source module: \texttt{Unadjusted\allowbreak{}Langevin/\allowbreak{}Chain\allowbreak{}Law.\allowbreak{}lean}.\par}


For measurable \(F(x,z)\), the recursion \(X_{n+1}=F(X_n,Z_n)\) has a law determined by the initial law and the full noise law when the initialization and noise are independent.

\subsection{\texorpdfstring{noise\allowbreak{}Driven\allowbreak{}Chain}{noiseDrivenChain}}


Define \(X_0=x\), \(X_{n+1}=F(X_n,w_n)\) for any update and noise path. No distribution assumptions are part of the recursion; they enter only when a law is assigned to its inputs.

\subsection{\texorpdfstring{noise\allowbreak{}Driven\allowbreak{}Chain zero}{noiseDrivenChain zero}}


The zero clause of the recursion evaluates to the supplied initial state. This definitional identity is the base case for all law and measurability inductions.

\subsection{\texorpdfstring{noise\allowbreak{}Driven\allowbreak{}Chain succ}{noiseDrivenChain succ}}


The successor clause unfolds to \(X_{n+1}=F(X_n,w_n)\). This exact recursion lets the generic construction later be identified with ULA by induction.

\subsection{\texorpdfstring{noise\allowbreak{}Driven\allowbreak{}Chain measurable}{noiseDrivenChain measurable}}


Induct on \(n\). The base map is the initial-state projection. At the successor, pair the measurable current state with the measurable coordinate evaluation \(w\mapsto w_n\), then compose with the measurable uncurried update. Thus each iterate is a measurable function of initialization and the full noise path.

\subsection{\texorpdfstring{chain\allowbreak{}Law}{chainLaw}}


Push \(\mu\otimes Q\) through the measurable recursion output at index \(n\). The product gives independent initialization and noise path. This canonical law depends only on the initial law, noise-path law, and update, not on a particular sample-space implementation.

\subsection{\texorpdfstring{chain\allowbreak{}Law zero}{chainLaw zero}}


At index zero the output map is first-coordinate projection. Since \(Q\) has mass one, the first marginal of \(\mu\otimes Q\) is exactly \(\mu\). This proves that the canonical law starts from the supplied initial law.

\subsection{\texorpdfstring{noise\allowbreak{}Driven\allowbreak{}Chain law of indep}{noiseDrivenChain law of indep}}


Independence gives the joint input law \(P_{(X,W)}=P_X\otimes P_W\). Push that identity through the measurable recursion output and use pushforward composition. The resulting output law is exactly \texttt{chain\allowbreak{}Law F (P.\allowbreak{}map X) (P.\allowbreak{}map W) n}.

\subsection{\texorpdfstring{noise\allowbreak{}Driven\allowbreak{}Chain law of iid}{noiseDrivenChain law of iid}}


Joint independence of the innovations and their common law \(\eta\) identify the entire noise-path law with \(\bigotimes_n\eta\). Substitute this identity into the independent-initialization law theorem. Countability of the time index supplies path-valued measurability.

\subsection{\texorpdfstring{noise\allowbreak{}Driven\allowbreak{}Chain law unique}{noiseDrivenChain law unique}}


Apply the canonical iid-law formula on both probability spaces. Equality of their initial laws and the shared innovation law makes the right sides identical, so their iterate laws agree. The theorem compares distributions without coupling the two implementations.

\subsection{\texorpdfstring{noise\allowbreak{}Driven\allowbreak{}Chain eq ula\allowbreak{}Chain}{noiseDrivenChain eq ulaChain}}


Induct on the index with \(F=\operatorname{ulaUpdate}\). The zero values agree, and the successor equations apply the same update to equal previous states and the same noise vector. This proves pointwise equality of the generic recursion and the repository's ULA chain.

\subsection{\texorpdfstring{ula\allowbreak{}Update measurable}{ulaUpdate measurable}}


Lipschitzness makes \(\nabla U\) continuous. The joint ULA update is obtained from coordinate projections and that gradient by scalar multiplication and addition, hence is measurable. This discharges the generic recursion's update-measurability premise.

\subsection{\texorpdfstring{ula\allowbreak{}Chain law of iid}{ulaChain law of iid}}


Use measurability of \texttt{ula\allowbreak{}Update} in the generic iid law theorem, then rewrite the generic recursion using \texttt{noise\allowbreak{}Driven\allowbreak{}Chain\_\allowbreak{}eq\_\allowbreak{}ula\allowbreak{}Chain}. The result identifies the actual ULA implementation with its canonical product-noise law.

\subsection{\texorpdfstring{lp\allowbreak{}Ula\allowbreak{}Chain law of iid}{lpUlaChain law of iid}}


The \(L^2\) chain representative agrees almost everywhere with the pointwise ULA recursion of the input representatives. Apply the pointwise law theorem and transfer pushforwards across that almost-everywhere equality. Quotient representative choices therefore do not change the algorithm law.

\section{\texorpdfstring{Realize every initial coupling independently of the noise}{Realize every initial coupling independently of the noise}}
\label{proof-module-UnadjustedLangevin-InitialExtension-lean}
{\small\color{gray}Source module: \texttt{Unadjusted\allowbreak{}Langevin/\allowbreak{}Initial\allowbreak{}Extension.\allowbreak{}lean}.\par}


Given any coupling \(\gamma\) of \(\mu\) and \(\pi\), extend it to \(\gamma\otimes Q\). This preserves its quadratic cost and supplies independent noise without requiring an optimal coupling.

\subsection{\texorpdfstring{Coupling.\allowbreak{}is\allowbreak{}Probability\allowbreak{}Measure}{Coupling.isProbabilityMeasure}}


The coupling's first marginal is the initial probability. Pushforward preserves total mass, so \(c(H\times H)=\mu(H)=1\). This constructs normalization of the joint measure using a marginal identity already stored in \texttt{Coupling}.

\subsection{\texorpdfstring{initial\allowbreak{}Extension}{initialExtension}}


Extend an initial coupling \(c(dx,dy)\) by an independent noise sample \(s\sim Q\), using \(c\otimes Q\) on \((H\times H)\times S\). Both algorithms retain the prescribed initial joint law while sharing fresh noise independent of the pair.

\subsection{\texorpdfstring{initial\allowbreak{}Extension is\allowbreak{}Probability\allowbreak{}Measure}{initialExtension isProbabilityMeasure}}


The coupling is a probability by its first marginal, and \(Q\) is a probability by hypothesis. Their product therefore has total mass one. This instance supports all subsequent \(L^2\) constructions on the extension.

\subsection{\texorpdfstring{initial\allowbreak{}Extension joint}{initialExtension joint}}


Projecting \(c\otimes Q\) to the first factor integrates out the probability noise coordinate. Since \(Q(S)=1\), the resulting measure is exactly \(c\), preserving the entire initial dependence structure.

\subsection{\texorpdfstring{initial\allowbreak{}Extension initial}{initialExtension initial}}


Compose the joint projection with its first state projection. The extension's joint marginal is \(c\), whose first state marginal is \(\mu\). Pushforward composition therefore yields the specified algorithm initialization law.

\subsection{\texorpdfstring{initial\allowbreak{}Extension reference}{initialExtension reference}}


Compose projection to the initial pair with its second state projection. The coupling's second marginal is \(\pi\), so the reference initialization on the extension has exactly law \(\pi\).

\subsection{\texorpdfstring{initial\allowbreak{}Extension noise}{initialExtension noise}}


The coupling measure has mass one. Projecting its product with \(Q\) to the noise factor therefore returns \(Q\). The \texttt{SFinite} hypothesis supplies the product marginal theorem's measure-theoretic condition.

\subsection{\texorpdfstring{initial\allowbreak{}Extension independent}{initialExtension independent}}


The two coordinate projections of a product probability measure are independent. Here those coordinates are the full coupled initial pair and the fresh noise sample. Independence of either individual initial state from any noise-derived random object follows by composition.

\subsection{\texorpdfstring{initial\allowbreak{}Extension initial noise}{initialExtension initial noise}}


Map \(((x,y),s)\) to \((x,s)\). Product pushforward and the first marginal identity of \(c\) give \(\mu\otimes Q\). Thus removing the reference coordinate leaves the same algorithm input law for every initial coupling.

\subsection{\texorpdfstring{initial\allowbreak{}Extension map}{initialExtension map}}


Factor the map through \(((x,y),s)\mapsto(x,s)\), whose law is \(\mu\otimes Q\). Pushforward composition proves that every measurable output depending only on algorithm initialization and noise has the coupling-independent law \(f_*(\mu\otimes Q)\).

\subsection{\texorpdfstring{initial\allowbreak{}Extension chain\allowbreak{}Law}{initialExtension chainLaw}}


Apply \texttt{initial\allowbreak{}Extension\_\allowbreak{}map} to the measurable recursion output. Its canonical pushforward is exactly \texttt{chain\allowbreak{}Law}, so the same output law is realized on every initial-coupling extension. This equality is derived rather than made a field of the coupling.

\subsection{\texorpdfstring{initial\allowbreak{}Extension iid noise\allowbreak{}Law}{initialExtension iid noiseLaw}}


First project the extension to its product-path noise law, then evaluate coordinate \(n\). The marginal of \(\bigotimes_k\eta\) at that coordinate is \(\eta\), giving the requested noise law on the extension.

\subsection{\texorpdfstring{initial\allowbreak{}Extension iid independent}{initialExtension iid independent}}


The iid path has independent coordinate evaluations. Pull this independence back along the extension's measure-preserving noise projection. The resulting innovation sequence remains jointly independent on the coupled initialization space.

\subsection{\texorpdfstring{initial\allowbreak{}Extension initial independent}{initialExtension initial independent}}


The algorithm's initial state is a projection of the initial pair. Compose the pair/\allowbreak{}noise independence theorem with this state projection. This proves independence from the entire noise sample, not just from each individual noise vector.

\subsection{\texorpdfstring{initial\allowbreak{}Extension cost}{initialExtension cost}}


The difference \(x-y\) depends only on the initial pair. Its \(L^2\) extended norm under \(c\otimes Q\) equals the same norm under \(c\), because the noise factor has mass one. The latter is exactly \texttt{coupling\allowbreak{}Cost c.\allowbreak{}joint}.

\subsection{\texorpdfstring{initial\allowbreak{}Extension mem\allowbreak{}Lp difference}{initialExtension memLp difference}}


The state-difference function is measurable. Its extended \(L^2\) norm equals the finite initial coupling cost by the previous theorem. These two facts prove \texttt{Mem\allowbreak{}Lp} for the difference on the extension.

\subsection{\texorpdfstring{initial\allowbreak{}Extension mem\allowbreak{}Lp reference}{initialExtension memLp reference}}


The reference coordinate has law \(\pi\). Pull its supplied \(L^2(\pi)\) identity moment back through this law to get square integrability on the extension. This establishes the reference initialization without choosing a separate realization.

\subsection{\texorpdfstring{initial\allowbreak{}Extension mem\allowbreak{}Lp initial}{initialExtension memLp initial}}


Write \(x=(x-y)+y\). The difference is in \(L^2\) by finite coupling cost and the reference coordinate is in \(L^2\) by its law. Closure of \(L^2\) under addition gives square integrability of the initial algorithm state; no independent moment premise on \(\mu\) is introduced.

\subsection{\texorpdfstring{initial\allowbreak{}Extension lp\allowbreak{}Initial}{initialExtension lpInitial}}


Use \texttt{initial\allowbreak{}Extension\_\allowbreak{}mem\allowbreak{}Lp\_\allowbreak{}initial} to convert the first state coordinate into an actual \(L^2\) class. The class is represented almost everywhere by \(((x,y),s)\mapsto x\), preserving its concrete relation to the coupling coordinates.

\subsection{\texorpdfstring{initial\allowbreak{}Extension lp\allowbreak{}Reference}{initialExtension lpReference}}


Convert the square-integrable second state coordinate into an \(L^2\) class. Its representative equals \(((x,y),s)\mapsto y\) almost everywhere, with law \(\pi\), and it shares the same product probability space as the algorithm initialization.

\subsection{\texorpdfstring{initial\allowbreak{}Extension lp\allowbreak{}Initial law}{initialExtension lpInitial law}}


The representative chosen by \texttt{to\allowbreak{}Lp} agrees almost everywhere with the first state coordinate. Pushforwards are unchanged by such an equality. Combine this with \texttt{initial\allowbreak{}Extension\_\allowbreak{}initial} to obtain the initial law \(\mu\).

\subsection{\texorpdfstring{initial\allowbreak{}Extension lp\allowbreak{}Reference law}{initialExtension lpReference law}}


Replace the selected \(L^2\) representative by the second state projection almost everywhere. Its already proved marginal is \(\pi\), hence the class used as reference initialization has exactly that law.

\subsection{\texorpdfstring{initial\allowbreak{}Extension lp cost}{initialExtension lp cost}}


The representative of the class difference equals \(x-y\) almost everywhere. The norm of that \(L^2\) difference is the real value of its extended norm; \texttt{initial\allowbreak{}Extension\_\allowbreak{}cost} identifies it with \((\operatorname{couplingCost}c)_{\mathbb R}\). The finite-cost premise justifies this real conversion.

\section{\texorpdfstring{A reference for every initial coupling with one output law}{A reference for every initial coupling with one output law}}
\label{proof-module-UnadjustedLangevin-ProductReference-lean}
{\small\color{gray}Source module: \texttt{Unadjusted\allowbreak{}Langevin/\allowbreak{}Product\allowbreak{}Reference.\allowbreak{}lean}.\par}


This construction turns the analytic reference estimate into an exact initial Wasserstein term by realizing each finite initial coupling on a product with one fixed noise source.

\subsection{\texorpdfstring{pointwise\allowbreak{}Update}{pointwiseUpdate}}


Define the ordinary Euler update \(F_h(x,z)=x-h\nabla U(x)+\sqrt{2h}z\) in an arbitrary Hilbert-space presentation. It is the pointwise function whose lifted composition gives the \(L^2\) numerical chain.

\subsection{\texorpdfstring{pointwise\allowbreak{}Update measurable}{pointwiseUpdate measurable}}


The Lipschitz gradient is continuous. Compose it with the first coordinate projection and combine the remaining continuous scalar and vector operations. The uncurried map \((x,z)\mapsto F_h(x,z)\) is therefore measurable.

\subsection{\texorpdfstring{product\allowbreak{}Noise mem\allowbreak{}Lp}{productNoise memLp}}


The noise projection from \(c\otimes Q\) to \(Q\) is measure preserving because the coupling is a probability. Pull back the supplied \(L^2(Q)\) noise function along this projection. Its square integrability is unchanged.

\subsection{\texorpdfstring{product\allowbreak{}Noise}{productNoise}}


Use the preceding square-integrability proof to form the class of \(((x,y),s)\mapsto W_n(s)\) on the concrete extension. Both the numerical chain and reference use this same \(L^2\) innovation.

\subsection{\texorpdfstring{product\allowbreak{}Noise coe}{productNoise coe}}


The representative selected by \texttt{to\allowbreak{}Lp} agrees almost everywhere with the pulled-back original noise \(W_n(s)\). This identity connects quotient-space noise arithmetic to the actual source random functions.

\subsection{\texorpdfstring{product\allowbreak{}Output\allowbreak{}Law}{productOutputLaw}}


Push \(\mu\otimes Q\) through the pointwise noise-driven chain starting at \(x\) with source path \(W(s)\). This definition omits the reference coupling entirely, making explicit the common algorithm output law for all coupling realizations.

\subsection{\texorpdfstring{product chain coe}{product chain coe}}


Induct through the lifted chain and use the representative identities for initial state, product noise, lifted drift, addition, and scalar multiplication. At each fixed index the result agrees almost everywhere with the ordinary recursion driven by \(x\) and \(W(s)\).

\subsection{\texorpdfstring{product chain output\allowbreak{}Law}{product chain outputLaw}}


Replace the \(L^2\) chain by its almost-everywhere pointwise recursion. That recursion depends only on the algorithm's initial coordinate and noise sample, so \texttt{initial\allowbreak{}Extension\_\allowbreak{}map} gives its coupling-independent law. The result is exactly \texttt{product\allowbreak{}Output\allowbreak{}Law}.

\subsection{\texorpdfstring{product\allowbreak{}Probability}{productProbability}}


A coupling of the initial probability with the reference measure is itself a probability, and its product with the fixed probability noise source is a probability. This local instance supplies normalization on every extension separately.

\subsection{\texorpdfstring{Product\allowbreak{}Reference}{ProductReference}}


For every finite-cost initial coupling, the structure supplies a concrete \texttt{L2\allowbreak{}Reference} on that coupling's product extension, with its initial state matching the second coordinate. Algorithm output laws and initial costs are proved outside the structure; no convergence estimate is a field.

\subsection{\texorpdfstring{Product\allowbreak{}Reference.\allowbreak{}output finite}{ProductReference.output finite}}


Finite initial Wasserstein distance supplies at least one finite-cost initial coupling. On its concrete extension, the numerical and reference variables belong to \(L^2\), giving a finite-cost output coupling. Their established laws identify its marginals with the common output law and \(\pi\), hence the output Wasserstein distance is finite.

\subsection{\texorpdfstring{dalalyan small step of product reference}{dalalyan small step of product reference}}


For each finite-cost initial coupling \(c\), apply the analytic \(L^2\) Dalalyan theorem on \(c\otimes Q\). The initial class-distance theorem identifies its initial error exactly with the real coupling cost, and the common output-law theorem identifies its numerical law.



The resulting bound holds for every finite \(c\), with coefficient \((1-mh)^n\ge0\). Infimizing their costs gives the genuine \(W_2(\mu,\pi)\), without selecting an optimal coupling. Output finiteness is proved first so the final real-valued inequality cannot conceal infinity.

\subsection{\texorpdfstring{dalalyan large step of product reference}{dalalyan large step of product reference}}


Repeat the couplingwise analytic argument for the stable large-step branch, whose contraction factor is \(Lh-1\in[0,1)\). The geometric error sum has denominator \(2-Lh>0\). The finite-coupling infimum theorem replaces the initial class distance by the actual initial Wasserstein distance and yields the displayed large-step bias term.

\section{\texorpdfstring{Gaussian and Brownian noise moments}{Gaussian and Brownian noise moments}}
\label{proof-module-UnadjustedLangevin-GaussianMoments-lean}
{\small\color{gray}Source module: \texttt{Unadjusted\allowbreak{}Langevin/\allowbreak{}Gaussian\allowbreak{}Moments.\allowbreak{}lean}.\par}


Fix a natural-number dimension \(d\in\mathbb N\), write \(E_d=\mathbb R^d\) with its Euclidean norm, and let \(P\) be a measure on \(\Omega\). A law statement \(X_i\sim\mathcal N(0,v)\) below means that the coordinate is almost everywhere measurable and has the stated pushforward measure. We prove the required \(L^2\) norms from these laws. Coordinate independence is unnecessary for second moments because \(\|X\|^2=\sum_iX_i^2\) has no cross terms.



For Brownian processes the input is mathlib’s actual scalar pre-Brownian law predicate, whose marginal and increment laws supply the Gaussian hypotheses. These facts establish the noise moments; existence and stationarity of a coupled Langevin reference require separate inputs.



The \(L^2\) norms and expectations below concern the actual measure \(P\), with \texttt{Mem\allowbreak{}Lp} giving square integrability before passing to a quotient class. Laws are pushforward identities; independence is used only for cross-term cancellation, while marginal Gaussian second moments supply the dimension factor.

\subsection{\texorpdfstring{Second moment of a centered Gaussian}{Second moment of a centered Gaussian}}


For nonnegative variance \(v\), the centered Gaussian law has second moment

\[\int_{\mathbb R}x^2\,d\mathcal N(0,v)(x)=v.\]

The known Gaussian variance is \(v\). Expand variance as \(\int(x-\mathbb E[x])^2\), and insert the Gaussian mean \(\mathbb E[x]=0\).



Simplify \(x-0=x\). The variance integral is now exactly the required second moment.

\subsection{\texorpdfstring{Gaussian coordinates give a square-\allowbreak{}integrable vector}{Gaussian coordinates give a square-integrable vector}}


Let \(X:\Omega\to E_d\), and suppose every coordinate satisfies \(X_i\sim\mathcal N(0,v)\) under \(P\), with \(v\ge0\). Then \(X\in L^2(P;E_d)\).



In finite-dimensional Euclidean space, square integrability is equivalent to square integrability of each coordinate. Each coordinate belongs to \(L^2\) because its Gaussian law has a finite second moment; transfer that membership along the equality of laws.

\subsection{\texorpdfstring{Sum the coordinate moments}{Sum the coordinate moments}}


Under the same coordinate Gaussian laws,

\[\int_\Omega\|X(\omega)\|^2\,dP(\omega)=dv.\]

For each coordinate \(i\), retain the proved square integrability \(X_i\in L^2(P)\). It will justify integrating the finite sum of squared coordinates.



Change variables using the law of \(X_i\), then apply the scalar Gaussian moment:

\[\int_\Omega X_i^2\,dP=\int_{\mathbb R}x^2\,d\mathcal N(0,v)=v.\]

Expand the squared Euclidean norm and interchange the finite sum and integral:

\[\int_\Omega\|X\|^2\,dP=\sum_{i=0}^{d-1}\int_\Omega X_i^2\,dP.\]

Every summand is integrable by its \(L^2\) membership.



Each of the \(d\) summands equals \(v\); hence their sum is \(dv\).

\subsection{\texorpdfstring{The Gaussian vector as an L\ensuremath{{}^{2}} class}{The Gaussian vector as an L² class}}


Let \([X]\) be the \(L^2\) equivalence class constructed from \(X\), whose coordinates have laws \(\mathcal N(0,v)\). Then

\[\|[X]\|_{L^2}^2=dv.\]

Express the squared \(L^2\) norm as the self-inner-product, then use the definition of the \(L^2\) inner product as an integral.



The chosen representative of \([X]\) equals \(X\) almost everywhere. Therefore the integrand \(\langle[X](\omega),[X](\omega)\rangle\) agrees almost everywhere with \(\|X(\omega)\|^2\), so the two integrals agree.



Substitute the integral identity and the vector second-moment calculation \(\int\|X\|^2\,dP=dv\).

\subsection{\texorpdfstring{The square-\allowbreak{}root Gaussian norm}{The square-root Gaussian norm}}


For the same Gaussian coordinate laws,

\[\|[X]\|_{L^2}=\sqrt{dv}.\]

This includes \(d=0\) and \(v=0\).



Replace \(dv\) by \(\|[X]\|^2\). Since norms are nonnegative, \(\sqrt{\|[X]\|^2}=\|[X]\|\).

\subsection{\texorpdfstring{Include the Langevin diffusion factor}{Include the Langevin diffusion factor}}


If every coordinate of \(X\) has law \(\mathcal N(0,v)\), then the Langevin-scaled vector satisfies

\[\|\sqrt2[X]\|_{L^2}=\sqrt{2d}\sqrt v.\]

Norm homogeneity gives \(\|\sqrt2[X]\|=\sqrt2\|[X]\|\), because \(\sqrt2\ge0\). Insert \(\|[X]\|=\sqrt{dv}\), then split the square roots of the nonnegative products.



Rearrange \(\sqrt2\sqrt d\sqrt v\) into \(\sqrt{2d}\sqrt v\).

\subsection{\texorpdfstring{Assemble Brownian coordinates}{Assemble Brownian coordinates}}


For scalar processes \(W_i:[0,\infty)\times\Omega\to\mathbb R\), define the Euclidean vector

\[W_t(\omega):=(W_0(t,\omega),\ldots,W_{d-1}(t,\omega))\in E_d.\]

Collect the coordinates at the fixed time \(t\) and outcome \(\omega\), using the finite-dimensional \(\ell^2\) norm.

\subsection{\texorpdfstring{Brownian marginals are Gaussian}{Brownian marginals are Gaussian}}


Assume each scalar \(W_i\) satisfies \(\mathrm{IsPreBrownianReal}\) under \(P\). For each \(t\ge0\) and coordinate \(i\),

\[W_t^{(i)}\sim\mathcal N(0,t).\]

Apply the scalar Brownian marginal-law theorem at time \(t\). The \(i\)-th coordinate of the assembled vector is definitionally the original scalar process.

\subsection{\texorpdfstring{Brownian vectors belong to L\ensuremath{{}^{2}}}{Brownian vectors belong to L²}}


Under the scalar pre-Brownian assumptions, for every \(t\ge0\),

\[W_t\in L^2(P;E_d).\]

Apply Gaussian vector square integrability with variance \(v=t\), using the just-proved marginal laws of all coordinates.

\subsection{\texorpdfstring{The exact Brownian noise norm}{The exact Brownian noise norm}}


For scalar pre-Brownian coordinates and \(t\ge0\),

\[\|\sqrt2[W_t]\|_{L^2}=\sqrt{2d}\sqrt t.\]

Specialize the scaled Gaussian vector norm to the Brownian vector, whose coordinate variance is \(t\). This derives the square-root noise scale from the Brownian laws.

\subsection{\texorpdfstring{Define an increment from any starting time}{Define an increment from any starting time}}


For \(a,t\ge0\), define the vector increment over the interval \([a,a+t]\) by

\[\Delta_aW_t(\omega):=W_{a+t}(\omega)-W_a(\omega).\]

Subtract the two Euclidean vectors at the interval endpoints. Each coordinate is the corresponding scalar Brownian increment.

\subsection{\texorpdfstring{Increment variance is elapsed time}{Increment variance is elapsed time}}


Under the scalar pre-Brownian assumptions, for every \(a,t\ge0\) and every coordinate \(i\),

\[W_i(a+t)-W_i(a)\sim\mathcal N(0,t).\]

Expand the vector increment to its scalar coordinate. Apply the Brownian difference-law theorem at times \(a+t\) and \(a\); its variance is the nonnegative distance between those times.



It remains to identify that variance. As real numbers,

\[|(a+t)-a|=|t|=t,\]

because \(t\ge0\). Thus the increment law is exactly \(\mathcal N(0,t)\).

\subsection{\texorpdfstring{Increments are square integrable}{Increments are square integrable}}


For pre-Brownian coordinates and any \(a,t\ge0\),

\[\Delta_aW_t\in L^2(P;E_d).\]

Each increment coordinate has the centered Gaussian law of variance \(t\), so the Gaussian vector integrability theorem applies.

\subsection{\texorpdfstring{The same noise norm on every reference segment}{The same noise norm on every reference segment}}


For scalar pre-Brownian coordinates and every \(a,t\ge0\),

\[\boxed{\|\sqrt2[W_{a+t}-W_a]\|_{L^2}=\sqrt{2d}\sqrt t.}\]

Apply the scaled Gaussian vector norm to the increment coordinate laws. The answer depends only on elapsed time \(t\), so every reference interval has the same moment bound.

\section{\texorpdfstring{Coordinate Gaussian moments and L\ensuremath{{}^{2}} Brownian continuity}{Coordinate Gaussian moments and L² Brownian continuity}}
\label{proof-module-UnadjustedLangevin-GaussianCoordinates-lean}
{\small\color{gray}Source module: \texttt{Unadjusted\allowbreak{}Langevin/\allowbreak{}Gaussian\allowbreak{}Coordinates.\allowbreak{}lean}.\par}


Coordinatewise Gaussian laws suffice for the second moment \(\mathbb E\|X\|^2=dv\). This module carefully obtains Brownian \(L^2\) continuity and increment norms without assuming independence between spatial coordinates.

\subsection{\texorpdfstring{gaussian coordinates mem\allowbreak{}Lp}{gaussian coordinates memLp}}


Each centered Gaussian coordinate has a finite second moment. In finite Euclidean dimension, \(\|X\|^2=\sum_iX_i^2\), so the finite sum is integrable. Coordinatewise measurability also assembles to vector measurability, giving \(X\in L^2\). Coordinate independence is unnecessary for this moment statement.

\subsection{\texorpdfstring{gaussian coordinates second moment}{gaussian coordinates second moment}}


Use \(\|X\|^2=\sum_iX_i^2\) and finite-sum integral linearity. The coordinate law \(X_i\sim\mathcal N(0,v)\) gives \(\mathbb EX_i^2=v\), hence \(\mathbb E\|X\|^2=|\iota|v\). The dimension factor is exactly the cardinality of the finite coordinate set.

\subsection{\texorpdfstring{gaussian coordinates Lp norm sq}{gaussian coordinates Lp norm sq}}


For the given \(L^2\) class, its squared norm is the integral of the squared norm of its representative. Apply the preceding Gaussian coordinate moment formula to obtain \(\|X\|_2^2=|\iota|v\).

\subsection{\texorpdfstring{gaussian coordinates Lp norm}{gaussian coordinates Lp norm}}


Take nonnegative square roots in \(\|X\|_2^2=|\iota|v\). Norm nonnegativity selects the positive root, giving \(\|X\|_2=\sqrt{|\iota|v}\), including zero variance and zero dimension.

\subsection{\texorpdfstring{gaussian coordinates to\allowbreak{}Lp coordinate\allowbreak{}Law}{gaussian coordinates toLp coordinateLaw}}


The representative of \texttt{Mem\allowbreak{}Lp.\allowbreak{}to\allowbreak{}Lp X} agrees almost everywhere with \(X\). Coordinate evaluation preserves that equality, so each coordinate law transfers to the canonical \(L^2\) representative by almost-everywhere invariance of pushforwards.

\subsection{\texorpdfstring{gaussian coordinates to\allowbreak{}Lp norm}{gaussian coordinates toLp norm}}


Apply the exact \(L^2\) Gaussian norm theorem to \texttt{to\allowbreak{}Lp X}, using its transferred coordinate laws. This gives \(\|[X]\|_2=\sqrt{|\iota|v}\) for the class constructed from the original random vector.

\subsection{\texorpdfstring{gaussian coordinates scaled l2 norm}{gaussian coordinates scaled l2 norm}}


The scalar-norm identity gives \(\|\sqrt2[X]\|_2=\sqrt2\sqrt{|\iota|v}\). Nonnegativity permits the square-root product rules, producing \(\sqrt{2|\iota|}\sqrt v\). This is the precise Langevin noise constant used in the local error estimate.

\subsection{\texorpdfstring{brownian\allowbreak{}Coordinates}{brownianCoordinates}}


Assemble the scalar coordinate processes into \(B_t(\omega)=(W_i(t,\omega))_{i\in\iota}\). Finite Euclidean coordinates determine the vector; no new random source or independence is introduced by this definition.

\subsection{\texorpdfstring{brownian\allowbreak{}Coordinates coordinate\allowbreak{}Law}{brownianCoordinates coordinateLaw}}


At a fixed time \(t\ge0\), each scalar pre-Brownian coordinate has law \(\mathcal N(0,t)\). Unfolding the assembled vector shows that its \(i\)-th coordinate is exactly that scalar variable, giving the required coordinate law.

\subsection{\texorpdfstring{brownian\allowbreak{}Coordinates\allowbreak{}Increment}{brownianCoordinatesIncrement}}


Define the vector increment \(\Delta_aB_t=B_{a+t}-B_a\) coordinate by coordinate. This records both the starting time and elapsed nonnegative time, matching the local noise input on an arbitrary Langevin segment.

\subsection{\texorpdfstring{brownian\allowbreak{}Coordinates\allowbreak{}Increment coordinate\allowbreak{}Law}{brownianCoordinatesIncrement coordinateLaw}}


The scalar Brownian increment law depends only on elapsed time: \(W_i(a+t)-W_i(a)\sim\mathcal N(0,t)\). Simplify \((a+t)-a=t\) in the variance and unfold the vector increment to obtain its coordinate law.

\subsection{\texorpdfstring{brownian\allowbreak{}Coordinates\allowbreak{}Increment mem\allowbreak{}Lp}{brownianCoordinatesIncrement memLp}}


Apply \texttt{gaussian\_\allowbreak{}coordinates\_\allowbreak{}mem\allowbreak{}Lp} to the increment's coordinate laws of variance \(t\). The assembled Brownian vector increment is square integrable even without independence between different coordinate processes.

\subsection{\texorpdfstring{brownian\allowbreak{}Coordinates\allowbreak{}Increment scaled l2 norm}{brownianCoordinatesIncrement scaled l2 norm}}


Apply the exact scaled Gaussian-coordinate norm theorem with variance \(t\). It yields \(\|\sqrt2[\Delta_aB_t]\|_2=\sqrt{2|\iota|}\sqrt t\), establishing the noise radius from Brownian laws alone.

\subsection{\texorpdfstring{brownian\allowbreak{}Coordinates\allowbreak{}Lp}{brownianCoordinatesLp}}


Convert the assembled Brownian vector at time \(t\) to its \(L^2\) class using its Gaussian-coordinate square-integrability proof. This places the time-indexed process in the Hilbert space where the deterministic integral equation is solved.

\subsection{\texorpdfstring{brownian\allowbreak{}Coordinates\allowbreak{}Lp coe}{brownianCoordinatesLp coe}}


The \texttt{to\allowbreak{}Lp} representative theorem gives \([B_t](\omega)=B_t(\omega)\) almost everywhere for each fixed time. This is the bridge from class-valued continuity statements to coordinate laws of the original Brownian process.

\subsection{\texorpdfstring{brownian\allowbreak{}Coordinates\allowbreak{}Lp sub coordinate\allowbreak{}Law}{brownianCoordinatesLp sub coordinateLaw}}


The difference of two selected representatives agrees almost everywhere with \(B_t-B_s\). Brownian increment laws, including the sign reversal when \(t<s\), give a centered Gaussian with variance \(|t-s|\); centered Gaussian symmetry handles the reversed order.

\subsection{\texorpdfstring{brownian\allowbreak{}Coordinates\allowbreak{}Lp norm sub}{brownianCoordinatesLp norm sub}}


Apply the Gaussian \(L^2\) norm formula to the difference class, whose coordinates have variance \(|t-s|\). Thus \(\|[B_t]-[B_s]\|_2=\sqrt{|\iota|\,|t-s|}\). This exact increment norm is stronger than the continuity statement derived from it.

\subsection{\texorpdfstring{brownian\allowbreak{}Coordinates\allowbreak{}Lp continuous}{brownianCoordinatesLp continuous}}


As \(s\to t\), \(|t-s|\to0\), hence the exact norm formula gives \(\|[B_t]-[B_s]\|_2\to0\). Lean uses continuity of the square root and the metric characterization of continuity. This needs finite-dimensional Brownian laws, with no sample-path continuity premise.

\subsection{\texorpdfstring{brownian\allowbreak{}Coordinates\allowbreak{}Lp increment scaled norm}{brownianCoordinatesLp increment scaled norm}}


Set the two times in the difference-norm identity to \(a+t\) and \(a\). Their distance is \(t\), and scalar multiplication by \(\sqrt2\) gives \(\sqrt{2|\iota|}\sqrt t\). This is the class-valued version of the earlier increment estimate.

\subsection{\texorpdfstring{brownian\allowbreak{}Coordinates\allowbreak{}Lp zero}{brownianCoordinatesLp zero}}


At time zero every coordinate has the degenerate centered Gaussian law of variance zero. Equivalently its \(L^2\) norm is zero by the moment formula. Norm-zero equivalence proves that the whole class is zero, irrespective of the chosen pointwise representative.

\subsection{\texorpdfstring{brownian\allowbreak{}Coordinates\allowbreak{}Real\allowbreak{}Lp}{brownianCoordinatesRealLp}}


Extend the nonnegative-time \(L^2\) process to real times by \(t\mapsto[B_{\max(t,0)}]\), implemented with \texttt{Real.\allowbreak{}to\allowbreak{}NNReal}. It is constant zero for negative times and agrees with the original class at nonnegative times.

\subsection{\texorpdfstring{brownian\allowbreak{}Coordinates\allowbreak{}Real\allowbreak{}Lp continuous}{brownianCoordinatesRealLp continuous}}


The clipping map \(t\mapsto\max(t,0)\) into nonnegative reals is continuous. Compose it with the established \(L^2\)-continuous Brownian map to obtain continuity of the real-time extension, including at zero.

\subsection{\texorpdfstring{brownian\allowbreak{}Coordinates\allowbreak{}Real\allowbreak{}Lp zero}{brownianCoordinatesRealLp zero}}


At real time zero the clipping map returns nonnegative zero. The previously proved zero-time \(L^2\) identity therefore simplifies the real-time extension to zero.

\subsection{\texorpdfstring{brownian\allowbreak{}Coordinates\allowbreak{}Real\allowbreak{}Lp increment scaled norm}{brownianCoordinatesRealLp increment scaled norm}}


For \(a,t\ge0\), clipping leaves \(a\) and \(a+t\) unchanged. Rewrite the real-time extension at both endpoints as the nonnegative-time Brownian classes, then apply their scaled increment formula. This supplies \(\sqrt{2d}\sqrt t\) in the real interval-integral formulation.

\section{\texorpdfstring{Existence of the additive Langevin integral equation}{Existence of the additive Langevin integral equation}}
\label{proof-module-UnadjustedLangevin-LangevinIntegral-lean}
{\small\color{gray}Source module: \texttt{Unadjusted\allowbreak{}Langevin/\allowbreak{}Langevin\allowbreak{}Integral.\allowbreak{}lean}.\par}


In a complete real normed space, a globally Lipschitz drift and a continuous additive path \(B\), \(B(0)=0\), determine a continuous solution of \(Y_t-x=-\int_0^tg(Y_s)ds+B_t\) on every finite nonnegative horizon. The construction is deterministic and later applies with the Banach space \(L^2(P)\).

\subsection{\texorpdfstring{Solve on a short interval}{Solve on a short interval}}


Remove the additive input by writing \(Y(t)=Z(t)+B(t)\). The new unknown satisfies \(Z(t)=x+\int_0^t f(s,Z(s))ds\), with \(f(s,z)=-g(z+B(s))\). This time-dependent drift is \(L\)-Lipschitz in \(z\) and continuous in time even though \(B\) need not be differentiable.



Bound \(\|B(t)\|\le C'\) on the compact interval and set \(G=\|g(x)\|+LC'\), radius \(a=2hG\), and drift bound \(M=G+La\). The condition \(Lh\le1/2\) proves \(Mh\le a\), giving the ball-invariance condition in \texttt{Is\allowbreak{}Picard\allowbreak{}Lindelof}. The proof checks these bounds and invokes \texttt{ODE.\allowbreak{}Fun\allowbreak{}Space.\allowbreak{}exists\_\allowbreak{}is\allowbreak{}Fixed\allowbreak{}Pt\_\allowbreak{}next}.



The chosen fixed point gives a continuous \(Z\) with its integral identity. Add back \(B\), use \(B(0)=0\), and rewrite the negative drift integral to obtain \(Y(t)-x=-\int_0^tg(Y(s))ds+B(t)\). The solution is represented as a real-time function, with the required identity asserted on \([0,h]\).

\subsection{\texorpdfstring{splice langevin integral}{splice langevin integral}}


Use \(Y(t)\) up to time \(a\), and \(Z(t-a)\) afterward. The equality \(Z(0)=Y(a)\) makes the pasted path continuous at the join. Split \(\int_0^t=\int_0^a+\int_a^t\), translate the second integral, and add the two supplied equations; the terms \(B(a)\) cancel, giving the same equation on \([0,a+b]\).

\subsection{\texorpdfstring{exists langevin integral}{exists langevin integral}}


Choose a positive segment length small enough that \(L\delta\le1/2\). Construct the first Picard segment, then repeatedly solve with initial state at the previous endpoint and shifted input \(B(a+s)-B(a)\). The splicing lemma combines finitely many segments covering the prescribed horizon. The local contraction restriction determines the construction mesh and imposes no restriction on the final horizon.

\subsection{\texorpdfstring{exists lp langevin integral small}{exists lp langevin integral small}}


Apply the Banach-space Picard theorem to \(H=L^2(P;K)\), \(g(X)=\operatorname{lpDrift}(\nabla U)(X)\). The lifted drift has the same Lipschitz constant as \(\nabla U\). Completeness of \(L^2\) and continuity of the supplied \(L^2\)-noise path are the exact fixed-point hypotheses.

\subsection{\texorpdfstring{exists lp langevin integral}{exists lp langevin integral}}


Use the finite-horizon Banach-space existence theorem with the lifted gradient drift. The conclusion is continuity in the \(L^2\) norm and an equality of \(L^2\) classes at every time in the horizon. The Bochner time integral therefore lives in the same Hilbert space; no pathwise differentiability of the random noise is asserted.

\section{\texorpdfstring{Construct the L\ensuremath{{}^{2}} Langevin equation from Brownian-\allowbreak{}coordinate laws}{Construct the L² Langevin equation from Brownian-coordinate laws}}
\label{proof-module-UnadjustedLangevin-BrownianLangevin-lean}
{\small\color{gray}Source module: \texttt{Unadjusted\allowbreak{}Langevin/\allowbreak{}Brownian\allowbreak{}Langevin.\allowbreak{}lean}.\par}


Coordinate pre-Brownian laws produce the continuous \(L^2\) additive noise needed by the deterministic integral-equation existence theorem. The exact noise norm is \(\sqrt{2d}\sqrt t\).

\subsection{\texorpdfstring{brownian\allowbreak{}Langevin\allowbreak{}Noise}{brownianLangevinNoise}}


Define \(N_a(t)=\sqrt2([B_{a+t}]-[B_a])\) in \(L^2\), using the real-time Brownian extension. The subtraction fixes the input at zero at local time zero, and the \(\sqrt2\) factor gives the overdamped Langevin diffusion normalization.

\subsection{\texorpdfstring{brownian\allowbreak{}Langevin\allowbreak{}Noise continuous}{brownianLangevinNoise continuous}}


The real-time Brownian map is continuous in \(L^2\) from its Gaussian increment moments. Composition with \(t\mapsto a+t\), subtraction of the fixed class \([B_a]\), and scalar multiplication preserve continuity. Thus the additive input satisfies the deterministic Picard theorem's continuity requirement.

\subsection{\texorpdfstring{brownian\allowbreak{}Langevin\allowbreak{}Noise zero}{brownianLangevinNoise zero}}


At local time zero both Brownian terms are \([B_a]\), so they cancel and their scalar multiple is zero. Lean unfolds the definition and uses additive identities; no exceptional-set statement is needed because this is equality in \(L^2\).

\subsection{\texorpdfstring{brownian\allowbreak{}Langevin\allowbreak{}Noise norm}{brownianLangevinNoise norm}}


For \(a,t\ge0\), apply the exact real-time Brownian increment moment formula. It gives \(\|N_a(t)\|_2=\sqrt{2d}\sqrt t\), with \(d=|\iota|\). The Gaussian laws prove this value; it is not an additional noise-bound hypothesis.

\subsection{\texorpdfstring{exists lp langevin integral of brownian small}{exists lp langevin integral of brownian small}}


Supply the continuous, zero-starting \(N_a\) to \texttt{exists\_\allowbreak{}lp\_\allowbreak{}langevin\_\allowbreak{}integral\_\allowbreak{}small}. It constructs \(Y_t-Y_0=-\int_0^t\operatorname{lpDrift}(\nabla U)(Y_s)ds+N_a(t)\) for \(Lh\le1/2\). Attach the exact noise norm theorem to the existential package. Neither the integral equation nor its noise estimate is assumed.

\subsection{\texorpdfstring{exists lp langevin integral of brownian}{exists lp langevin integral of brownian}}


Use the unrestricted finite-horizon \(L^2\) existence theorem with the same constructed Brownian input. The resulting solution is continuous in \(L^2\), starts at the prescribed random variable, and satisfies the actual gradient integral equation. The exact \(\sqrt{2d}\sqrt t\) noise formula is included as a proved property.

\section{\texorpdfstring{Transfer the Gibbs score estimate to an L\ensuremath{{}^{2}} reference}{Transfer the Gibbs score estimate to an L² reference}}
\label{proof-module-UnadjustedLangevin-StationaryMoments-lean}
{\small\color{gray}Source module: \texttt{Unadjusted\allowbreak{}Langevin/\allowbreak{}Stationary\allowbreak{}Moments.\allowbreak{}lean}.\par}


A reference variable \(Y\) whose law is \(\pi\) inherits \(\|\nabla U(Y)\|_{L^2}\le\sqrt{Ld}\). The assertion that its law is Gibbs is a displayed premise here.

\subsection{\texorpdfstring{lp\allowbreak{}Drift norm le of gibbs law}{lpDrift norm le of gibbs law}}


The representative of \texttt{lp\allowbreak{}Drift} equals \(\nabla U(Y)\) almost everywhere. Change variables using \(P\circ Y^{-1}=\pi_U\), so its squared \(L^2(P)\) norm is \(\int\|\nabla U\|^2d\pi_U\le Ld\). Nonnegativity of the norm and of \(Ld\) then gives \(\|\nabla U(Y)\|_2\le\sqrt{Ld}\). The law identity, rather than independence, is the input to this transfer.

\subsection{\texorpdfstring{lp\allowbreak{}Drift norm le of strong\allowbreak{}Convex gibbs law}{lpDrift norm le of strongConvex gibbs law}}


Use strong convexity to prove \(Z_U<\infty\), and apply the previous law-transfer theorem with that proof. The conclusion is the same \(\sqrt{Ld}\) bound for the lifted drift of any \(L^2\) variable with Gibbs law. There is no separate assumed drift-moment bound hidden in this specialization.

\section{\texorpdfstring{Assemble continuous Brownian reference segments}{Assemble continuous Brownian reference segments}}
\label{proof-module-UnadjustedLangevin-BrownianReference-lean}
{\small\color{gray}Source module: \texttt{Unadjusted\allowbreak{}Langevin/\allowbreak{}Brownian\allowbreak{}Reference.\allowbreak{}lean}.\par}


Construct consecutive \(L^2\) Langevin segments from actual Brownian-coordinate noise and package them as an analytic reference. At this intermediate stage Gibbs stationarity remains an explicit premise, later discharged in CanonicalULA.

\subsection{\texorpdfstring{brownian\allowbreak{}Segment\allowbreak{}From}{brownianSegmentFrom}}


Choose the integral solution supplied by \texttt{exists\_\allowbreak{}lp\_\allowbreak{}langevin\_\allowbreak{}integral\_\allowbreak{}of\_\allowbreak{}brownian} for the segment starting at absolute time \(nh\), with the supplied initial \(L^2\) variable. Classical choice selects data from a proved existence theorem; its defining equation and continuity are recovered by the specification theorem.

\subsection{\texorpdfstring{brownian\allowbreak{}Segment\allowbreak{}From spec}{brownianSegmentFrom spec}}


Apply the choice specification to obtain continuity, the time-zero initial value, the gradient integral equation on \([0,h]\), and the exact noise norm. These four properties describe the selected segment independently of its noncomputable implementation.

\subsection{\texorpdfstring{brownian\allowbreak{}Knots}{brownianKnots}}


Set the zeroth knot to the prescribed initial variable. Define the next knot to be the endpoint at local time \(h\) of the segment started at the current knot. This recursion ensures consecutive segments meet exactly as \(L^2\) classes.

\subsection{\texorpdfstring{brownian\allowbreak{}Segments}{brownianSegments}}


At index \(n\), select \texttt{brownian\allowbreak{}Segment\allowbreak{}From} with the recursively constructed \(n\)-th knot as initial state. The absolute noise start is \(nh\), so all segments are driven by increments of the same Brownian process.

\subsection{\texorpdfstring{brownian\allowbreak{}Segments start}{brownianSegments start}}


The chosen-segment specification states its value at zero equals its supplied initial state. Here that state is the \(n\)-th knot, so \(Y^{(n)}_0=K_n\) follows directly by specialization.

\subsection{\texorpdfstring{brownian\allowbreak{}Segments initial}{brownianSegments initial}}


Apply the segment-start identity at \(n=0\), then unfold the zeroth knot. The actual constructed reference therefore begins at exactly the prescribed random variable, preserving its chosen joint realization with the numerical initialization.

\subsection{\texorpdfstring{brownian\allowbreak{}Euler\allowbreak{}Noise}{brownianEulerNoise}}


Normalize the actual segment endpoint noise by \(1/\sqrt{2h}\). This defines the \(L^2\) innovation used by the Euler chain, so it shares the exact Brownian increment driving the reference segment.

\subsection{\texorpdfstring{brownian\allowbreak{}Euler\allowbreak{}Noise endpoint}{brownianEulerNoise endpoint}}


For \(h>0\), \(\sqrt{2h}\ne0\). Multiply the normalized noise class by \(\sqrt{2h}\) and cancel the inverse scalar. The resulting equality identifies the Euler noise contribution with the same Brownian endpoint input in the integral equation.

\subsection{\texorpdfstring{brownian\allowbreak{}Reference}{brownianReference}}


Populate \texttt{L2\allowbreak{}Reference} with the constructed segments. Their specification proves continuity and the integral equation; recursive knots prove endpoint compatibility; normalized noise proves endpoint matching; Gibbs-score moment transfer proves the drift bound. The supplied all-time Gibbs-law premise remains visible and is later discharged by the stationarity construction.

\subsection{\texorpdfstring{brownian\allowbreak{}Reference of strong\allowbreak{}Convex}{brownianReference of strongConvex}}


Use strong convexity to derive finiteness of the Gibbs partition, then call \texttt{brownian\allowbreak{}Reference}. All analytic, normalization, moment, and noise fields are proved internally. The all-time law premise is still explicit in this constructor and must be provided by the stationary-law theorem.

\section{\texorpdfstring{Strong convexity forces m \ensuremath{\le} L}{Strong convexity forces m ≤ L}}
\label{proof-module-UnadjustedLangevin-PotentialAssumptions-lean}
{\small\color{gray}Source module: \texttt{Unadjusted\allowbreak{}Langevin/\allowbreak{}Potential\allowbreak{}Assumptions.\allowbreak{}lean}.\par}


On a nontrivial Hilbert space the numerical compatibility \(m\le L\) follows from strong convexity and a Lipschitz gradient of the same differentiable potential.

\subsection{\texorpdfstring{Compare monotonicity with Cauchy–Schwarz}{Compare monotonicity with Cauchy–Schwarz}}


Choose \(x\ne0\), possible because the Hilbert space is nontrivial. Strong monotonicity gives \(m\|x\|^2\le\langle\nabla U(x)-\nabla U(0),x\rangle\), while Cauchy–Schwarz and Lipschitzness bound the same quantity by \(L\|x\|^2\). Cancel the strictly positive \(\|x\|^2\) to prove \(m\le L\). No positivity premise on \(m\) is needed here.

\section{\texorpdfstring{Normalized increments are independent standard Gaussians}{Normalized increments are independent standard Gaussians}}
\label{proof-module-UnadjustedLangevin-BrownianInnovations-lean}
{\small\color{gray}Source module: \texttt{Unadjusted\allowbreak{}Langevin/\allowbreak{}Brownian\allowbreak{}Innovations.\allowbreak{}lean}.\par}


For a pre-Brownian coordinate and positive grid spacing \(h\), \(Z_n=(B_{(n+1)h}-B_{nh})/\sqrt h\) has standard normal law. Across independent coordinate paths, the entire coordinate-time family is independent.

\subsection{\texorpdfstring{normalized\allowbreak{}Brownian\allowbreak{}Increment}{normalizedBrownianIncrement}}


Define \(\xi_n=(B_{(n+1)h}-B_{nh})/\sqrt h\). The denominator is interpreted as real scalar division; positive-step law results below explicitly require \(h>0\), so that the variance rescaling is legitimate.

\subsection{\texorpdfstring{normalized\allowbreak{}Brownian\allowbreak{}Increment has\allowbreak{}Law}{normalizedBrownianIncrement hasLaw}}


A Brownian increment over an interval of length \(h\) is \(\mathcal N(0,h)\). Scaling by \(1/\sqrt h\) multiplies its variance by \(1/h\), producing \(\mathcal N(0,1)\). The proof rewrites \((n+1)h-nh=h\) and uses \((\sqrt h)^2=h\), justified by \(h>0\).

\subsection{\texorpdfstring{normalized\allowbreak{}Brownian\allowbreak{}Increment mem\allowbreak{}Lp}{normalizedBrownianIncrement memLp}}


The exact standard Gaussian law provides every finite moment, in particular the second. Pull this square integrability back through \texttt{normalized\allowbreak{}Brownian\allowbreak{}Increment\_\allowbreak{}has\allowbreak{}Law} to prove \(\xi_n\in L^2(P)\).

\subsection{\texorpdfstring{normalized\allowbreak{}Brownian\allowbreak{}Increment second moment}{normalizedBrownianIncrement second moment}}


Transfer the integral of \(x^2\) along the standard Gaussian law of \(\xi_n\). Its mean is zero and variance one, hence \(\mathbb E\xi_n^2=1\). This is an exact second-moment identity for the normalized increment.

\subsection{\texorpdfstring{normalized\allowbreak{}Brownian\allowbreak{}Increment independent}{normalizedBrownianIncrement independent}}


Brownian increments on the consecutive grid intervals are independent. Normalizing each by the same deterministic scalar is a coordinatewise measurable transformation and preserves independence. The formulation permits \(h=0\), where every increment is the constant zero and the independence assertion remains valid.

\subsection{\texorpdfstring{i\allowbreak{}Indep\allowbreak{}Fun uncurry nat of aemeasurable}{iIndepFun uncurry nat of aemeasurable}}


Assume the row-valued random maps \((X_{i,n})_n\) are independent and each row is internally independent. Replace the almost-everywhere measurable coordinates by measurable versions; countability of \(\mathbb N\) lets all replacements in a row hold on one full-measure set. Apply independence of independent families, then transfer the result back across the coordinatewise almost-everywhere equalities.

\subsection{\texorpdfstring{normalized\allowbreak{}Brownian\allowbreak{}Increments independent}{normalizedBrownianIncrements independent}}


Each grid increment depends measurably on two entries of its coordinate grid path. Independence of those paths gives independence across coordinate rows, while Brownian independent increments gives independence within each row. Apply \texttt{i\allowbreak{}Indep\allowbreak{}Fun\_\allowbreak{}uncurry\_\allowbreak{}nat\_\allowbreak{}of\_\allowbreak{}aemeasurable} to obtain joint independence of all \((i,n)\) innovations.

\subsection{\texorpdfstring{normalized\allowbreak{}Brownian\allowbreak{}Increments independent of paths}{normalizedBrownianIncrements independent of paths}}


Restrict each independent full coordinate path to the grid \(nh\), preserving independence under composition. The preceding theorem then yields joint coordinate-time independence of the normalized increments. This avoids requiring a separate measurable map into an uncountable path space.

\subsection{\texorpdfstring{The entire innovation law}{The entire innovation law}}


Every coordinate-time innovation has law \(\mathcal N(0,1)\), and the family is jointly independent. The countable product-law theorem therefore identifies the whole innovation map with \(\bigotimes_{(i,n)}\mathcal N(0,1)\). Countability supplies measurability of the path-valued map from coordinatewise almost-everywhere measurability.

\section{\texorpdfstring{Pull Brownian noise onto every coupled initialization}{Pull Brownian noise onto every coupled initialization}}
\label{proof-module-UnadjustedLangevin-ProductBrownian-lean}
{\small\color{gray}Source module: \texttt{Unadjusted\allowbreak{}Langevin/\allowbreak{}Product\allowbreak{}Brownian.\allowbreak{}lean}.\par}


Brownian-coordinate laws, innovation independence and their L\ensuremath{{}^{2}} realizations survive passage to the product \(\gamma\otimes Q\). This builds all reference fields except the explicitly displayed stationarity premise, discharged by the later canonical assembly.

\subsection{\texorpdfstring{pre\allowbreak{}Brownian comp measure\allowbreak{}Preserving}{preBrownian comp measurePreserving}}


A measure-preserving map leaves the law of every measurable random variable unchanged after pullback. Apply this to every finite-dimensional Brownian vector. Their Gaussian laws and covariance data are therefore preserved, giving a pre-Brownian process on the new sample space.

\subsection{\texorpdfstring{brownian\allowbreak{}Coordinates\allowbreak{}Lp pullback}{brownianCoordinatesLp pullback}}


Both sides represent the vector \((B_i(t,f(\omega)))_i\) almost everywhere. Pull back the source representative equality using measure preservation and combine it with the target \texttt{to\allowbreak{}Lp} equality. \texttt{Lp.\allowbreak{}ext} identifies the two quotient classes.

\subsection{\texorpdfstring{brownian\allowbreak{}Euler\allowbreak{}Noise pullback}{brownianEulerNoise pullback}}


The \(L^2\) pullback is linear, hence commutes with subtraction and scalar multiplication. Apply the preceding Brownian-vector pullback identity to the two endpoints in the normalized increment. This identifies the pulled-back Euler noise with the noise constructed after pullback.

\subsection{\texorpdfstring{source\allowbreak{}Euler\allowbreak{}Noise}{sourceEulerNoise}}


Take the actual selected representative of the previously constructed \(L^2\) Brownian Euler noise on the source space. This is the random function supplied to the generic product algorithm before choosing any initial coupling.

\subsection{\texorpdfstring{source\allowbreak{}Euler\allowbreak{}Noise mem\allowbreak{}Lp}{sourceEulerNoise memLp}}


Every representative of a given \(L^2\) class is square integrable by \texttt{Lp.\allowbreak{}mem\allowbreak{}Lp}. Apply that fact to the Brownian Euler noise class. This avoids assuming an extra noise-moment condition in the product construction.

\subsection{\texorpdfstring{source\allowbreak{}Euler\allowbreak{}Noise coe}{sourceEulerNoise coe}}


Unfold the nested class-valued noise definition and apply representative equalities for scalar multiplication, subtraction, and the Brownian vector classes. Almost everywhere the result is \((\sqrt{2h})^{-1}\sqrt2(B_{nh+h}-B_{nh})\), with real-time clipping as in the original definition.

\subsection{\texorpdfstring{source\allowbreak{}Euler\allowbreak{}Noise coordinate coe}{sourceEulerNoise coordinate coe}}


Evaluate the vector representative formula in coordinate \(i\). For \(h>0\), use \(\sqrt{2h}=\sqrt2\sqrt h\) and cancel the nonzero \(\sqrt2\), obtaining \((B_i((n+1)h)-B_i(nh))/\sqrt h\).

\subsection{\texorpdfstring{source\allowbreak{}Euler\allowbreak{}Noise coordinate coe nnreal}{sourceEulerNoise coordinate coe nnreal}}


For nonnegative mesh times, \texttt{Real.\allowbreak{}to\allowbreak{}NNReal} leaves their real values unchanged. Rewrite both clipped grid endpoints into the corresponding nonnegative-real products. The coordinate formula becomes the ordinary Brownian grid increment divided by \(\sqrt h\).

\subsection{\texorpdfstring{source\allowbreak{}Euler\allowbreak{}Noise normalized coordinate}{sourceEulerNoise normalized coordinate}}


The nonnegative-grid formula is exactly the definition of \texttt{normalized\allowbreak{}Brownian\allowbreak{}Increment}. After converting the positive real step to its nonnegative-real form, the two random functions agree almost everywhere. This gives a direct bridge to the established innovation-law theorems.

\subsection{\texorpdfstring{source\allowbreak{}Euler\allowbreak{}Noise coordinate\allowbreak{}Law}{sourceEulerNoise coordinateLaw}}


A normalized Brownian increment has law \(\mathcal N(0,1)\). Transfer this law across \texttt{source\allowbreak{}Euler\allowbreak{}Noise\_\allowbreak{}normalized\_\allowbreak{}coordinate} using almost-everywhere invariance of laws. Hence the actual source noise representative has the required Gaussian coordinate law.

\subsection{\texorpdfstring{source\allowbreak{}Euler\allowbreak{}Noise independent}{sourceEulerNoise independent}}


The Brownian grid-path theorem gives independence of all normalized coordinates indexed by \((i,n)\). Reindex by the injective coordinate swap to \((n,i)\), then transfer across the almost-everywhere coordinate identities. This matches the index order expected by the vector innovation theorem.

\subsection{\texorpdfstring{product\allowbreak{}Pre\allowbreak{}Brownian}{productPreBrownian}}


Projection from an initial-coupling product extension to the noise source is measure preserving. Pull each scalar Brownian coordinate back along it using \texttt{pre\allowbreak{}Brownian\_\allowbreak{}comp\_\allowbreak{}measure\allowbreak{}Preserving}. The actual coordinate processes on the extension retain their Brownian finite-dimensional laws.

\subsection{\texorpdfstring{product\allowbreak{}Noise eq brownian\allowbreak{}Euler\allowbreak{}Noise}{productNoise eq brownianEulerNoise}}


The lifted source noise and the extension's constructed Brownian noise are both the measure-preserving pullback of the source \(L^2\) Brownian Euler noise. Use representative extensionality to identify the former with this pullback, then apply \texttt{brownian\allowbreak{}Euler\allowbreak{}Noise\_\allowbreak{}pullback} to identify the latter.

\subsection{\texorpdfstring{product\allowbreak{}Euler\allowbreak{}Noise coordinate\allowbreak{}Law}{productEulerNoise coordinateLaw}}


Pull the source standard-Gaussian coordinate law through the extension's noise projection. The lifted product-noise representative agrees almost everywhere with the raw pulled-back source noise. Transfer the law across that equality to obtain the coordinate law for the exact classes used by the numerical chain.

\subsection{\texorpdfstring{product\allowbreak{}Euler\allowbreak{}Noise independent}{productEulerNoise independent}}


Identify the source coordinate-time path law as the product of standard Gaussians. Pull that law to the extension; its coordinate family is independent by the product-law independence criterion. Finally replace raw source pullbacks by the almost-everywhere equal product-noise representatives.

\subsection{\texorpdfstring{product\allowbreak{}Euler\allowbreak{}Noise independent\allowbreak{}Initial}{productEulerNoise independentInitial}}


The raw first initial coordinate depends only on the coupling factor, while the whole source noise sequence depends only on the noise factor. Product independence therefore applies. Countability of the sequence lets representative equalities hold simultaneously at all indices; transport independence to the selected \(L^2\) initialization and noise representatives.

\subsection{\texorpdfstring{product\allowbreak{}Euler\allowbreak{}Noise independent\allowbreak{}Reference}{productEulerNoise independentReference}}


Use the same product independence argument with the second initial coordinate. Both initialization and whole noise-path representative equalities are transferred almost everywhere, proving the actual reference class is independent of all fresh Euler noise.

\subsection{\texorpdfstring{brownian\allowbreak{}Product\allowbreak{}Probability}{brownianProductProbability}}


The coupling is a probability by its marginal identity and the Brownian source measure is a probability. Their product extension is normalized. This local instance supplies the probability-space hypotheses of the Brownian pullback and \(L^2\) reference constructors.

\subsection{\texorpdfstring{product\allowbreak{}Gibbs\allowbreak{}Mem\allowbreak{}Lp}{productGibbsMemLp}}


Name the already proved finite-second-moment theorem for the strongly convex Gibbs law. The abbreviation supplies exactly the \(L^2(\pi_U)\) identity moment needed to construct the reference initialization on every coupling extension.

\subsection{\texorpdfstring{brownian\allowbreak{}Product\allowbreak{}L2\allowbreak{}Reference}{brownianProductL2Reference}}


Construct the Brownian reference on the extension using the supplied stationarity of its selected segments. Replace its Brownian-defined Euler noise by the common lifted source noise using \texttt{product\allowbreak{}Noise\_\allowbreak{}eq\_\allowbreak{}brownian\allowbreak{}Euler\allowbreak{}Noise}. All remaining continuity, integral, score, and noise fields come from the established constructors.

\subsection{\texorpdfstring{product\allowbreak{}Reference of constructed brownian}{productReference of constructed brownian}}


For each finite initial coupling, use \texttt{brownian\allowbreak{}Product\allowbreak{}L2\allowbreak{}Reference}. Its initial reference equals the second-coordinate \(L^2\) class by the segment-initial identity. These constructions fill the \texttt{Product\allowbreak{}Reference} fields. Only the explicitly quantified segment-stationarity premise remains an input at this intermediate stage.

\subsection{\texorpdfstring{dalalyan of constructed product brownian}{dalalyan of constructed product brownian}}


Derive \(m\le L\) from strong convexity and smoothness, construct the product reference from the stated stationary segment laws, and apply the small-step product-reference theorem. This intermediate endpoint has the genuine initial \(W_2\) term; the later independent Brownian construction discharges its stationarity premise.

\section{\texorpdfstring{Compatibility exports for the independent product realization}{Compatibility exports for the independent product realization}}
\label{proof-module-IPLA-ProductReference-lean}
{\small\color{gray}Source module: \texttt{IPLA/\allowbreak{}Product\allowbreak{}Reference.\allowbreak{}lean}.\par}


This file re-exports the shared product realization and its law-identification theorem. A product of an initial coupling and a noise source realizes every needed coupling while fixing a common algorithm output law.

\subsection{\texorpdfstring{Separate initialization from the noise source}{Separate initialization from the noise source}}


On \((H\times H)\times S\) with law \(c\otimes Q\), let the two initial variables be the first and second coordinates of \(c\), and let every innovation depend only on \(S\). Product independence makes both initial variables independent of the noise source. The first chain evolves by \(x\mapsto x-h\nabla V(x)+\sqrt{2h}\xi_k\), and its output law is the pushforward of \(\mu_0\otimes Q\) by the measurable recursion; hence it depends only on \(\mu_0\), not on the chosen coupling \(c\). This verifies the common-output obligation needed to take the coupling infimum. Finite coupling cost and Gibbs moments supply the \(L^2\) initial realization; the exported output-finiteness statement guards subsequent real transport bounds.



The \texttt{export} commands expose the already proved declarations under the IPLA namespace. Their definitions and proofs are unchanged, so the mathematical argument is the one in the imported source modules; the dependency diagram below opens those full proofs. This file itself makes no additional analytic assumption and contains no replacement proof.

\section{\texorpdfstring{Canonical Gaussian path law and implementation independence}{Canonical Gaussian path law and implementation independence}}
\label{proof-module-IPLA-AlgorithmLaw-lean}
{\small\color{gray}Source module: \texttt{IPLA/\allowbreak{}Algorithm\allowbreak{}Law.\allowbreak{}lean}.\par}


The scalar Gaussian coordinate hypotheses uniquely determine the entire countable noise path. Independence of initialization then determines every output law, even across different underlying probability spaces.

\subsection{\texorpdfstring{Construct the independent Gaussian noise path}{Construct the independent Gaussian noise path}}


A noise coordinate is a pair \((k,a)\) consisting of iteration and state coordinate. Put the countable product law \(\bigotimes_{(k,a)}\mathcal N(0,1)\) on scalar paths \(w\). Packing sets \((\operatorname{pack}w)_k(a)=w(k,a)\); its measurability follows because every output coordinate is a measurable product projection, followed by the measurable Euclidean coordinate identification. The canonical vector path law is the pushforward of that scalar product law by packing.



The two probability instances are proved as well: each scalar Gaussian factor has mass one, hence the infinite product is a probability measure; pushing it through measurable packing preserves total mass. These are \texttt{canonical\allowbreak{}Flat\allowbreak{}Noise\allowbreak{}Law\_\allowbreak{}is\allowbreak{}Probability\allowbreak{}Measure} and \texttt{canonical\allowbreak{}Noise\allowbreak{}Path\allowbreak{}Law\_\allowbreak{}is\allowbreak{}Probability\allowbreak{}Measure}, whose source appears in the complete-module panel.




\textbf{Correspondence with the displayed declarations.}


\begin{itemize}

\item \texttt{Noise\allowbreak{}Coordinate}: A coordinate records the iteration and the parameter or particle block.
\item \texttt{canonical\allowbreak{}Flat\allowbreak{}Noise\allowbreak{}Law}: The canonical scalar noise source: independent standard real Gaussians.
\item \texttt{pack\allowbreak{}Noise\allowbreak{}Path}: Group scalar coordinates into the Euclidean noise vector at each time.
\item \texttt{pack\allowbreak{}Noise\allowbreak{}Path\_\allowbreak{}measurable}: Packing the countable coordinate source is measurable.
\item \texttt{canonical\allowbreak{}Noise\allowbreak{}Path\allowbreak{}Law}: The canonical vector-valued Gaussian noise-path law.
\end{itemize}
\subsection{\texorpdfstring{The algorithm hypotheses identify the entire path law}{The algorithm hypotheses identify the entire path law}}


Mutual independence of the scalar family implies that its joint map has law equal to the infinite product of its individual coordinate laws. Since each law is \(\mathcal N(0,1)\),

\[\mathcal L((\xi_{k,a})_{k,a})=\bigotimes_{k,a}\mathcal N(0,1).\]

Composing this map with measurable packing and using the composition rule for pushforwards gives \(\mathcal L((\xi_k)_k)=\operatorname{canonicalNoisePathLaw}\). The argument proves independence of vector blocks from scalar independence and never assumes a grouping principle without justification.




The canonical flat source is the product of one scalar Gaussian for each time/\allowbreak{}coordinate pair. The packing map merely groups these coordinates into a vector at each time; coordinate evaluation proves its measurability. The supplied innovation structure has the same finite-coordinate Gaussian product laws by mutual independence, so the product-law uniqueness theorem identifies its whole flat path law. Applying the measurable packing map gives equality of the full vector path law, including all cross-time relations.




\textbf{Correspondence with the displayed declarations.}


\begin{itemize}

\item \texttt{Gaussian\allowbreak{}Innovations.\allowbreak{}flat\allowbreak{}Noise\allowbreak{}Path\_\allowbreak{}law}: The Gaussian coordinate hypotheses identify the complete flattened path law.
\item \texttt{Gaussian\allowbreak{}Innovations.\allowbreak{}noise\allowbreak{}Path\_\allowbreak{}law}: The vector-valued path law is derived from the scalar Gaussian hypotheses, including every coordinate and every iteration.
\end{itemize}
\subsection{\texorpdfstring{Identify the common update and recursion}{Identify the common update and recursion}}


The generic pointwise ULA update for \(V_N\) and step \(\gamma/N\) is \(z-(\gamma/N)\nabla V_N(z)+\sqrt{2(\gamma/N)}\xi\). Reassociating multiplication inside the square root makes it the scaled IPLA update. The generic noise-driven chain and \texttt{scaled\allowbreak{}Chain} therefore agree by induction: their initial values match and each successor uses the same map on the same preceding state and noise vector.




\textbf{Correspondence with the displayed declarations.}


\begin{itemize}

\item \texttt{pointwise\allowbreak{}Update\_\allowbreak{}eq\_\allowbreak{}scaled\allowbreak{}Update}: The generic ULA update for \texttt{U\_\allowbreak{}N} is exactly the scaled IPLA update.
\item \texttt{noise\allowbreak{}Driven\allowbreak{}Chain\_\allowbreak{}eq\_\allowbreak{}scaled\allowbreak{}Chain}: Exact identification of the generic noise recursion with the IPLA recursion.
\end{itemize}
\subsection{\texorpdfstring{A common initial law gives common output laws}{A common initial law gives common output laws}}


Let \(F_k(z_0,\xi)\) be the measurable map that iterates the update \(k\) times. Initialization–path independence gives the joint input law \(\mu_0\otimes\mathcal G\), where \(\mathcal G\) is the canonical Gaussian path law. Hence

\[\rho_k=(F_k)_\#(\mu_0\otimes\mathcal G).\]

The pointwise/\allowbreak{}L\ensuremath{{}^{2}} identity identifies this measure with \texttt{laws}. The right-hand side depends only on \(U,\gamma,\mu_0\), so two Gaussian implementations with the same initial law give identical output laws. This identity needs no stationary diffusion or reference coupling.




The exact update identity matches the generic noise-driven recursion to the scaled IPLA recursion. Independence then makes the joint initialization/\allowbreak{}path law the product of the initial pushforward and the canonical noise path law. Push this common source through the measurable time-\(k\) recursion. Two implementations on different probability spaces therefore have the same output law whenever their initial laws agree; the conclusion is about the actual entire algorithm law at the specified time.




\textbf{Correspondence with the displayed declarations.}


\begin{itemize}

\item \texttt{laws\_\allowbreak{}eq\_\allowbreak{}canonical\allowbreak{}Chain\allowbreak{}Law}: The law in the IPLA convergence statement is the canonical law of the actual Gaussian algorithm. No reference process is involved in this result.
\item \texttt{laws\_\allowbreak{}eq\_\allowbreak{}of\_\allowbreak{}initial\allowbreak{}Law\_\allowbreak{}eq}: Two Gaussian implementations of IPLA with the same initial law have the same output law, even when they use different underlying probability spaces.
\end{itemize}
\section{\texorpdfstring{Euler approximations converge to the integral solution}{Euler approximations converge to the integral solution}}
\label{proof-module-UnadjustedLangevin-EulerLimit-lean}
{\small\color{gray}Source module: \texttt{Unadjusted\allowbreak{}Langevin/\allowbreak{}Euler\allowbreak{}Limit.\allowbreak{}lean}.\par}


This module proves a deterministic convergence theorem for additive Euler discretization. With mesh \(h=T/n\), its terminal value converges to the continuous integral solution, including when the state space itself is \(L^2(P)\).

\subsection{\texorpdfstring{additive\allowbreak{}Euler\allowbreak{}Chain}{additiveEulerChain}}


Define \(X_0=x\) and \(X_{k+1}=X_k-hg(X_k)+B((k+1)h)-B(kh)\). The additive input is one fixed function evaluated at mesh endpoints; consequently the numerical chain and the integral solution can be compared with exactly matching noise increments.

\subsection{\texorpdfstring{additive\allowbreak{}Euler\allowbreak{}Chain zero}{additiveEulerChain zero}}


Evaluating the recursive definition at index zero gives \(X_0=x\). Lean proves this by definitional equality; no continuity of the input or positivity of the mesh is required for this initialization identity.

\subsection{\texorpdfstring{additive\allowbreak{}Euler\allowbreak{}Chain succ}{additiveEulerChain succ}}


Unfold the successor clause once to expose \(X_{k+1}=X_k-hg(X_k)+\Delta B_k\). This definitional identity is the equation used when subtracting the exact solution's increment and deriving the error recurrence.

\subsection{\texorpdfstring{langevin integral increment}{langevin integral increment}}


Subtract the integral equations at \(b\) and \(a\). Continuity of \(g\circ Y\) makes it interval integrable, so interval-integral additivity gives \(\int_0^b-\int_0^a=\int_a^b\). Thus \(Y(b)-Y(a)=-\int_a^b g(Y(s))ds+B(b)-B(a)\), including oriented intervals when \(b<a\).

\subsection{\texorpdfstring{finite euler recurrence bound}{finite euler recurrence bound}}


From \(e_{k+1}\le qe_k+b\), \(e_0=0\), \(q\ge1\), and \(b\ge0\), induction proves \(e_n\le nbq^n\). In the successor step bound \(q\cdot kbq^k+b\) by \((k+1)bq^{k+1}\), using \(q^{k+1}\ge1\). This coarse bound avoids dividing by \(q-1\), so it also covers \(q=1\).

\subsection{\texorpdfstring{explicit\allowbreak{}Euler norm sub le}{explicitEuler norm sub le}}


Write the difference as \((x-y)-h(g(x)-g(y))\). The triangle inequality and Lipschitzness give \(\|\cdot\|\le\|x-y\|+hL\|x-y\|=(1+Lh)\|x-y\|\). The hypothesis \(h\ge0\) removes the absolute value in the scalar norm.

\subsection{\texorpdfstring{additive\allowbreak{}Euler\allowbreak{}Chain error le of drift modulus}{additiveEulerChain error le of drift modulus}}


At one step the common noise cancels. The Euler map contributes \((1+Lh)e_k\), and the exact integral remainder contributes at most \(h\varepsilon\), because the drift changes by at most \(\varepsilon\) within one mesh interval. Apply the finite recurrence bound and \((1+Lh)^n\le e^{Lnh}\), then use \(nh=T\) to obtain \(e_n\le T\varepsilon e^{LT}\).

\subsection{\texorpdfstring{additive\allowbreak{}Euler\allowbreak{}Chain tendsto}{additiveEulerChain tendsto}}


The continuous drift path \(g\circ Y\) is uniformly continuous on \([0,T]\). Given any positive error tolerance, choose a drift modulus \(\varepsilon\) making \(T\varepsilon e^{LT}\) small, then take the mesh \(T/(n+1)\) below its continuity radius. The previous error bound gives endpoint convergence. The case \(T=0\) is separated so no division by the horizon is needed.

\subsection{\texorpdfstring{lp langevin integral euler tendsto}{lp langevin integral euler tendsto}}


Apply the deterministic matching-input Euler convergence theorem in \(L^2(P;H)\), using the Lipschitz lifted gradient. The numerical endpoint converges in the \(L^2\) norm to the constructed solution at \(T\). This is the topology later used to pass smooth bounded tests and identify laws.

\section{\texorpdfstring{Identify recursive sampling laws with iterated Markov transitions}{Identify recursive sampling laws with iterated Markov transitions}}
\label{proof-module-UnadjustedLangevin-MarkovLaw-lean}
{\small\color{gray}Source module: \texttt{Unadjusted\allowbreak{}Langevin/\allowbreak{}Markov\allowbreak{}Law.\allowbreak{}lean}.\par}


The path-space law of an iid-noise recursion agrees with the repeated one-step transition \(\mu\mapsto F_\#(\mu\otimes\eta)\). The proof explicitly establishes independence of the current state and the fresh innovation.

\subsection{\texorpdfstring{finite\allowbreak{}Noise\allowbreak{}Chain}{finiteNoiseChain}}


Evaluate the same recursion using only a finite prefix \((w_0,\ldots,w_{n-1})\). At the successor, recurse on the first \(n\) coordinates and use the final coordinate as the next innovation. This makes finite dependence explicit for the fresh-noise independence proof.

\subsection{\texorpdfstring{finite\allowbreak{}Noise\allowbreak{}Chain eq}{finiteNoiseChain eq}}


Induct on prefix length. Restricting an infinite path to its first \(n\) entries gives the same previous state, and the final finite coordinate is \(w_n\). The successor updates thus agree, identifying finite-prefix and infinite-path implementations pointwise.

\subsection{\texorpdfstring{finite\allowbreak{}Noise\allowbreak{}Chain measurable}{finiteNoiseChain measurable}}


The zero map is the initial projection. At the successor, restriction to the earlier finite coordinates and evaluation of the last coordinate are measurable; combine them with the measurable update and the induction hypothesis. This proves the state depends measurably on a finite prefix.

\subsection{\texorpdfstring{iid prefix independent fresh}{iid prefix independent fresh}}


The index sets \(\{0,\ldots,n-1\}\) and \(\{n\}\) are disjoint. Independence of the coordinate family on the iid product source therefore gives independence of the whole prefix vector and the next coordinate. This is stronger than pairwise independence of individual past entries.

\subsection{\texorpdfstring{markov\allowbreak{}Step}{markovStep}}


Define \(M_\eta(\mu)=F_*(\mu\otimes\eta)\). Its product measure means the innovation is fresh and independent of the current state. The subsequent theorem derives that independence for the canonical recursion, rather than assuming it at each step.

\subsection{\texorpdfstring{markov\allowbreak{}Laws}{markovLaws}}


Set \(\mu_0=\mu\), \(\mu_{n+1}=M_\eta(\mu_n)\). These are iterated independent-noise transition laws. The recursion alone does not yet assert that a previously defined sample-path algorithm realizes them.

\subsection{\texorpdfstring{chain\allowbreak{}Law joint fresh}{chainLaw joint fresh}}


The current state is a measurable function of initialization and the first \(n\) noises. Product-source independence and prefix/\allowbreak{}fresh independence make it independent of \(w_n\). Its marginal is \texttt{chain\allowbreak{}Law}, and the innovation marginal is \(\eta\); their joint law is therefore the product of those actual marginals.

\subsection{\texorpdfstring{chain\allowbreak{}Law succ iid}{chainLaw succ iid}}


The pointwise successor is \(F(X_n,w_n)\). Push the proved product joint law of \((X_n,w_n)\) through \(F\); this yields \(\mu_{n+1}=M_\eta(\mu_n)\) for the canonical recursion laws. No later-time state/\allowbreak{}noise independence is assumed.

\subsection{\texorpdfstring{chain\allowbreak{}Law eq markov\allowbreak{}Laws}{chainLaw eq markovLaws}}


The canonical path law and the iterated Markov law have the same initial value and the same successor equation. Induction on \(n\) identifies them at every index. This closes the connection between sample-path and transition-law definitions.

\subsection{\texorpdfstring{markov\allowbreak{}Step is\allowbreak{}Probability}{markovStep isProbability}}


The independent product of two probability measures has mass one. A measurable pushforward preserves that total mass, so \texttt{markov\allowbreak{}Step} is again a probability measure.

\subsection{\texorpdfstring{markov\allowbreak{}Laws is\allowbreak{}Probability}{markovLaws isProbability}}


Induct from the initial probability and apply \texttt{markov\allowbreak{}Step\_\allowbreak{}is\allowbreak{}Probability} at each successor. This proves normalization of every transition iterate and provides the probability instances used in later law comparisons.

\subsection{\texorpdfstring{noise\allowbreak{}Driven\allowbreak{}Chain law eq markov\allowbreak{}Laws}{noiseDrivenChain law eq markovLaws}}


First identify the independent iid realization with the canonical product-path output law. Then apply \texttt{chain\allowbreak{}Law\_\allowbreak{}eq\_\allowbreak{}markov\allowbreak{}Laws}. Every realization satisfying the original initialization/\allowbreak{}noise assumptions consequently has the same proved Markov iterate law.

\section{\texorpdfstring{A finite O(h\ensuremath{{}^{3/2}}) rejected-\allowbreak{}displacement bound}{A finite O(h³ᐟ²) rejected-displacement bound}}
\label{proof-module-UnadjustedLangevin-MALAError-lean}
{\small\color{gray}Source module: \texttt{Unadjusted\allowbreak{}Langevin/\allowbreak{}MALAError.\allowbreak{}lean}.\par}


The Metropolis-adjusted Langevin proposal provides an exactly invariant approximation to the diffusion. This module proves that its mean rejected displacement under Gibbs initialization is at most \(Ch\sqrt h\), with a finite constant derived from the potential assumptions.

\subsection{\texorpdfstring{mala\allowbreak{}Proposal}{malaProposal}}


Define \(T_h(x,\xi)=x-h\nabla U(x)+\sqrt{2h}\xi\). This is the same Euler proposal used by ULA; MALA applies an accept/\allowbreak{}reject rule to this proposed state.

\subsection{\texorpdfstring{mala\allowbreak{}Log\allowbreak{}Ratio}{malaLogRatio}}


The log acceptance ratio is \(R_h(x,y)=U(x)-U(y)+(\|y-x+h\nabla U(x)\|^2-\|x-y+h\nabla U(y)\|^2)/(4h)\). The common Gaussian normalizers cancel in the reverse/\allowbreak{}forward proposal density ratio.

\subsection{\texorpdfstring{mala\allowbreak{}Log\allowbreak{}Ratio expansion}{malaLogRatio expansion}}


Expand both squared norms using the real inner-product norm identity. The \(\|y-x\|^2\) terms cancel, and division by \(4h\) gives \(R_h=U(x)-U(y)+\frac12\langle\nabla U(x)+\nabla U(y),y-x\rangle-\frac h4(\|\nabla U(y)\|^2-\|\nabla U(x)\|^2)\). The premise \(h\ne0\) justifies denominator cancellation.

\subsection{\texorpdfstring{gradient norm sq difference bound}{gradient norm sq difference bound}}


Factor \(|\|g_y\|^2-\|g_x\|^2|=|\|g_y\|-\|g_x\||(\|g_y\|+\|g_x\|)\). The reverse triangle inequality bounds the first factor by \(\|g_y-g_x\|\), then Lipschitzness bounds it by \(L\|y-x\|\). The other factor is nonnegative, so multiplication preserves the estimate.

\subsection{\texorpdfstring{potential trapezoid bound}{potential trapezoid bound}}


Use the smooth upper and lower first-order models for the potential increment at the two endpoints. Compare the increment with the average endpoint directional derivative, and use the Lipschitz bound on the derivative difference. The resulting trapezoid error is at most \(L\|y-x\|^2\); no second derivative of \(U\) is assumed.

\subsection{\texorpdfstring{mala\allowbreak{}Log\allowbreak{}Ratio abs bound}{malaLogRatio abs bound}}


Insert the exact log-ratio expansion and use the triangle inequality. Bound its potential trapezoid term by \(L\|y-x\|^2\) and its squared-gradient difference by \(L\|y-x\|(\|g_x\|+\|g_y\|)\). Positivity of \(h\) gives the displayed coefficient \(hL/4\) without an absolute value on the step.

\subsection{\texorpdfstring{metropolis rejection le abs}{metropolis rejection le abs}}


If \(r\ge0\), \(e^r\ge1\), so rejection is zero. If \(r<0\), \(1-e^r\le-r=|r|\) by \(e^r\ge1+r\). Both cases also give \(0\le1-\min(1,e^r)\), establishing a nonnegative rejection probability and its log-ratio bound.

\subsection{\texorpdfstring{rejected displacement le}{rejected displacement le}}


Multiply \(1-\min(1,e^r)\le|r|\) by the nonnegative displacement norm \(\|y-x\|\). This yields the rejection-weighted coupling error bound with no condition on the states.

\subsection{\texorpdfstring{mala\allowbreak{}Amplitude}{malaAmplitude}}


Use the step-independent amplitude \(A(x,\xi)=\|\nabla U(x)\|+\sqrt2\|\xi\|\). For \(0\le h\le1\), it bounds the proposal displacement after extracting \(\sqrt h\); its cubic integrability will give a uniform rejection-error constant.

\subsection{\texorpdfstring{mala\allowbreak{}Rejection\allowbreak{}Displacement}{malaRejectionDisplacement}}


Define \(D_h(x,\xi)=\|T_h(x,\xi)-x\|(1-\min(1,e^{R_h(x,T_h(x,\xi))}))\). This is the acceptance-variable-averaged distance lost on rejection, so no additional uniform random variable appears in its integral.

\subsection{\texorpdfstring{mala\allowbreak{}Amplitude nonneg}{malaAmplitude nonneg}}


Both score and noise norms are nonnegative, and \(\sqrt2\ge0\). Their sum defining \(A\) is therefore nonnegative. This sign permits later powers, products, and integral comparisons without absolute-value ambiguity.

\subsection{\texorpdfstring{mala\allowbreak{}Rejection\allowbreak{}Displacement nonneg}{malaRejectionDisplacement nonneg}}


The proposal displacement norm is nonnegative and the rejection factor is nonnegative by \texttt{metropolis\_\allowbreak{}rejection\_\allowbreak{}le\_\allowbreak{}abs}. Their product is nonnegative, permitting use of extended integrals and domination by an integrable envelope.

\subsection{\texorpdfstring{mala\allowbreak{}Proposal displacement le}{malaProposal displacement le}}


The triangle inequality gives \(\|T_h-x\|\le h\|g_x\|+\sqrt{2h}\|\xi\|\). For \(0\le h\le1\), \(h\le\sqrt h\) and \(\sqrt{2h}=\sqrt2\sqrt h\). Factor out \(\sqrt h\) to obtain \(\|T_h-x\|\le\sqrt h A(x,\xi)\).

\subsection{\texorpdfstring{mala\allowbreak{}Rejection\allowbreak{}Displacement pointwise}{malaRejectionDisplacement pointwise}}


Use rejection \(\le|R_h|\), the log-ratio estimate, and \(\|g_{T_h}\|\le\|g_x\|+L\|T_h-x\|\). Replace the displacement by \(\sqrt hA\), and absorb the remaining step factors with \(h\le1\). The result is \(D_h\le h\sqrt h\,L(2+L)A^3\), whose envelope has no step-size dependence.

\subsection{\texorpdfstring{integrable norm fourth mul exp neg mul norm sq}{integrable norm fourth mul exp neg mul norm sq}}


A polynomial of degree four can be absorbed by part of Gaussian decay: \(\|x\|^4e^{-a\|x\|^2}\le C e^{-(a/2)\|x\|^2}\) for \(a>0\). Gaussian integrability and measurable domination therefore prove the fourth weighted moment is finite.

\subsection{\texorpdfstring{integrable norm fourth mul exp neg of strong\allowbreak{}Convex}{integrable norm fourth mul exp neg of strongConvex}}


The quadratic strong-convexity lower bound gives \(e^{-U(x)}\le C e^{-(m/4)\|x\|^2}\). Multiply by \(\|x\|^4\) and invoke the fourth Gaussian weighted-moment theorem. This derives the unnormalized Gibbs fourth moment from the original potential assumptions.

\subsection{\texorpdfstring{integrable norm fourth gibbs\allowbreak{}Measure}{integrable norm fourth gibbsMeasure}}


Change measure to the normalized Gibbs density. Integrability of \(\|x\|^4\) under \(\pi_U\) reduces to the just-proved integrability of \(\|x\|^4e^{-U(x)}\) times a finite normalization factor. No fourth-moment premise is added.

\subsection{\texorpdfstring{mem\allowbreak{}Lp three id gibbs\allowbreak{}Measure}{memLp three id gibbsMeasure}}


The fourth-moment theorem gives the identity variable in \(L^4(\pi_U)\). On a probability space, \(L^4\subset L^3\), so the identity belongs to \(L^3\). This supplies the cubic moment needed for the rejection envelope.

\subsection{\texorpdfstring{mala\allowbreak{}Pair\allowbreak{}Measure}{malaPairMeasure}}


Take \(\pi_U\otimes\mathcal N(0,I)\), representing a Gibbs state and an independent standard Gaussian proposal innovation. This is the joint law under which the local rejection error is averaged.

\subsection{\texorpdfstring{mala\allowbreak{}Amplitude mem\allowbreak{}Lp three}{malaAmplitude memLp three}}


The Gibbs state is in \(L^3\); Lipschitz growth \(\|\nabla U(x)\|\le\|\nabla U(0)\|+L\|x\|\) puts its score in \(L^3\). Gaussian noise also belongs to \(L^3\). Lift both to their independent product and add their norm envelopes to obtain \(A\in L^3\).

\subsection{\texorpdfstring{mala\allowbreak{}Amplitude cube integrable}{malaAmplitude cube integrable}}


For nonnegative \(A\), its third power is \(\|A\|^3\). The \(L^3\) characterization gives integrability of this function. This is the finite ordinary integral used to define the rejection constant.

\subsection{\texorpdfstring{mala\allowbreak{}Rejection\allowbreak{}Constant}{malaRejectionConstant}}


Set \(C_{U,L}=L(2+L)\int A(x,\xi)^3d(\pi_U\otimes\gamma)\). The preceding declaration proves the integral is finite under the potential assumptions. All dependence on the positive step size has been removed from this constant.

\subsection{\texorpdfstring{mala\allowbreak{}Rejection\allowbreak{}Constant nonneg}{malaRejectionConstant nonneg}}


The cubic amplitude integrand is nonnegative and its integral is nonnegative. Since \(L\ge0\), the factor \(L(2+L)\) is nonnegative as well. Their product gives \(C_{U,L}\ge0\).

\subsection{\texorpdfstring{mala\allowbreak{}Rejection\allowbreak{}Displacement integrable}{malaRejectionDisplacement integrable}}


Continuity of the potential and gradient makes the rejection-displacement function measurable. Its absolute value equals itself by nonnegativity, and its pointwise bound is the integrable function \(h\sqrt h\,L(2+L)A^3\). Apply measurable domination to prove ordinary integrability.

\subsection{\texorpdfstring{mala\allowbreak{}Rejection\allowbreak{}Displacement integral bound}{malaRejectionDisplacement integral bound}}


Integrate the pointwise rejection envelope using the established integrability of both sides. Pull the step-independent factors out of the integral and recognize \(C_{U,L}\). This gives \(\mathbb ED_h\le C_{U,L}h\sqrt h\) for \(0<h\le1\).

\subsection{\texorpdfstring{exists mala\allowbreak{}Rejection\allowbreak{}Constant}{exists malaRejectionConstant}}


Choose the explicitly defined \texttt{mala\allowbreak{}Rejection\allowbreak{}Constant U L}. Its nonnegativity and the integrated envelope theorem verify the existential statement for every \(0<h\le1\) simultaneously. The same constant works throughout the small-step limit.

\section{\texorpdfstring{Detailed balance through accepted and rejected flows}{Detailed balance through accepted and rejected flows}}
\label{proof-module-UnadjustedLangevin-Metropolis-lean}
{\small\color{gray}Source module: \texttt{Unadjusted\allowbreak{}Langevin/\allowbreak{}Metropolis.\allowbreak{}lean}.\par}


The Metropolis correction preserves the initial marginal whenever the accepted proposal flow is invariant under swapping its two states. A positive proposal-pair density makes this symmetry explicit.

\subsection{\texorpdfstring{accepted\allowbreak{}Flow}{acceptedFlow}}


Weight the current/\allowbreak{}proposal joint measure \(\nu(dx,dy)\) by acceptance probability \(a(x,y)\). Thus \(A(dx,dy)=a(x,y)\nu(dx,dy)\), with \texttt{of\allowbreak{}Real} embedding its real nonnegative weight. The bounds \(0\le a\le1\) are supplied in the conservation theorems below.

\subsection{\texorpdfstring{rejected\allowbreak{}Flow}{rejectedFlow}}


The rejected pair measure is \(R(dx,dy)=(1-a(x,y))\nu(dx,dy)\). It retains the proposed state as an auxiliary coordinate; the eventual output maps it to the original state \(x\), preserving all rejection mass.

\subsection{\texorpdfstring{metropolis\allowbreak{}Output}{metropolisOutput}}


Define the output as \(\operatorname{snd}_*A+\operatorname{fst}_*R\). Accepted moves output \(y\), while rejected moves output \(x\). This is the complete transition measure and requires no renormalization of accepted moves.

\subsection{\texorpdfstring{accepted\allowbreak{}Flow add rejected\allowbreak{}Flow}{acceptedFlow add rejectedFlow}}


Since \(0\le a\le1\), the two nonnegative weights add to one after embedding in extended scalars. Additivity of with-density measures then gives \(A+R=\nu\). The proof checks both nonnegativity guards before combining \texttt{of\allowbreak{}Real} values.

\subsection{\texorpdfstring{marginals eq of swap invariant}{marginals eq of swap invariant}}


The identity \(\operatorname{fst}\circ\operatorname{swap}=\operatorname{snd}\) and pushforward composition give \(\operatorname{snd}_*\nu=\operatorname{fst}_*(\operatorname{swap}_*\nu)\). If the pair measure is swap invariant, this is \(\operatorname{fst}_*\nu\), proving equality of its marginals.

\subsection{\texorpdfstring{metropolis\allowbreak{}Output eq initial}{metropolisOutput eq initial}}


Detailed balance makes the two marginals of the accepted flow equal. Replace its second marginal by its first, then combine the accepted and rejected first marginals using \(A+R=\nu\). The output is therefore \(\operatorname{fst}_*\nu\), exactly the original state law.

\subsection{\texorpdfstring{map weighted comp}{map weighted comp}}


For a measurable \(f\) and weight \(w\), prove \(f_*((w\circ f)\mu)=w(f_*\mu)\). On each measurable set, the pushforward definition and nonnegative change-of-variables formula turn both sides into \(\int\mathbf1_A(f(x))w(f(x))d\mu(x)\). This is the weight-transport identity used by the actual MALA proposal map.

\subsection{\texorpdfstring{map weighted comp equiv}{map weighted comp equiv}}


Specialize the preceding pushforward/\allowbreak{}weight identity to a measurable equivalence. Its bundled measurability supplies the required hypothesis. The application below uses the coordinate swap equivalence to prove symmetry of an accepted density.

\subsection{\texorpdfstring{metropolis\allowbreak{}Acceptance}{metropolisAcceptance}}


For positive joint density \(w(x,y)\), define \(a(x,y)=\min(1,w(y,x)/w(x,y))\). The ratio is reverse flow divided by forward flow; positivity ensures both its meaning and nonnegative acceptance.

\subsection{\texorpdfstring{weight mul metropolis\allowbreak{}Acceptance}{weight mul metropolisAcceptance}}


Distribute the positive factor \(w(x,y)\) through the minimum. Cancellation of the nonzero denominator gives \(w(x,y)a(x,y)=\min(w(x,y),w(y,x))\). This algebra is the reason the accepted joint density is symmetric.

\subsection{\texorpdfstring{accepted\allowbreak{}Flow density}{acceptedFlow density}}


Two successive changes of density multiply their weights. Combine the original positive density \(w\) and acceptance \(a\), then apply the weight/\allowbreak{}minimum identity. The accepted flow becomes the base measure weighted by \(\min(w,w\circ\operatorname{swap})\), a visibly symmetric density.

\subsection{\texorpdfstring{metropolis\allowbreak{}Output density eq initial}{metropolisOutput density eq initial}}


The base measure is swap invariant, and the accepted minimum density is unchanged by swapping its arguments. The weight-transport theorem therefore proves accepted-flow detailed balance. The general accepted/\allowbreak{}rejected conservation theorem then identifies the output with the initial marginal.

\section{\texorpdfstring{Couple ULA to an invariant Metropolis chain}{Couple ULA to an invariant Metropolis chain}}
\label{proof-module-UnadjustedLangevin-MALACoupling-lean}
{\small\color{gray}Source module: \texttt{Unadjusted\allowbreak{}Langevin/\allowbreak{}MALACoupling.\allowbreak{}lean}.\par}


A common Gaussian proposal couples unadjusted ULA with MALA. Under a stable small step, their expected distance grows by at most the Gibbs-average rejected displacement at each step. This is an auxiliary approximation for stationarity.

\subsection{\texorpdfstring{mala\allowbreak{}Acceptance}{malaAcceptance}}


Define \(a_h(x,y)=\min(1,e^{R_h(x,y)})\), with the proved target/\allowbreak{}proposal log ratio. The minimum caps acceptance at one and the exponential ensures nonnegativity; these bounds justify the accepted/\allowbreak{}rejected measure split.

\subsection{\texorpdfstring{mala\allowbreak{}Proposal\allowbreak{}Acceptance}{malaProposalAcceptance}}


Compose \(a_h\) with \((x,\xi)\mapsto(x,T_h(x,\xi))\). The resulting \(a_h(x,T_h(x,\xi))\) is the acceptance weight on the original state/\allowbreak{}noise product space, prior to pushing to state/\allowbreak{}proposal pairs.

\subsection{\texorpdfstring{mala\allowbreak{}Acceptance nonneg}{malaAcceptance nonneg}}


Both entries of the minimum, one and the exponential, are nonnegative. Their minimum is nonnegative. This direct order proof is used when combining real weights through \texttt{ENNReal.\allowbreak{}of\allowbreak{}Real}.

\subsection{\texorpdfstring{mala\allowbreak{}Acceptance le one}{malaAcceptance le one}}


The minimum \(\min(1,e^{R_h})\) is at most its first entry. This proves \(a_h\le1\), so the rejection weight \(1-a_h\) is nonnegative.

\subsection{\texorpdfstring{mala\allowbreak{}Proposal measurable}{malaProposal measurable}}


The gradient is continuous by Lipschitzness. Projections, scalar multiplication, subtraction, and addition preserve continuity, so the joint proposal map \((x,\xi)\mapsto x-h\nabla U(x)+\sqrt{2h}\xi\) is measurable.

\subsection{\texorpdfstring{mala\allowbreak{}Proposal\allowbreak{}Acceptance measurable}{malaProposalAcceptance measurable}}


Differentiability makes \(U\) continuous, and Lipschitzness makes its gradient continuous. The log-ratio formula is assembled from continuous norms, inner products, and arithmetic; exponentiation and minimum preserve measurability. Composing with the measurable proposal gives the state/\allowbreak{}noise acceptance weight.

\subsection{\texorpdfstring{ula\allowbreak{}Law\allowbreak{}Step}{ulaLawStep}}


Define \(P_h\mu\) as the pushforward of \(\mu\otimes\gamma\) by the Euler proposal \(T_h\). The product measure encodes that the fresh Gaussian is independent of the current state. This transition is the actual independent-noise ULA step on laws.

\subsection{\texorpdfstring{mala\allowbreak{}One\allowbreak{}Step}{malaOneStep}}


Split \(\mu\otimes\gamma\) into acceptance- and rejection-weighted measures. Map the accepted part to \(T_h(x,\xi)\), and the rejected part to \(x\), then add them. This defines one full MALA transition with all rejection mass retained.

\subsection{\texorpdfstring{mala\allowbreak{}Coupled\allowbreak{}Step}{malaCoupledStep}}


For a pair \((x,y)\), draw one shared Gaussian \(\xi\). ULA always moves to \(T_h(x,\xi)\); MALA moves to \(T_h(y,\xi)\) with weight \(a_h(y,T_h(y,\xi))\), otherwise stays at \(y\). The two weighted pushforwards define the explicit coupling measure.

\subsection{\texorpdfstring{weighted accept reject add}{weighted accept reject add}}


For a measurable \(a\in[0,1]\), combine the two with-density measures using \(a+(1-a)=1\). Nonnegativity of both terms justifies the \texttt{of\allowbreak{}Real} addition identity, so their sum is exactly the original measure.

\subsection{\texorpdfstring{mala\allowbreak{}Coupled\allowbreak{}Step fst}{malaCoupledStep fst}}


The first-coordinate update is identical on acceptance and rejection. Combine those two weights to one before pushing forward. The remaining product marginal is the original first marginal times a fresh Gaussian, so the first output law is \(P_h(\operatorname{fst}_*\rho)\).

\subsection{\texorpdfstring{mala\allowbreak{}Coupled\allowbreak{}Step snd}{malaCoupledStep snd}}


Project both branches of the explicit coupling to the second state. They are exactly the accepted proposal and rejected current-state branches of MALA. Product marginal identities replace \(\rho\) by its second marginal, giving \(\operatorname{snd}_*\rho'=M_h(\operatorname{snd}_*\rho)\).

\subsection{\texorpdfstring{mala\allowbreak{}Coupled\allowbreak{}Step is\allowbreak{}Probability}{malaCoupledStep isProbability}}


The first marginal is a measurable pushforward of a probability product by \texttt{mala\allowbreak{}Coupled\allowbreak{}Step\_\allowbreak{}fst}. It has mass one, and a measure and its first-coordinate pushforward have the same total mass. Therefore the joint coupling also has total mass one.

\subsection{\texorpdfstring{mala\allowbreak{}Coupled\allowbreak{}Chain}{malaCoupledChain}}


Initialize on the diagonal \((X_0,Y_0)=(X,X)\) with \(X\sim\mu\), then recursively apply the explicit coupled step. Equal starting states give zero initial transport cost, while subsequent first and second marginals follow ULA and MALA respectively.

\subsection{\texorpdfstring{mala\allowbreak{}Coupled\allowbreak{}Chain is\allowbreak{}Probability}{malaCoupledChain isProbability}}


The diagonal pushforward of the initial probability is a probability. Induct on the number of transitions, using probability preservation of \texttt{mala\allowbreak{}Coupled\allowbreak{}Step} at each successor. This produces the instance needed for product noise and finite-cost arguments at every iteration.

\subsection{\texorpdfstring{mala\allowbreak{}Coupled\allowbreak{}Chain fst}{malaCoupledChain fst}}


At zero, diagonal initialization has first marginal \(\mu\). In the successor step apply \texttt{mala\allowbreak{}Coupled\allowbreak{}Step\_\allowbreak{}fst} and the induction hypothesis. The resulting recursion is precisely the iteration \(P_h^n\mu\).

\subsection{\texorpdfstring{mala\allowbreak{}Coupled\allowbreak{}Chain snd}{malaCoupledChain snd}}


At zero the second marginal is \(\pi\). The successor marginal theorem gives \(M_h\) applied to the current second marginal; substitute the induction hypothesis and the supplied one-step invariance \(M_h\pi=\pi\). Thus the second coordinate remains stationary through the coupling.

\subsection{\texorpdfstring{mala\allowbreak{}Proposal nonexpansive}{malaProposal nonexpansive}}


Shared noise cancels in \(T_h(x,\xi)-T_h(y,\xi)\). The sharp strong-convexity/\allowbreak{}smoothness contraction theorem for the deterministic gradient map bounds the remaining norm by \(\|x-y\|\) when \(0\le h\le2/(m+L)\). No Gaussian moment is used in this pointwise comparison.

\subsection{\texorpdfstring{mala\allowbreak{}Coupled\allowbreak{}Step pointwise}{malaCoupledStep pointwise}}


The accepted distance is at most \(\|x-y\|\) by nonexpansiveness. On rejection, insert \(T_h(y,\xi)\) and use the triangle inequality to bound the distance by \(\|x-y\|+\|T_h(y,\xi)-y\|\). Weight the two branches; their common old-distance coefficients add to one, and the remaining term is the rejection displacement.

\subsection{\texorpdfstring{mala\allowbreak{}L1\allowbreak{}Cost}{malaL1Cost}}


Define \(\mathcal C(\rho)=\int\|x-y\|d\rho\) as a nonnegative extended integral. Keeping it extended avoids assuming integrability before the cost recurrence proves it. Subsequent theorems establish finiteness and then convert it to an ordinary expected distance.

\subsection{\texorpdfstring{lintegral weighted map}{lintegral weighted map}}


Pull an extended nonnegative integral through a measurable pushforward, then through a with-density measure. Multiplication of the two nonnegative real factors becomes the embedded product \(a(x)b(f(x))\). This gives an exact integration formula for each branch of the coupling.

\subsection{\texorpdfstring{mala\allowbreak{}Coupled\allowbreak{}Step cost formula}{malaCoupledStep cost formula}}


Apply \texttt{lintegral\_\allowbreak{}weighted\_\allowbreak{}map} to accepted and rejected branches of the coupled measure. Their cost is the integral of the acceptance-weighted proposal distance plus rejection-weighted stay distance over \(\rho\otimes\gamma\). The formula exposes exactly the pointwise expression bounded by \texttt{mala\allowbreak{}Coupled\allowbreak{}Step\_\allowbreak{}pointwise}.

\subsection{\texorpdfstring{mala\allowbreak{}Coupled\allowbreak{}Step cost le}{malaCoupledStep cost le}}


Integrate the pointwise coupling bound. The old distance integrates to \(\mathcal C(\rho)\); the error depends only on the second state and innovation. Its second marginal is Gibbs by hypothesis, so that integral is bounded by \(C_{U,L}h\sqrt h\). Thus \(\mathcal C(\rho')\le\mathcal C(\rho)+C_{U,L}h\sqrt h\), valid initially as an extended inequality.

\subsection{\texorpdfstring{mala\allowbreak{}Coupled\allowbreak{}Chain cost le}{malaCoupledChain cost le}}


Induct from zero diagonal cost. The second marginal remains Gibbs by the supplied one-step MALA invariance, so the one-step cost bound applies at every iteration. Summing the same local error gives \(\mathcal C(\rho_n)\le nC_{U,L}h\sqrt h\).

\subsection{\texorpdfstring{mala\allowbreak{}Coupled\allowbreak{}Chain distance integrable}{malaCoupledChain distance integrable}}


The quantitative extended cost estimate bounds \(\int\|x-y\|d\rho_n\) by a finite embedded real number. The distance is measurable and nonnegative, so finite extended integral is equivalent to ordinary integrability. This discharges the moment condition rather than assuming it.

\subsection{\texorpdfstring{mala\allowbreak{}Coupled\allowbreak{}Chain distance integral le}{malaCoupledChain distance integral le}}


Convert the finite extended cost inequality to real scalars using \texttt{ENNReal.\allowbreak{}to\allowbreak{}Real}. The newly proved integrability identifies the real value of the extended integral with the ordinary expected distance. Consequently \(\mathbb E_{\rho_n}\|X-Y\|\le nC_{U,L}h\sqrt h\).

\section{\texorpdfstring{Identify Gaussian proposals with their radial density}{Identify Gaussian proposals with their radial density}}
\label{proof-module-UnadjustedLangevin-GaussianProposalDensity-lean}
{\small\color{gray}Source module: \texttt{Unadjusted\allowbreak{}Langevin/\allowbreak{}Gaussian\allowbreak{}Proposal\allowbreak{}Density.\allowbreak{}lean}.\par}


For \(h>0\), the affine standard Gaussian proposal \(c+\sqrt{2h}\,Z\) has density proportional to \(e^{-\|y-c\|^2/(4h)}\). The proof identifies probability measures by characteristic functions and does not require a preselected coordinate formula for the Gaussian normalizer.

\subsection{\texorpdfstring{radial\allowbreak{}Gaussian\allowbreak{}Integral}{radialGaussianIntegral}}


Set \(J(a)=\int_H e^{-a\|z\|^2}dz\) as an ordinary real integral. The positive-precision results below prove its integrability and positivity before using \(J(a)^{-1}\) as a normalizing constant.

\subsection{\texorpdfstring{radial\allowbreak{}Gaussian\allowbreak{}Normalizer}{radialGaussianNormalizer}}


Define \(C(a)=J(a)^{-1}\). This is an actual integral normalizer on the given finite-dimensional Hilbert space, avoiding a coordinate-dependent choice of a formula for Lebesgue density.

\subsection{\texorpdfstring{radial\allowbreak{}Gaussian\allowbreak{}Integral pos}{radialGaussianIntegral pos}}


For \(a>0\), Gaussian integrability was established in \texttt{integrable\_\allowbreak{}exp\_\allowbreak{}neg\_\allowbreak{}mul\_\allowbreak{}norm\_\allowbreak{}sq}. The integrand is strictly positive, so its ordinary integral is strictly positive by \texttt{integral\_\allowbreak{}exp\_\allowbreak{}pos}.

\subsection{\texorpdfstring{radial\allowbreak{}Gaussian\allowbreak{}Normalizer pos}{radialGaussianNormalizer pos}}


The inverse of a strictly positive real number is strictly positive. Apply this to \(J(a)\), whose positivity was proved immediately above. This also proves the normalizer is nonzero for later cancellation.

\subsection{\texorpdfstring{radial\allowbreak{}Gaussian\allowbreak{}Density}{radialGaussianDensity}}


Define \(q_{a,c}(z)=C(a)e^{-a\|z-c\|^2}\), embedded in extended nonnegative scalars. The center changes the translated argument, while the normalizer depends only on precision and dimension.

\subsection{\texorpdfstring{integrable radial\allowbreak{}Gaussian\allowbreak{}Density real}{integrable radialGaussianDensity real}}


Translation by \(c\) preserves Lebesgue measure, so \(e^{-a\|z-c\|^2}\) is integrable whenever \(a>0\). Multiply by the finite real constant \(C(a)\) to obtain integrability of the displayed real density.

\subsection{\texorpdfstring{integral radial\allowbreak{}Gaussian\allowbreak{}Density real}{integral radialGaussianDensity real}}


Translate \(z\mapsto z-c\) in the integral and pull out \(C(a)\). The result is \(C(a)J(a)=1\), justified by \(J(a)>0\). Thus normalization is independent of the center.

\subsection{\texorpdfstring{radial\allowbreak{}Gaussian\allowbreak{}Density is\allowbreak{}Probability}{radialGaussianDensity isProbability}}


The total mass of the with-density measure is the extended integral of its density. Nonnegativity and the proved real integrability identify it with \texttt{of\allowbreak{}Real} of the real integral, which is one. This constructs probability normalization, including finiteness.

\subsection{\texorpdfstring{integral complex radial\allowbreak{}Gaussian}{integral complex radialGaussian}}


For real \(a\), \(e^{-(a:\mathbb C)\|z\|^2}\) is the complex embedding of the real Gaussian exponential. Rewrite the integrand with \texttt{Complex.\allowbreak{}of\allowbreak{}Real\_\allowbreak{}exp}; the continuous real embedding commutes with integration, giving the complex embedding of \(J(a)\).

\subsection{\texorpdfstring{integral radial\allowbreak{}Gaussian character}{integral radialGaussian character}}


Apply the finite-dimensional complex Gaussian Fourier integral with imaginary linear coefficient. Combining the two exponentials and using \(i^2=-1\) gives \(\int e^{-a\|z\|^2}e^{i\langle z,t\rangle}dz=J(a)e^{-\|t\|^2/(4a)}\). The zero-frequency Gaussian formula identifies the normalization factor with the same real integral \(J(a)\).

\subsection{\texorpdfstring{char\allowbreak{}Fun radial\allowbreak{}Gaussian\allowbreak{}Density}{charFun radialGaussianDensity}}


Change from the density measure to a weighted Lebesgue integral. Translate \(z=u+c\); the character splits into \(e^{i\langle u,t\rangle}e^{i\langle c,t\rangle}\). The Fourier formula gives a factor \(J(a)\), canceled by \(C(a)\), leaving \(e^{-\|t\|^2/(4a)}e^{i\langle c,t\rangle}\).

\subsection{\texorpdfstring{gaussian\allowbreak{}Proposal\allowbreak{}Normalizer}{gaussianProposalNormalizer}}


Specialize the precision to \(a=1/(4h)\). The resulting constant normalizes a Gaussian with covariance \(2hI\). Its independence of the current state is essential when the same factor cancels in the forward and reverse Metropolis ratio.

\subsection{\texorpdfstring{gaussian\allowbreak{}Proposal\allowbreak{}Density}{gaussianProposalDensity}}


Define \(q_h(c,y)=C(1/(4h))e^{-\|y-c\|^2/(4h)}\). For \(h>0\), this is the density of a covariance-\(2hI\) Gaussian centered at \(c\). The density identification is proved later rather than built into the definition.

\subsection{\texorpdfstring{gaussian\allowbreak{}Proposal\allowbreak{}Normalizer pos}{gaussianProposalNormalizer pos}}


For \(h>0\), the precision \(1/(4h)\) is positive. Apply positivity of the radial Gaussian normalizer at that precision. This nonzero constant is used in detailed-balance cancellation.

\subsection{\texorpdfstring{radial\allowbreak{}Gaussian\allowbreak{}Density eq proposal\allowbreak{}Density}{radialGaussianDensity eq proposalDensity}}


Unfold the two density definitions and the proposal normalizer. Their exponents agree by \(-(1/(4h))\|y-c\|^2=-\|y-c\|^2/(4h)\). Function extensionality gives equality of the density functions, with no probabilistic argument needed.

\subsection{\texorpdfstring{gaussian proposal density}{gaussian proposal density}}


Compute characteristic functions on both sides. The affine image \(c+\sqrt{2h}Z\) has character \(e^{-h\|t\|^2}e^{i\langle c,t\rangle}\); the radial-density formula with \(a=1/(4h)\) gives the same expression. Both measures are probabilities, so \texttt{Measure.\allowbreak{}ext\_\allowbreak{}of\_\allowbreak{}char\allowbreak{}Fun} identifies them. The square-root identity uses \(h>0\).

\subsection{\texorpdfstring{gaussian proposal radial density}{gaussian proposal radial density}}


Rewrite the proposal density using \texttt{radial\allowbreak{}Gaussian\allowbreak{}Density\_\allowbreak{}eq\_\allowbreak{}proposal\allowbreak{}Density} in the established affine-Gaussian law equality. This yields the equivalent precision-parameter form used in the detailed-balance proof.

\subsection{\texorpdfstring{ula proposal radial density}{ula proposal radial density}}


Set the affine Gaussian center to \(c=x-h\nabla U(x)\). The fixed-state density theorem then gives \(x-h\nabla U(x)+\sqrt{2h}Z\) its exact radial density. No regularity of \(U\) is required because \(x\) is fixed throughout this statement.

\subsection{\texorpdfstring{ula proposal density}{ula proposal density}}


Expand the radial density at center \(x-h\nabla U(x)\). The exponent becomes \(-\|y-x+h\nabla U(x)\|^2/(4h)\), with the same positive normalizer. This is the explicit ULA transition density used in the Metropolis acceptance ratio.

\section{\texorpdfstring{The joint initial-\allowbreak{}state and proposal density}{The joint initial-state and proposal density}}
\label{proof-module-UnadjustedLangevin-GaussianProposalPairDensity-lean}
{\small\color{gray}Source module: \texttt{Unadjusted\allowbreak{}Langevin/\allowbreak{}Gaussian\allowbreak{}Proposal\allowbreak{}Pair\allowbreak{}Density.\allowbreak{}lean}.\par}


For a measurable center \(c(x)\), independent \(X\sim\mu\) and \(Z\sim\mathcal N(0,I)\) produce the joint law \(\mu(dx)q_h(c(x),y)dy\).

\subsection{\texorpdfstring{Disintegrate the Gaussian change of variables}{Disintegrate the Gaussian change of variables}}


For \(X\sim\mu\) and independent \(Z\sim\mathcal N(0,I)\), fix \(x\) and apply the proved proposal-density identity to \(c(x)+\sqrt{2h}Z\). It gives conditional integration against \(q_h(c(x),y)dy\).



To obtain equality of joint measures, evaluate a measurable set \(A\subset H\times H\). Tonelli writes both sides as integrals over \(x\) of their \(y\)-sections. The fixed-center density theorem identifies each section, hence the measures agree. Measurability of \(c\) supplies measurability of the pair update and joint density; no probability or moment condition on \(\mu\) is added.

\section{\texorpdfstring{Exact Gibbs invariance of the Metropolis-\allowbreak{}adjusted step}{Exact Gibbs invariance of the Metropolis-adjusted step}}
\label{proof-module-UnadjustedLangevin-MALAInvariance-lean}
{\small\color{gray}Source module: \texttt{Unadjusted\allowbreak{}Langevin/\allowbreak{}MALAInvariance.\allowbreak{}lean}.\par}


The Gaussian proposal density identifies the MALA pair flow. Its reversed-to-forward density ratio is precisely the exponential of the implemented log ratio, so the Metropolis balance theorem proves exact Gibbs invariance.

\subsection{\texorpdfstring{mala\allowbreak{}Proposal\allowbreak{}Joint}{malaProposalJoint}}


Push \(\mu\otimes\gamma\) through \((x,\xi)\mapsto(x,T_h(x,\xi))\). Its first coordinate retains the current state, and its second records the actual Gaussian proposal. This joint measure is the input to the accepted-flow detailed-balance theorem.

\subsection{\texorpdfstring{mala\allowbreak{}Joint\allowbreak{}Weight}{malaJointWeight}}


Set \(w_h(x,y)=\exp(-U(x)-\|y-x+h\nabla U(x)\|^2/(4h))\). The Gibbs and Gaussian normalization constants are kept in a separate symmetric base measure, so \(w_h\) is strictly positive and its reverse/\allowbreak{}forward ratio is easy to compute.

\subsection{\texorpdfstring{mala\allowbreak{}Joint\allowbreak{}Weight ratio}{malaJointWeight ratio}}


Divide the two positive exponentials and use \(e^a/e^b=e^{a-b}\). Expanding the definition of \(w_h\) shows the exponent difference is exactly \texttt{mala\allowbreak{}Log\allowbreak{}Ratio}. Thus \(w_h(y,x)/w_h(x,y)=e^{R_h(x,y)}\).

\subsection{\texorpdfstring{mala\allowbreak{}One\allowbreak{}Step eq metropolis\allowbreak{}Output}{malaOneStep eq metropolisOutput}}


Transport acceptance and rejection weights through the measurable pair map using \texttt{map\_\allowbreak{}weighted\_\allowbreak{}comp}. Projecting the accepted pair to its second coordinate gives the proposal; projecting rejected pairs to the first gives the unchanged state. These are precisely the two branches defining \texttt{mala\allowbreak{}One\allowbreak{}Step}.

\subsection{\texorpdfstring{mala\allowbreak{}Proposal\allowbreak{}Joint gibbs density}{malaProposalJoint gibbs density}}


Apply the Gaussian joint-density theorem with center \(x-h\nabla U(x)\), then multiply its density by the Gibbs density. Combine the exponential factors into \(w_h(x,y)\). Pull the two state-independent normalizers into the base measure \(Z_U^{-1}C_h(dx\,dy)\), which is invariant under coordinate swap.

\subsection{\texorpdfstring{mala\allowbreak{}One\allowbreak{}Step gibbs invariant}{malaOneStep gibbs invariant}}


The joint proposal measure has a positive density \(w_h\) against a symmetric base. Its acceptance rule is \(\min(1,w_h(y,x)/w_h(x,y))\) by \texttt{mala\allowbreak{}Joint\allowbreak{}Weight\_\allowbreak{}ratio}. Hence its accepted density is the symmetric minimum \(\min(w_h(x,y),w_h(y,x))\).



Apply \texttt{metropolis\allowbreak{}Output\_\allowbreak{}density\_\allowbreak{}eq\_\allowbreak{}initial} and the exact step/\allowbreak{}output identity. The initial marginal of the proposal pair is Gibbs, since projecting \(\pi_U\otimes\gamma\) to its state returns \(\pi_U\). Thus \(M_h\pi_U=\pi_U\) is proved with no stationarity input.

\section{\texorpdfstring{Characteristic functions are continuous in L\ensuremath{{}^{2}}}{Characteristic functions are continuous in L²}}
\label{proof-module-UnadjustedLangevin-LpCharacteristic-lean}
{\small\color{gray}Source module: \texttt{Unadjusted\allowbreak{}Langevin/\allowbreak{}Lp\allowbreak{}Characteristic.\allowbreak{}lean}.\par}


At each fixed frequency, the characteristic function of the law of a real \(L^2\) variable depends continuously on that variable. Consequently centered Gaussian laws are closed under \(L^2\) convergence, with variance given by the squared norm.

\subsection{\texorpdfstring{centered\allowbreak{}Characteristic}{centeredCharacteristic}}


Set \(c_t(x)=e^{itx}-1\), represented by the unit-radius circle map centered at \(-1\). Subtracting one makes \(c_t(0)=0\), which is the condition needed by the Lipschitz composition construction on \(L^2\).

\subsection{\texorpdfstring{centered\allowbreak{}Characteristic lipschitz}{centeredCharacteristic lipschitz}}


Multiplication by \(t\) is \(|t|\)-Lipschitz, and the unit circle parametrization is \(1\)-Lipschitz. Composing gives \(|c_t(x)-c_t(y)|\le|t||x-y|\). The proof identifies the nonnegative Lipschitz constants before applying the composition theorem.

\subsection{\texorpdfstring{centered\allowbreak{}Characteristic zero}{centeredCharacteristic zero}}


At zero the exponential is one, hence \(c_t(0)=e^0-1=0\). Unfolding the circle-map definition makes this a scalar identity; this vanishing proof permits \texttt{comp\allowbreak{}Lp} to lift the map without an integrability assumption on a nonzero constant.

\subsection{\texorpdfstring{characteristic\allowbreak{}Lp}{characteristicLp}}


Lift \(c_t\) to \(L^2\) by Lipschitz composition, then add the constant class one. Thus \(\operatorname{characteristicLp}(t,X)\) represents \(e^{itX}\). The probability assumption makes the constant function square integrable.

\subsection{\texorpdfstring{characteristic\allowbreak{}Lp coe}{characteristicLp coe}}


Combine the almost-everywhere representative identities for lifted composition, \(L^2\) addition, and the constant class. Their intersection has full measure; there the representative is \((e^{itX}-1)+1=e^{itX}\). This connects quotient-space operations with the random function used in the characteristic integral.

\subsection{\texorpdfstring{char\allowbreak{}Fun map eq inner characteristic\allowbreak{}Lp}{charFun map eq inner characteristicLp}}


The pushforward integral theorem gives \(\widehat{P_X}(t)=\int e^{itX}dP\). In complex \(L^2\), inner product with the constant one has exactly this integral because its conjugate is one. The representative identity above therefore yields \(\widehat{P_X}(t)=\langle1,\operatorname{characteristicLp}(t,X)\rangle\).

\subsection{\texorpdfstring{continuous char\allowbreak{}Fun map Lp}{continuous charFun map Lp}}


Lipschitz composition is continuous on \(L^2\); adding the constant one preserves continuity. Pairing with one is a continuous linear functional. The preceding integral identity rewrites their composition as \(X\mapsto\widehat{P_X}(t)\), proving continuity of the actual-law characteristic function at every fixed frequency.

\subsection{\texorpdfstring{continuous gaussian\allowbreak{}Real char\allowbreak{}Fun Lp}{continuous gaussianReal charFun Lp}}


The centered Gaussian characteristic function is \(e^{-t^2\|X\|_2^2/2}\). Norm, squaring, conversion of a nonnegative real to \texttt{\ensuremath{\mathbb R}\ensuremath{\ge}0}, and the exponential are continuous. Lean unfolds the Gaussian formula and composes these continuity facts; the variance is the actual squared \(L^2\) norm.

\subsection{\texorpdfstring{has\allowbreak{}Law gaussian\allowbreak{}Real of tendsto Lp eventually}{hasLaw gaussianReal of tendsto Lp eventually}}


At each frequency, \(L^2\) convergence makes the actual characteristic functions converge. It also makes the centered Gaussian characteristic functions with variances \(\|X_i\|_2^2\) converge to variance \(\|Y\|_2^2\). Eventual Gaussianity identifies the two nets eventually, so uniqueness of limits identifies their limits. \texttt{Measure.\allowbreak{}ext\_\allowbreak{}of\_\allowbreak{}char\allowbreak{}Fun} then gives the claimed Gaussian law of \(Y\); a nontrivial filter is required for uniqueness.

\subsection{\texorpdfstring{has\allowbreak{}Law gaussian\allowbreak{}Real of tendsto Lp}{hasLaw gaussianReal of tendsto Lp}}


If every approximant has the stated Gaussian law, then the law identity holds eventually for any filter. Feed \texttt{Eventually.\allowbreak{}of\_\allowbreak{}forall} to the preceding closure theorem. This version is suited to finite orthogonal sums converging to an infinite Gaussian series.

\section{\texorpdfstring{A vanishing coupling distance determines the limiting law}{A vanishing coupling distance determines the limiting law}}
\label{proof-module-UnadjustedLangevin-CouplingLimit-lean}
{\small\color{gray}Source module: \texttt{Unadjusted\allowbreak{}Langevin/\allowbreak{}Coupling\allowbreak{}Limit.\allowbreak{}lean}.\par}


If \(X_n\to X\) in \(L^2\) and each \(\operatorname{Law}(X_n)\) couples to a fixed probability \(\pi\) with mean distance tending to zero, then \(\operatorname{Law}(X)=\pi\).

\subsection{\texorpdfstring{characteristic\allowbreak{}Test norm sub le}{characteristicTest norm sub le}}


The unit-circle map is \(1\)-Lipschitz, so the difference of the two characters is at most \(|\langle x-y,t\rangle|\). Cauchy–Schwarz then bounds this by \(\|t\|\|x-y\|\). Thus characteristic tests have an explicit Lipschitz constant at each fixed frequency.

\subsection{\texorpdfstring{char\allowbreak{}Fun marginals norm sub le}{charFun marginals norm sub le}}


Pull both marginal characteristic integrals back to the coupling space and subtract. Each character is bounded by one, hence integrable under the probability coupling. The integral norm inequality and the pointwise Lipschitz estimate give \(|\widehat\mu(t)-\widehat\pi(t)|\le\|t\|\mathbb E\|X-Y\|\), with the distance's integrability explicitly supplied.

\subsection{\texorpdfstring{continuous char\allowbreak{}Fun map Lp hilbert}{continuous charFun map Lp hilbert}}


Apply the continuous linear functional \(x\mapsto\langle t,x\rangle\) to an \(H\)-valued \(L^2\) variable. Its lifted map is continuous on \(L^2\). The Hilbert-valued characteristic function at \(t\) is the scalar projected characteristic function at frequency one, so the scalar continuity theorem proves the claim.

\subsection{\texorpdfstring{measure map Lp eq of coupling cost tendsto}{measure map Lp eq of coupling cost tendsto}}


For each \(t\), the coupling inequality forces the approximating characteristic functions toward \(\widehat\pi(t)\), since the expected distance tends to zero. Independently, \(L^2\) convergence forces them toward \(\widehat{P_X}(t)\). Uniqueness of limits along a nontrivial filter identifies these values, and \texttt{Measure.\allowbreak{}ext\_\allowbreak{}of\_\allowbreak{}char\allowbreak{}Fun} yields \(P_X=\pi\).

\section{\texorpdfstring{From scalar coordinate laws to independent Gaussian vectors}{From scalar coordinate laws to independent Gaussian vectors}}
\label{proof-module-UnadjustedLangevin-GaussianInnovationVectors-lean}
{\small\color{gray}Source module: \texttt{Unadjusted\allowbreak{}Langevin/\allowbreak{}Gaussian\allowbreak{}Innovation\allowbreak{}Vectors.\allowbreak{}lean}.\par}


For finitely many spatial coordinates, independent standard normal scalar innovations assemble into independent standard Gaussian vectors. Scalar normal laws without joint independence would not suffice for this conclusion.

\subsection{\texorpdfstring{gaussian\allowbreak{}Innovations path\allowbreak{}Law}{gaussianInnovations pathLaw}}


Regroup the jointly independent scalar coordinates indexed by \((n,i)\) into vector blocks \(W_n\). A finite block with independent standard Gaussian coordinates has the standard Euclidean Gaussian law. Product-measure identification and reindexing then show that the full vector path has law \(\bigotimes_n\mathcal N(0,I)\).

\subsection{\texorpdfstring{gaussian\allowbreak{}Innovations vector\allowbreak{}Law}{gaussianInnovations vectorLaw}}


Project the established whole-path product law to its \(n\)-th vector coordinate. Product marginals give \(P\circ W_n^{-1}=\mathcal N(0,I)\). This derives the multivariate law from the scalar coordinate assumptions rather than adding it as an algorithm hypothesis.

\subsection{\texorpdfstring{gaussian\allowbreak{}Innovations vector\allowbreak{}Independent}{gaussianInnovations vectorIndependent}}


The coordinate maps of the vector product law are independent. Pull that independence back through the actual innovation-path map using \texttt{gaussian\allowbreak{}Innovations\_\allowbreak{}path\allowbreak{}Law}. Thus the vector-valued \(W_n\) are independent over time, not just uncorrelated.

\section{\texorpdfstring{Prove Gibbs stationarity by refining invariant Metropolis chains}{Prove Gibbs stationarity by refining invariant Metropolis chains}}
\label{proof-module-UnadjustedLangevin-GibbsStationarity-lean}
{\small\color{gray}Source module: \texttt{Unadjusted\allowbreak{}Langevin/\allowbreak{}Gibbs\allowbreak{}Stationarity.\allowbreak{}lean}.\par}


A Brownian Langevin integral solution started from Gibbs stays Gibbs. The proof uses convergent Euler approximations and exactly invariant MALA chains, with no weak-PDE uniqueness assumption.

\subsection{\texorpdfstring{hilbert lp\allowbreak{}Chain coe}{hilbert lpChain coe}}


Induct over the chain recursion. The representative of the lifted gradient equals the pointwise gradient almost everywhere, and the representative formulas for addition and scalar multiplication give the same Euler successor. Intersect the finite set of null-set complements needed at each index to identify the \(L^2\) chain with the ordinary recursion.

\subsection{\texorpdfstring{hilbert lp\allowbreak{}Chain law eq ula\allowbreak{}Law\allowbreak{}Step}{hilbert lpChain law eq ulaLawStep}}


Use the almost-everywhere pointwise-chain identification, the independent-iid realization law theorem, and the canonical path/\allowbreak{}Markov-law identity. With innovation law \(\mathcal N(0,I)\), the generic independent-noise step is exactly \texttt{ula\allowbreak{}Law\allowbreak{}Step}. Thus the \(L^2\) implementation has law \(P_h^n(P_{X_0})\).

\subsection{\texorpdfstring{additive\allowbreak{}Euler\allowbreak{}Chain eq hilbert lp\allowbreak{}Chain}{additiveEulerChain eq hilbert lpChain}}


Induct on the mesh index. Both chains start at the same \(L^2\) variable. Their successor equations agree because the supplied additive increment equals \(\sqrt{2h}W_k\), making the noise terms identical as classes.

\subsection{\texorpdfstring{brownian\allowbreak{}Langevin\allowbreak{}Noise grid increment}{brownianLangevinNoise grid increment}}


Apply the Brownian input cocycle at \(a=kh\), \(t=h\). The elementary identity \((k+1)h=kh+h\) matches its endpoints to the grid notation. This supplies the exact additive increment used by the normalized Euler noise.

\subsection{\texorpdfstring{brownian additive\allowbreak{}Euler\allowbreak{}Chain eq lp\allowbreak{}Chain}{brownian additiveEulerChain eq lpChain}}


The Brownian grid-increment identity and \texttt{brownian\allowbreak{}Euler\allowbreak{}Noise\_\allowbreak{}endpoint} verify the increment premise of the generic chain equality. Therefore the deterministic matching-noise Euler approximation is exactly the already defined Brownian-normalized \(L^2\) Euler chain.

\subsection{\texorpdfstring{brownian\allowbreak{}Euler\allowbreak{}Noise vector\allowbreak{}Gaussian}{brownianEulerNoise vectorGaussian}}


Each normalized coordinate is standard Gaussian. The supplied joint scalar-coordinate independence lets \texttt{gaussian\allowbreak{}Innovations\_\allowbreak{}vector\allowbreak{}Law} and \texttt{gaussian\allowbreak{}Innovations\_\allowbreak{}vector\allowbreak{}Independent} assemble those coordinates into iid standard multivariate Gaussian vectors. These are the exact law hypotheses needed by the Markov-recursion theorem.

\subsection{\texorpdfstring{Pass the invariant Gibbs marginal to the continuous limit}{Pass the invariant Gibbs marginal to the continuous limit}}


Fix \(T>0\) and choose meshes \(h_n=T/N_n\to0\) small enough for the nonexpansive MALA coupling. Matching-noise Euler approximations converge in \(L^2\) to \(Y_T\). Their laws are exactly \(P_{h_n}^{N_n}\pi_U\), by the proved iid Brownian innovation laws and initialization independence.



Couple each Euler law to \(N_n\) MALA steps started from Gibbs. Exact MALA invariance makes the second marginal \(\pi_U\). The rejection bound gives expected coupling distance at most \(N_nC h_n\sqrt{h_n}=TC\sqrt{h_n}\to0\); integrability of this distance is proved from the finite extended-cost estimate.



Apply \texttt{measure\_\allowbreak{}map\_\allowbreak{}Lp\_\allowbreak{}eq\_\allowbreak{}of\_\allowbreak{}coupling\_\allowbreak{}cost\_\allowbreak{}tendsto}: characteristic functions of the Euler endpoints converge to those of \(Y_T\) by \(L^2\) convergence, and to those of \(\pi_U\) by the vanishing coupling cost. Equality of characteristic functions proves \(P_{Y_T}=\pi_U\). The zero-horizon case follows directly from the integral equation.

\section{\texorpdfstring{Identify actual Gaussian IPLA with the shared ULA construction}{Identify actual Gaussian IPLA with the shared ULA construction}}
\label{proof-module-IPLA-ULAAdapter-lean}
{\small\color{gray}Source module: \texttt{IPLA/\allowbreak{}ULAAdapter.\allowbreak{}lean}.\par}


This adapter translates the particle coordinates into shared ULA probability interfaces. It proves the needed Gaussian and law identities, while all Brownian construction and stationarity arguments remain in their shared modules.

\subsection{\texorpdfstring{Brownian increments satisfy the complete algorithm noise specification}{Brownian increments satisfy the complete algorithm noise specification}}


On the independent product of an initial coupling and a Brownian source, define \(\xi_{k,a}=(B_a((k+1)h)-B_a(kh))/\sqrt h\). Each Brownian coordinate has centered Gaussian increments of variance \(h\), so each normalized coordinate has law \(\mathcal N(0,1)\). Independent Brownian coordinates and disjoint increments give mutual independence across coordinates and times. The product construction makes the complete noise path independent of the initial coupling. The shared product-noise theorems prove these three assertions, which are assembled into \texttt{Gaussian\allowbreak{}Innovations}.



The local instance \texttt{algorithm\allowbreak{}Product\allowbreak{}Probability} follows because a coupling of probability measures has total mass one and the independent source law has mass one; their product therefore has total mass one. This justifies the probability-space hypotheses of the innovation and L\ensuremath{{}^{2}} law theorems.




\textbf{Correspondence with the displayed declarations.}


\begin{itemize}

\item \texttt{product\allowbreak{}Gaussian\allowbreak{}Innovations}: The concrete product realization satisfies every distributional field of the actual IPLA algorithm. Standard Gaussian laws and all independence claims are conclusions from the independent Brownian coordinate source.
\end{itemize}
\subsection{\texorpdfstring{Transfer the product realization to any implementation}{Transfer the product realization to any implementation}}


The Brownian product's initial random variable has the prescribed law \(P\circ Z_0^{-1}\), and its innovations satisfy the Gaussian structure just proved. Implementation independence therefore identifies its IPLA law with any supplied Gaussian implementation having the same initial law. The shared product-chain law identity then gives

\[\rho_k=\operatorname{productOutputLaw}(V_N,\gamma/N,\mu_0,k).\]

Finite coupling cost and the target second moment are used here to construct the product initial state in \(L^2\). Neither stationarity nor a reference diffusion is assumed for this law identity.




Canonical Brownian Euler increments are proved standard Gaussian in every coordinate, mutually independent across times and coordinates, and independent of the coupled initialization. These are exactly the three fields of \texttt{Gaussian\allowbreak{}Innovations}. Their initial law agrees with the supplied initialization law on the product extension, so canonical-law uniqueness identifies the resulting algorithm output laws. The final Markov-step identity then shows that the actual supplied chain iterates the same ULA kernel at step \(\gamma/N\).




\textbf{Correspondence with the displayed declarations.}


\begin{itemize}

\item \texttt{laws\_\allowbreak{}eq\_\allowbreak{}product\allowbreak{}Brownian\allowbreak{}Output}: The law produced by the Brownian product construction is exactly the law of any supplied Gaussian IPLA implementation with the same initialization. This bridge has no stationarity or reference-process hypothesis.
\end{itemize}
\subsection{\texorpdfstring{The actual chain iterates the ULA Markov transition}{The actual chain iterates the ULA Markov transition}}


The innovation hypotheses imply independent vector-valued Gaussian increments by grouping the scalar coordinate family. Applying the shared Hilbert-space law theorem to \(V_N\), step \(\gamma/N\), and smoothness \(NL\) gives

\[\rho_k=\operatorname{ulaLawStep}(V_N,\gamma/N)^{[k]}(\mu_0).\]

The continuous \(L^2\) representatives are measurable; the initial state is independent of the full noise path by hypothesis. This exact transition identity transfers both convergence and finiteness statements from ULA to the actual IPLA laws.




\textbf{Correspondence with the displayed declarations.}


\begin{itemize}

\item \texttt{laws\_\allowbreak{}eq\_\allowbreak{}ula\allowbreak{}Law\allowbreak{}Step}: The actual IPLA law is precisely ULA for the rescaled lifted potential. All vector Gaussian and independence facts follow from the displayed scalar innovation specification of the algorithm.
\end{itemize}
\section{\texorpdfstring{Construct an isonormal Gaussian family from product coordinates}{Construct an isonormal Gaussian family from product coordinates}}
\label{proof-module-UnadjustedLangevin-GaussianHilbert-lean}
{\small\color{gray}Source module: \texttt{Unadjusted\allowbreak{}Langevin/\allowbreak{}Gaussian\allowbreak{}Hilbert.\allowbreak{}lean}.\par}


An explicit product of standard normal coordinates realizes a linear isometry \(I:H\to L^2\) with \(I(x)\sim\mathcal N(0,\|x\|^2)\) and covariance \(\mathbb E[I(x)I(y)]=\langle x,y\rangle\).

\subsection{\texorpdfstring{source}{source}}


Use the product probability \(Q=\bigotimes_{i\in\iota}\mathcal N(0,1)\) on real coordinate families. The construction uses \texttt{Measure.\allowbreak{}infinite\allowbreak{}Pi}; independence and coordinate marginals are properties of this concrete source rather than extra assumptions on an unspecified Gaussian process.

\subsection{\texorpdfstring{coordinate has\allowbreak{}Law}{coordinate hasLaw}}


The coordinate projection \(\omega\mapsto\omega_i\) pushes the product measure to its \(i\)-th factor \(\mathcal N(0,1)\). The product marginal theorem proves both almost-everywhere measurability and the pushforward identity, packaged as \texttt{Has\allowbreak{}Law}.

\subsection{\texorpdfstring{coordinate\allowbreak{}Lp}{coordinateLp}}


A standard Gaussian has finite second moment, so the coordinate projection belongs to \(L^2(Q)\). Apply its \texttt{Mem\allowbreak{}Lp.\allowbreak{}to\allowbreak{}Lp} construction to obtain the quotient class \(Z_i\); the definition stores the original coordinate function together with its integrability proof.

\subsection{\texorpdfstring{coordinate\allowbreak{}Lp coe}{coordinateLp coe}}


The representative selected for \(Z_i\) agrees with \(\omega_i\) almost everywhere by the defining property of \texttt{to\allowbreak{}Lp}. This equality is the bridge used to transfer coordinate laws, means, and independence calculations to the \(L^2\) classes.

\subsection{\texorpdfstring{coordinates independent}{coordinates independent}}


The coordinate projections of an infinite product measure are jointly independent. Invoke that product theorem for the scalar Gaussian factors. This is independence of every finite set of coordinates, which is precisely what finite Gaussian linear combinations require.

\subsection{\texorpdfstring{coordinate mean}{coordinate mean}}


Change variables through the coordinate law and use the centered Gaussian mean: \(\int\omega_i\,dQ=\int x\,d\mathcal N(0,1)=0\). Gaussian integrability makes this an ordinary expectation identity.

\subsection{\texorpdfstring{coordinate\allowbreak{}Lp inner}{coordinateLp inner}}


For \(i=j\), the \(L^2\) inner product is \(\mathbb EZ_i^2=1\). For \(i\ne j\), independence factors \(\mathbb EZ_iZ_j=(\mathbb EZ_i)(\mathbb EZ_j)=0\). Almost-everywhere representative identities move between \(Z_i\) and the product coordinates. Thus \(\langle Z_i,Z_j\rangle=\mathbf1_{i=j}\).

\subsection{\texorpdfstring{coordinates orthonormal}{coordinates orthonormal}}


The preceding inner-product formula is the orthonormality criterion for a family in a real Hilbert space. Lean converts the diagonal value one and off-diagonal value zero into \texttt{Orthonormal \ensuremath{\mathbb R} coordinate\allowbreak{}Lp}, the input for the orthogonal-series extension.

\subsection{\texorpdfstring{source is\allowbreak{}Probability\allowbreak{}Measure}{source isProbabilityMeasure}}


Each scalar Gaussian factor is a probability measure. The product-probability instance therefore gives \(Q(\Omega)=1\). This instance is what permits constant functions and integrable functions to be treated in the usual probability-space \(L^2\) theory below.

\subsection{\texorpdfstring{synthesis}{synthesis}}


Extend the orthonormal family \(Z_i\) to a linear isometry \(S:\ell^2(\iota)\to L^2(Q)\). The orthogonal-family extension theorem constructs the limit of finite sums and preserves squared norm, \(\|S(a)\|_2^2=\sum_i a_i^2\). It constructs an \(L^2\) limit rather than assuming pointwise convergence of a random series.

\subsection{\texorpdfstring{synthesis has\allowbreak{}Sum}{synthesis hasSum}}


The synthesis extension theorem states \(\sum_i a_iZ_i=S(a)\), with convergence of finite partial sums in \(L^2\). In Lean \texttt{Has\allowbreak{}Sum} is convergence along finite subsets, so the identity does not require choosing an enumeration of the basis index set.

\subsection{\texorpdfstring{isonormal}{isonormal}}


For a Hilbert basis \(b\), compose its coordinate isometry \(b.\mathrm{repr}:H\to\ell^2(\iota)\) with \texttt{synthesis}. The resulting linear isometry \(I_b\) sends \(x\) to the Gaussian series with coefficients \(\langle b_i,x\rangle\), and preserves all inner products.

\subsection{\texorpdfstring{isonormal has\allowbreak{}Sum}{isonormal hasSum}}


Apply \texttt{synthesis\_\allowbreak{}has\allowbreak{}Sum} to the coefficient family \(b.\mathrm{repr}(x)\). Unfolding the composition gives \(I_b(x)=\sum_i\langle b_i,x\rangle Z_i\) in \(L^2\). This is the exact series represented by the bundled isonormal isometry.

\subsection{\texorpdfstring{integral eq inner one}{integral eq inner one}}


On a probability space, the constant one is in \(L^2\). The definition of the real \(L^2\) inner product gives \(\mathbb EX=\langle1,X\rangle\), after replacing the selected constant representative by one almost everywhere. This exhibits expectation as a continuous functional.

\subsection{\texorpdfstring{coordinate\allowbreak{}Lp mean}{coordinateLp mean}}


Use \(Z_i=\omega_i\) almost everywhere to replace the \(L^2\) representative inside the integral. The already established centered coordinate expectation then gives \(\mathbb EZ_i=0\).

\subsection{\texorpdfstring{synthesis mean}{synthesis mean}}


Each finite partial sum \(\sum_{i\in F}a_iZ_i\) has mean zero by linearity and the coordinate means. Pair the \(L^2\)-convergent series with the constant one; continuity of the inner product passes the zero expectations to the limit, giving \(\mathbb ES(a)=0\).

\subsection{\texorpdfstring{isonormal mean}{isonormal mean}}


The isonormal variable is \texttt{synthesis} applied to the Hilbert-basis coefficients of \(x\). Specialize the synthesis mean theorem to these coefficients to get \(\mathbb EI_b(x)=0\), without redoing the infinite-series limit.

\subsection{\texorpdfstring{variance eq norm sq}{variance eq norm sq}}


For an \(L^2\) random variable with mean zero, the variance formula reduces to \(\operatorname{Var}(X)=\int X^2dQ\). The \(L^2\) norm identity identifies this with \(\|X\|_2^2\). The explicit mean premise is used to remove the centering term.

\subsection{\texorpdfstring{partial sum coe}{partial sum coe}}


Induct over the finite set. The zero sum has zero representative almost everywhere; addition and scalar multiplication of \(L^2\) classes agree almost everywhere with the corresponding pointwise operations. Intersect the finitely many full-measure sets to obtain \(\sum_{i\in F}a_iZ_i=\sum_{i\in F}a_i\omega_i\) almost everywhere.

\subsection{\texorpdfstring{partial sum has\allowbreak{}Gaussian\allowbreak{}Law}{partial sum hasGaussianLaw}}


Independent scalar standard Gaussians remain Gaussian after a finite linear combination. Apply that finite-sum Gaussian theorem to \(a_i\omega_i\), and transfer the law across \texttt{partial\_\allowbreak{}sum\_\allowbreak{}coe}. This proves Gaussianity of the actual \(L^2\) partial sum, not merely of a formally similar expression.

\subsection{\texorpdfstring{partial sum mean}{partial sum mean}}


The integral of a finite \(L^2\) sum is the sum of its integrals, because every summand is integrable. Each coordinate mean is zero, hence \(\mathbb E\sum_{i\in F}a_iZ_i=0\). This supplies the centering parameter needed to identify its Gaussian law.

\subsection{\texorpdfstring{partial sum has\allowbreak{}Law}{partial sum hasLaw}}


A scalar Gaussian law is determined by its mean and variance. The finite sum is Gaussian and centered; its variance is its squared \(L^2\) norm. Therefore its law is \(\mathcal N(0,\|\sum_{i\in F}a_iZ_i\|_2^2)\), with the variance converted to a nonnegative real.

\subsection{\texorpdfstring{synthesis has\allowbreak{}Law}{synthesis hasLaw}}


Finite partial sums converge in \(L^2\) to \(S(a)\), and each has the centered Gaussian law with its own squared norm as variance. Apply \texttt{has\allowbreak{}Law\_\allowbreak{}gaussian\allowbreak{}Real\_\allowbreak{}of\_\allowbreak{}tendsto\_\allowbreak{}Lp}; the isometry then rewrites the limiting variance as \(\|a\|_{\ell^2}^2\). This proves Gaussianity of the completed infinite series.

\subsection{\texorpdfstring{isonormal has\allowbreak{}Law}{isonormal hasLaw}}


Apply the synthesis law theorem to \(b.\mathrm{repr}(x)\). The Hilbert-basis coordinate map preserves norm, so the variance becomes \(\|x\|^2\). Hence \(I_b(x)\sim\mathcal N(0,\|x\|^2)\), including the degenerate law for \(x=0\).

\subsection{\texorpdfstring{isonormal has\allowbreak{}Gaussian\allowbreak{}Law}{isonormal hasGaussianLaw}}


The exact law theorem identifies \(I_b(x)\) with a scalar Gaussian probability. Forgetting its explicit mean and variance parameters yields the bundled \texttt{Has\allowbreak{}Gaussian\allowbreak{}Law} predicate used by the finite-dimensional Gaussian-process construction.

\subsection{\texorpdfstring{isonormal covariance}{isonormal covariance}}


Both variables are centered, so their covariance is \(\mathbb E[I_b(x)I_b(y)]\), the \(L^2\) inner product. The linear isometry preserves that inner product, giving \(\operatorname{Cov}(I_b(x),I_b(y))=\langle x,y\rangle_H\). Lean checks integrability through their \(L^2\) membership.

\section{\texorpdfstring{Scalar Gaussian linear maps give jointly Gaussian processes}{Scalar Gaussian linear maps give jointly Gaussian processes}}
\label{proof-module-UnadjustedLangevin-GaussianLinearProcess-lean}
{\small\color{gray}Source module: \texttt{Unadjusted\allowbreak{}Langevin/\allowbreak{}Gaussian\allowbreak{}Linear\allowbreak{}Process.\allowbreak{}lean}.\par}


If a linear map \(T:H\to L^2(P;\mathbb R)\) sends every input to a Gaussian variable, all finite collections of its evaluations are jointly Gaussian.

\subsection{\texorpdfstring{has\allowbreak{}Gaussian\allowbreak{}Law pi of linear\allowbreak{}Map Lp}{hasGaussianLaw pi of linearMap Lp}}


Use the finite-dimensional Gaussian criterion: every scalar linear functional of \((T(f_i))_i\) must be Gaussian. A finite linear combination equals \(T(\sum_i a_if_i)\) as an \(L^2\) class by linearity, hence agrees almost everywhere as a random variable. The assumed Gaussianity of every \(T(x)\) proves the criterion.

\subsection{\texorpdfstring{is\allowbreak{}Gaussian\allowbreak{}Process of linear\allowbreak{}Map Lp}{isGaussianProcess of linearMap Lp}}


For each finite set of times, apply the preceding joint-Gaussian theorem to the corresponding finite family of vectors \(f(t)\). This verifies the finite-dimensional-law definition of \texttt{Is\allowbreak{}Gaussian\allowbreak{}Process}. The index set may be uncountable; this statement alone asserts no sample-path continuity.

\section{\texorpdfstring{An explicit Gaussian construction of pre-\allowbreak{}Brownian coordinates}{An explicit Gaussian construction of pre-Brownian coordinates}}
\label{proof-module-UnadjustedLangevin-CanonicalBrownian-lean}
{\small\color{gray}Source module: \texttt{Unadjusted\allowbreak{}Langevin/\allowbreak{}Canonical\allowbreak{}Brownian.\allowbreak{}lean}.\par}


The isonormal family on the time Hilbert space constructs a Gaussian process with covariance \(\min(s,t)\). Independent product copies supply the Brownian coordinates required by the canonical ULA theorem. The construction proves pre-Brownian laws; continuous sample paths are not asserted or needed, because continuity is established in \(L^2\).

\subsection{\texorpdfstring{Time\allowbreak{}Hilbert}{TimeHilbert}}


Take the deterministic Hilbert space \(L^2(\mathbb R,dt)\). Its vectors will encode time intervals before an isonormal map turns their inner products into covariances. It is distinct from the random-variable \(L^2\) space on the Gaussian source.

\subsection{\texorpdfstring{time\allowbreak{}Indicator}{timeIndicator}}


Let \(e_t=\mathbf1_{(0,t]}\) as a class in \(L^2(dt)\), for \(t\ge0\). Lean constructs it with \texttt{indicator\allowbreak{}Const\allowbreak{}Lp} and the finite Lebesgue measure of \texttt{Ioc 0 t}. Endpoints do not affect this Lebesgue \(L^2\) class, but this half-open convention makes interval intersection identities direct.

\subsection{\texorpdfstring{time\allowbreak{}Indicator inner}{timeIndicator inner}}


The product of the two indicators is the indicator of \([0,\min(s,t)]\) almost everywhere. Integrating it gives \(\langle e_s,e_t\rangle=\min(s,t)\). The proof evaluates the Lebesgue measure of that finite interval, producing the Brownian covariance before introducing randomness.

\subsection{\texorpdfstring{basis\allowbreak{}Index}{basisIndex}}


Choose the index type supplied by a Hilbert basis of the time Hilbert space. It is a set of Hilbert-space vectors used as a type; no explicit countable enumeration is required by the orthogonal-family construction.

\subsection{\texorpdfstring{basis}{basis}}


Select a Hilbert basis for \texttt{Time\allowbreak{}Hilbert} by Hilbert-space basis existence. This basis supplies the coordinate isometry for the explicit isonormal Gaussian family; all subsequent covariance identities follow from isometry and do not depend on which basis was selected.

\subsection{\texorpdfstring{Sample}{Sample}}


A sample is a real family indexed by the chosen time-Hilbert basis. This concrete product space stores the independent standard Gaussian coefficients used to synthesize all Brownian times.

\subsection{\texorpdfstring{probability}{probability}}


Use \texttt{Gaussian\allowbreak{}Hilbert.\allowbreak{}source} on the basis index type, namely the product of standard scalar Gaussians. Its established probability instance provides normalization and independent coordinate laws for the Brownian construction.

\subsection{\texorpdfstring{instance at 53}{instance at 53}}


The scalar source is the Gaussian coordinate product probability. Reuse \texttt{Gaussian\allowbreak{}Hilbert.\allowbreak{}source\_\allowbreak{}is\allowbreak{}Probability\allowbreak{}Measure} to provide normalization for the chosen basis index. This instance enables the \(L^2\) expectation and covariance theorems for the constructed Brownian process.

\subsection{\texorpdfstring{value\allowbreak{}Lp}{valueLp}}


Apply the isonormal isometry to \(e_t\), obtaining \(B_t=I(e_t)\) as an \(L^2\) random variable. The construction immediately gives a centered Gaussian with variance \(\|e_t\|^2=t\); joint Gaussianity follows from linearity.

\subsection{\texorpdfstring{process}{process}}


At each time select the strongly measurable representative of \texttt{value\allowbreak{}Lp}. This defines an actual function \(B_t(\omega)\) for use in \texttt{Has\allowbreak{}Law} and finite-dimensional distributions. Equalities of \(L^2\) classes still translate to almost-everywhere equalities of these representatives.

\subsection{\texorpdfstring{process measurable}{process measurable}}


The representative selected for an \(L^2\) class is measurable in this setting. Apply the representative measurability theorem to \texttt{value\allowbreak{}Lp t}, obtaining measurability of the actual time-\(t\) random variable.

\subsection{\texorpdfstring{process mean}{process mean}}


The isonormal mean theorem gives \(\mathbb EI(e_t)=0\). Since \texttt{process t} is precisely the selected representative of this class, this is the expectation of the displayed random function \(B_t\).

\subsection{\texorpdfstring{process covariance}{process covariance}}


The isonormal covariance formula gives \(\operatorname{Cov}(I(e_s),I(e_t))=\langle e_s,e_t\rangle\). Substitute \texttt{time\allowbreak{}Indicator\_\allowbreak{}inner} to obtain \(\min(s,t)\). This calculation establishes the Brownian covariance for the concrete source and representatives.

\subsection{\texorpdfstring{is\allowbreak{}Gaussian\allowbreak{}Process}{isGaussianProcess}}


The isonormal isometry is linear and every output is scalar Gaussian. Apply \texttt{is\allowbreak{}Gaussian\allowbreak{}Process\_\allowbreak{}of\_\allowbreak{}linear\allowbreak{}Map\_\allowbreak{}Lp} to the family \(t\mapsto e_t\). Every finite collection of the selected random variables is therefore jointly Gaussian.

\subsection{\texorpdfstring{is\allowbreak{}Pre\allowbreak{}Brownian}{isPreBrownian}}


Combine centered Gaussian finite-dimensional laws with covariance \(\min(s,t)\). The Gaussian-process Brownian characterization yields \texttt{Is\allowbreak{}Pre\allowbreak{}Brownian\allowbreak{}Real}. This proves Brownian finite-dimensional laws; the prefix “pre” records that sample-path continuity is not part of this predicate.

\subsection{\texorpdfstring{Vector\allowbreak{}Sample}{VectorSample}}


For a coordinate index type \(I\), take one scalar Brownian sample for each \(i\in I\). This is the product sample type \(I\to\mathrm{Sample}\), allowing independent entire coordinate processes rather than only independent values at one time.

\subsection{\texorpdfstring{vector\allowbreak{}Probability}{vectorProbability}}


Equip the vector sample space with the infinite product of the scalar Brownian source probabilities. The coordinate processes inherit independence from this actual product construction.

\subsection{\texorpdfstring{coordinate}{coordinate}}


Define the \(i\)-th coordinate process by \(B_i(t,\omega)=B(t,\omega_i)\). It first projects to the \(i\)-th scalar sample, then evaluates the already constructed scalar process at \(t\).

\subsection{\texorpdfstring{coordinate is\allowbreak{}Pre\allowbreak{}Brownian}{coordinate isPreBrownian}}


The product coordinate projection has the scalar Brownian source as its law. Pull the scalar pre-Brownian process back along this projection; its finite-dimensional distributions, mean, and covariance are preserved. Thus each vector coordinate has the established scalar Brownian laws.

\subsection{\texorpdfstring{coordinate paths independent}{coordinate paths independent}}


The product coordinate projections are independent as sample-valued maps. Compose each projection with the map sending a scalar sample to its entire time-indexed process. Independence is preserved under these coordinatewise functions, giving independence of the whole coordinate paths.

\subsection{\texorpdfstring{coordinate grid paths independent}{coordinate grid paths independent}}


Restrict each independent coordinate path to the deterministic grid \(nh\). This is a coordinatewise function of the path, so independence is preserved. The result supplies exactly the countable grid-path hypothesis used in the innovation theorems.

\subsection{\texorpdfstring{instance at 89}{instance at 89}}


Each scalar Brownian sample source has mass one. The infinite product over vector coordinates therefore also has mass one. This instance concerns the actual vector sample measure, not an assumed normalization of an abstract process law.

\section{\texorpdfstring{Stability and uniqueness for additive-\allowbreak{}noise solutions}{Stability and uniqueness for additive-noise solutions}}
\label{proof-module-UnadjustedLangevin-LangevinUniqueness-lean}
{\small\color{gray}Source module: \texttt{Unadjusted\allowbreak{}Langevin/\allowbreak{}Langevin\allowbreak{}Uniqueness.\allowbreak{}lean}.\par}


Two solutions driven by the same continuous additive noise satisfy \(\|Y_t-Z_t\|\le e^{Lt}\|x-y\|\) on the prescribed interval. Equality of initial values gives uniqueness. Only global Lipschitzness is used here; the sharper contractive estimate uses strong convexity elsewhere.

\subsection{\texorpdfstring{langevin integral stability}{langevin integral stability}}


Subtract the two integral equations so that their common additive input cancels. The integral norm inequality and Lipschitz drift give \(u(t)\le\|x-y\|+L\int_0^t u(s)ds\), where \(u(t)=\|Y(t)-Z(t)\|\) is continuous. Grönwall's inequality yields \(u(t)\le\|x-y\|e^{Lt}\). All time integrals are justified by continuity on the compact horizon.

\subsection{\texorpdfstring{langevin integral unique}{langevin integral unique}}


Set the two initial states equal in the stability estimate. The right side becomes \(0e^{Lt}=0\), forcing \(\|Y(t)-Z(t)\|=0\) for each \(t\in[0,h]\). Norm-zero equivalence gives equality of the vectors, packaged in Lean as \texttt{Eq\allowbreak{}On} on the interval.

\subsection{\texorpdfstring{lp langevin integral stability}{lp langevin integral stability}}


Specialize Banach-space stability to the actual lifted gradient on \(L^2(P;K)\). Its Lipschitz theorem supplies the required constant, so \(\|Y(t)-Z(t)\|_2\le e^{Lt}\|Y(0)-Z(0)\|_2\). This compares solutions with a common \(L^2\)-noise input and possibly different initial random variables.

\subsection{\texorpdfstring{lp langevin integral unique}{lp langevin integral unique}}


Apply the same-input uniqueness theorem in \(L^2(P;K)\). Equality here is equality of \(L^2\) classes for each time, hence almost-everywhere equality of representatives at each fixed time. The conclusion does not by itself select one exceptional null set valid simultaneously at every real time.

\section{\texorpdfstring{Identify chosen reference segments with one global solution}{Identify chosen reference segments with one global solution}}
\label{proof-module-UnadjustedLangevin-BrownianSegmentsGlobal-lean}
{\small\color{gray}Source module: \texttt{Unadjusted\allowbreak{}Langevin/\allowbreak{}Brownian\allowbreak{}Segments\allowbreak{}Global.\allowbreak{}lean}.\par}


The segment construction is compatible with the unique global integral solution. This lets stationarity of a global Brownian Langevin solution prove stationarity of every chosen reference segment.

\subsection{\texorpdfstring{brownian\allowbreak{}Langevin\allowbreak{}Noise increment}{brownianLangevinNoise increment}}


Expand the two increments from zero and subtract them. Their common \([B_0]\) term cancels, leaving \(\sqrt2([B_{a+t}]-[B_a])=N_a(t)\). This is an algebraic cocycle identity in \(L^2\), valid for the real-time extension as defined.

\subsection{\texorpdfstring{lp brownian integral shift}{lp brownian integral shift}}


Subtract the global integral equations at \(a+t\) and \(a\). Integral additivity gives the drift integral over \([a,a+t]\), and time translation turns it into \(\int_0^t\nabla U(Y_{a+s})ds\). The noise cocycle gives \(N_a(t)\), yielding the local segment equation with initial state \(Y_a\).

\subsection{\texorpdfstring{brownian\allowbreak{}Segment\allowbreak{}From eq global}{brownianSegmentFrom eq global}}


The shifted global path and the selected local segment are continuous solutions of the same \(L^2\) integral equation with the same initial value and noise input. Apply \(L^2\) uniqueness to identify them on \([0,h]\). The horizon premise ensures the shifted interval lies inside the global solution interval.

\subsection{\texorpdfstring{brownian\allowbreak{}Knots eq global}{brownianKnots eq global}}


Induct on the knot index. At zero, the global equation and zero-starting noise give the prescribed initial state. If the current knot equals \(Y_{nh}\), segment/\allowbreak{}global uniqueness makes its endpoint equal \(Y_{(n+1)h}\), establishing the next knot identity.

\subsection{\texorpdfstring{brownian\allowbreak{}Segments eq global}{brownianSegments eq global}}


The knot/\allowbreak{}global theorem supplies the matching initial value for the \(n\)-th segment. Apply \texttt{brownian\allowbreak{}Segment\allowbreak{}From\_\allowbreak{}eq\_\allowbreak{}global} to conclude \(Y^{(n)}_t=Y_{nh+t}\) on \([0,h]\). This is a deterministic identification of the constructions, with no stationarity premise.

\section{\texorpdfstring{Dalalyan’s theorem for the actual Gaussian ULA transition}{Dalalyan’s theorem for the actual Gaussian ULA transition}}
\label{proof-module-UnadjustedLangevin-CanonicalULA-lean}
{\small\color{gray}Source module: \texttt{Unadjusted\allowbreak{}Langevin/\allowbreak{}Canonical\allowbreak{}ULA.\allowbreak{}lean}.\par}


Let \(U:\mathbb R^d\to\mathbb R\) be differentiable and \(m\)-strongly convex with \(m>0\), with \(L\)-Lipschitz gradient. The canonical Brownian source derives all stochastic reference obligations and proves the Wasserstein bound for the iterated Gaussian transition \(P_h\), with finite initial \(W_2\).

\subsection{\texorpdfstring{Identify the canonical output law}{Identify the canonical output law}}


The Brownian source provides jointly independent standard Gaussian innovations; the product extension makes initialization independent of the entire sequence. The realization/\allowbreak{}Markov-law theorem identifies its output with \(P_h^k\mu\). The product output-law identity removes the chosen finite coupling from this law equality.

\subsection{\texorpdfstring{product\allowbreak{}Reference of independent brownian}{productReference of independent brownian}}


On each finite-coupling extension, pull back the independent Brownian coordinates and use the product independence theorems for the reference initialization. Construct a global \(L^2\) integral solution over the required finite horizon.



The proved MALA/\allowbreak{}Euler stationarity theorem gives Gibbs law for that global solution at every needed absolute time. Segment/\allowbreak{}global uniqueness identifies each recursively constructed reference segment with that solution. Thus the segment-stationarity premise of \texttt{product\allowbreak{}Reference\_\allowbreak{}of\_\allowbreak{}constructed\_\allowbreak{}brownian} is discharged, filling every reference field from the stated source and potential assumptions.

\subsection{\texorpdfstring{canonical\allowbreak{}Euler\allowbreak{}Noise}{canonicalEulerNoise}}


Specialize the source Euler noise to the explicitly constructed independent coordinate Brownian source. This abbreviation fixes one concrete implementation of the standard Gaussian innovations for the unconditional law theorem.

\subsection{\texorpdfstring{canonical\allowbreak{}Product\allowbreak{}Reference}{canonicalProductReference}}


The canonical Brownian construction supplies the coordinate Brownian laws and independence of their grid paths for every positive step. Apply \texttt{product\allowbreak{}Reference\_\allowbreak{}of\_\allowbreak{}independent\_\allowbreak{}brownian}. No stochastic reference structure or stationary-law assertion remains an external premise.

\subsection{\texorpdfstring{canonical\allowbreak{}Output eq ula\allowbreak{}Law\allowbreak{}Step}{canonicalOutput eq ulaLawStep}}


Finite initial transport provides a finite-cost coupling, and Gibbs has a finite second moment. Apply the general Brownian output-law identity on that extension with the canonical independent Brownian source. It proves the canonical product implementation has exactly the ordinary ULA law \(P_h^k\mu\).

\subsection{\texorpdfstring{dalalyan output finite}{dalalyan output finite}}


The canonical product reference supplies a finite-cost output coupling for every iteration whenever initial transport is finite. Its output law is \(P_h^k\mu\). Substituting this law proves \(W_2(P_h^k\mu,\pi_U)<\infty\), legitimizing the real Wasserstein notation used in the final estimates.

\subsection{\texorpdfstring{dalalyan}{dalalyan}}


Construct \texttt{canonical\allowbreak{}Product\allowbreak{}Reference} from differentiability, \(m\)-strong convexity, and \(L\)-Lipschitz gradient. Nontriviality of the finite coordinate space proves \(m\le L\); the finite-coupling product theorem gives the small-step bound with bias \((33/20)(L/m)\sqrt{hd}\).



Finally replace the canonical product output by \(P_h^k\mu\) using \texttt{canonical\allowbreak{}Output\_\allowbreak{}eq\_\allowbreak{}ula\allowbreak{}Law\allowbreak{}Step}. The resulting inequality is about the actual Gaussian ULA laws, and output finiteness has already been proved. The independent Brownian construction and stationary reference are internal parts of this proof.

\subsection{\texorpdfstring{dalalyan large step}{dalalyan large step}}


Use the same fully constructed canonical reference in the large-step product theorem. The step assumptions give \(0\le Lh-1<1\) and \(2-Lh>0\), so the bias is \((33/20)Lh\sqrt{hd}/(2-Lh)\). Rewrite its output law as the actual iterated Gaussian ULA law.

\section{\texorpdfstring{Dalalyan bounds for an implemented L\ensuremath{{}^{2}} Gaussian recursion}{Dalalyan bounds for an implemented L² Gaussian recursion}}
\label{proof-module-UnadjustedLangevin-GaussianULA-lean}
{\small\color{gray}Source module: \texttt{Unadjusted\allowbreak{}Langevin/\allowbreak{}Gaussian\allowbreak{}ULA.\allowbreak{}lean}.\par}


This endpoint applies the canonical-law theorem to an actual chain on a supplied probability space. Its inputs are the original potential assumptions, a square-integrable initialization, independent standard Gaussian scalar innovations and independence of the initial variable from the innovation sequence.

\subsection{\texorpdfstring{initial\allowbreak{}Law probability}{initialLaw probability}}


The representative of the \(L^2\) initialization is measurable. Its pushforward of a probability measure therefore has total mass one, giving the probability instance required by the law-level canonical ULA theorem.

\subsection{\texorpdfstring{gaussian\allowbreak{}ULA initial transport finite}{gaussianULA initial transport finite}}


The actual initialization is in \(L^2\), so its law has a finite second moment. Strong convexity proves the same for Gibbs and proves its normalization. The independent product coupling of these two laws has finite quadratic cost, hence their Wasserstein distance is finite.

\subsection{\texorpdfstring{dalalyan of gaussian innovations}{dalalyan of gaussian innovations}}


Assemble the jointly independent standard scalar Gaussian coordinates into iid multivariate standard innovations. Initialization independence and the \(L^2\) representative theorem identify the implementation's law with \(P_h^k(P_{X_0})\). Derive finite initial transport from its moments, then apply the unconditional canonical small-step theorem and rewrite the law back to the implementation.

\subsection{\texorpdfstring{dalalyan large step of gaussian innovations}{dalalyan large step of gaussian innovations}}


Use the same implementation/\allowbreak{}Markov-law identification and derived finite initial transport, then apply the canonical large-step theorem. The numerical implementation therefore inherits contraction factor \(Lh-1\) and bias \((33/20)Lh\sqrt{hd}/(2-Lh)\) throughout the stated stable large-step regime.

\section{\texorpdfstring{The Langevin generator annihilates Gibbs averages}{The Langevin generator annihilates Gibbs averages}}
\label{proof-module-UnadjustedLangevin-Generator-lean}
{\small\color{gray}Source module: \texttt{Unadjusted\allowbreak{}Langevin/\allowbreak{}Generator.\allowbreak{}lean}.\par}


For compactly supported \(C^2\) test functions, \(\mathcal L\varphi=\Delta\varphi-\langle\nabla U,\nabla\varphi\rangle\) has zero Gibbs mean. These identities establish infinitesimal invariance; they alone do not assert stationarity of a stochastic process.

\subsection{\texorpdfstring{directional\allowbreak{}Generator}{directionalGenerator}}


For a fixed direction \(v\), define \(L_v\varphi(x)=D(D\varphi(\cdot)[v])(x)[v]-\langle\nabla U(x),v\rangle D\varphi(x)[v]\). The first term is a directional second derivative of the test; the second is the directional gradient drift. No Hessian of \(U\) is part of this definition.

\subsection{\texorpdfstring{langevin\allowbreak{}Generator}{langevinGenerator}}


Choose the standard orthonormal basis and set \(L\varphi=\sum_iL_{e_i}\varphi\). The finite sum makes the construction meaningful in dimension zero and gives a direct path to directional integration by parts. A subsequent identity proves that this is the usual basis-independent \(\Delta\varphi-\langle\nabla U,\nabla\varphi\rangle\).

\subsection{\texorpdfstring{integral gibbs eq weighted}{integral gibbs eq weighted}}


Unfold the normalized density and use the integral change-of-density theorem: \(\int f\,d\pi_U=(Z_U^{-1})_{\mathbb R}\int e^{-U(x)}f(x)dx\). Continuity of \(U\) makes the density measurable. The factor is written \texttt{(gibbs\allowbreak{}Partition U)\ensuremath{{}^{-1}}.\allowbreak{}to\allowbreak{}Real} because the partition was defined in extended nonnegative scalars.

\subsection{\texorpdfstring{integrable gibbs of continuous has\allowbreak{}Compact\allowbreak{}Support}{integrable gibbs of continuous hasCompactSupport}}


A continuous compactly supported \(f\) gives a continuous compactly supported product \(e^{-U}f\), hence a Lebesgue-integrable function. The with-density integrability criterion and multiplication by the real normalization factor transfer this fact to \(\pi_U\). This explicitly supplies the integrability needed by subsequent linearity arguments.

\subsection{\texorpdfstring{fderiv exp neg apply}{fderiv exp neg apply}}


Apply the chain rule to \(x\mapsto e^{-U(x)}\): \(D(e^{-U})(x)[v]=-e^{-U(x)}DU(x)[v]\). The real Hilbert gradient identity rewrites \(DU(x)[v]\) as \(\langle\nabla U(x),v\rangle\). Only differentiability of \(U\) is used.

\subsection{\texorpdfstring{integral weighted directional\allowbreak{}Generator eq zero}{integral weighted directionalGenerator eq zero}}


Differentiate \(e^{-U(x)}D\varphi(x)[v]\) in direction \(v\). The product rule gives exactly \(e^{-U(x)}L_v\varphi(x)\). Since derivatives of a compactly supported \(C^2\) test remain compactly supported, this weighted first derivative is integrable and has zero integral by integration by parts. The proof establishes continuity and compact support before invoking that theorem.

\subsection{\texorpdfstring{integral directional\allowbreak{}Generator gibbs eq zero}{integral directionalGenerator gibbs eq zero}}


Apply the weighted-to-Gibbs integral formula to \(L_v\varphi\). Its weighted Lebesgue integral is zero, so multiplication by the normalization factor still gives zero. This is infinitesimal invariance for a single direction and does not yet identify the law of a time-evolving process.

\subsection{\texorpdfstring{continuous directional\allowbreak{}Generator}{continuous directionalGenerator}}


The \(C^2\) hypothesis makes the first and second derivatives of \(\varphi\) continuous. The gradient of \(U\) is continuous by its Lipschitz hypothesis. Evaluating continuous linear maps in the fixed direction \(v\), multiplying the scalar factors, and subtracting therefore gives continuity of \(L_v\varphi\).

\subsection{\texorpdfstring{has\allowbreak{}Compact\allowbreak{}Support directional\allowbreak{}Generator}{hasCompactSupport directionalGenerator}}


Outside the compact support of \(\varphi\), the test is locally zero, so its first and second derivatives vanish there. Both terms of \(L_v\varphi\) therefore vanish. The support inclusion puts the support of the generator inside the same compact set; no support property of \(U\) is required.

\subsection{\texorpdfstring{integrable directional\allowbreak{}Generator gibbs\allowbreak{}Measure}{integrable directionalGenerator gibbsMeasure}}


Combine continuity of \(L_v\varphi\) with preservation of compact support. The established compact-support Gibbs integrability theorem then gives absolute integrability. This permits ordinary integral algebra, even though the ambient generator is not bounded uniformly for arbitrary tests.

\subsection{\texorpdfstring{integrable langevin\allowbreak{}Generator gibbs\allowbreak{}Measure}{integrable langevinGenerator gibbsMeasure}}


The full generator is the finite sum of its directional pieces. Each piece is integrable under \(\pi_U\), so the finite-sum integrability theorem applies. Finiteness of the basis is the reason no exchange of an infinite series and integral is required.

\subsection{\texorpdfstring{integral weighted langevin\allowbreak{}Generator eq zero}{integral weighted langevinGenerator eq zero}}


Expand the finite basis sum defining \(L\varphi\) and distribute the Gibbs weight over it. Compact support supplies integrability of every summand, allowing \(\int\sum_i e^{-U}L_{e_i}\varphi=\sum_i\int e^{-U}L_{e_i}\varphi\). Every directional integral vanishes, hence so does their sum.

\subsection{\texorpdfstring{integral langevin\allowbreak{}Generator gibbs eq zero}{integral langevinGenerator gibbs eq zero}}


Normalization changes \(\int L\varphi\,d\pi_U\) into the real inverse partition times its weighted Lebesgue integral. The latter was proved to be zero. This establishes the full Gibbs infinitesimal invariance identity for compactly supported \(C^2\) tests.

\subsection{\texorpdfstring{langevin\allowbreak{}Generator eq laplacian sub inner}{langevinGenerator eq laplacian sub inner}}


The sum of the directional second derivatives over an orthonormal basis is the Laplacian. For the drift sum, \(D\varphi(x)[e_i]=\langle\nabla\varphi(x),e_i\rangle\); Parseval turns the sum of products with \(\langle\nabla U(x),e_i\rangle\) into \(\langle\nabla U(x),\nabla\varphi(x)\rangle\). The proof rewrites the basis definition into this conventional formula pointwise.

\subsection{\texorpdfstring{integral laplacian sub gradient gibbs eq zero}{integral laplacian sub gradient gibbs eq zero}}


Replace the conventional expression \(\Delta\varphi-\langle\nabla U,\nabla\varphi\rangle\) by the equal basis-defined generator inside the integral. The established zero-generator integral then closes the result. The rewrite makes the displayed mathematical identity refer to the same function that the preceding Lean proof integrated.

\section{\texorpdfstring{Physical-\allowbreak{}time generator and Gibbs integration identity}{Physical-time generator and Gibbs integration identity}}
\label{proof-module-IPLA-Generator-lean}
{\small\color{gray}Source module: \texttt{IPLA/\allowbreak{}Generator.\allowbreak{}lean}.\par}


For scaled state \(z\), physical IPLA time is \(N\) times the ordinary lifted Langevin time. The generator is consequently \(N^{-1}(\Delta-\nabla V_N\cdot\nabla)\).

\subsection{\texorpdfstring{Identify the generator on a twice differentiable test function}{Identify the generator on a twice differentiable test function}}


Define \(\mathcal A_N\varphi=N^{-1}\mathcal L_{V_N}\varphi\), where \(\mathcal L_V\) is the shared Langevin generator. For \(C^2\) test functions the shared coordinate definition equals the Laplacian minus the gradient inner product, hence

\[\mathcal A_N\varphi(z)=\frac1N\left(\Delta\varphi(z)-\langle\nabla V_N(z),\nabla\varphi(z)\rangle\right).\]

This is a direct identity of operators on the scaled state; no invariant-measure conclusion is built into its definition.




\textbf{Correspondence with the displayed declarations.}


\begin{itemize}

\item \texttt{diffusion\allowbreak{}Generator}: Generator in the physical IPLA time variable, acting on the rescaled state.
\item \texttt{diffusion\allowbreak{}Generator\_\allowbreak{}eq}: The exact Laplacian and drift formula for the IPLA generator.
\end{itemize}
\subsection{\texorpdfstring{Derive integrability and infinitesimal invariance}{Derive integrability and infinitesimal invariance}}


For compactly supported \(C^2\) \(\varphi\), the shared Gibbs integration-by-parts theorem proves that \(\mathcal L_{V_N}\varphi\) is integrable and has Gibbs integral zero. Its potential hypotheses follow from differentiability and Lipschitzness of the joint gradient via the geometry module. Multiplication by the constant \(N^{-1}\) preserves integrability and gives

\[\int\mathcal A_N\varphi\,d\Pi_N=\frac1N\int\mathcal L_{V_N}\varphi\,d\Pi_N=0.\]

The theorem also returns the probability property, derived from positive strong convexity. Explicit integrability rules out a zero caused by the default integral value of a nonintegrable function. Full time-evolution stationarity is supplied separately by the shared Langevin construction.




For a \(C^2\) compactly supported test function, the shared Gibbs integration identity gives \(\int[\Delta\varphi-\langle\nabla V_N,\nabla\varphi\rangle]d\Pi_N=0\). Physical IPLA time multiplies this generator by \(N^{-1}\), so linearity of the integral preserves zero. The theorem also proves integrability of the generator before asserting this identity, using the compact support and the proved lifted score regularity. This distinguishes an actual invariance calculation from a formal integral expression whose integrability has not been justified.




\textbf{Correspondence with the displayed declarations.}


\begin{itemize}

\item \texttt{gibbs\allowbreak{}Target\_\allowbreak{}infinitesimal\allowbreak{}Invariant}: Gibbs infinitesimal invariance for IPLA under only the original joint potential hypotheses. Integrability is a conclusion, so the zero integral cannot result from Lean's default value for a nonintegrable function.
\end{itemize}
\section{\texorpdfstring{Zero-\allowbreak{}dimensional state spaces have zero transport cost}{Zero-dimensional state spaces have zero transport cost}}
\label{proof-module-UnadjustedLangevin-Degenerate-lean}
{\small\color{gray}Source module: \texttt{Unadjusted\allowbreak{}Langevin/\allowbreak{}Degenerate.\allowbreak{}lean}.\par}


On a subsingleton state space every two states coincide, so every coupling has zero quadratic difference cost. Probability laws on such a space therefore have zero Wasserstein distance.

\subsection{\texorpdfstring{coupling\allowbreak{}Cost eq zero of subsingleton}{couplingCost eq zero of subsingleton}}


In a subsingleton state space every pair satisfies \(x=y\), so the displacement function is identically zero. Its extended \(L^2\) norm is zero under every pair measure, hence every coupling cost vanishes.

\subsection{\texorpdfstring{wasserstein2 eq zero of subsingleton}{wasserstein2 eq zero of subsingleton}}


The independent product of two probabilities is a coupling, so the coupling type is nonempty. Every coupling has zero cost by the preceding result; taking their infimum gives \(W_2(\mu,\pi)=0\).

\subsection{\texorpdfstring{wasserstein2\allowbreak{}Real eq zero of subsingleton}{wasserstein2Real eq zero of subsingleton}}


Rewrite the extended Wasserstein distance as zero using the subsingleton theorem. Its real conversion is also zero. This covers dimension-zero spaces without invoking the nontrivial-space comparison \(m\le L\).

\section{\texorpdfstring{Zero-\allowbreak{}dimensional parameter and state spaces}{Zero-dimensional parameter and state spaces}}
\label{proof-module-IPLA-Degenerate-lean}
{\small\color{gray}Source module: \texttt{IPLA/\allowbreak{}Degenerate.\allowbreak{}lean}.\par}


Euclidean space over an empty coordinate type contains a single vector. These elementary cases allow the public theorem to retain independent dimensions that may both be zero.

\subsection{\texorpdfstring{The parameter error vanishes when d\ensuremath{\theta} = 0}{The parameter error vanishes when dθ = 0}}


If \(d_\theta=0\), every parameter vector equals every other one. In particular \(p_\theta z=\theta_*\), so \(\|p_\theta z-\theta_*\|^2=0\) pointwise. Its integral under any measure, and the square root of that integral, are zero. No probability, potential or noise assumption is required.




\textbf{Correspondence with the displayed declarations.}


\begin{itemize}

\item \texttt{parameter\allowbreak{}RMS\_\allowbreak{}zero\_\allowbreak{}dimension}: There is no parameter error when the parameter dimension is zero. This holds for every latent dimension and every measure.
\end{itemize}
\subsection{\texorpdfstring{Every term of the full zero-\allowbreak{}dimensional bound vanishes}{Every term of the full zero-dimensional bound vanishes}}


If also \(d_x=0\), the whole state is a singleton. Any two probability measures on it are equal and their \(W_2\) distance is zero. The equilibrium radius \(\sqrt{2d_\theta/(Nm)}\) and discretization factor \(\sqrt{d_x+d_\theta/N}\) are also zero. Substituting these facts reduces the full claimed inequality to \(0\le0\), irrespective of the values of \(m,\gamma,L\). This handles the case in which the nonempty-coordinate analytic theorem cannot be instantiated.




When \(d_\theta=0\), the parameter space is a singleton, hence every parameter error vector and its square vanish. If also \(d_x=0\), the whole state space is a singleton. Any probability law there is the same Dirac law, so its Wasserstein distance to the Gibbs probability is zero. The dimension factors in the remaining terms are also zero, giving the public bound in this case directly, without invoking a theorem requiring a nonempty coordinate index.




\textbf{Correspondence with the displayed declarations.}


\begin{itemize}

\item \texttt{zero\_\allowbreak{}dimensions\_\allowbreak{}bound}: The complete displayed error bound in the zero-dimensional state. No regularity of a potential, noise law, or initialization is needed here.
\end{itemize}
\section{\texorpdfstring{Gibbs Dirichlet form, symmetry and dissipation}{Gibbs Dirichlet form, symmetry and dissipation}}
\label{proof-module-UnadjustedLangevin-GeneratorDirichlet-lean}
{\small\color{gray}Source module: \texttt{Unadjusted\allowbreak{}Langevin/\allowbreak{}Generator\allowbreak{}Dirichlet.\allowbreak{}lean}.\par}


For compactly supported \(C^2\) tests \(\varphi,\psi\), the generator obeys \(\int\varphi\mathcal L\psi\,d\pi=-\int\langle\nabla\varphi,\nabla\psi\rangle d\pi\). This provides the symmetric Dirichlet form at the test-function level.

\subsection{\texorpdfstring{fderiv exp neg mul apply}{fderiv exp neg mul apply}}


Differentiate the product \(e^{-U}\varphi\): \(D(e^{-U}\varphi)(x)[v]=e^{-U(x)}(D\varphi(x)[v]-\langle\nabla U(x),v\rangle\varphi(x))\). Lean combines the product rule with the already proved derivative of the Gibbs weight, then collects the common exponential factor algebraically.

\subsection{\texorpdfstring{integral weighted directional dirichlet}{integral weighted directional dirichlet}}


Integrate by parts the derivative of \(e^{-U}\varphi\) against \(D\psi[v]\). The weighted product rule separates a \(D\varphi[v]D\psi[v]\) term from the drift term, giving \(\int e^{-U}\varphi L_v\psi=-\int e^{-U}D\varphi[v]D\psi[v]\). Compact support and \(C^2\) regularity justify every derivative and integral involved.

\subsection{\texorpdfstring{integral directional\allowbreak{}Generator dirichlet gibbs}{integral directionalGenerator dirichlet gibbs}}


Convert both sides of the weighted directional identity to Gibbs integrals. The same real normalization factor multiplies both integrals, producing \(\int\varphi L_v\psi\,d\pi_U=-\int D\varphi[v]D\psi[v]d\pi_U\). This is a change of measure, with no new integration-by-parts step.

\subsection{\texorpdfstring{integrable gradient inner gibbs\allowbreak{}Measure}{integrable gradient inner gibbsMeasure}}


The gradients are continuous; the gradient of the compactly supported test \(\varphi\) vanishes outside its compact support. Thus \(x\mapsto\langle\nabla\varphi(x),\nabla\psi(x)\rangle\) is continuous with compact support. The general compact-test Gibbs integrability theorem applies even though \(\psi\) itself need not have compact support in this statement.

\subsection{\texorpdfstring{integrable test mul langevin\allowbreak{}Generator gibbs\allowbreak{}Measure}{integrable test mul langevinGenerator gibbsMeasure}}


The generator of the \(C^2\) test \(\psi\) is continuous under the Lipschitz-gradient assumption on \(U\). Multiplication by the continuous compactly supported \(\varphi\) makes the product compactly supported, hence Gibbs integrable. This is the exact product needed on the left side of the Dirichlet identity.

\subsection{\texorpdfstring{integral langevin\allowbreak{}Generator dirichlet gibbs}{integral langevinGenerator dirichlet gibbs}}


Sum the directional Dirichlet identities over the orthonormal basis. Finite-sum integrability allows the integral to pass through the sum, and Parseval identifies \(\sum_iD\varphi[e_i]D\psi[e_i]=\langle\nabla\varphi,\nabla\psi\rangle\). Therefore \(\int\varphi L\psi\,d\pi_U=-\int\langle\nabla\varphi,\nabla\psi\rangle d\pi_U\).

\subsection{\texorpdfstring{integral langevin\allowbreak{}Generator symmetric gibbs}{integral langevinGenerator symmetric gibbs}}


Apply the Dirichlet identity in both orders \((\varphi,\psi)\) and \((\psi,\varphi)\). Symmetry of the real inner product makes the resulting energy integrands equal pointwise. Hence \(\int\varphi L\psi\,d\pi_U=\int\psi L\varphi\,d\pi_U\).

\subsection{\texorpdfstring{integral langevin\allowbreak{}Generator self gibbs}{integral langevinGenerator self gibbs}}


Specialize the bilinear identity to \(\psi=\varphi\), then rewrite \(\langle\nabla\varphi,\nabla\varphi\rangle=\|\nabla\varphi\|^2\). This identifies the quadratic generator form exactly: \(\int\varphi L\varphi\,d\pi_U=-\int\|\nabla\varphi\|^2d\pi_U\).

\subsection{\texorpdfstring{integral langevin\allowbreak{}Generator dissipative gibbs}{integral langevinGenerator dissipative gibbs}}


The self-pairing identity expresses the generator quadratic form as minus an integral of a nonnegative square. The integral is nonnegative, so its negative is at most zero. The proof uses the actual integrable energy established earlier, giving dissipativity on compactly supported \(C^2\) tests.

\section{\texorpdfstring{A weak forward-\allowbreak{}equation interface and a conditional stationarity route}{A weak forward-equation interface and a conditional stationarity route}}
\label{proof-module-UnadjustedLangevin-FokkerPlanck-lean}
{\small\color{gray}Source module: \texttt{Unadjusted\allowbreak{}Langevin/\allowbreak{}Fokker\allowbreak{}Planck.\allowbreak{}lean}.\par}


The weak Fokker–Planck equation is formulated by integrating \(\mathcal L\varphi\) against time-dependent probability measures. The final theorem in this module assumes uniqueness of this weak equation explicitly; the main ULA development proves stationarity by the separate Metropolis approximation route.

\subsection{\texorpdfstring{Is\allowbreak{}Weak\allowbreak{}Fokker\allowbreak{}Planck\allowbreak{}Solution}{IsWeakFokkerPlanckSolution}}


A curve \(\rho_t\) is in this weak solution class when it is probability valued, has the stated moment and integrability properties, and satisfies the integrated generator equation against every compactly supported \(C^2\) test. The structure fields keep the analytic conditions visible, preventing a formal real integral from substituting for an unproved integrable quantity.

\subsection{\texorpdfstring{constant gibbs is\allowbreak{}Weak\allowbreak{}Fokker\allowbreak{}Planck\allowbreak{}Solution}{constant gibbs isWeakFokkerPlanckSolution}}


Use the proved Gibbs normalization and second moment for the probability and moment fields. Compact-support generator integrability supplies the test-integral fields. Finally \(\int L\varphi\,d\pi_U=0\) makes the time integral vanish, while the constant curve has zero change in \(\int\varphi\,d\rho_t\). These verify each field of the weak equation for the stationary candidate.

\subsection{\texorpdfstring{Weak\allowbreak{}Fokker\allowbreak{}Planck\allowbreak{}Unique}{WeakFokkerPlanckUnique}}


This proposition quantifies over two curves in the stated weak solution class with the same initial measure and requires their equality throughout \([0,T]\). It is a named analytic uniqueness obligation. Defining the proposition does not prove it, and no axiom or automatic instance makes it available.

\subsection{\texorpdfstring{gibbs stationary of weak\allowbreak{}Fokker\allowbreak{}Planck\allowbreak{}Uniqueness}{gibbs stationary of weakFokkerPlanckUniqueness}}


Compare the supplied weak solution \(\rho\) with the constant Gibbs weak solution constructed above. Their time-zero measures agree by hypothesis. The explicit \texttt{hunique} premise therefore identifies them at every time. Applying this conditional bridge to a stochastic process still requires its law curve to satisfy the weak equation and requires this uniqueness premise.

\section{\texorpdfstring{The standalone strongly convex IPLA error bound}{The standalone strongly convex IPLA error bound}}
\label{proof-module-IPLA-Theorem-lean}
{\small\color{gray}Source module: \texttt{IPLA/\allowbreak{}Theorem.\allowbreak{}lean}.\par}


For the fixed observed-data slice \(U_y\), assume differentiability, joint \(m\)-strong convexity with \(m>0\), \(L\)-Lipschitz gradient, \(N>0\), and \(0<\gamma\le2/(m+L)\). Initialization is square integrable, and innovations have the stated independent standard Gaussian coordinate laws and initialization independence. Let \(\theta_*\) maximize the actual latent marginal. The theorem applies to every pair of independent natural dimensions, including zero.

\subsection{\texorpdfstring{Combine lifted transport with proved equilibrium concentration}{Combine lifted transport with proved equilibrium concentration}}


Suppose the state has at least one coordinate. The lifted geometry gives \(Nm\)-strong convexity and \(NL\)-smoothness in dimension \(D=d_\theta+Nd_x\). The Gibbs theorem normalizes \(\Pi_N\), while square-integrable initialization and Gibbs moments imply finite initial transport. Apply the shared Gaussian ULA theorem at step \(h=\gamma/N\), using the supplied Gaussian hypotheses. The step restriction rescales exactly and gives

\[W_2(\rho_k,\Pi_N)\le(1-m\gamma)^kW_2(\mu_0,\Pi_N)+\frac{33}{20}\frac Lm\sqrt{d_x+d_\theta/N}\sqrt\gamma.\]

The ULA output-finiteness result and the proved law identity ensure this is finite actual transport, before any real conversion. The equilibrium concentration theorem gives \(\operatorname{RMS}_{\Pi_N}(p_\theta-\theta_*)\le\sqrt{2d_\theta/(Nm)}\). For any coupling of \(\rho_k\) and \(\Pi_N\), Minkowski and \(\|p_\theta\|\le1\) imply the sampled parameter RMS is at most coupling RMS plus that equilibrium radius. Taking the coupling infimum and using finiteness yields

\[\left(\int\|p_\theta z-\theta_*\|^2d\rho_k\right)^{1/2}\le\sqrt{\frac{2d_\theta}{Nm}}+(1-m\gamma)^kW_2(\mu_0,\Pi_N)+\frac{33}{20}\frac Lm\sqrt{d_x+\frac{d_\theta}{N}}\sqrt\gamma.\]

Stationarity, finite transport and concentration are derived here through the shared and particle theorems; they are not additional theorem hypotheses.




The proof first constructs target probability and initial finite transport from the lifted strong-convexity assumptions and \(L^2\) initialization. Gaussian ULA then gives the state-space transport rate, with step \(\gamma/N\), modulus \(Nm\), score constant \(NL\), and dimension \(d_\theta+Nd_x\). The law adapter identifies its output with the actual supplied IPLA recursion and proves the output transport distance is finite.



The projection inequality gives \(R_{\theta_*}(\rho_k)\le W_2(\rho_k,\Pi_N)+R_{\theta_*}(\Pi_N)\); equilibrium concentration supplies the second term. The facts \texttt{houtput} and \texttt{hmoment} ensure finite values before projected RMS is rewritten as a square-root integral. Finally, the exact particle/\allowbreak{}time scaling simplifies the transport terms. Thus normalization, stationary reference and parameter concentration are derived inputs, while joint strong convexity, smoothness and the Gaussian innovation laws remain the stated assumptions.




\textbf{Correspondence with the displayed declarations.}


\begin{itemize}

\item \texttt{nonempty\_\allowbreak{}dimension\_\allowbreak{}bound}: For fixed observations \texttt{y}, write \texttt{U y (pair theta x) = -\allowbreak{}log p\_\allowbreak{}theta(x,y)}. The general-dimension counterpart of equation (34), with independent parameter dimension \texttt{dtheta} and latent dimension \texttt{dx}, for the actual Gaussian IPLA recursion and its defined Gibbs target. Every hypothesis is displayed here or is one of the Gaussian-law and independence fields of \texttt{Gaussian\allowbreak{}Innovations}, whose definition is included in the standalone file. Normalization, all finite transport moments, a stationary Langevin reference, and the equilibrium parameter concentration are derived in the proof.
\end{itemize}
\subsection{\texorpdfstring{Include every zero-\allowbreak{}dimensional case in the public theorem}{Include every zero-dimensional case in the public theorem}}


The public theorem splits on whether the state coordinate index type is nonempty. In the nonempty case, instantiate the preceding theorem directly, retaining arbitrary independent \(d_\theta,d_x\). If the index type is empty, its parameter part forces \(d_\theta=0\); since \(N>0\), a positive \(d_x\) would supply a latent coordinate, so \(d_x=0\) as well. Normalize the Gibbs target by the same strong-convexity theorem and apply the singleton-state bound, where every error, transport distance and dimension term vanishes. This completes the displayed estimate in all dimensions. Existence and uniqueness of a permissible \(\theta_*\) are established separately by \texttt{exists\allowbreak{}Unique\_\allowbreak{}marginal\allowbreak{}Maximizer}, so the maximizer binder does not demand an unproved existence assumption.




\textbf{Correspondence with the displayed declarations.}


\begin{itemize}

\item \texttt{theorem1}: Strongly convex IPLA for independent, arbitrary parameter and latent dimensions. All stationary-reference and concentration inputs are derived. The zero-dimensional cases are included without extra assumptions.
\end{itemize}
\section{\texorpdfstring{Exact entropy scaling with an explicit lifted LSI constant}{Exact entropy scaling with an explicit lifted LSI constant}}
\label{proof-module-IPLA-LSI-Scaling-lean}
{\small\color{gray}Source module: \texttt{IPLA/\allowbreak{}LSI/\allowbreak{}Scaling.\allowbreak{}lean}.\par}


Keep the LSI constant \(\beta\) of the actual lifted target explicit. Substitute \(h=\gamma/N\), \(\ell=NL\), \(D=d_\theta+Nd_x\) into the shared ULA entropy argument. The regime \(\beta=N\alpha\) is an additional lifted-target assumption, with its own scalar specializations.

\subsection{\texorpdfstring{Translate the numerical step and compatibility conditions}{Translate the numerical step and compatibility conditions}}


For \(N>0\), \(L>0\), factor the denominator and cancel positive \(N\) to get

\[\frac\gamma N\le\frac\beta{4(NL)^2}\quad\Longleftrightarrow\quad\gamma\le\frac\beta{4NL^2}.\]

If \(\beta=N\alpha\), multiplication by positive \(N\) gives \(N\alpha\le NL\Longleftrightarrow\alpha\le L\), and cancellation gives \(N\alpha/(4NL^2)=\alpha/(4L^2)\). The forward compatibility lemma simply uses the equivalence. These are exact ordered-field identities; they do not derive an LSI for the lifted model from any pair-model LSI.




The general lifted theorem uses \(h=\gamma/N\), score constant \(L_N=NL\), dimension \(D=d_\theta+Nd_x\), and an LSI constant \(\beta\) for the actual lifted Gibbs target. Thus \(h\le\beta/(4L_N^2)\) becomes \(\gamma\le\beta/(4NL^2)\). Only after additionally assuming \(\beta=N\alpha\) can the step restriction simplify to \(\gamma\le\alpha/(4L^2)\). Positive \(N\) justifies these cancellations; no product tensorization is used, since the pair potentials share a parameter.




\textbf{Correspondence with the displayed declarations.}


\begin{itemize}

\item \texttt{lifted\_\allowbreak{}lsi\_\allowbreak{}step\_\allowbreak{}iff}: The ULA restriction for the lifted step is exactly this IPLA restriction.
\item \texttt{linear\_\allowbreak{}lsi\_\allowbreak{}compatibility\_\allowbreak{}iff}: The scaled compatibility condition is equivalent to the base comparison.
\item \texttt{linear\_\allowbreak{}lsi\_\allowbreak{}step\_\allowbreak{}iff}: When the actual lifted LSI constant is \texttt{N * alpha}, the IPLA step restriction has no remaining particle-number factor.
\end{itemize}
\subsection{\texorpdfstring{Rewrite the contraction exponent}{Rewrite the contraction exponent}}


Associativity of real multiplication and division gives \(-\beta(\gamma/N)k=-\beta\gamma k/N\), so the corresponding exponentials agree. When \(\beta=N\alpha\) and \(N>0\), cancellation yields

\[e^{-\beta\gamma k/N}=e^{-\alpha\gamma k},\qquad e^{-(N\alpha)(\gamma/N)k}=e^{-\alpha\gamma k}.\]

The second version composes the first two identities. The rate becomes independent of \(N\) only under this explicitly supplied linear scaling of the actual lifted LSI constant.




\textbf{Correspondence with the displayed declarations.}


\begin{itemize}

\item \texttt{lifted\_\allowbreak{}entropy\_\allowbreak{}rate\_\allowbreak{}eq}: The general decay rate retains the actual lifted LSI constant.
\item \texttt{linear\_\allowbreak{}entropy\_\allowbreak{}rate\_\allowbreak{}eq}: The rewritten general rate simplifies under the additional linear LSI scaling assumption.
\item \texttt{linear\_\allowbreak{}ula\_\allowbreak{}entropy\_\allowbreak{}rate\_\allowbreak{}eq}: The same simplification directly from the ULA step notation.
\end{itemize}
\subsection{\texorpdfstring{Rewrite full and normalized entropy biases}{Rewrite full and normalized entropy biases}}


The shared ULA bias is \(8hD\ell^2/\beta\). Substitute the lifted values and cancel one particle factor:

\[\frac{8(\gamma/N)D(NL)^2}{\beta}=\frac{8\gamma ND L^2}{\beta},\qquad \frac1N\frac{8\gamma ND L^2}{\beta}=\frac{8\gamma DL^2}{\beta}.\]

The dimension identity \(D/N=d_x+d_\theta/N\) follows by distributing division over the sum. If \(\beta=N\alpha\), the normalized bias becomes

\[\frac{8\gamma DL^2}{N\alpha}=\frac{8\gamma L^2}{\alpha}\left(d_x+\frac{d_\theta}{N}\right).\]

The formal algebra handles \(\alpha=0\) separately because real division is total; applications use \(\alpha>0\). For any bound \(H\le qH_0+8\gamma ND L^2/\beta\), divide by positive \(N\), distribute, and use the second identity to obtain the normalized entropy bound.




The full-state bias is \(8hDL_N^2/\beta=8\gamma NDL^2/\beta\); dividing entropy by \(N\) removes exactly one particle factor. Under \(\beta=N\alpha\), its intensive form becomes \((8\gamma L^2/\alpha)(d_x+d_\theta/N)\). The initial entropy is divided by \(N\) at the same time. The one-step forcing is different: \(6h^2DL_N^2=6\gamma^2DL^2\), so its normalized form is \(6\gamma^2L^2(d_x+d_\theta/N)\). Keeping these full and intensive quantities separate avoids losing a factor of \(N\).




\textbf{Correspondence with the displayed declarations.}


\begin{itemize}

\item \texttt{lifted\_\allowbreak{}entropy\_\allowbreak{}bias\_\allowbreak{}eq}: The unnormalized entropy bias after substituting the lifted step and smoothness. The dimension \texttt{D} is kept as a scalar parameter.
\item \texttt{normalized\_\allowbreak{}entropy\_\allowbreak{}bias\_\allowbreak{}eq}: Dividing the entropy by the particle number cancels one bias factor.
\item \texttt{normalized\_\allowbreak{}dimension\_\allowbreak{}eq}: The intensive dimension appearing in normalized entropy.
\item \texttt{linear\_\allowbreak{}normalized\_\allowbreak{}entropy\_\allowbreak{}bias\_\allowbreak{}eq}: The normalized bias under an LSI constant \texttt{N * alpha} for the actual lift.
\item \texttt{normalize\_\allowbreak{}entropy\_\allowbreak{}bound}: A general entropy bound can be divided by the positive particle number, with its initial contribution and bias both normalized exactly.
\end{itemize}
\subsection{\texorpdfstring{Rewrite the one-\allowbreak{}step forcing and normalized recurrence}{Rewrite the one-step forcing and normalized recurrence}}


The squared step cancels both particle factors in the squared smoothness:

\[6(\gamma/N)^2D(NL)^2=6\gamma^2DL^2.\]

Dividing by \(N\) and substituting the intensive dimension gives \(6\gamma^2DL^2/N=6\gamma^2L^2(d_x+d_\theta/N)\). Thus any supplied lifted recurrence normalizes to

\[\frac{H_{k+1}}N\le e^{-\beta\gamma/N}\frac{H_k}N+\frac{6\gamma^2DL^2}{N}.\]

This scalar helper explicitly assumes its recurrence. Its proof divides by positive \(N\), rewrites the exponential argument, and applies the forcing identity. Under \(\beta=N\alpha\), the rate and intensive-dimension identities give the specialized recurrence with rate \(e^{-\alpha\gamma}\) and forcing \(6\gamma^2L^2(d_x+d_\theta/N)\). The analytic theorem derives the recurrence separately from interpolation data.




\textbf{Correspondence with the displayed declarations.}


\begin{itemize}

\item \texttt{lifted\_\allowbreak{}entropy\_\allowbreak{}forcing\_\allowbreak{}eq}: The unnormalized one-step forcing loses both particle factors from the squared step.
\item \texttt{normalized\_\allowbreak{}entropy\_\allowbreak{}forcing\_\allowbreak{}eq}: The normalized one-step forcing has the intensive dimension.
\item \texttt{normalize\_\allowbreak{}lifted\_\allowbreak{}entropy\_\allowbreak{}recurrence}: Normalizing the lifted one-step recurrence. Its analytic recurrence is an explicit premise; this theorem performs only scalar scaling.
\item \texttt{linear\_\allowbreak{}normalized\_\allowbreak{}entropy\_\allowbreak{}recurrence}: The normalized recurrence under the additional linear lifted-LSI scale.
\end{itemize}
\section{\texorpdfstring{Actual Gaussian ULA laws and their frozen-\allowbreak{}drift interpolation}{Actual Gaussian ULA laws and their frozen-drift interpolation}}
\label{proof-module-UnadjustedLangevin-LSI-ULALaw-lean}
{\small\color{gray}Source module: \texttt{Unadjusted\allowbreak{}Langevin/\allowbreak{}LSI/\allowbreak{}ULALaw.\allowbreak{}lean}.\par}


Define the frozen-drift law \(\rho_{k,t}=\operatorname{Law}(X_k-t\nabla U(X_k)+\sqrt{2t}Z)\), with independent standard Gaussian \(Z\), and \(\rho_{k+1}=\rho_{k,h}\). The endpoint and probability identities connect the entropy argument to the actual algorithm.

\subsection{\texorpdfstring{ula\allowbreak{}Entropy}{ulaEntropy}}


Set \(H_k=\operatorname{KL}(\rho_k\mid\nu)_{\mathbb R}\). The definition uses the actual extended KL of the ULA law, followed by \texttt{to\allowbreak{}Real}. Separate theorems prove finite KL wherever this is used as ordinary real entropy; the definition alone would not exclude infinity.

\subsection{\texorpdfstring{ula\allowbreak{}Entropy zero}{ulaEntropy zero}}


Substitute the definitional identity \(\rho_0=\mu\) into the entropy definition. This gives \(H_0=\operatorname{KL}(\mu\mid\nu)_{\mathbb R}\), the initial term appearing in the final exponential bound.

\subsection{\texorpdfstring{interpolation\allowbreak{}Law}{interpolationLaw}}


Define \(\rho_t=(T_t)_*(\mu\otimes\gamma)\), where \(T_t(x,z)=x-t\nabla U(x)+\sqrt{2t}z\). The product encodes independence. This is an alias of the actual \texttt{ula\allowbreak{}Law\allowbreak{}Step}; its frozen-drift Gaussian interpolation interpretation applies to \(t\ge0\).

\subsection{\texorpdfstring{ula\allowbreak{}Laws}{ulaLaws}}


Iterate the actual independent-Gaussian transition at fixed step \(h\): \(\rho_k=P_h^k\mu\). The definition is on measures, while the separate chain-law theorems identify independent iid implementations with these same laws.

\subsection{\texorpdfstring{interpolation\allowbreak{}Law endpoint}{interpolationLaw endpoint}}


At elapsed time \(h\), the interpolation from \(\rho_k\) is \(P_h\rho_k\). The preceding successor identity therefore proves \(\rho_{k,h}=\rho_{k+1}\). This is the law equality used to turn a continuous interpolation bound into the discrete entropy step.

\subsection{\texorpdfstring{ula\allowbreak{}Laws is\allowbreak{}Probability}{ulaLaws isProbability}}


Induct from the initial probability \(\mu\). At the successor rewrite the iterate as one interpolation step at time \(h\), then apply probability preservation. Every algorithm law is therefore a probability before applying KL and LSI identities.

\subsection{\texorpdfstring{interpolation\allowbreak{}Law zero}{interpolationLaw zero}}


At \(t=0\), both the drift increment and Gaussian increment vanish, so the update map is first-coordinate projection. Integrating out the standard Gaussian probability returns \(\mu\). This endpoint identity needs no potential regularity.

\subsection{\texorpdfstring{ula\allowbreak{}Laws zero}{ulaLaws zero}}


The zeroth iterate of a function is the identity, hence \(\rho_0=\mu\). Lean uses definitional equality, which makes the initial-entropy rewrite exact rather than an asymptotic endpoint argument.

\subsection{\texorpdfstring{ula\allowbreak{}Laws succ}{ulaLaws succ}}


The successor-iterate identity gives \(\rho_{k+1}=P_h(\rho_k)\). Lean invokes \texttt{Function.\allowbreak{}iterate\_\allowbreak{}succ\_\allowbreak{}apply'}, whose order matches applying one more transition to the current law.

\subsection{\texorpdfstring{interpolation\allowbreak{}Law is\allowbreak{}Probability}{interpolationLaw isProbability}}


Lipschitzness makes the gradient continuous and the joint update measurable. The product \(\mu\otimes\gamma\) is a probability, and its pushforward has the same total mass. Evaluating the mapped measure on the whole space proves normalization of \(\rho_t\).

\section{\texorpdfstring{The relative-\allowbreak{}density log-\allowbreak{}Sobolev convention}{The relative-density log-Sobolev convention}}
\label{proof-module-UnadjustedLangevin-LSI-FunctionalInequality-lean}
{\small\color{gray}Source module: \texttt{Unadjusted\allowbreak{}Langevin/\allowbreak{}LSI/\allowbreak{}Functional\allowbreak{}Inequality.\allowbreak{}lean}.\par}


Write \(\rho=r\nu\), \(H_\nu(\rho)=\int r\log r\,d\nu\), and \(J_\nu(\rho)=\int\|\nabla\log r\|^2d\rho\). This module defines the explicit smooth positive-density LSI convention \(2\alpha H\le J\) and guards real-valued conversions by finite divergence and Fisher information.

\subsection{\texorpdfstring{tilted\allowbreak{}Measure}{tiltedMeasure}}


For a real relative density \(r\), define \(\rho_r(dx)=r(x)\nu(dx)\) using \texttt{with\allowbreak{}Density (ENNReal.\allowbreak{}of\allowbreak{}Real \ensuremath{\circ} r)}. Positivity and normalization are later hypotheses, rather than part of this raw measure definition. A normalized positive \(r\) is the density of the current law relative to the target.

\subsection{\texorpdfstring{relative\allowbreak{}Fisher}{relativeFisher}}


Define \(I_\nu(r)=\int\|\nabla\log r\|^2d\rho_r\) as an extended nonnegative integral. The explicit density is differentiated pointwise, avoiding differentiation of an arbitrary almost-everywhere Radon–Nikodym representative. Later real-valued estimates prove finiteness before taking \texttt{to\allowbreak{}Real}.

\subsection{\texorpdfstring{Is\allowbreak{}Log\allowbreak{}Sobolev}{IsLogSobolev}}


The target is a probability and satisfies \(2\alpha\operatorname{KL}(\rho_r\mid\nu)\le I_\nu(r)\) for every strictly positive \(C^1\) normalized relative density. The factor \(2\alpha\) fixes the convention used throughout the proof. This test class and its normalization are explicit; the definition contains no algorithm or entropy-decay claim.

\subsection{\texorpdfstring{Why finite KL and Fisher information are separate obligations}{Why finite KL and Fisher information are separate obligations}}


Apply the extended LSI to the supplied positive \(C^1\) probability density. Both KL and Fisher are assumed finite, so \texttt{ENNReal.\allowbreak{}to\allowbreak{}Real\_\allowbreak{}le\_\allowbreak{}to\allowbreak{}Real} preserves the inequality. Nonnegativity of \(2\alpha\) identifies the real value of its embedding, giving \(2\alpha H_\nu(r)\le I_\nu(r)_{\mathbb R}\).

\section{\texorpdfstring{Integrate a forced entropy-\allowbreak{}dissipation inequality}{Integrate a forced entropy-dissipation inequality}}
\label{proof-module-UnadjustedLangevin-LSI-EntropyDecay-lean}
{\small\color{gray}Source module: \texttt{Unadjusted\allowbreak{}Langevin/\allowbreak{}LSI/\allowbreak{}Entropy\allowbreak{}Decay.\allowbreak{}lean}.\par}


An integrating factor solves \(H\prime\le-aH+b\) as \(H(t)\le e^{-at}H(0)+(b/a)(1-e^{-at})\), under the displayed positive-rate, continuity and one-sided derivative assumptions. The entropy wrapper supplies the appropriate rate from LSI and dissipation.

\subsection{\texorpdfstring{A scalar affine comparison from right derivatives}{A scalar affine comparison from right derivatives}}


For continuous \(H\) with right derivative bounded by \(-aH+b\), \(a>0\), compare it to the exact scalar ODE solution. The one-sided Grönwall theorem accepts the right derivative through its right-slope liminf inequality. Simplifying its explicit solution gives \(H(t)\le e^{-at}H(0)+(b/a)(1-e^{-at})\). Neither \(H\) nor \(b\) needs to be nonnegative.

\subsection{\texorpdfstring{Scalar LSI and dissipation give rate 2\ensuremath{\alpha}}{Scalar LSI and dissipation give rate 2α}}


The supplied dissipation identity is \(H'=-I\), and the scalar LSI premise is \(2\alpha H\le I\). Therefore \(H'\le-2\alpha H\). Apply the affine comparison theorem with \(a=2\alpha\) and \(b=0\), obtaining \(H(t)\le e^{-2\alpha t}H(0)\). This scalar theorem does not itself establish the dissipation identity for a law-valued evolution.

\section{\texorpdfstring{Relative entropy along an identified density path}{Relative entropy along an identified density path}}
\label{proof-module-UnadjustedLangevin-LSI-EntropyFlow-lean}
{\small\color{gray}Source module: \texttt{Unadjusted\allowbreak{}Langevin/\allowbreak{}LSI/\allowbreak{}Entropy\allowbreak{}Flow.\allowbreak{}lean}.\par}


Density entropy is the real conversion of the actual KL divergence of the tilted measure. With finite-entropy guards, LSI and a supplied dissipation inequality give exponential entropy decay along this path.

\subsection{\texorpdfstring{Define actual density entropy with its finiteness guard}{Define actual density entropy with its finiteness guard}}


Define the real entropy of a specified relative density as \(H_\nu(r)=\operatorname{KL}(r\nu\mid\nu)_{\mathbb R}\). It uses the same \texttt{tilted\allowbreak{}Measure} and actual Mathlib KL as the LSI definition. Applications must separately establish finite KL before using it as an ordinary entropy value.

\subsection{\texorpdfstring{Apply the finite real LSI at every time}{Apply the finite real LSI at every time}}


At each interpolation time, positivity, \(C^1\) regularity, normalization, and finite KL/\allowbreak{}Fisher allow \texttt{Is\allowbreak{}Log\allowbreak{}Sobolev.\allowbreak{}entropy\_\allowbreak{}le\_\allowbreak{}fisher} to give \(2\alpha H(t)\le I(t)\). Combine this with the explicit right-dissipation identity and entropy continuity in \texttt{entropy\_\allowbreak{}decay\_\allowbreak{}of\_\allowbreak{}dissipation}. The conclusion is decay at rate \(2\alpha\) for that admissible density path; identifying an actual evolution with the path remains an application obligation.

\section{\texorpdfstring{The conditional one-\allowbreak{}step KL contraction with bias 6}{The conditional one-step KL contraction with bias 6}}
\label{proof-module-UnadjustedLangevin-LSI-OneStep-lean}
{\small\color{gray}Source module: \texttt{Unadjusted\allowbreak{}Langevin/\allowbreak{}LSI/\allowbreak{}One\allowbreak{}Step.\allowbreak{}lean}.\par}


Assume the displayed LSI, gradient-moment and interpolation differential bounds. For \(0<h\le\alpha/(4L^2)\), a scalar integrating-factor argument proves \(H_{k+1}\le e^{-\alpha h}H_k+6h^2dL^2\).

\subsection{\texorpdfstring{Absorb a small initial-\allowbreak{}entropy forcing into the exponential gap}{Absorb a small initial-entropy forcing into the exponential gap}}


For \(0\le x\le1/4\), the tangent bound for the exponential gives \(e^{-3x/2}\ge1/2\). Multiply \(e^{x/2}\ge1+x/2\) by \(e^{-3x/2}\); then \(e^{-x}\ge e^{-3x/2}+(x/2)e^{-3x/2}\ge e^{-3x/2}+x/4\). This quantitative rate gap absorbs the initial-entropy forcing.

\subsection{\texorpdfstring{Uniform forcing, affine comparison and the coefficient 6}{Uniform forcing, affine comparison and the coefficient 6}}


Combine \(2\alpha H\le I\) and \(H'\le-3I/4+t^2L^2G_0+2tdL^2\). The gradient bound \(G_0\le(4L^2/\alpha)H(0)+2dL\) and \(h\le\alpha/(4L^2)\) give a uniform inequality \(H'\le-(3\alpha/2)H+b\), with \(b=(\alpha/4)H(0)+(5/2)hdL^2\). The proof establishes \(hL\le1/4\), \(\alpha h\le1/4\), and the coefficient bound \(h^2L^2(4L^2/\alpha)\le\alpha/4\) before this substitution.



Apply \texttt{affine\_\allowbreak{}decay\_\allowbreak{}of\_\allowbreak{}deriv\_\allowbreak{}right\_\allowbreak{}le} and bound its forcing integral by \(bh\), using \(1-e^{-ah}\le ah\). The exponential-gap lemma absorbs \((\alpha h/4)H(0)\), changing the decay factor from \(e^{-3\alpha h/2}\) to \(e^{-\alpha h}\). Nonnegativity of \(H(0)\) justifies multiplying that scalar bound. Finally enlarge the remaining nonnegative \((5/2)h^2dL^2\) term to \(6h^2dL^2\), obtaining the stated one-step estimate.

\section{\texorpdfstring{Iterate entropy contraction and obtain the bias constant 8}{Iterate entropy contraction and obtain the bias constant 8}}
\label{proof-module-UnadjustedLangevin-LSI-DiscreteEntropy-lean}
{\small\color{gray}Source module: \texttt{Unadjusted\allowbreak{}Langevin/\allowbreak{}LSI/\allowbreak{}Discrete\allowbreak{}Entropy.\allowbreak{}lean}.\par}


Starting from \(H_{k+1}\le e^{-\alpha h}H_k+6h^2dL^2\), induction retains the exact geometric sum. The small-step denominator estimate \(1-e^{-\alpha h}\ge\frac34\alpha h\) gives \(H_k\le e^{-\alpha hk}H_0+8hdL^2/\alpha\).

\subsection{\texorpdfstring{Keep the exact geometric sum before simplifying it}{Keep the exact geometric sum before simplifying it}}


Set \(q=e^{-\alpha h}\in(0,1)\), \(b=6h^2dL^2\). The affine recurrence theorem sums \(H_{k+1}\le qH_k+b\) into \(H_k\le q^kH_0+b(1-q^k)/(1-q)\). Rewrite \(q^k=e^{-\alpha hk}\) using the exponential natural-multiple identity. No sign assumption on \(H_0\) is needed.

\subsection{\texorpdfstring{three quarters mul le one sub exp neg}{three quarters mul le one sub exp neg}}


The tangent bound \(e^x\ge1+x\), multiplied by \(e^{-x}\), gives \(1-e^{-x}\ge x/(1+x)\). For \(0\le x\le1/4\), elementary positive-denominator algebra shows \(x/(1+x)\ge3x/4\). This lower bound controls the geometric-series denominator in the bias floor.

\subsection{\texorpdfstring{discrete entropy recurrence bound}{discrete entropy recurrence bound}}


The nonnegative forcing gives \(H_k\le e^{-\alpha hk}H_0+6h^2dL^2/(1-e^{-\alpha h})\). For \(\alpha h\le1/4\), the denominator is at least \(3\alpha h/4>0\). Divide and cancel to obtain the bias floor \(8hdL^2/\alpha\). Lean proves positivity of the denominator before multiplying across it.

\section{\texorpdfstring{Conditional LSI theorem for the actual ULA laws}{Conditional LSI theorem for the actual ULA laws}}
\label{proof-module-UnadjustedLangevin-LSI-Theorem-lean}
{\small\color{gray}Source module: \texttt{Unadjusted\allowbreak{}Langevin/\allowbreak{}LSI/\allowbreak{}Theorem.\allowbreak{}lean}.\par}


The actual iterated Gaussian ULA laws satisfy \(H_k\le e^{-\alpha hk}H_0+8hdL^2/\alpha\), once the displayed interpolation identification, finite KL, entropy/\allowbreak{}Fisher dissipation, LSI, and moment hypotheses are supplied. This endpoint does not construct all those analytic data from LSI and smoothness alone.



The full actual Gaussian ULA endpoint in UnconditionalTheorem.lean derives all necessary analytic estimates from smoothness, LSI, and finite initial KL. This module retains a reusable interface for explicitly specified interpolation data.

\subsection{\texorpdfstring{The target estimate and the conditional Lean statement}{The target estimate and the conditional Lean statement}}


For every index, derive the one-step recurrence from the preceding interpolation theorem. The numerical conditions \(\alpha\le L\), \(h\le\alpha/(4L^2)\) imply \(\alpha h\le1/4\). Apply \texttt{discrete\_\allowbreak{}entropy\_\allowbreak{}recurrence\_\allowbreak{}bound} and rewrite the initial value with \texttt{ula\allowbreak{}Entropy\_\allowbreak{}zero}. This gives \(H_k\le e^{-\alpha hk}\operatorname{KL}(\mu\mid\pi_U)_{\mathbb R}+8hdL^2/\alpha\), conditional on the displayed analytic path hypotheses.

\subsection{\texorpdfstring{Why finite KL and Fisher information are separate obligations}{Why finite KL and Fisher information are separate obligations}}


At time zero the supplied relative-density law identity gives \(r_{k,0}\nu=\rho_k\), because \texttt{interpolation\allowbreak{}Law\_\allowbreak{}zero} returns its input law. Substitute this equality into finite KL of the time-zero density path. It yields finite actual KL of the algorithm law at every index.

\subsection{\texorpdfstring{Combine production, LSI, and the initial gradient moment}{Combine production, LSI, and the initial gradient moment}}


Define \(G_k=\int\|\nabla U(x)\|^2d\rho_k(x)\) for the actual ULA law. Since Lean's ordinary integral is totalized, later hypotheses explicitly prove integrability of the squared gradient before treating \(G_k\) as an energy. The definition itself contains no entropy-moment estimate.

\subsection{\texorpdfstring{The direct endpoint constructs the recurrence from the path hypotheses}{The direct endpoint constructs the recurrence from the path hypotheses}}


Identify the entropy path's endpoints using \texttt{interpolation\allowbreak{}Law\_\allowbreak{}zero} and \texttt{interpolation\allowbreak{}Law\_\allowbreak{}endpoint}: its time-zero entropy is \(H_k\), and its time-\(h\) entropy is \(H_{k+1}\). The law identity and probability preservation prove that every intermediate relative density is normalized.



Apply the real LSI inequality using the explicit positive smooth density and finite KL/\allowbreak{}Fisher guards. The scalar \texttt{one\_\allowbreak{}step\_\allowbreak{}entropy\_\allowbreak{}of\_\allowbreak{}interpolation} then combines it with the supplied entropy derivative, production, and gradient-moment estimates. Rewriting the two endpoints yields \(H_{k+1}\le e^{-\alpha h}H_k+6h^2dL^2\). Those analytic interpolation premises remain explicit in this theorem.

\section{\texorpdfstring{Differentiate the entropy integrand under domination}{Differentiate the entropy integrand under domination}}
\label{proof-module-UnadjustedLangevin-LSI-EntropyDerivative-lean}
{\small\color{gray}Source module: \texttt{Unadjusted\allowbreak{}Langevin/\allowbreak{}LSI/\allowbreak{}Entropy\allowbreak{}Derivative.\allowbreak{}lean}.\par}


The scalar identity \((r\log r)\prime=r\prime(\log r+1)\), for \(r>0\), lifts to an integral derivative under the explicitly supplied common local integrable bound. Domination is a premise, not a consequence of Gaussian convolution in this module.

\subsection{\texorpdfstring{has\allowbreak{}Deriv\allowbreak{}At mul log}{hasDerivAt mul log}}


If \(r(t)>0\), the product and logarithm chain rules give \((r\log r)'=\dot r\log r+r\dot r/r\). Positivity supplies \(r(t)\ne0\), so cancellation gives \(\dot r(\log r+1)\). Lean proves that cancellation explicitly before simplifying the derivative expression.

\subsection{\texorpdfstring{has\allowbreak{}Deriv\allowbreak{}At integral mul log of dominated}{hasDerivAt integral mul log of dominated}}


For almost every spatial point, apply \texttt{has\allowbreak{}Deriv\allowbreak{}At\_\allowbreak{}mul\_\allowbreak{}log} throughout a common time neighborhood. Its derivative is bounded in norm by the supplied integrable spatial envelope, uniformly over that neighborhood.



The dominated differentiation theorem uses this bound, integrability of \(r_t\log r_t\), and measurability of both integrands. It concludes both integrability of \(\dot r_t(\log r_t+1)\) and \(\frac d{dt}\int r_t\log r_t d\nu=\int\dot r_t(\log r_t+1)d\nu\). The derivative of the integral is a conclusion, not a hypothesis.

\section{\texorpdfstring{Identify density entropy and Fisher integrals}{Identify density entropy and Fisher integrals}}
\label{proof-module-UnadjustedLangevin-LSI-KLOfDensity-lean}
{\small\color{gray}Source module: \texttt{Unadjusted\allowbreak{}Langevin/\allowbreak{}LSI/\allowbreak{}KLOf\allowbreak{}Density.\allowbreak{}lean}.\par}


For a strictly positive relative density, with the stated differentiability and integrability, the actual likelihood ratio equals \(\log r\). This identifies the formal KL divergence with \(\int r\log r\,d\nu\), proves its finiteness and identifies the real Fisher value with its score-energy integral.

\subsection{\texorpdfstring{llr tilted\allowbreak{}Measure ae eq log}{llr tiltedMeasure ae eq log}}


The with-density measure is absolutely continuous with respect to \(\nu\), and its Radon–Nikodym derivative equals \texttt{of\allowbreak{}Real r} \(\nu\)-almost everywhere. Absolute continuity transfers that equality to \(\rho_r\)-almost everywhere. Taking the real logarithm gives the actual log-likelihood ratio \(\log r\).

\subsection{\texorpdfstring{kl\allowbreak{}Div tilted\allowbreak{}Measure ne top}{klDiv tiltedMeasure ne top}}


Absolute continuity follows from the with-density construction. Integrability of \(r\log r\) gives integrability of \(\log r\) under \(\rho_r\), and the log-likelihood identity identifies that with the KL integrand. The standard finite-KL criterion therefore proves \(\operatorname{KL}(\rho_r\mid\nu)<\infty\).

\subsection{\texorpdfstring{density\allowbreak{}Entropy eq integral}{densityEntropy eq integral}}


For normalized \(\nu\) and \(\rho_r\), the finite-KL integral formula has mass correction \(\nu(E)-\rho_r(E)=0\). Replace the log-likelihood ratio by \(\log r\), then change density to obtain \(H_\nu(r)=\int r\log r\,d\nu\). The integrability premise is used before invoking that formula.

\subsection{\texorpdfstring{relative\allowbreak{}Fisher to\allowbreak{}Real eq integral}{relativeFisher toReal eq integral}}


Nonnegativity and integrability identify the real value of the extended Fisher integral with the ordinary integral under \(\rho_r\). The weighted-measure formula then rewrites it as \(I_\nu(r)_{\mathbb R}=\int r\|\nabla\log r\|^2d\nu\). Both representations refer to the same supplied density.

\subsection{\texorpdfstring{integral tilted\allowbreak{}Measure}{integral tiltedMeasure}}


For measurable \(r\ge0\), the Bochner with-density formula gives \(\int g\,d\rho_r=\int r(x)g(x)d\nu\). The embedded density is finite pointwise and has real value \(r\), so the theorem's extended-to-real scalar action simplifies to ordinary multiplication.

\subsection{\texorpdfstring{integrable log tilted\allowbreak{}Measure iff}{integrable log tiltedMeasure iff}}


Apply the with-density integrability criterion to \(g=\log r\). Since the density is measurable, finite, and nonnegative, it yields \(\log r\in L^1(\rho_r)\iff r\log r\in L^1(\nu)\). This is the exact entropy-integrability bridge used before converting KL to real values.

\subsection{\texorpdfstring{relative\allowbreak{}Fisher ne top}{relativeFisher ne top}}


The squared score is nonnegative. Its ordinary integrability implies finiteness of its extended nonnegative integral, by the \texttt{has\allowbreak{}Finite\allowbreak{}Integral\_\allowbreak{}iff\_\allowbreak{}of\allowbreak{}Real} criterion. That extended integral is exactly \texttt{relative\allowbreak{}Fisher}.

\section{\texorpdfstring{Relative scores turn the weak equation into entropy production}{Relative scores turn the weak equation into entropy production}}
\label{proof-module-UnadjustedLangevin-LSI-ScoreCalculus-lean}
{\small\color{gray}Source module: \texttt{Unadjusted\allowbreak{}Langevin/\allowbreak{}LSI/\allowbreak{}Score\allowbreak{}Calculus.\allowbreak{}lean}.\par}


The chain rule \(\nabla\log r=\nabla r/r\) converts the entropy test \(\log r+1\) into the relative score. The weak identity supplied for this test then gives Fisher dissipation, with an extra score/\allowbreak{}error pairing when drift forcing is present.

\subsection{\texorpdfstring{gradient log of pos}{gradient log of pos}}


The logarithm chain rule gives \(D\log r(x)=r(x)^{-1}Dr(x)\) when \(r(x)>0\). The real Hilbert-space gradient is the Riesz representative of the derivative; linearity of that identification gives \(\nabla\log r(x)=r(x)^{-1}\nabla r(x)\).

\subsection{\texorpdfstring{gradient log add one of pos}{gradient log add one of pos}}


Adding a constant to a scalar function leaves its Fréchet derivative unchanged. Apply the gradient's derivative-based definition to obtain \(\nabla(\log r+1)=\nabla\log r\). The explicit differentiability/\allowbreak{}positivity arguments align the statement with the entropy-test calculus.

\subsection{\texorpdfstring{inner gradient entropy test}{inner gradient entropy test}}


Rearrange the positive-density score identity to \(\nabla r=r\nabla\log r\). Pair this with \(\nabla(\log r+1)=\nabla\log r\) and use \(\langle s,s\rangle=\|s\|^2\). The Dirichlet integrand is therefore \(\langle\nabla r,\nabla(\log r+1)\rangle=r\|\nabla\log r\|^2\).

\subsection{\texorpdfstring{entropy production of weak identity}{entropy production of weak identity}}


The strictly positive \(C^1\) density makes \(\psi=\log r+1\) a \(C^1\) test. The entropy derivative integrability is supplied, and the pointwise gradient identity converts its Dirichlet integrability to the supplied weighted Fisher integrability. Substitute this admissible test into the explicit weak equation to obtain \(\int\dot r(\log r+1)d\nu=-\int r\|\nabla\log r\|^2d\nu\).

\subsection{\texorpdfstring{Test the weak equation with a noncompact function}{Test the weak equation with a noncompact function}}


Use the same \(C^1\) entropy test in the forced weak equation. Rewrite the Dirichlet term by \texttt{inner\_\allowbreak{}gradient\_\allowbreak{}entropy\_\allowbreak{}test} and the forcing gradient by \texttt{gradient\_\allowbreak{}log\_\allowbreak{}add\_\allowbreak{}one\_\allowbreak{}of\_\allowbreak{}pos}. The three explicit integrability conditions justify the substitution, giving entropy production \(-I+\int r\langle e,\nabla\log r\rangle d\nu\). Admissibility of this noncompact test remains part of the weak-equation premise.

\section{\texorpdfstring{Spend one quarter of Fisher information on drift error}{Spend one quarter of Fisher information on drift error}}
\label{proof-module-UnadjustedLangevin-LSI-DriftError-lean}
{\small\color{gray}Source module: \texttt{Unadjusted\allowbreak{}Langevin/\allowbreak{}LSI/\allowbreak{}Drift\allowbreak{}Error.\allowbreak{}lean}.\par}


Young’s inequality \(\langle a,b\rangle\le\frac14\|a\|^2+\|b\|^2\) bounds the error pairing in entropy production. Weighted energy integrability yields an actual integral inequality \(H\prime\le-\frac34J+\mathcal E\).

\subsection{\texorpdfstring{inner le quarter norm sq add norm sq}{inner le quarter norm sq add norm sq}}


Cauchy–Schwarz gives \(\langle e,s\rangle\le\|e\|\|s\|\). Expand the nonnegative square \((\|e\|-\|s\|/2)^2\ge0\) to get \(\|e\|\|s\|\le\|e\|^2+\|s\|^2/4\). This is the Young coefficient that preserves three quarters of Fisher information.

\subsection{\texorpdfstring{integrable weighted inner of energy}{integrable weighted inner of energy}}


Multiply the absolute Young bound by the nonnegative weight \(r\). It dominates \(|r\langle e,s\rangle|\) by \(r\|s\|^2/4+r\|e\|^2\). Both energy functions are integrable, so the assumed measurability and \texttt{Integrable.\allowbreak{}mono'} prove integrability of the cross term.

\subsection{\texorpdfstring{entropy production le of forced weak identity}{entropy production le of forced weak identity}}


First derive cross-term integrability from the two finite energies and its measurability. The forced weak identity tested by \(\log r+1\) gives production \(-I+\int r\langle e,s\rangle\). Insert the integrated Young bound to obtain \(-\frac34I+\int r\|e\|^2\). The entropy differential inequality is therefore proved from the stated spatial weak identity.

\subsection{\texorpdfstring{abs inner le quarter norm sq add norm sq}{abs inner le quarter norm sq add norm sq}}


Use the absolute Cauchy–Schwarz inequality \(|\langle e,s\rangle|\le\|e\|\|s\|\), followed by the same square expansion. The absolute-value version controls integrability of the signed cross term before its integral is estimated.

\subsection{\texorpdfstring{integral weighted inner le quarter fisher add energy}{integral weighted inner le quarter fisher add energy}}


Multiply the pointwise Young inequality by \(r\ge0\), then integrate. The supplied integrability permits splitting the upper integral and pulling out \(1/4\). Thus \(\int r\langle e,s\rangle\le\frac14\int r\|s\|^2+\int r\|e\|^2\). No probability normalization of the weight is needed.

\section{\texorpdfstring{Local regularity needed for actual entropy calculus}{Local regularity needed for actual entropy calculus}}
\label{proof-module-UnadjustedLangevin-LSI-Regularity-lean}
{\small\color{gray}Source module: \texttt{Unadjusted\allowbreak{}Langevin/\allowbreak{}LSI/\allowbreak{}Regularity.\allowbreak{}lean}.\par}


A LocalEntropyRegularity package supplies positive densities, time derivatives, and local integrable domination. Its consequences prove finite KL, continuity, derivatives and a production inequality for the identified density entropy; existence of these packages is retained as a premise.

\subsection{\texorpdfstring{Local\allowbreak{}Entropy\allowbreak{}Regularity}{LocalEntropyRegularity}}


At a fixed time \(t\), store a common time neighborhood, an integrable spatial envelope for \(\dot r_u(\log r_u+1)\), local pointwise positivity/\allowbreak{}differentiability, and measurability/\allowbreak{}integrability of the entropy integrands. These are sufficient data for dominated differentiation. The structure contains no derivative of the entropy integral and no entropy-production estimate.

\subsection{\texorpdfstring{Local\allowbreak{}Entropy\allowbreak{}Regularity.\allowbreak{}has\allowbreak{}Deriv\allowbreak{}At entropy}{LocalEntropyRegularity.hasDerivAt entropy}}


Pass the stored neighborhood, envelope, and pointwise derivative fields to \texttt{has\allowbreak{}Deriv\allowbreak{}At\_\allowbreak{}integral\_\allowbreak{}mul\_\allowbreak{}log\_\allowbreak{}of\_\allowbreak{}dominated}. It concludes both integrability of the differentiated integrand and the derivative formula for \(\int r_t\log r_t d\nu\). This is the calculus theorem supplied by the local data, not a field assumed in them.

\subsection{\texorpdfstring{Test an admissible forced weak equation}{Test an admissible forced weak equation}}


Use local dominated differentiation to compute \(H'(t)=\int\dot r_t(\log r_t+1)d\nu\). The forced weak equation, weighted Fisher integrability, and drift-energy integrability then permit the entropy test and integrated Young inequality. The result is \(H'(t)\le-3I(t)/4+E(t)\), with \(E(t)=\int r_t\|e_t\|^2d\nu\). No numerical bound on \(E(t)\) is assumed in this declaration.

\subsection{\texorpdfstring{kl\allowbreak{}Div finite of local\allowbreak{}Regularity}{klDiv finite of localRegularity}}


The local regularity data include integrability of \(r_t\log r_t\). Positivity and \(C^1\) regularity give a measurable nonnegative density. Apply \texttt{kl\allowbreak{}Div\_\allowbreak{}tilted\allowbreak{}Measure\_\allowbreak{}ne\_\allowbreak{}top} to obtain finite actual KL before interpreting its \texttt{to\allowbreak{}Real} value as entropy.

\subsection{\texorpdfstring{density\allowbreak{}Entropy continuous\allowbreak{}On of regular}{densityEntropy continuousOn of regular}}


At every point of the closed interval, dominated differentiation makes the ordinary entropy integral differentiable and hence continuous. Normalization and the finite-entropy density formula identify that integral with actual KL entropy at every interval point. \texttt{Continuous\allowbreak{}On.\allowbreak{}congr} transfers continuity to \(t\mapsto\operatorname{KL}(\rho_t\mid\nu)_{\mathbb R}\).

\subsection{\texorpdfstring{density\allowbreak{}Entropy has\allowbreak{}Deriv\allowbreak{}Within\allowbreak{}At of regular}{densityEntropy hasDerivWithinAt of regular}}


The ordinary entropy integral has the derivative computed by local regularity. Its equality with actual KL is known only for normalized densities in \([0,T]\), so the proof works with right derivatives at \(t\in[0,T)\).



In a sufficiently small right neighborhood, \(u\ge t\ge0\) and \(u<T\), so the density formula identifies both functions there and at \(t\). \texttt{congr\_\allowbreak{}of\_\allowbreak{}eventually\allowbreak{}Eq} transfers the derivative to actual KL. This avoids requiring path normalization outside the interpolation interval.

\subsection{\texorpdfstring{Express the production bound using relative Fisher information}{Express the production bound using relative Fisher information}}


Convert squared-score integrability under \(\rho_t\) into weighted Fisher integrability under \(\nu\) by the with-density criterion. Apply the regular weak-equation derivative estimate, and identify its weighted integral with \texttt{relative\allowbreak{}Fisher}'s real value. Pair that estimate with the actual-KL right derivative from the preceding theorem to get \(H'(t)\le-3I(t)/4+E(t)\) for the identified density law.

\section{\texorpdfstring{Derive entropy calculus from regular interpolation hypotheses}{Derive entropy calculus from regular interpolation hypotheses}}
\label{proof-module-UnadjustedLangevin-LSI-RegularizedTheorem-lean}
{\small\color{gray}Source module: \texttt{Unadjusted\allowbreak{}Langevin/\allowbreak{}LSI/\allowbreak{}Regularized\allowbreak{}Theorem.\allowbreak{}lean}.\par}


This wrapper replaces a supplied entropy derivative estimate with explicit regular density data, a forced weak identity and an error-energy bound. It then invokes the conditional ULA entropy endpoint. Constructing those regular interpolation data remains a displayed premise.



The full actual Gaussian ULA endpoint in UnconditionalTheorem.lean derives all necessary analytic estimates from smoothness, LSI, and finite initial KL. This module retains a reusable interface for explicitly specified interpolation data.

\subsection{\texorpdfstring{Construct the entropy-\allowbreak{}calculus hypotheses and invoke the actual-\allowbreak{}law theorem}{Construct the entropy-calculus hypotheses and invoke the actual-law theorem}}


First use the density/\allowbreak{}law identity to prove probability normalization along each interpolation. Local regularity proves finite actual KL, score integrability proves finite Fisher, and \texttt{density\allowbreak{}Entropy\_\allowbreak{}continuous\allowbreak{}On\_\allowbreak{}of\_\allowbreak{}regular} gives continuity of actual entropy.



Define the candidate derivative \(H'_k(t)=\int\dot r_{k,t}(\log r_{k,t}+1)d\pi_U\). The regular weak-equation theorem proves that this is the actual right derivative and that it is bounded by \(-3I/4+E\). Insert the supplied drift-energy bound \(E\le t^2L^2G_k+2tdL^2\), then invoke \texttt{ula\_\allowbreak{}kl\_\allowbreak{}bound\_\allowbreak{}of\_\allowbreak{}interpolation}.



The final estimate is therefore derived without assuming an entropy derivative, a production inequality, or a discrete recurrence. This wrapper still explicitly requires the regular density construction, its admissible weak equation, its drift-energy estimate, and the entropy-to-gradient moment bound; the wrapper does not itself construct them from LSI.

\section{\texorpdfstring{A reusable regular-\allowbreak{}interpolation interface}{A reusable regular-interpolation interface}}
\label{proof-module-UnadjustedLangevin-LSI-InterpolationData-lean}
{\small\color{gray}Source module: \texttt{Unadjusted\allowbreak{}Langevin/\allowbreak{}LSI/\allowbreak{}Interpolation\allowbreak{}Data.\allowbreak{}lean}.\par}


The RegularInterpolationData structure bundles density paths identifying the actual ULA interpolation, regularity, weak-equation and error controls. Its methods derive path probability, finite KL, the one-step recurrence and the iterated entropy bound. Existence of this package is not asserted by its definition.



The full actual Gaussian ULA endpoint in UnconditionalTheorem.lean derives all necessary analytic estimates from smoothness, LSI, and finite initial KL. This module retains a reusable interface for explicitly specified interpolation data.

\subsection{\texorpdfstring{Specify the actual laws, density regularity and energy data}{Specify the actual laws, density regularity and energy data}}


Package the relative densities \(r_{k,t}\), their time derivatives, and drift fields together with the exact law identity, positivity, spatial regularity, local entropy domination, finite score/\allowbreak{}error energies, admissible weak equation, drift-energy bound, and gradient-moment bound. These remain explicit mathematical inputs. No entropy derivative, entropy recurrence, or convergence estimate is a structure field.

\subsection{\texorpdfstring{Regular\allowbreak{}Interpolation\allowbreak{}Data.\allowbreak{}path probability}{RegularInterpolationData.path probability}}


The data's law field identifies the tilted density with the actual interpolation from \(\rho_k\). Probability preservation proves that interpolation is normalized. Rewrite by the law identity to obtain normalization of each density along the supplied path.

\subsection{\texorpdfstring{Regular\allowbreak{}Interpolation\allowbreak{}Data.\allowbreak{}kl finite}{RegularInterpolationData.kl finite}}


Use local entropy regularity to prove finite KL for each relative density, then apply the time-zero law identification in \texttt{ula\_\allowbreak{}kl\_\allowbreak{}finite\_\allowbreak{}of\_\allowbreak{}interpolation}. The target probability instance is explicit; in the final theorem it is supplied by the LSI hypothesis.

\subsection{\texorpdfstring{Regular\allowbreak{}Interpolation\allowbreak{}Data.\allowbreak{}one step}{RegularInterpolationData.one step}}


Define \(H'_k(t)\) by the differentiated entropy integral. The regular weak-equation theorem proves its right-derivative identity and Fisher/\allowbreak{}drift production bound. Data fields supply probability, KL/\allowbreak{}Fisher finiteness, continuity, numerical error-energy and gradient bounds. Feed these proved ingredients to \texttt{ula\_\allowbreak{}kl\_\allowbreak{}one\_\allowbreak{}step\_\allowbreak{}of\_\allowbreak{}interpolation} to derive the recurrence rather than assuming it in the data.

\subsection{\texorpdfstring{Regular\allowbreak{}Interpolation\allowbreak{}Data.\allowbreak{}kl bound}{RegularInterpolationData.kl bound}}


Unpack every analytic field into \texttt{ula\_\allowbreak{}kl\_\allowbreak{}bound\_\allowbreak{}of\_\allowbreak{}regular\_\allowbreak{}interpolation}, retaining LSI, potential smoothness, and the step regime as separate arguments. Its already proved differentiation, one-step estimate, and geometric iteration give the full displayed KL bound for the actual ULA laws represented by the data.

\section{\texorpdfstring{Conditional entropy convergence for the actual Gaussian particle laws}{Conditional entropy convergence for the actual Gaussian particle laws}}
\label{proof-module-IPLA-LSI-Theorem-lean}
{\small\color{gray}Source module: \texttt{IPLA/\allowbreak{}LSI/\allowbreak{}Theorem.\allowbreak{}lean}.\par}


Let \(V_N(z)=\sum_iU(\theta,\sqrt Nz_i)\), \(D=d_\theta+Nd_x\), and \(\Pi_N\propto e^{-V_N}\). Assume differentiability and \(L\)-Lipschitz joint gradient, an LSI of constant \(\beta>0\) for this actual target, \(\beta\le NL\), and \(0<\gamma\le\beta/(4NL^2)\). Take an initial probability law \(\mu_0\). The endpoint also assumes \texttt{Regular\allowbreak{}Interpolation\allowbreak{}Data} for \((V_N,\beta,\gamma/N,NL,\mu_0)\). These density, weak-equation, energy and moment obligations remain conditional; they have not been constructed from LSI and smoothness alone.

\subsection{\texorpdfstring{The canonical particle laws are the actual Gaussian algorithm laws}{The canonical particle laws are the actual Gaussian algorithm laws}}


Define \(\rho_k=\operatorname{ulaLawStep}(V_N,\gamma/N)^{[k]}\mu_0\); at \(k=0\) this is \(\mu_0\) by the iterate definition. The lifted gradient is \(NL\)-Lipschitz, so the measurable Gaussian transition preserves probability, and induction through the shared ULA law theorem proves \(\rho_k\) is a probability for every \(k\). If \(\mu_0\) is the initial law of the explicit independent-Gaussian IPLA recursion, \texttt{laws\_\allowbreak{}eq\_\allowbreak{}ula\allowbreak{}Law\allowbreak{}Step} identifies it with this same iterate. Thus the entropy result concerns the actual algorithm rather than a substituted reference process.




\textbf{Correspondence with the displayed declarations.}


\begin{itemize}

\item \texttt{particle\allowbreak{}Laws}: The canonical Markov laws of IPLA, using its proved lifted-ULA transition.
\item \texttt{particle\allowbreak{}Laws\_\allowbreak{}is\allowbreak{}Probability}: The actual lifted IPLA transition preserves probability mass.
\item \texttt{particle\allowbreak{}Laws\_\allowbreak{}eq\_\allowbreak{}laws}: These Markov laws agree with the existing noise-driven Gaussian IPLA implementation, including its scalar coordinate laws and independence.
\end{itemize}
\subsection{\texorpdfstring{Density regularity proves finite actual KL before real conversion}{Density regularity proves finite actual KL before real conversion}}


The LSI assumption includes that \(\Pi_N\) is a probability measure. The interpolation data give positive relative densities \(r_k(t)\) with \(r_k(t)\Pi_N\) equal to the actual frozen-drift interpolation law, together with local entropy regularity and spatial smoothness. The shared finiteness theorem applies with \(h=\gamma/N\ge0\) and proves

\[\operatorname{KL}(\rho_k\Vert\Pi_N)<\infty\quad\text{for every }k.\]

The time-zero interpolation identity transfers this density-level KL finiteness to the actual iterate. In the probability-law application, the integrable entropy density also equals that actual finite KL. This guard precedes every use of \texttt{to\allowbreak{}Real}; infinite divergence is never turned into an artificial zero. The proof requires the supplied regularity, not a strong-convexity premise.




\textbf{Correspondence with the displayed declarations.}


\begin{itemize}

\item \texttt{particle\_\allowbreak{}kl\_\allowbreak{}finite}: The identified regular interpolation gives finite actual IPLA relative entropy. Target normalization follows from the lifted LSI assumption.
\end{itemize}
\subsection{\texorpdfstring{Derive entropy production and the one-\allowbreak{}step inequality}{Derive entropy production and the one-step inequality}}


Put \(h=\gamma/N\), \(\ell=NL\), and \(H_k=\operatorname{KL}(\rho_k\Vert\Pi_N)\). The data provide the weak interpolation equation, differentiation and integration regularity, an error field \(e_k(t)\), and the bounds

\[\int r_k(t)\|e_k(t)\|^2d\Pi_N\le t^2\ell^2G_k+2tD\ell^2,\qquad G_k\le\frac{4\ell^2}{\beta}H_k+2D\ell.\]

The shared calculus tests the weak equation against the logarithmic density and uses Young's inequality to derive \(H'(t)\le-3J(t)/4+\int r\|e\|^2d\Pi_N\). LSI gives \(J(t)\ge2\beta H(t)\). Since \(h\le\beta/(4\ell^2)\) and \(\beta\le\ell\), one has \(h\ell\le1/4\) and \(\beta h\le1/4\). The error bounds then imply

\[H'(t)\le-\frac32\beta H(t)+\frac\beta4H_k+\frac52hD\ell^2.\]

The shared affine differential comparison and \(e^{-3x/2}+x/4\le e^{-x}\) for \(0\le x\le1/4\) give \(H_{k+1}\le e^{-\beta h}H_k+6h^2D\ell^2\). The IPLA proof checks positivity of \(L\) from \(0<\beta\le NL\), translates the step restriction, and substitutes the exact lifted dimension and forcing. Therefore

\[H_{k+1}\le e^{-\beta\gamma/N}H_k+6\gamma^2DL^2.\]

Dividing by positive \(N\) gives the normalized recurrence with forcing \(6\gamma^2L^2(d_x+d_\theta/N)\). Neither recurrence is a field of the interpolation structure.




Let \(H_k=\operatorname{KL}(\rho_k\Vert\Pi_N)\). The canonical laws are those of the actual lifted Gaussian Euler kernel, and the law identity relates them to the supplied implementation. The regular interpolation data connect the density at both endpoints of every Euler interval to these same laws, and supply the analytic weak-equation and moment regularity used by the shared entropy theorem. Finite KL is proved before \texttt{to\allowbreak{}Real} is used.



Apply that theorem with \(h=\gamma/N\), \(L_N=NL\), dimension \(D\), and actual lifted LSI constant \(\beta\). The compatibility and positive-step facts are established by \texttt{h\allowbreak{}L}, \texttt{h\allowbreak{}Ls}, \texttt{h\ensuremath{\beta}\allowbreak{}Ls} and \texttt{hscaled}. The result is \[H_{k+1}\le e^{-\beta\gamma/N}H_k+6\gamma^2DL^2.\] The normalized declaration divides this proved inequality by positive \(N\) and substitutes \(D/N=d_x+d_\theta/N\).




\textbf{Correspondence with the displayed declarations.}


\begin{itemize}

\item \texttt{particle\_\allowbreak{}kl\_\allowbreak{}one\_\allowbreak{}step}: The one-step entropy recurrence is derived from the regular weak interpolation equation. The forcing is \texttt{6 * gamma\ensuremath{{}^{2}} * D * L\ensuremath{{}^{2}}} in the full state.
\item \texttt{particle\_\allowbreak{}kl\_\allowbreak{}one\_\allowbreak{}step\_\allowbreak{}normalized}: The actual one-step recurrence per particle has the intensive dimension \texttt{dx + dtheta /\allowbreak{} N}.
\end{itemize}
\subsection{\texorpdfstring{Iterate and simplify the finite-\allowbreak{}entropy bound}{Iterate and simplify the finite-entropy bound}}


Let \(q=e^{-\beta h}\). Repeated substitution in the one-step estimate gives

\[H_k\le q^kH_0+6h^2D\ell^2\sum_{j=0}^{k-1}q^j.\]

Because \(0<\beta h\le1/4\), the elementary bound \(1-e^{-\beta h}\ge3\beta h/4\) bounds the geometric sum by \(4/(3\beta h)\). Thus the shared theorem proves \(H_k\le e^{-\beta hk}H_0+8hD\ell^2/\beta\). Substitute \(h=\gamma/N\), \(\ell=NL\), and \(D=d_\theta+Nd_x\) to obtain

\[H_k\le e^{-\beta\gamma k/N}H_0+\frac{8\gamma ND L^2}{\beta}.\]

The proof invokes the shared bound after deriving every lifted numerical and regularity input, then uses the exact scaling identities. Dividing the entire inequality by positive \(N\) yields

\[\frac{H_k}{N}\le e^{-\beta\gamma k/N}\frac{H_0}{N}+\frac{8\gamma DL^2}{\beta}.\]


Iterating the finite-entropy inequality uses the shared ULA entropy bound, whose geometric sum is controlled with the positive exponent \(\beta h\). After the exact lifted substitutions it yields \[H_k\le e^{-\beta\gamma k/N}H_0+\frac{8\gamma NDL^2}{\beta}.\] The normalized result divides both the initial term and bias by \(N\). The linear-LSI specialization then substitutes the additional lifted hypothesis \(\operatorname{LSI}(\Pi_N,N\alpha)\), producing rate \(e^{-\alpha\gamma k}\) and intensive dimension \(d_x+d_\theta/N\). The Gaussian implementation corollaries use equality of laws; they introduce no new entropy inequality.




\textbf{Correspondence with the displayed declarations.}


\begin{itemize}

\item \texttt{particle\_\allowbreak{}kl\_\allowbreak{}bound}: The shared ULA entropy argument in IPLA coordinates. The LSI constant \texttt{beta} belongs to the full lifted target; in particular it is not assumed to equal \texttt{N} times a pair-target constant.
\item \texttt{particle\_\allowbreak{}kl\_\allowbreak{}bound\_\allowbreak{}normalized}: Dividing the full-state relative entropy by the particle number gives its natural per-particle normalization.
\end{itemize}
\subsection{\texorpdfstring{Specialize only when the lifted target supplies \ensuremath{\beta} = N\ensuremath{\alpha}}{Specialize only when the lifted target supplies β = Nα}}


Assume the actual lifted target has the supplied LSI constant \(N\alpha\), with \(0<\alpha\le L\), and \(0<\gamma\le\alpha/(4L^2)\). Then \(N\alpha>0\), compatibility \(N\alpha\le NL\) holds, and the step equivalence supplies the general theorem's restriction. Its rate cancels to \(e^{-\alpha\gamma k}\); the normalized bias uses \(D/N=d_x+d_\theta/N\). Consequently

\[\frac{H_k}{N}\le e^{-\alpha\gamma k}\frac{H_0}{N}+\frac{8\gamma L^2}{\alpha}\left(d_x+\frac{d_\theta}{N}\right).\]

This is a specialization of a lifted-target LSI already given in the hypotheses. The shared parameter makes the lifted model nonproduct, so no pair-target LSI tensorization is inferred.




\textbf{Correspondence with the displayed declarations.}


\begin{itemize}

\item \texttt{particle\_\allowbreak{}kl\_\allowbreak{}bound\_\allowbreak{}of\_\allowbreak{}linear\_\allowbreak{}lsi}: If the lifted target actually has LSI constant \texttt{N * alpha}, the rate is independent of \texttt{N} and the normalized dimension is \texttt{dx + dtheta /\allowbreak{} N}. This is a specialization of a supplied lifted-target LSI, rather than a tensorization assertion about the original pair model.
\end{itemize}
\subsection{\texorpdfstring{Transfer the canonical bounds to the supplied Gaussian implementation}{Transfer the canonical bounds to the supplied Gaussian implementation}}


For a square-integrable initialization on \(P\) and noises satisfying all scalar Gaussian and independence fields, the law identity gives \(\rho_k=\operatorname{laws}(P,U,\gamma,k)\). Rewriting the normalized general bound by this equality proves the first Gaussian implementation corollary. Rewriting the specialized \(N\alpha\)-LSI bound proves the second one, with all algorithm and interpolation assumptions retained. The entropy on the right is the full initial joint divergence relative to \(\Pi_N\); normalization divides both it and the subsequent joint divergence by \(N\). These conclusions by themselves do not assert parameter concentration about a mode.




\textbf{Correspondence with the displayed declarations.}


\begin{itemize}

\item \texttt{gaussian\_\allowbreak{}ipla\_\allowbreak{}kl\_\allowbreak{}bound\_\allowbreak{}normalized}: The normalized bound for the original independent-Gaussian IPLA implementation. It is obtained by identifying its laws with \texttt{particle\allowbreak{}Laws}.
\item \texttt{gaussian\_\allowbreak{}ipla\_\allowbreak{}kl\_\allowbreak{}bound\_\allowbreak{}of\_\allowbreak{}linear\_\allowbreak{}lsi}: The \texttt{N * alpha} lifted-LSI specialization for the actual Gaussian IPLA implementation, with every algorithm-law assumption retained.
\end{itemize}
\section{\texorpdfstring{Parameter KL by data processing and guarded real conversion}{Parameter KL by data processing and guarded real conversion}}
\label{proof-module-IPLA-LSI-Marginal-lean}
{\small\color{gray}Source module: \texttt{IPLA/\allowbreak{}LSI/\allowbreak{}Marginal.\allowbreak{}lean}.\par}


The parameter reference law is exactly the projection of the actual lifted Gibbs target. Measurable projection decreases extended KL, and finite joint KL makes the real bound transfer valid. This measure-theoretic argument needs no marginal-density or strong-convexity formula.

\subsection{\texorpdfstring{Define the projected target and prove its probability property}{Define the projected target and prove its probability property}}


With \(p=p_\theta\), define \(\Pi_N^\theta=p_\#\Pi_N\). The projection is a continuous linear map, hence measurable. For any probability measure \(\mu\), \((p_\#\mu)(\mathbb R^{d_\theta})=\mu(p^{-1}(\mathbb R^{d_\theta}))=\mu(\mathbb R^D)=1\). This proves probability preservation for any state law and then for the named target whenever \(\Pi_N\) is a probability. In the conditional LSI application that normalization comes from the LSI assumption itself.




\textbf{Correspondence with the displayed declarations.}


\begin{itemize}

\item \texttt{parameter\allowbreak{}Target}: The parameter target is the actual lifted Gibbs law projected to its parameter coordinates. Probability normalization remains an explicit premise.
\item \texttt{parameter\allowbreak{}Law\_\allowbreak{}is\allowbreak{}Probability}: Projecting a probability-valued state law preserves probability mass.
\item \texttt{parameter\allowbreak{}Target\_\allowbreak{}is\allowbreak{}Probability}: The Gibbs parameter target is normalized whenever the actual joint target is normalized. No concentration or convexity hypothesis is introduced.
\end{itemize}
\subsection{\texorpdfstring{Extended KL decreases under the parameter projection}{Extended KL decreases under the parameter projection}}


The shared data-processing theorem for a measurable map yields

\[\operatorname{KL}(p_\#\mu\Vert p_\#\nu)\le\operatorname{KL}(\mu\Vert\nu).\]

At the density level this is convexity of \(u\log u\) under averaging over forgotten coordinates; the measure theorem does not require a constructed conditional density. If the right-hand divergence is finite, the inequality immediately excludes infinity on the left. Since both extended quantities are finite, order is preserved when taking their real values. This proves the real data-processing inequality with its necessary finiteness guard.




The parameter projection is a continuous linear map and hence measurable. Data processing gives the extended inequality \(\operatorname{KL}(p_\#\mu\Vert p_\#\nu)\le\operatorname{KL}(\mu\Vert\nu)\) before any real conversion. A finite joint KL therefore proves finite marginal KL, and only then can \texttt{ENNReal.\allowbreak{}to\allowbreak{}Real\_\allowbreak{}le\_\allowbreak{}to\allowbreak{}Real} transfer the order. This is a divergence statement about the actual pushed-forward target, with no density formula or optimizer concentration assumed.




\textbf{Correspondence with the displayed declarations.}


\begin{itemize}

\item \texttt{parameter\_\allowbreak{}kl\_\allowbreak{}le\_\allowbreak{}joint}: Extended-valued data processing for the common parameter projection.
\item \texttt{parameter\_\allowbreak{}kl\_\allowbreak{}ne\_\allowbreak{}top\_\allowbreak{}of\_\allowbreak{}joint}: Finite joint KL implies finite parameter KL.
\item \texttt{parameter\_\allowbreak{}kl\allowbreak{}Real\_\allowbreak{}le\_\allowbreak{}joint}: The real data-processing inequality, with finiteness of both divergences established before taking their \texttt{to\allowbreak{}Real} values.
\end{itemize}
\subsection{\texorpdfstring{Transfer a full-\allowbreak{}joint or per-\allowbreak{}particle bound with the correct factor}{Transfer a full-joint or per-particle bound with the correct factor}}


If finite full joint KL has real bound \(H\le B\), transitivity with real data processing gives parameter KL at most \(B\). If instead \(H/N\le B\) and \(N>0\), multiplication by the positive count yields \(H\le NB\), hence

\[\operatorname{KL}(p_\#\mu\Vert p_\#\nu)\le NB.\]

The proof uses the positive-denominator order equivalence before applying the full-joint transfer lemma. This factor is necessary because normalized joint entropy is not the same quantity as parameter-marginal entropy.




If the available estimate is intensive, \(H/N\le B\), positive \(N\) gives the full joint bound \(H\le NB\), and data processing then gives marginal KL at most \(NB\). It does not give marginal KL at most \(B\). The initial entropy in the final bound is still full joint entropy: projection decreases initial entropy, so replacing that term by initial parameter entropy would make the right-hand side smaller without a supporting argument.




\textbf{Correspondence with the displayed declarations.}


\begin{itemize}

\item \texttt{parameter\_\allowbreak{}kl\_\allowbreak{}bound\_\allowbreak{}of\_\allowbreak{}joint\_\allowbreak{}bound}: Any finite real bound on full joint KL is also a parameter-marginal bound.
\item \texttt{parameter\_\allowbreak{}kl\_\allowbreak{}bound\_\allowbreak{}of\_\allowbreak{}per\_\allowbreak{}particle\_\allowbreak{}bound}: A per-particle joint entropy bound controls unscaled parameter KL after multiplication by the particle count.
\end{itemize}
\subsection{\texorpdfstring{Specialize all transfers to the actual Gibbs parameter target}{Specialize all transfers to the actual Gibbs parameter target}}


Set \(\nu=\Pi_N\) in the three generic statements. By the definition of \(\Pi_N^\theta\), this yields extended data processing \(\operatorname{KL}(p_\#\mu\Vert\Pi_N^\theta)\le\operatorname{KL}(\mu\Vert\Pi_N)\), finiteness of parameter KL whenever joint KL is finite, and any real upper-bound transfer. The proofs are direct specializations and retain normalization of both laws and the finite-joint guard. No additional concentration or optimizer estimate appears in this measure projection step.




\textbf{Correspondence with the displayed declarations.}


\begin{itemize}

\item \texttt{parameter\_\allowbreak{}kl\_\allowbreak{}le\_\allowbreak{}gibbs\_\allowbreak{}joint}: Extended KL data processing with the actual Gibbs parameter target.
\item \texttt{parameter\_\allowbreak{}kl\_\allowbreak{}ne\_\allowbreak{}top\_\allowbreak{}of\_\allowbreak{}gibbs\_\allowbreak{}joint}: Finite actual Gibbs-relative joint KL gives finite parameter KL.
\item \texttt{parameter\_\allowbreak{}kl\_\allowbreak{}bound\_\allowbreak{}of\_\allowbreak{}gibbs\_\allowbreak{}joint\_\allowbreak{}bound}: A joint KL estimate transfers to the actual Gibbs parameter marginal, without constructing a formula for its marginal density.
\end{itemize}
\section{\texorpdfstring{Why the LSI constant cannot exceed smoothness}{Why the LSI constant cannot exceed smoothness}}
\label{proof-module-UnadjustedLangevin-LSI-Compatibility-lean}
{\small\color{gray}Source module: \texttt{Unadjusted\allowbreak{}Langevin/\allowbreak{}LSI/\allowbreak{}Compatibility.\allowbreak{}lean}.\par}


Write \(\nu(dx)=Z^{-1}e^{-U(x)}dx\), with \(\nabla U\) globally \(L\)-Lipschitz, in a nonzero finite-dimensional real Hilbert space. The original positive \(C^1\) density formulation of LSI implies \(\alpha\le L\). The proof first derives a variance inequality for bounded smooth tests, then combines bounded sinusoidal tests with Gibbs integration by parts and the previously proved score bound. No position moment, transport inequality, Hessian, or strong convexity is assumed.

\subsection{\texorpdfstring{integrable exp mul bounded}{integrable exp mul bounded}}


If \(|f|\le B\), then \(tf\le |t|B\), so \(0<e^{tf}\le e^{|t|B}\). The right side is an integrable constant under every finite measure. Measurability follows from that of \(f\); \texttt{mono'} then proves actual integrability of the exponential. This establishes the whole real line as the domain of the moment generating function used below.

\subsection{\texorpdfstring{exponential test}{exponential test}}


Let \(M(t)=\int e^{tf}\,d\nu\), \(\psi(t)=\log M(t)\), and \(r_t=e^{tf}/M(t)\). Boundedness of \(f\) proves \(0<M(t)<\infty\); the normalized tilt is therefore a probability measure. The proof constructs strict positivity and \(C^1\) regularity of \texttt{r} before testing the original LSI definition.



The logarithmic identity \texttt{hl} is \(\log r_t=tf-\psi(t)\). Integration against the tilted law and Mathlib's identity \(\int f\,r_t\,d\nu=\psi'(t)\) give the actual KL value \(H(r_t\nu\mid\nu)=t\psi'(t)-\psi(t)\); boundedness proves its finiteness first. The gradient identity \texttt{hgrad} gives \(\nabla\log r_t=t\nabla f\). If \(\|\nabla f\|\le D\), its squared norm is at most \(t^2D^2\). This is integrable under the probability tilt, so Fisher information is finite. LSI now yields \(2\alpha(t\psi'(t)-\psi(t))\le t^2D^2\). Both finiteness guards are retained when converting the extended-valued inequality to real numbers.

\subsection{\texorpdfstring{tendsto entropy quotient}{tendsto entropy quotient}}


For a \(C^3\) scalar function \(\psi\) with \(\psi(0)=0\), both numerator and denominator of \((t\psi'(t)-\psi(t))/t^2\) tend to zero as \(t\downarrow0\). The product rule cancels the two copies of \(\psi'(t)\), leaving numerator derivative \(t\psi''(t)\); the denominator derivative is \(2t\). On the punctured right neighborhood, their quotient equals \(\psi''(t)/2\). Continuity of the second derivative and one application of l'Hopital's rule prove the limit \(\psi''(0)/2\). The proof explicitly supplies the nonzero denominator derivative and both zero limits.

\subsection{\texorpdfstring{variance le of bounded}{variance le of bounded}}


For bounded \(f\), its moment and cumulant generating functions are analytic on all real parameters. Mathlib's second derivative formula at zero is \(\psi''(0)=\int f^2\,d\nu-(\int f\,d\nu)^2\), since \(M(0)=1\). Thus no unproved assertion that LSI implies a spectral gap is inserted.



Divide the exponential-test inequality by \(t^2>0\), for \(t>0\), and use \texttt{tendsto\_\allowbreak{}entropy\_\allowbreak{}quotient}. The left side tends to \(\alpha\operatorname{Var}_\nu(f)\), while the right side remains \(D^2\). Passing to the limit gives \(\alpha\operatorname{Var}_\nu(f)\le D^2\). The boundedness and gradient-bound hypotheses here describe this test function; they are not additional assumptions on the target.

\subsection{\texorpdfstring{integral gibbs weighted}{integral gibbs weighted}}


Unfold \(\nu=Z^{-1}(e^{-U}dx)\). Scalar multiplication of a measure multiplies its real Bochner integral by \((Z^{-1})^{\mathrm{real}}\). The with-density integral formula then gives \(\int g\,d\nu=(Z^{-1})^{\mathrm{real}}\int e^{-U}g\,dx\). Continuity of \(U\) makes the weight measurable, and its pointwise value is finite. This equality concerns total integrals; integrability for its later uses is established separately.

\subsection{\texorpdfstring{integrable gibbs weighted}{integrable gibbs weighted}}


Strict positivity of the Gibbs partition follows from differentiability of \(U\); the hypothesis here also excludes an infinite partition. Consequently multiplication of the weighted Lebesgue measure by \(Z^{-1}\) preserves integrability in both directions. The with-density integrability equivalence then gives \(g\in L^1(\nu)\) exactly when \(e^{-U}g\in L^1(dx)\). This is the justification for moving the bounded-test integration-by-parts calculation to Lebesgue measure.

\subsection{\texorpdfstring{gibbs stein bounded}{gibbs stein bounded}}


For bounded \(C^1\) \(\psi\) with bounded gradient and any fixed \(v\), prove \(\mathbb E_\nu[\psi\langle\nabla U,v\rangle]=\mathbb E_\nu\langle\nabla\psi,v\rangle\). The already established score MGF implies integrability of \(\langle\nabla U,v\rangle\). Multiplying it by bounded \(\psi\) preserves integrability. Bounded \(\nabla\psi\) also makes its directional derivative integrable.



The three weighted functions required by Mathlib's global integration-by-parts theorem are therefore integrable: \(\psi e^{-U}\), \((D_v\psi)e^{-U}\), and \(\psi D_v(e^{-U})\). The explicit derivative \texttt{hw} is \(D_v(e^{-U})=-e^{-U}\langle\nabla U,v\rangle\). Integration by parts cancels the two minus signs; multiplying by the same Gibbs normalizer on both sides gives the identity. Boundary vanishing and coordinate moments are not separate hypotheses.

\subsection{\texorpdfstring{integral mul sq le}{integral mul sq le}}


For real \(L^2\) functions \(f,g\), Holder's inequality proves \(fg\in L^1\); their squares are integrable as well. For every real \(c\), integrate \((cf+g)^2\ge0\) to obtain the quadratic \((\int f^2)c^2+2(\int fg)c+\int g^2\ge0\). Its discriminant is nonpositive. Dividing the resulting algebraic inequality by four gives \((\int fg)^2\le(\int f^2)(\int g^2)\). This elementary squared Cauchy-Schwarz form avoids square roots in the compatibility argument.

\subsection{\texorpdfstring{log\allowbreak{}Sobolev constant le smoothness}{logSobolev constant le smoothness}}


Choose \(\|v\|=1\), possible because the space is nontrivial. Put \(s(x)=\langle\nabla U(x),v\rangle\). The score MGF gives \(\mathbb E s^2\le L\); the bounded Stein identity with \(\psi=1\) gives \(\mathbb E s=0\).



For \(t>0\), use the bounded test \(f_t(x)=\sin(t\langle v,x\rangle)\). Its gradient is \(t\cos(t\langle v,x\rangle)v\), so its norm is at most \(|t|\). The bounded-test variance theorem gives \(\alpha\operatorname{Var}(f_t)\le t^2\). Stein's identity and score centering give \(\mathbb E[s(f_t-\mathbb Ef_t)]=t\mathbb E\cos(t\langle v,X\rangle)\). Squared Cauchy-Schwarz now yields \(t^2(\mathbb E\cos(t\langle v,X\rangle))^2\le L\operatorname{Var}(f_t)\).



Multiply by \(\alpha\ge0\), insert the variance bound, and cancel \(t^2>0\). This proves \(\alpha(\mathbb E\cos(t\langle v,X\rangle))^2\le L\). Dominated convergence applies because cosine is bounded by one and tends pointwise to one as \(t\downarrow0\). Its expectation tends to one, proving \(\alpha\le L\). No integrability of \(X\) itself is needed.

\subsection{\texorpdfstring{constant le smoothness}{constant le smoothness}}


This wrapper removes two temporary restrictions from the previous proof. If \(\alpha<0\), then \(\alpha\le0\le L\) immediately. Otherwise, LSI includes the assertion that the Gibbs law is a probability measure. An infinite partition would make its normalizing scalar zero and hence make the Gibbs measure zero, contradicting mass one. Thus the partition is finite internally, and the nonnegative-constant compatibility theorem applies. The final hypotheses are precisely LSI, differentiability, Lipschitz gradient, finite dimension, and a nontrivial state space.

\section{\texorpdfstring{Finite entropy supplies the gradient moment directly}{Finite entropy supplies the gradient moment directly}}
\label{proof-module-UnadjustedLangevin-LSI-EntropyMoment-lean}
{\small\color{gray}Source module: \texttt{Unadjusted\allowbreak{}Langevin/\allowbreak{}LSI/\allowbreak{}Entropy\allowbreak{}Moment.\allowbreak{}lean}.\par}


A finite initial KL divergence already controls the squared Gibbs score when \(\nabla U\) is Lipschitz. Gaussian averaging of the directional score MGF gives \(\mathbb E_\nu e^{\|\nabla U\|^2/(4L)}\le e^{d/2}\). The entropy inequality then proves both integrability and \(G_\mu\le4L H(\mu\mid\nu)+2dL\). This direct argument closes the moment input to ULA without assuming an entropy-cost coupling or a transport inequality.

\subsection{\texorpdfstring{mul le kl\allowbreak{}Fun add exp sub one}{mul le klFun add exp sub one}}


Let \(\Phi(a)=a\log a-a+1\), with its continuous value \(\Phi(0)=1\). For \(a>0\), the elementary inequality \(e^u\ge1+u\), applied to \(u=b-\log a\), gives \(e^b\ge a+ab-a\log a\). Rearranging proves \(ab\le\Phi(a)+e^b-1\). The zero-density case is handled separately: its left side is zero, and its right side is \(e^b\ge0\). Thus this entropy/\allowbreak{}exponential Young inequality remains valid at zero Radon-Nikodym density.

\subsection{\texorpdfstring{integrable of finite kl of integrable exp}{integrable of finite kl of integrable exp}}


Write \(a=d\mu/d\nu\). Finite KL supplies absolute continuity \(\mu\ll\nu\) and integrability of the actual log likelihood; equivalently \(\Phi(a)\) is integrable under \(\nu\). For a measurable nonnegative observable \(f\), the Young inequality gives \(0\le af\le\Phi(a)+e^f\). Both terms on the right are integrable by the finite-KL and exponential-moment hypotheses.



Domination proves \(af\in L^1(\nu)\). The Radon-Nikodym weighted-integrability equivalence then gives \(f\in L^1(\mu)\). In particular, the moment conclusion is derived before any ordinary expectation is manipulated; it is not an extra assumption hidden in the entropy inequality.

\subsection{\texorpdfstring{integral le kl add of integral exp le}{integral le kl add of integral exp le}}


Assume \(\int e^f d\nu\le e^c\). First derive \(f\in L^1(\mu)\) with the preceding theorem. Apply the Young inequality to \(a=d\mu/d\nu\) and \(f-c\), and integrate: \(\int(f-c)d\mu\le\int\Phi(a)d\nu+\int e^{f-c}d\nu-1\).



Probability normalization identifies \(\int\Phi(a)d\nu\) with the genuine finite KL divergence and \(\int(f-c)d\mu=\int f d\mu-c\). The MGF hypothesis makes \(\int e^{f-c}d\nu\le1\), canceling the final two terms. Thus \(\int f d\mu\le H(\mu\mid\nu)+c\), together with the integrability witness used in its derivation. The proof explicitly justifies each weighted integral and subtraction.

\subsection{\texorpdfstring{integrable radial gaussian}{integrable radial gaussian}}


For \(b>0\), the real radial Gaussian \(e^{-b\|v\|^2}\) is integrable on every finite-dimensional real Hilbert space. The proof specializes the complex Gaussian Fourier integrability theorem to real positive precision and takes its norm; the complex norm is exactly the displayed real exponential. This supplies the finite integral needed in the Gaussian linearization argument.

\subsection{\texorpdfstring{lintegral radial gaussian}{lintegral radial gaussian}}


Because the radial Gaussian is nonnegative and integrable, its nonnegative Lebesgue integral equals the embedded ordinary integral. The Gaussian Fourier formula evaluates the latter as \((\pi/b)^{d/2}\). Hence \(\int^- e^{-b\|v\|^2}dv=\operatorname{ofReal}((\pi/b)^{d/2})\), a positive finite normalization factor that can safely be canceled later.

\subsection{\texorpdfstring{lintegral radial gaussian mul exp inner}{lintegral radial gaussian mul exp inner}}


Complete the square in \(-b\|v\|^2+\langle g,v\rangle\): it equals \(\|g\|^2/(4b)-b\|v-g/(2b)\|^2\). Exponentiating separates the constant factor \(e^{\|g\|^2/(4b)}\) from a translated radial Gaussian. Translation invariance of Lebesgue measure and the radial normalization give \(\int^- e^{-b\|v\|^2}e^{\langle g,v\rangle}dv=e^{\|g\|^2/(4b)}(\pi/b)^{d/2}\). The identity is proved for nonnegative integrals before any integrability of a random score is assumed.

\subsection{\texorpdfstring{lintegral exp norm sq le of directional mgf}{lintegral exp norm sq le of directional mgf}}


Suppose \(\mathbb E e^{\langle g,v\rangle}\le e^{L\|v\|^2/2}\) for every direction \(v\). Multiply by \(e^{-L\|v\|^2}\) and integrate in \(v\). Tonelli applies to the nonnegative jointly measurable integrand, so the order of integration can be exchanged without assuming the desired exponential moment in advance.



Gaussian linearization identifies the left side with \((\pi/L)^{d/2}\mathbb E e^{\|g\|^2/(4L)}\). The directional MGF bound controls the right side by \(\int e^{-(L/2)\|v\|^2}dv=(2\pi/L)^{d/2}\). Cancel the positive finite factor \((\pi/L)^{d/2}\) in the extended nonnegative reals to obtain \(\mathbb E e^{\|g\|^2/(4L)}\le2^{d/2}\).

\subsection{\texorpdfstring{gibbs score exponential moment}{gibbs score exponential moment}}


The smooth Gibbs score has the previously proved directional MGF bound \(\mathbb E_\nu e^{\langle\nabla U,v\rangle}\le e^{L\|v\|^2/2}\), obtained from translation invariance and first-order smoothness. Apply Gaussian averaging with \(g=\nabla U\). This gives a finite nonnegative exponential integral bounded by \(2^{d/2}\); measurability and this finiteness prove ordinary integrability.



Convert the nonnegative integral back to its real expectation. Since \(\log2\le1\), \(2^{d/2}\le e^{d/2}\), giving the convenient bound \(\mathbb E_\nu e^{\|\nabla U\|^2/(4L)}\le e^{d/2}\). Only differentiability, Lipschitz gradient, positive \(L\), and finite Gibbs normalization are used; neither convexity nor LSI enters this exponential-moment theorem.

\subsection{\texorpdfstring{gibbs gradient moment bound of finite kl}{gibbs gradient moment bound of finite kl}}


Apply the entropy inequality with \(f(x)=\|\nabla U(x)\|^2/(4L)\) and \(c=d/2\). The Gibbs exponential-moment theorem provides both exponential integrability and the MGF bound; finite initial KL provides the absolute-continuity and entropy inputs. Therefore \(\mathbb E_\mu[\|\nabla U\|^2/(4L)]\le H(\mu\mid\nu)+d/2\), with this observable integrable under \(\mu\).



Multiply by \(4L>0\), proving \(G_\mu\le4L H(\mu\mid\nu)+2dL\) and actual square integrability of the score. The proof also derives finite Gibbs partition from its assumed probability normalization. This replaces the previously missing LSI-to-transport route by a stronger direct entropy/\allowbreak{}score estimate.

\subsection{\texorpdfstring{gibbs gradient moment bound of finite kl of alpha le}{gibbs gradient moment bound of finite kl of alpha le}}


If \(0<\alpha\le L\), then \(4L\le4L^2/\alpha\). Finite KL is nonnegative, so replacing the direct estimate's coefficient \(4L\) by \(4L^2/\alpha\) preserves the inequality. Thus \(G_\mu\le(4L^2/\alpha)H(\mu\mid\nu)+2dL\), the coefficient used in the ULA one-step calculation. Integrability is inherited from the stronger direct theorem. Compatibility \(\alpha\le L\) is proved from LSI in a separate module.

\section{\texorpdfstring{Construct the positive Gaussian mixture and Gibbs-\allowbreak{}relative density}{Construct the positive Gaussian mixture and Gibbs-relative density}}
\label{proof-module-UnadjustedLangevin-LSI-InterpolationDensity-lean}
{\small\color{gray}Source module: \texttt{Unadjusted\allowbreak{}Langevin/\allowbreak{}LSI/\allowbreak{}Interpolation\allowbreak{}Density.\allowbreak{}lean}.\par}


At positive time the actual update has an explicit Gaussian mixture density, even if its initial law is singular. Let \(c_t(x)=x-t\nabla U(x)\) and let \(C_t>0\) be the Gaussian normalizer. The Lebesgue density is \(q_t(y)=\int C_t\exp(-\|y-c_t(x)\|^2/(4t))\,d\mu(x)\). This module proves its law identity, everywhere positivity, pointwise finiteness and spatial continuity. For the normalized Gibbs target, \(r_t(y)=Z e^{U(y)}q_t(y)\) is then a positive continuous relative probability density. GaussianEntropy, GaussianTimeDerivative, GaussianWeakEquation and GaussianEntropyIdentity extend this construction to the full positive-time entropy calculus used by the final theorem.

\subsection{\texorpdfstring{Define the explicit Gaussian mixture before taking real values}{Define the explicit Gaussian mixture before taking real values}}


For the state-dependent center \(c_t(x)=x-t\nabla U(x)\), let

\[K_t(x,y)=C_t\exp\!\left(-\frac{\|y-c_t(x)\|^2}{4t}\right),\qquad q_t^{\mathrm{ext}}(y)=\int K_t(x,y)\,d\mu(x).\]

\texttt{interpolation\allowbreak{}Density} is this mixture as an extended nonnegative integral: Lean integrates \texttt{ULA.\allowbreak{}gaussian\allowbreak{}Proposal\allowbreak{}Density t (x -\allowbreak{} t \ensuremath{\bullet} gradient U x) y}. The normalizer \(C_t\) is the inverse of the actual radial Gaussian integral, already identified with the standard-noise proposal law in \texttt{Gaussian\allowbreak{}Proposal\allowbreak{}Density.\allowbreak{}lean}.



The definition is total in time and potential, but the Gaussian density interpretation is used only for \(t>0\). Keeping the integral in \(\mathbb R_{\ge0}\cup\{\infty\}\) makes nonnegativity immediate and preserves a genuine finiteness obligation before taking a real density. The later normalizer bound discharges that obligation at every proposed state.

\subsection{\texorpdfstring{Apply measurable parameter integration to the mixture}{Apply measurable parameter integration to the mixture}}


If \(\nabla U\) is Lipschitz, its state-dependent center \(c_t\) is continuous and measurable. The joint Gaussian integrand \((x,y)\mapsto K_t(x,y)\) is therefore measurable: subtraction, scalar multiplication, squared norm and exponential preserve the required measurability.



The theorem \texttt{Measurable.\allowbreak{}lintegral\_\allowbreak{}prod\_\allowbreak{}left} now proves that \(y\mapsto q_t^{\mathrm{ext}}(y)\) is measurable. Lean unfolds only the Gaussian proposal density, records gradient measurability as \texttt{hg}, and lets \texttt{fun\_\allowbreak{}prop} assemble the elementary operations. The mixing probability measure supplies the finite, hence sigma-finite, hypothesis for parameter integration. No positive-time assumption is needed for this measurability statement.

\subsection{\texorpdfstring{Prove the Gaussian mixture is the density of the actual update}{Prove the Gaussian mixture is the density of the actual update}}


For every \(t>0\),

\[\rho_t=(Y_t)_\#(\mu\otimes\mathcal N(0,I))=q_t^{\mathrm{ext}}\,dy.\]

To prove equality of measures, fix a measurable set \(A\). The pushforward formula and Tonelli express the actual proposal probability as

\[\rho_t(A)=\int \mathbb P\bigl(c_t(x)+\sqrt{2t}Z\in A\bigr)\,d\mu(x).\]

For each fixed \(x\), \texttt{ULA.\allowbreak{}gaussian\_\allowbreak{}proposal\_\allowbreak{}density} already identifies the affine Gaussian proposal with density \(K_t(x,\cdot)\). The local identity \texttt{hinner} therefore changes the inner probability to \(\int_AK_t(x,y)\,dy\). The joint integrand is measurable, as proved by \texttt{hd}, so \texttt{lintegral\_\allowbreak{}lintegral\_\allowbreak{}swap} interchanges the two nonnegative integrals:

\[\rho_t(A)=\int_A\left(\int K_t(x,y)\,d\mu(x)\right)dy=\int_Aq_t^{\mathrm{ext}}(y)\,dy.\]

The Lean proof mirrors these three steps with \texttt{Measure.\allowbreak{}map\_\allowbreak{}apply}, \texttt{Measure.\allowbreak{}prod\_\allowbreak{}apply}, the fixed-state density identity, and Tonelli under the restricted volume measure on \(A\). It assumes no density for \(\mu\).

\subsection{\texorpdfstring{Use positivity of every Gaussian component}{Use positivity of every Gaussian component}}


At every positive time and every point \(y\), the mixture density is strictly positive:

\[q_t^{\mathrm{ext}}(y)>0.\]

The normalizer is positive by \texttt{ULA.\allowbreak{}gaussian\allowbreak{}Proposal\allowbreak{}Normalizer\_\allowbreak{}pos}, and the real exponential is positive. Thus the component density \(K_t(x,y)\) is positive for every initial point \(x\). Lean proves this as \texttt{hqpos}, using positivity of \texttt{ENNReal.\allowbreak{}of\allowbreak{}Real}.



The support of \(x\mapsto K_t(x,y)\) is consequently all of \(E\); this is \texttt{hsupp}. Its measurability is \texttt{hq}. The characterization \texttt{lintegral\_\allowbreak{}pos\_\allowbreak{}iff\_\allowbreak{}support} reduces positivity of its integral to positive mass of that support. Since \(\mu(E)=1\), the integral is positive. The assertion holds pointwise in \(y\), not only almost everywhere.

\subsection{\texorpdfstring{Bound every mixture value by the common Gaussian peak}{Bound every mixture value by the common Gaussian peak}}


For \(t>0\), the Gaussian exponent is nonpositive, hence

\[0\le K_t(x,y)\le C_t,\qquad q_t^{\mathrm{ext}}(y)\le C_t.\]

Lean proves the exponential bound with \texttt{Real.\allowbreak{}exp\_\allowbreak{}le\_\allowbreak{}one\_\allowbreak{}iff}: the numerator \(-\|y-c_t(x)\|^2\) is nonpositive and the denominator \(4t\) is positive. Multiplication by the nonnegative normalizer preserves the inequality, as does \texttt{ENNReal.\allowbreak{}of\allowbreak{}Real}.



\texttt{lintegral\_\allowbreak{}mono} integrates this pointwise bound. The integral of the constant \(C_t\) equals \(C_t\mu(E)=C_t\) because \(\mu\) is a probability measure. This theorem requires no regularity of the gradient: its bound concerns the nonnegative integrand at every point, and the upper nonnegative integral remains monotone even before measurability is established.

\subsection{\texorpdfstring{Derive pointwise finiteness before real conversion}{Derive pointwise finiteness before real conversion}}


The preceding bound gives \(q_t^{\mathrm{ext}}(y)\le\operatorname{ofReal}(C_t)<\infty\) for each \(t>0\) and \(y\). The extended density therefore never equals infinity.



The Lean proof is the order argument \texttt{ne\_\allowbreak{}top\_\allowbreak{}of\_\allowbreak{}le\_\allowbreak{}ne\_\allowbreak{}top}, applied to the universal fact \texttt{ENNReal.\allowbreak{}of\allowbreak{}Real\_\allowbreak{}ne\_\allowbreak{}top} and \texttt{interpolation\allowbreak{}Density\_\allowbreak{}le\_\allowbreak{}normalizer}. This small theorem is the explicit finiteness witness used by every later conversion from the extended mixture density to an ordinary real function. It prevents the total operation \texttt{to\allowbreak{}Real} from silently turning an infinite density into zero.

\subsection{\texorpdfstring{Choose the finite real representative of the mixture}{Choose the finite real representative of the mixture}}


The real density is defined by

\[q_t(y)=\bigl(q_t^{\mathrm{ext}}(y)\bigr).\mathrm{toReal}.\]

This is \texttt{interpolation\allowbreak{}Lebesgue\allowbreak{}Density U \ensuremath{\mu} t y}. The definition itself is total; its positive-time interpretation uses the pointwise finiteness theorem above. Once \(q_t^{\mathrm{ext}}(y)\ne\infty\) is supplied, converting \(q_t(y)\) back with \texttt{ENNReal.\allowbreak{}of\allowbreak{}Real} recovers the original extended density exactly. The subsequent law identity, positivity argument, ordinary integral formula and continuity theorem all retain the positive-time hypotheses that justify this interpretation.

\subsection{\texorpdfstring{Identify the actual law with the finite real density}{Identify the actual law with the finite real density}}


For \(t>0\), the actual proposal law satisfies

\[\rho_t(A)=\int_Aq_t(y)\,dy\quad\text{for every measurable }A.\]

The extended-density law identity is already proved. The present proof changes only its integrand representative: \texttt{ENNReal.\allowbreak{}of\allowbreak{}Real\_\allowbreak{}to\allowbreak{}Real} and \texttt{interpolation\allowbreak{}Density\_\allowbreak{}ne\_\allowbreak{}top} prove \(\operatorname{ofReal}(q_t(y))=q_t^{\mathrm{ext}}(y)\) for every \(y\). Function extensionality then gives equality of the two density functions and hence of their \texttt{with\allowbreak{}Density} measures.



The separate finiteness theorem is invoked at the exact point where it is needed. This is why the real-valued representative still denotes the genuine Gaussian mixture, including from an initial law that has no Lebesgue density.

\subsection{\texorpdfstring{Deduce absolute continuity for every positive time}{Deduce absolute continuity for every positive time}}


The law identity \(\rho_t=q_t^{\mathrm{ext}}dy\) immediately implies \(\rho_t\ll dy\): every Lebesgue-null set has zero integral against the density and thus zero proposal probability.



Lean rewrites the actual \texttt{interpolation\allowbreak{}Law} using \texttt{interpolation\allowbreak{}Law\_\allowbreak{}eq\_\allowbreak{}with\allowbreak{}Density}, then applies the general theorem \texttt{with\allowbreak{}Density\_\allowbreak{}absolutely\allowbreak{}Continuous}. The result holds for any initial probability law when the gradient is Lipschitz. In particular, positive-time Gaussian smoothing produces absolute continuity from a Dirac or other singular initial law. This conclusion alone does not imply finite KL relative to a particular target, nor entropy continuity at the initial time.

\subsection{\texorpdfstring{Convert the nonnegative mixture to its ordinary real integral}{Convert the nonnegative mixture to its ordinary real integral}}


The real density has the expected ordinary integral formula:

\[q_t(y)=\int C_t\exp\!\left(-\frac{\|y-c_t(x)\|^2}{4t}\right)d\mu(x).\]

The proof records nonnegativity of the real integrand as \texttt{hnonneg} and the uniform bound by \(C_t\) as \texttt{hbound}. Because \(\mu\) is a probability measure, the constant majorant is integrable. Continuity of the Lipschitz gradient supplies measurability, and \texttt{Integrable.\allowbreak{}mono'} proves \texttt{hint}, integrability of the displayed real Gaussian integrand.



\texttt{of\allowbreak{}Real\_\allowbreak{}integral\_\allowbreak{}eq\_\allowbreak{}lintegral\_\allowbreak{}of\allowbreak{}Real} can now identify its ordinary integral with the extended mixture integral after applying \texttt{of\allowbreak{}Real}. Taking \texttt{to\allowbreak{}Real} reverses this conversion because the real integral is nonnegative. This proves the formula without treating a nonintegrable Bochner integral as an ordinary expectation.

\subsection{\texorpdfstring{Pass spatial continuity through the probability mixture}{Pass spatial continuity through the probability mixture}}


At a fixed positive time, \(q_t:E\to\mathbb R\) is continuous. For every fixed initial state \(x\), the Gaussian kernel \(y\mapsto K_t(x,y)\) is continuous, and its absolute value is bounded by the same constant \(C_t\), independent of \(x\) and \(y\).



The proof first replaces \texttt{interpolation\allowbreak{}Lebesgue\allowbreak{}Density} by the ordinary integral formula from the previous theorem. It applies \texttt{continuous\_\allowbreak{}of\_\allowbreak{}dominated} with constant majorant \(x\mapsto C_t\). Measurability in \(x\) follows from gradient continuity. The nonpositive Gaussian exponent proves \(\|K_t(x,y)\|=K_t(x,y)\le C_t\); probability normalization makes this majorant integrable. Pointwise continuity in \(y\) is then assembled from subtraction, norm, square and exponential.



These are exactly the four inputs of dominated convergence for continuous parameters. The result proves spatial continuity of the actual mixture, while GaussianEntropy.lean and GaussianTimeDerivative.lean prove the corresponding spatial and time derivatives.

\subsection{\texorpdfstring{Derive finite partition from target normalization}{Derive finite partition from target normalization}}


Recall \(\nu=Z^{-1}(e^{-U}dy)\) with the extended partition function \(Z=\int e^{-U}dy\). If this measure is a probability measure, then \(Z\ne\infty\).



Indeed, \(Z=\infty\) would give \(Z^{-1}=0\), and hence \(\nu=0\). Its mass on the whole space would then be zero, contradicting \(\nu(E)=1\). Lean names the hypothetical zero-measure equality \texttt{hzero}, rewrites the probability mass identity \texttt{hmass}, and closes the contradiction by arithmetic.



No strong convexity or LSI estimate is used. This matters because LSI already includes probability normalization of its target, so its finite partition requirement can be recovered from that normalization rather than imported through a stronger convexity hypothesis.

\subsection{\texorpdfstring{Write the explicit density relative to Gibbs}{Write the explicit density relative to Gibbs}}


The normalized Gibbs target has Lebesgue density \(Z^{-1}e^{-U(y)}\). Therefore the relative density corresponding to the mixture \(q_t(y)dy\) is

\[r_t(y)=Z e^{U(y)}q_t(y).\]

\texttt{interpolation\allowbreak{}Relative\allowbreak{}Density} is exactly this product, with \(Z\) represented by \texttt{(ULA.\allowbreak{}gibbs\allowbreak{}Partition U).\allowbreak{}to\allowbreak{}Real} and \(q_t\) by \texttt{interpolation\allowbreak{}Lebesgue\allowbreak{}Density}. The later theorems separately prove finite partition, positivity and the density-law identity before using this expression as a normalized relative density.



The definition makes the coordinate conversion explicit: the ratio to the target is obtained by multiplying the Lebesgue density by the reciprocal target density. It does not package an entropy derivative, convergence bound or weak PDE as part of the meaning of a density.

\subsection{\texorpdfstring{Multiply measurable factors to obtain relative-\allowbreak{}density measurability}{Multiply measurable factors to obtain relative-density measurability}}


Differentiability of \(U\) implies continuity and measurability of \(e^U\). The extended Gaussian mixture is measurable by \texttt{interpolation\allowbreak{}Density\_\allowbreak{}measurable}, and its real conversion remains measurable by \texttt{ennreal\_\allowbreak{}to\allowbreak{}Real}. Multiplication by the constant \(Z^{\mathrm{toReal}}\) therefore proves measurability of \(r_t\).



The Lean expression follows these factors in order: \texttt{measurable\_\allowbreak{}const.\allowbreak{}mul} combines the partition scalar with \texttt{Real.\allowbreak{}measurable\_\allowbreak{}exp.\allowbreak{}comp h\allowbreak{}U.\allowbreak{}continuous.\allowbreak{}measurable}, then the final \texttt{.\allowbreak{}mul} inserts the real Gaussian mixture. This statement concerns measurability only, so it does not need a positive-time assumption or a finiteness conversion.

\subsection{\texorpdfstring{Obtain relative-\allowbreak{}density continuity at positive time}{Obtain relative-density continuity at positive time}}


For \(t>0\), spatial continuity of \(q_t\) and continuity of \(e^U\) give continuity of \(r_t(y)=Z e^{U(y)}q_t(y)\).



The proof assembles three continuous factors directly: the constant partition scalar, \texttt{Real.\allowbreak{}continuous\_\allowbreak{}exp.\allowbreak{}comp h\allowbreak{}U.\allowbreak{}continuous}, and \texttt{interpolation\allowbreak{}Lebesgue\allowbreak{}Density\_\allowbreak{}continuous}. Its two applications of \texttt{Continuous.\allowbreak{}mul} mirror the two products in the definition. The conclusion is a property of the explicit actual-law density. It supplies continuity, and GaussianEntropy.lean proves the C\ensuremath{{}^{1}} regularity needed to apply the declared LSI.

\subsection{\texorpdfstring{Retain all finiteness guards when proving strict positivity}{Retain all finiteness guards when proving strict positivity}}


When the target Gibbs measure is a probability measure and \(t>0\), all three factors of \(r_t(y)=Z e^{U(y)}q_t(y)\) are strictly positive. The partition is positive by \texttt{ULA.\allowbreak{}gibbs\allowbreak{}Partition\_\allowbreak{}pos} and finite by \texttt{gibbs\allowbreak{}Partition\_\allowbreak{}ne\_\allowbreak{}top\_\allowbreak{}of\_\allowbreak{}probability}. The mixture is positive by \texttt{interpolation\allowbreak{}Density\_\allowbreak{}pos} and finite by \texttt{interpolation\allowbreak{}Density\_\allowbreak{}ne\_\allowbreak{}top}.



For each of these two extended-valued quantities, \texttt{ENNReal.\allowbreak{}to\allowbreak{}Real\_\allowbreak{}pos} requires both nonzero value and finiteness. The Lean witnesses \texttt{h\allowbreak{}Zpos}, \texttt{h\allowbreak{}Zfin}, \texttt{hqpos} and \texttt{hqfin} provide exactly those hypotheses. The ordinary exponential is positive independently. Two applications of \texttt{mul\_\allowbreak{}pos} prove strict positivity of the relative density at every spatial point.

\subsection{\texorpdfstring{Cancel the target density and identify the exact relative law}{Cancel the target density and identify the exact relative law}}


The constructed relative density satisfies

\[r_t\nu=q_tdy=\rho_t\qquad(t>0).\]

The proof first records positive finite partition and measurability of the Gibbs weight and relative density. The actual-law mixture theorem changes \(\rho_t\) to its Lebesgue density. On the left, \texttt{with\allowbreak{}Density\_\allowbreak{}smul\_\allowbreak{}measure}, \texttt{with\allowbreak{}Density\_\allowbreak{}mul} and \texttt{with\allowbreak{}Density\_\allowbreak{}smul'} combine multiplication by the relative density with multiplication by the Gibbs density into one density with respect to volume.



The remaining equality is pointwise. The real-to-extended conversion \texttt{hrel} proves

\[\operatorname{ofReal}(r_t(y))=Z\,\operatorname{ofReal}(e^{U(y)})\,q_t^{\mathrm{ext}}(y).\]

Its uses of \texttt{of\allowbreak{}Real\_\allowbreak{}to\allowbreak{}Real} invoke the finite partition and pointwise finite mixture bounds. Next \texttt{hexp} proves \(e^{-U(y)}e^{U(y)}=1\) in the extended nonnegative representation. The full density product is rearranged into \((Z^{-1}Z)(e^{-U(y)}e^{U(y)})q_t^{\mathrm{ext}}(y)\). Both cancellation steps are justified explicitly, leaving exactly the mixture density.



This theorem constructs law identification at positive times; no density is supplied as an assumption. Spatial \(C^1\) regularity, time derivatives, entropy domination and the entropy weak equation are not conclusions of this algebraic density identity.

\subsection{\texorpdfstring{Transfer normalization from the actual proposal law}{Transfer normalization from the actual proposal law}}


The measure induced by the constructed relative density is a probability measure:

\[\int r_t\,d\nu=1\qquad(t>0).\]

Lean expresses this as \texttt{Is\allowbreak{}Probability\allowbreak{}Measure (tilted\allowbreak{}Measure \ensuremath{\nu} r\_\allowbreak{}t)}. The proof rewrites that measure with \texttt{tilted\allowbreak{}Measure\_\allowbreak{}interpolation\allowbreak{}Relative\allowbreak{}Density}, reducing normalization to the already proved probability preservation of the actual independent-Gaussian update. \texttt{interpolation\allowbreak{}Law\_\allowbreak{}is\allowbreak{}Probability} supplies the result from probability of the initial law and measurability of the Lipschitz-gradient proposal.



Thus positivity, continuity, law identification and normalization all refer to the same explicitly constructed density. Normalization alone does not prove finite entropy or Fisher information, and GaussianEntropy.lean derives those finiteness properties for the actual positive-time interpolation.

\section{\texorpdfstring{Gaussian displacement and frozen-\allowbreak{}gradient error energy}{Gaussian displacement and frozen-gradient error energy}}
\label{proof-module-UnadjustedLangevin-LSI-GaussianDrift-lean}
{\small\color{gray}Source module: \texttt{Unadjusted\allowbreak{}Langevin/\allowbreak{}LSI/\allowbreak{}Gaussian\allowbreak{}Drift.\allowbreak{}lean}.\par}


For \(Y_t=X-t\nabla U(X)+\sqrt{2t}Z\), independent Gaussian noise makes the cross term vanish, giving \(\mathbb E\|Y_t-X\|^2=t^2\mathbb E\|\nabla U(X)\|^2+2td\). Gradient Lipschitzness bounds the corresponding frozen-gradient error by \(L^2\) times this moment.

\subsection{\texorpdfstring{integral norm sq std\allowbreak{}Gaussian}{integral norm sq stdGaussian}}


Choose a finite orthonormal basis and use \(\|Z\|^2=\sum_i\langle e_i,Z\rangle^2\). Each standard Gaussian coordinate has second moment one, proved using its dual variance formula and zero mean. Gaussian \(L^2\) membership justifies integrating the finite sum, giving \(\mathbb E\|Z\|^2=d\), also when \(d=0\).

\subsection{\texorpdfstring{integral norm sq affine std\allowbreak{}Gaussian}{integral norm sq affine stdGaussian}}


Expand \(\|ag(X)+bZ\|^2=a^2\|g(X)\|^2+2ab\langle g(X),Z\rangle+b^2\|Z\|^2\). The \(L^2\) hypothesis on \(g\) and Gaussian \(L^2\) membership prove all three terms integrable on the independent product.



The product bilinear-integral theorem factors the cross expectation as \(\langle\mathbb Eg(X),\mathbb EZ\rangle=0\). The product marginals identify the first energy, and \texttt{integral\_\allowbreak{}norm\_\allowbreak{}sq\_\allowbreak{}std\allowbreak{}Gaussian} identifies the noise energy. Hence the integral is \(a^2\mathbb E\|g\|^2+b^2d\).

\subsection{\texorpdfstring{mala\allowbreak{}Proposal sub initial}{malaProposal sub initial}}


Subtract \(x\) from \(T_t(x,z)=x-t\nabla U(x)+\sqrt{2t}z\). Additive cancellation leaves \(-t\nabla U(x)+\sqrt{2t}z\). This algebraic identity shows why a state second moment is unnecessary when only displacement energy is estimated.

\subsection{\texorpdfstring{frozen displacement second moment}{frozen displacement second moment}}


Apply the affine Gaussian energy formula with \(a=-t\), \(b=\sqrt{2t}\), and \(g=\nabla U\). For \(t\ge0\), \(b^2=2t\), yielding \(\mathbb E\|T_t-X\|^2=t^2\mathbb E\|\nabla U(X)\|^2+2td\). Independence and centering have already discharged the cross term.

\subsection{\texorpdfstring{frozen gradient error energy}{frozen gradient error energy}}


Lipschitzness gives \(\|\nabla U(T_t)-\nabla U(X)\|^2\le L^2\|T_t-X\|^2\) pointwise. Convert integrability of the squared gradient to \(L^2\) membership, so the displacement has an integrable squared norm.



Continuity of the proposal and gradient proves measurability of the squared discrepancy. Domination by the integrable displacement energy proves its integrability; \texttt{integral\_\allowbreak{}mono} and the exact displacement moment then give \(L^2(t^2G_0+2td)\). This is a joint-law estimate; a state-space conditional drift is constructed separately.

\subsection{\texorpdfstring{frozen displacement mem\allowbreak{}Lp}{frozen displacement memLp}}


Pull the square-integrable initial gradient to the first product coordinate and the square-integrable Gaussian identity to the second. Their scalar multiples and sum lie in \(L^2\). The displacement identity rewrites this sum as \(T_t(X,Z)-X\).

\section{\texorpdfstring{Construct the conditional Gaussian drift and transfer its energy}{Construct the conditional Gaussian drift and transfer its energy}}
\label{proof-module-UnadjustedLangevin-LSI-ConditionalDrift-lean}
{\small\color{gray}Source module: \texttt{Unadjusted\allowbreak{}Langevin/\allowbreak{}LSI/\allowbreak{}Conditional\allowbreak{}Drift.\allowbreak{}lean}.\par}


The actual frozen Gaussian interpolation is \(Y_t=X-t\nabla U(X)+\sqrt{2t}Z\), where \((X,Z)\) has the independent law \(P=\mu\otimes\mathcal N(0,I)\). Write \(D_t=\nabla U(Y_t)-\nabla U(X)\). This module constructs a measurable field \(e_t\) satisfying \(e_t(Y_t)=\mathbb E[D_t\mid Y_t]\), proves its flux identities, and derives its squared-energy bound by conditional Jensen. Thus the state-dependent error is a consequence of the joint Gaussian construction. The entropy weak equation is a separate analytic obligation.

\subsection{\texorpdfstring{Construct a conditional mean on the observed state space}{Construct a conditional mean on the observed state space}}


Let \(Y:\Omega\to S\) be measurable, let \(P\) be a probability measure, and let the Banach-valued vector \(D\) satisfy \(\mathbb E_P\|D\|^2<\infty\). The theorem constructs a strongly measurable \(e:S\to F\) such that

\[e(Y)=\mathbb E_P[D\mid\sigma(Y)]\quad P\text{-almost surely},\qquad \int_S\|e(y)\|^2\,d(Y_\#P)(y)\le\int_\Omega\|D\|^2\,dP.\]

The proof first applies Doob–Dynkin factorization to the strongly measurable conditional expectation. In Lean, \texttt{strongly\allowbreak{}Measurable\_\allowbreak{}cond\allowbreak{}Exp} and \texttt{exists\_\allowbreak{}eq\_\allowbreak{}measurable\_\allowbreak{}comp} produce \texttt{e}, its strong measurability \texttt{he}, and the exact factorization \texttt{hfactor}. The equality \texttt{hcomp} writes that factorization in the pointwise function form needed for later rewrites.



Next \texttt{h\allowbreak{}D.\allowbreak{}cond\allowbreak{}Exp} proves that the conditional expectation is in \(L^2(P)\). Its squared norm is therefore integrable. The pushforward integrability theorem converts this to integrability of \(\|e\|^2\) under \(Y_\#P\); this is \texttt{hint}. Finally, \texttt{integral\_\allowbreak{}map\_\allowbreak{}of\_\allowbreak{}strongly\allowbreak{}Measurable} identifies the left energy with \(\int\|\mathbb E[D\mid\sigma(Y)]\|^2dP\). The conditional Jensen theorem \texttt{integral\_\allowbreak{}norm\_\allowbreak{}cond\allowbreak{}Exp\_\allowbreak{}rpow\_\allowbreak{}le}, with exponent \(2\), bounds that integral by \(\int\|D\|^2dP\). The real-power notation in that library theorem is converted to the ordinary square by \texttt{Real.\allowbreak{}rpow\_\allowbreak{}two}.

\subsection{\texorpdfstring{Identify the vector flux on every measurable set}{Identify the vector flux on every measurable set}}


Suppose a strongly measurable field \(e\) satisfies the conditional-mean identity \(e(Y)=\mathbb E[D\mid\sigma(Y)]\), with \(D\) integrable. For every measurable \(A\subseteq S\),

\[\int_A e(y)\,d(Y_\#P)(y)=\int_{Y^{-1}(A)}D(\omega)\,dP(\omega).\]

The first step, \texttt{set\allowbreak{}Integral\_\allowbreak{}map}, pulls the left integral back to \(\int_{Y^{-1}(A)}e(Y)dP\). The almost-sure conditional-mean identity changes its integrand to the conditional expectation. Because \(Y^{-1}(A)\) belongs to the pulled-back sigma algebra, \texttt{set\allowbreak{}Integral\_\allowbreak{}cond\allowbreak{}Exp} replaces that expectation by \(D\) in the set integral. The probability assumption guarantees the finite-measure hypotheses for conditional expectation. This proves a vector-valued identification of the flux, extending beyond the energy estimate.

\subsection{\texorpdfstring{Transfer integrable inner-\allowbreak{}product tests through the conditional mean}{Transfer integrable inner-product tests through the conditional mean}}


Let \(\varphi:S\to F\) be strongly measurable and suppose \(\langle D,\varphi(Y)\rangle\) is integrable. Then

\[\int_S\langle e(y),\varphi(y)\rangle\,d(Y_\#P)(y)=\int_\Omega\langle D,\varphi(Y)\rangle\,dP.\]

This is the flux identity needed when the test depends on the current interpolated state. The proof gives \(Y\) its pulled-back sigma algebra and observes that \(\varphi\circ Y\) is strongly measurable for that algebra; the Lean witnesses are \texttt{h\allowbreak{}Ypull} and \texttt{h\ensuremath{\varphi}.\allowbreak{}comp\_\allowbreak{}measurable h\allowbreak{}Ypull}. Conditional expectation can therefore pull this factor out of the continuous bilinear inner product:

\[\mathbb E[\langle D,\varphi(Y)\rangle\mid Y]=\langle\mathbb E[D\mid Y],\varphi(Y)\rangle.\]

\texttt{cond\allowbreak{}Exp\_\allowbreak{}bilin\_\allowbreak{}of\_\allowbreak{}strongly\allowbreak{}Measurable\_\allowbreak{}right} proves this equality as \texttt{hpull}, using \texttt{inner\allowbreak{}SL \ensuremath{\mathbb R}} as the continuous bilinear map. The pushforward integral is pulled back, \texttt{hcond} substitutes the conditional expectation, and integrating \texttt{hpull} removes the conditioning. The statement allows noncompact tests whenever its displayed product is integrable; it does not prove a Fokker–Planck equation or the integrability of an arbitrary entropy test.

\subsection{\texorpdfstring{Specify the discrepancy on the actual independent joint law}{Specify the discrepancy on the actual independent joint law}}


For \(p=(x,z)\), define

\[D_t(p)=\nabla U\!\left(x-t\nabla U(x)+\sqrt{2t}\,z\right)-\nabla U(x).\]

This is exactly \texttt{frozen\allowbreak{}Gradient\allowbreak{}Discrepancy U t p}: \texttt{ULA.\allowbreak{}mala\allowbreak{}Proposal} supplies the interpolated state and the two gradient terms are subtracted in the displayed order. Its sign agrees with the additional relative-density drift \(e_t(y)=\nabla U(y)-\mathbb E[\nabla U(X)\mid Y_t=y]\). No conditional expectation, density, or error bound is part of this definition. For \(t\ge0\), the later energy theorem uses the actual independent Gaussian product law.

\subsection{\texorpdfstring{Record precisely what is observed at interpolation time}{Record precisely what is observed at interpolation time}}


The observed variable is \(Y_t(x,z)=x-t\nabla U(x)+\sqrt{2t}z\). The sigma algebra \(\sigma(Y_t)\) consists of preimages \(Y_t^{-1}(A)\) of measurable subsets of the state space. Lean writes it as \texttt{Measurable\allowbreak{}Space.\allowbreak{}comap} of the proposal map.



The abbreviation \texttt{frozen\allowbreak{}Drift\allowbreak{}Sigma\allowbreak{}Algebra U t} makes the conditional expectation statements refer to this exact observation. Conditioning retains the interpolated position and averages out information about the starting point and Gaussian innovation that the position alone does not determine. Lipschitzness of the gradient later proves that this pulled-back sigma algebra lies inside the ambient measurable structure.

\subsection{\texorpdfstring{Construct the actual interpolation error and prove its energy bound}{Construct the actual interpolation error and prove its energy bound}}


Assume \(\nabla U\) is \(L\)-Lipschitz and \(G_0=\int\|\nabla U\|^2d\mu<\infty\). For each \(t\ge0\), the theorem constructs a strongly measurable \(e_t:E\to E\) with

\[e_t(Y_t)=\mathbb E[D_t\mid Y_t],\qquad \int\|e_t(y)\|^2\,d\rho_t(y)\le t^2L^2G_0+2tnL^2,\]\[\rho_t=(Y_t)_\#(\mu\otimes\mathcal N(0,I)),\qquad n=\dim E.\]

The proof names the actual proposal map \texttt{Y} and discrepancy \texttt{D}. \texttt{ULA.\allowbreak{}mala\allowbreak{}Proposal\_\allowbreak{}measurable} proves measurability of \(Y\). \texttt{frozen\_\allowbreak{}gradient\_\allowbreak{}error\_\allowbreak{}energy} gives square-integrability of \(D\) and the joint bound \(\mathbb E\|D\|^2\le L^2(t^2G_0+2tn)\). The continuity of the Lipschitz gradient gives \texttt{h\allowbreak{}Dmeas}, so \texttt{mem\allowbreak{}Lp\_\allowbreak{}two\_\allowbreak{}iff\_\allowbreak{}integrable\_\allowbreak{}sq\_\allowbreak{}norm} turns that bound into the \(L^2\) input \texttt{h\allowbreak{}D}.



The general conditional-mean theorem now constructs \texttt{err}. Its pushforward measure is definitionally the actual \texttt{interpolation\allowbreak{}Law}, so the field is integrable on the required marginal. Conditional Jensen followed by the joint Gaussian estimate gives the result; the final \texttt{ring} step expands \(L^2(t^2G_0+2tn)\) into the exact coefficient order used by the entropy interface. Neither the error field nor its numerical energy estimate is a hypothesis.

\subsection{\texorpdfstring{Express the constructed error energy in relative-\allowbreak{}density coordinates}{Express the constructed error energy in relative-density coordinates}}


Suppose a measurable nonnegative \(r\) identifies the actual interpolation law as \(\rho_t=r\nu\). The preceding construction gives a field \(e_t\) satisfying both the conditional-mean identity and

\[\int r(y)\|e_t(y)\|^2\,d\nu(y)\le t^2L^2G_0+2tnL^2.\]

The weighted integrand is also proved integrable. In Lean, \texttt{exists\_\allowbreak{}frozen\_\allowbreak{}drift\_\allowbreak{}error} first supplies the state-space field. The law identity \texttt{hlaw} changes its integrability and energy conclusions from the interpolation law to \texttt{tilted\allowbreak{}Measure \ensuremath{\nu} r}. \texttt{integrable\_\allowbreak{}with\allowbreak{}Density\_\allowbreak{}iff\_\allowbreak{}integrable\_\allowbreak{}smul'} then converts integrability under \(r\nu\) into integrability of \(r\|e_t\|^2\) under \(\nu\). Nonnegativity of \(r\) makes the real conversion of \texttt{ENNReal.\allowbreak{}of\allowbreak{}Real r} equal to \(r\).



Finally \texttt{integral\_\allowbreak{}tilted\allowbreak{}Measure} converts the energy integral itself. The density-law identity is a displayed input to this corollary; at positive times \texttt{Interpolation\allowbreak{}Density.\allowbreak{}lean} constructs such a relative density explicitly for a normalized Gibbs target. This corollary does not assume entropy production or prove the remaining weak differential equation.

\section{\texorpdfstring{Construct spatial smoothness, finite Fisher information and finite actual entropy}{Construct spatial smoothness, finite Fisher information and finite actual entropy}}
\label{proof-module-UnadjustedLangevin-LSI-GaussianEntropy-lean}
{\small\color{gray}Source module: \texttt{Unadjusted\allowbreak{}Langevin/\allowbreak{}LSI/\allowbreak{}Gaussian\allowbreak{}Entropy.\allowbreak{}lean}.\par}


For \(t>0\), the actual frozen Gaussian update \(Y_t=X-t\nabla U(X)+\sqrt{2t}\,Z\) has the positive mixture density constructed in \texttt{Interpolation\allowbreak{}Density.\allowbreak{}lean}. This module proves its spatial \(C^1\) regularity by differentiating the displayed Gaussian integral. Bayes' posterior formula and Jensen's inequality then give the universal Lebesgue Fisher bound \(\int\|\nabla\log p_t\|^2p_t\le d/(2t)\). If the initial gradient has finite energy, the relative score \(\nabla\log r_t=\nabla U+\nabla\log p_t\) also has finite energy. The target LSI consequently proves finite actual KL at every positive time. No initial density or second moment of the state is assumed.

\subsection{\texorpdfstring{Bound the quadratic Gaussian envelope uniformly}{Bound the quadratic Gaussian envelope uniformly}}


For \(a>0\), put \(u=a\|z\|^2\). The scalar inequality \(u e^{-u}\le e^{-1}\le1\) gives

\[\|z\|^2e^{-a\|z\|^2}\le a^{-1}.\]

Lean multiplies the desired inequality by the positive number \(a\), rewrites its left side as \(u e^{-u}\), and applies \texttt{Real.\allowbreak{}mul\_\allowbreak{}exp\_\allowbreak{}neg\_\allowbreak{}le\_\allowbreak{}exp\_\allowbreak{}neg\_\allowbreak{}one}. The final cancellation uses \texttt{mul\_\allowbreak{}inv\_\allowbreak{}cancel\ensuremath{{}_{0}}} with the explicit nonzero witness from \texttt{ha}. This global bound later controls the conditional displacement variance even when the initial distribution has no moments.

\subsection{\texorpdfstring{Bound the first Gaussian envelope without an initial moment}{Bound the first Gaussian envelope without an initial moment}}


The elementary estimate \(r\le1+r^2\), together with \(e^{-ar^2}\le1\) and the quadratic envelope bound, gives

\[\|z\|e^{-a\|z\|^2}\le1+a^{-1}.\]

The proof records the exponential bound as \texttt{hexp}, the quadratic bound as \texttt{hsq}, and \(r\le1+r^2\) as \texttt{hnorm}. Multiplication by the nonnegative exponential and addition then yield the displayed result. The bound is intentionally simple rather than optimal: a finite constant independent of the proposed state and mixture component is what dominated differentiation requires.

\subsection{\texorpdfstring{Specify the real Gaussian kernel and its derivative map}{Specify the real Gaussian kernel and its derivative map}}


With precision \(a\), center \(c\), and the already normalized Gaussian constant \(N_a\), define

\[K_a(c,y)=N_a e^{-a\|y-c\|^2},\qquad D_yK_a(c,y)[v]=-2aK_a(c,y)\langle y-c,v\rangle.\]

\texttt{real\allowbreak{}Gaussian\allowbreak{}Kernel} is the ordinary real density. \texttt{real\allowbreak{}Gaussian\allowbreak{}Kernel\allowbreak{}Deriv} is its proposed Fréchet derivative as a continuous linear functional: the scalar \texttt{real\allowbreak{}Gaussian\allowbreak{}Kernel a c y * (-\allowbreak{}2 * a)} multiplies \texttt{inner\allowbreak{}SL \ensuremath{\mathbb R} (y -\allowbreak{} c)}. The precision parameter keeps this helper reusable; the ULA specialization is \(a=1/(4t)\). The next theorem proves that this proposed derivative is the actual derivative.

\subsection{\texorpdfstring{Differentiate the kernel by the squared-\allowbreak{}norm and exponential chain rules}{Differentiate the kernel by the squared-norm and exponential chain rules}}


The map \(y\mapsto y-c\) has derivative the identity. The squared-norm derivative is \(v\mapsto2\langle y-c,v\rangle\). Multiplication by \(-a\), exponentiation, and multiplication by \(N_a\) therefore produce exactly

\[D_yK_a(c,y)[v]=-2aK_a(c,y)\langle y-c,v\rangle.\]

The local Lean proof \texttt{hd} is this sequence of \texttt{sub\_\allowbreak{}const}, \texttt{norm\_\allowbreak{}sq}, \texttt{const\_\allowbreak{}mul}, \texttt{exp}, and a final \texttt{const\_\allowbreak{}mul}. The final extensionality argument compares the two continuous linear maps at an arbitrary vector \texttt{v}; elementary scalar algebra proves they agree. No regularity of the potential is involved in this fixed-center calculation.

\subsection{\texorpdfstring{Find a common integrable bound for all component derivatives}{Find a common integrable bound for all component derivatives}}


The norm of the inner-product functional \(v\mapsto\langle y-c,v\rangle\) equals \(\|y-c\|\). Hence

\[\|D_yK_a(c,y)\|=2aN_a\|y-c\|e^{-a\|y-c\|^2}\le2aN_a(1+a^{-1}).\]

Lean uses \texttt{norm\_\allowbreak{}smul} and \texttt{inner\allowbreak{}SL\_\allowbreak{}apply\_\allowbreak{}norm}, removes absolute values using positivity of \(N_a\) and the exponential, and applies the first Gaussian envelope bound. The resulting constant is independent of both \(y\) and \(c\). A probability mixture therefore integrates this bound without requiring any moment of its center map.

\subsection{\texorpdfstring{Define the actual proposed mixture derivative}{Define the actual proposed mixture derivative}}


For \(c_t(x)=x-t\nabla U(x)\), define the continuous-linear-map-valued integral

\[P'_t(y)=\int D_yK_{1/(4t)}(c_t(x),y)\,d\mu(x).\]

This is \texttt{interpolation\allowbreak{}Lebesgue\allowbreak{}Density\allowbreak{}Deriv}. The symbol here denotes a spatial Fréchet derivative, not the time derivative of the density. The following domination lemmas establish its integrability, continuity, and equality with the actual spatial derivative of the explicit mixture. Keeping the definition as an integral of linear maps makes each directional evaluation correspond directly to the component kernel derivative.

\subsection{\texorpdfstring{Integrate each Gaussian component over the mixing probability law}{Integrate each Gaussian component over the mixing probability law}}


For a measurable center map \(c\), the function \(x\mapsto K_a(c(x),y)\) is measurable and satisfies \(0\le K_a(c(x),y)\le N_a\). Since \(\mu\) is a probability measure, the constant \(N_a\) is integrable, and so is the kernel.



Lean starts from \texttt{integrable\_\allowbreak{}const} and uses \texttt{Integrable.\allowbreak{}mono'}. Its first obligation is measurability, assembled from the center's measurability by \texttt{fun\_\allowbreak{}prop}. Its second obligation is the norm bound; positivity removes the absolute value, and \texttt{Real.\allowbreak{}exp\_\allowbreak{}le\_\allowbreak{}one\_\allowbreak{}iff} supplies the Gaussian peak bound.

\subsection{\texorpdfstring{Integrate the continuous-\allowbreak{}linear-\allowbreak{}map-\allowbreak{}valued component derivative}{Integrate the continuous-linear-map-valued component derivative}}


The common derivative bound is a finite constant \(2aN_a(1+a^{-1})\). Therefore

\[\int\|D_yK_a(c(x),y)\|\,d\mu(x)<\infty.\]

The Lean proof again applies \texttt{Integrable.\allowbreak{}mono'} to an integrable constant, now with values in the Banach space of continuous linear functionals. Measurability follows from the explicit derivative formula; the second obligation is exactly \texttt{norm\_\allowbreak{}real\allowbreak{}Gaussian\allowbreak{}Kernel\allowbreak{}Deriv\_\allowbreak{}le}. This supplies the Bochner-integrability premise needed to identify the derivative of the mixture.

\subsection{\texorpdfstring{Differentiate the mixture under the integral sign}{Differentiate the mixture under the integral sign}}


For each proposed state \(y\),

\[D_y\!\left(\int K_a(c(x),y)\,d\mu(x)\right)=\int D_yK_a(c(x),y)\,d\mu(x).\]

Lean applies \texttt{has\allowbreak{}FDeriv\allowbreak{}At\_\allowbreak{}integral\_\allowbreak{}of\_\allowbreak{}dominated\_\allowbreak{}of\_\allowbreak{}fderiv\_\allowbreak{}le} with the whole spatial domain as the common neighborhood. Kernel integrability supplies the value at the base point, derivative-kernel integrability supplies derivative measurability, the uniform Gaussian bound supplies domination, and \texttt{has\allowbreak{}FDeriv\allowbreak{}At\_\allowbreak{}real\allowbreak{}Gaussian\allowbreak{}Kernel} supplies pointwise differentiation. All these premises were proved from the actual kernel. In particular, a derivative of the mixture is not among the hypotheses.

\subsection{\texorpdfstring{Prove continuity of the mixture derivative with the same domination}{Prove continuity of the mixture derivative with the same domination}}


For every fixed initial point \(x\), the derivative kernel \(y\mapsto D_yK_a(c(x),y)\) is continuous. Its operator norm is bounded by the same integrable constant for every \(y\). Dominated continuity therefore proves continuity of

\[y\longmapsto\int D_yK_a(c(x),y)\,d\mu(x).\]

The Lean theorem \texttt{continuous\_\allowbreak{}of\_\allowbreak{}dominated} requires measurability of every section, a common integrable bound, and almost-everywhere continuity of the component map. The preceding lemmas provide the first three facts, while the explicit formula and \texttt{fun\_\allowbreak{}prop} prove continuity. This continuity is what strengthens mere differentiability to \(C^1\).

\subsection{\texorpdfstring{Combine the derivative and its continuity into a C1 theorem}{Combine the derivative and its continuity into a C1 theorem}}


A function is \(C^1\) exactly when it has a continuous Fréchet derivative. Apply this criterion to the derivative field \(y\mapsto\int D_yK_a(c(x),y)\,d\mu(x)\).



\texttt{cont\allowbreak{}Diff\_\allowbreak{}one\_\allowbreak{}iff\_\allowbreak{}has\allowbreak{}FDeriv\allowbreak{}At} reduces the Lean proof to the two preceding results: \texttt{continuous\_\allowbreak{}integral\_\allowbreak{}real\allowbreak{}Gaussian\allowbreak{}Kernel\allowbreak{}Deriv} gives continuity of this field, and \texttt{has\allowbreak{}FDeriv\allowbreak{}At\_\allowbreak{}integral\_\allowbreak{}real\allowbreak{}Gaussian\allowbreak{}Kernel} identifies it as the derivative at every point. The conclusion applies to every probability mixing measure, including singular ones.

\subsection{\texorpdfstring{Apply Gaussian smoothing to the actual frozen-\allowbreak{}drift law}{Apply Gaussian smoothing to the actual frozen-drift law}}


The explicit real density of the update satisfies

\[p_t(y)=\int K_{1/(4t)}(x-t\nabla U(x),y)\,d\mu(x).\]

The proof first establishes this equality as \texttt{heq}, using the integral formula from \texttt{Interpolation\allowbreak{}Density.\allowbreak{}lean} and the scalar identity \(-\|z\|^2/(4t)=-(1/(4t))\|z\|^2\). The center map is measurable because the gradient is Lipschitz. The general Gaussian-mixture \(C^1\) theorem then applies with positive precision \(1/(4t)\). Thus spatial regularity is a consequence of the actual transition law, not an assumption about a chosen interpolation density.

\subsection{\texorpdfstring{Identify the exact spatial derivative of the actual density}{Identify the exact spatial derivative of the actual density}}


The derivative field defined earlier is now certified:

\[D_yp_t(y)=P'_t(y)=\int D_yK_{1/(4t)}(x-t\nabla U(x),y)\,d\mu(x).\]

Lean applies the general dominated-differentiation theorem to the measurable frozen center and rewrites its mixture back to \texttt{interpolation\allowbreak{}Lebesgue\allowbreak{}Density}. The final \texttt{rfl} checks that the derivative integral is exactly \texttt{interpolation\allowbreak{}Lebesgue\allowbreak{}Density\allowbreak{}Deriv}. This identity is later used to define a jointly measurable score field while retaining its agreement with the actual spatial derivative.

\subsection{\texorpdfstring{Recover C1 regularity of the potential from its Lipschitz gradient}{Recover C1 regularity of the potential from its Lipschitz gradient}}


The assumptions are differentiability of \(U\) and Lipschitz continuity of \(\nabla U\). The Riesz isometry identifies

\[DU(x)=\operatorname{Riesz}(\nabla U(x)).\]

Since the Riesz map and the gradient are continuous, \(x\mapsto DU(x)\) is continuous. The criterion \texttt{cont\allowbreak{}Diff\_\allowbreak{}one\_\allowbreak{}iff\_\allowbreak{}fderiv} now yields \(U\in C^1\). Lean rewrites the derivative field by \texttt{to\allowbreak{}Dual\_\allowbreak{}comp\_\allowbreak{}gradient} and composes the two continuous maps. No Hessian or twice-continuously-differentiable potential is assumed.

\subsection{\texorpdfstring{Establish spatial admissibility of the actual Gibbs-\allowbreak{}relative density}{Establish spatial admissibility of the actual Gibbs-relative density}}


The relative density is \(r_t(y)=Z e^{U(y)}p_t(y)\). The partition factor \(Z\) is constant in space, \(e^U\) is \(C^1\) because \(U\) is \(C^1\), and the actual Gaussian mixture \(p_t\) is \(C^1\). Therefore \(r_t\) is \(C^1\) at every positive time.



The Lean proof follows this product exactly using \texttt{cont\allowbreak{}Diff\_\allowbreak{}const.\allowbreak{}mul}, the potential theorem followed by \texttt{exp}, and \texttt{interpolation\allowbreak{}Lebesgue\allowbreak{}Density\_\allowbreak{}cont\allowbreak{}Diff}. Together with pointwise positivity and normalization from the density-construction module, this discharges the smooth test-class premises of \texttt{Is\allowbreak{}Log\allowbreak{}Sobolev}.

\subsection{\texorpdfstring{Apply the target LSI directly to the constructed interpolation law}{Apply the target LSI directly to the constructed interpolation law}}


Use the actual positive normalized \(C^1\) relative density as the LSI test. This gives

\[2\alpha\,\mathrm{KL}(\rho_t\mid\nu)\le I(\rho_t\mid\nu).\]

This theorem retains extended nonnegative values, so neither finiteness is silently assumed. Lean obtains the Gibbs probability instance from \texttt{h\allowbreak{}LSI.\allowbreak{}1}; \texttt{h\allowbreak{}LSI.\allowbreak{}2} is applied with the already proved positivity, \(C^1\) regularity, and relative probability normalization. Finally \texttt{tilted\allowbreak{}Measure\_\allowbreak{}interpolation\allowbreak{}Relative\allowbreak{}Density} replaces the induced tilted measure with the actual Gaussian interpolation law. The only functional-inequality hypothesis is the target's LSI.

\subsection{\texorpdfstring{Record the square-\allowbreak{}norm Jensen inequality for probability measures}{Record the square-norm Jensen inequality for probability measures}}


For a square-integrable random vector \(F\) under a probability measure, convexity of squared norm gives

\[\left\|\int F\,d\mu\right\|^2\le\int\|F\|^2\,d\mu.\]

Lean obtains convexity by applying \texttt{Convex\allowbreak{}On.\allowbreak{}pow} to convexity and nonnegativity of the norm, with exponent two. \texttt{map\_\allowbreak{}integral\_\allowbreak{}le} then applies Jensen on the whole vector space. \texttt{Mem\allowbreak{}Lp} supplies both integrability of \(F\) and of \(\|F\|^2\). This general lemma is used under the normalized Gaussian posterior at each proposed state.

\subsection{\texorpdfstring{Define Bayes posterior on the initial mixture component}{Define Bayes posterior on the initial mixture component}}


At a fixed proposed point \(y\), let \(p(y)=\int K_a(c(x),y)\,d\mu(x)\). The posterior law of the initial component is

\[\lambda_y(dx)=\frac{K_a(c(x),y)}{p(y)}\,\mu(dx).\]

\texttt{gaussian\allowbreak{}Posterior} constructs this measure with \texttt{with\allowbreak{}Density}. The definition is total, while its intended probability interpretation uses \(a>0\), measurability of the center, and positivity of the mixing probability law. The following lemmas establish positivity of the denominator and normalization explicitly, before using Jensen.

\subsection{\texorpdfstring{Prove the real mixture denominator is everywhere positive}{Prove the real mixture denominator is everywhere positive}}


Every component \(K_a(c(x),y)\) is positive because its normalizer and exponential are positive. Its support as a function of \(x\) is therefore the whole initial-state space, which has probability one.



The Lean proof uses \texttt{integral\_\allowbreak{}pos\_\allowbreak{}iff\_\allowbreak{}support\_\allowbreak{}of\_\allowbreak{}nonneg}, with integrability from the common Gaussian peak bound. The local support equality \texttt{hs} converts the condition to positive mass of the whole space. Thus \(p(y)>0\) at every point, allowing Bayes normalization and the logarithmic score to be defined pointwise.

\subsection{\texorpdfstring{Check posterior normalization by cancellation of the actual mixture integral}{Check posterior normalization by cancellation of the actual mixture integral}}


The total posterior mass equals

\[\int\frac{K_a(c(x),y)}{p(y)}\,d\mu(x)=\frac{p(y)}{p(y)}=1.\]

Lean first records positive denominator \texttt{hp} and integrable kernel \texttt{hq}. The definition of \texttt{with\allowbreak{}Density} expresses mass as a nonnegative integral. \texttt{of\allowbreak{}Real\_\allowbreak{}integral\_\allowbreak{}eq\_\allowbreak{}lintegral\_\allowbreak{}of\allowbreak{}Real} converts it to the ordinary integral, justified by integrability and nonnegativity of the normalized kernel. \texttt{integral\_\allowbreak{}div} pulls out the denominator and \texttt{div\_\allowbreak{}self hp.\allowbreak{}ne'} performs the normalization.

\subsection{\texorpdfstring{Prove posterior displacement has finite conditional second moment}{Prove posterior displacement has finite conditional second moment}}


The posterior may arise from an initial law without moments. Nevertheless, at each fixed \(y\),

\[\int\|y-c(x)\|^2\,d\lambda_y(x)=\frac{1}{p(y)}\int K_a(c(x),y)\|y-c(x)\|^2\,d\mu(x)<\infty.\]

The quadratic Gaussian envelope bounds the numerator's integrand by \(N_a/a\). Lean's \texttt{integrable\_\allowbreak{}with\allowbreak{}Density\_\allowbreak{}iff\_\allowbreak{}integrable\_\allowbreak{}smul'} changes posterior integrability into the displayed weighted integral under \(\mu\). The proof \texttt{hsquare} dominates it by the integrable constant \(N_a/(ap(y))\). Measurability follows from the measurable center and absolute continuity of the posterior. The standard \(L^2\) characterization then supplies \texttt{Mem\allowbreak{}Lp} for the displacement.

\subsection{\texorpdfstring{Identify the mixture score as the posterior component-\allowbreak{}score mean}{Identify the mixture score as the posterior component-score mean}}


Dividing the mixture derivative by its positive value gives

\[\nabla\log p(y)=-2a\int(y-c(x))\,d\lambda_y(x).\]

Lean differentiates \(\log p\) using the actual mixture derivative and the nonzero denominator. The Riesz inverse commutes with the Bochner integral; \texttt{hgrad} identifies the inverse image of each component derivative with \(-2aK_a(c(x),y)(y-c(x))\). Expanding the posterior via \texttt{integral\_\allowbreak{}with\allowbreak{}Density\_\allowbreak{}eq\_\allowbreak{}integral\_\allowbreak{}to\allowbreak{}Real\_\allowbreak{}smul} gives the same integral divided by \(p(y)\). The final scalar cancellations establish the equality pointwise, rather than merely choosing an almost-everywhere score representative.

\subsection{\texorpdfstring{Use posterior Jensen to bound the pointwise mixture score energy}{Use posterior Jensen to bound the pointwise mixture score energy}}


Applying square-norm Jensen to the posterior score identity yields

\[p(y)\|\nabla\log p(y)\|^2\le4a^2\int K_a(c(x),y)\|y-c(x)\|^2\,d\mu(x).\]

The posterior is a probability measure, and its displacement is square integrable by the preceding lemmas; these are the exact hypotheses needed for Jensen. Lean records that inequality as \texttt{hj}, scales it by \(4a^2p(y)\), expands the posterior integral, and cancels \(p(y)\) using positivity. This gives a usable unnormalized density inequality whose right side can be integrated over the proposed state by Fubini.

\subsection{\texorpdfstring{Compute the exact Gaussian component displacement energy}{Compute the exact Gaussian component displacement energy}}


At precision \(1/(4t)\), the component density is the law of \(c+\sqrt{2t}Z\). Therefore

\[\int K_{1/(4t)}(c,y)\|y-c\|^2\,dy=\mathbb E\|\sqrt{2t}Z\|^2=2td.\]

The proof includes integrability, not only an integral identity. \texttt{gaussian\_\allowbreak{}proposal\_\allowbreak{}radial\_\allowbreak{}density} identifies the normalized density with the actual affine Gaussian law. Gaussian square integrability proves \texttt{hnoise}; \texttt{integrable\_\allowbreak{}map\_\allowbreak{}measure} transfers it to the proposed state, and the density-integrability equivalence gives the weighted Lebesgue integrand. The integral identity then follows from the map formula, \((\sqrt{2t})^2=2t\), and the already proved standard Gaussian moment \(\mathbb E\|Z\|^2=d\).

\subsection{\texorpdfstring{Integrate posterior Jensen to obtain the universal smoothing Fisher bound}{Integrate posterior Jensen to obtain the universal smoothing Fisher bound}}


Integrate the pointwise score inequality, use Fubini, and insert the component displacement moment:

\[\int p(y)\|\nabla\log p(y)\|^2\,dy\le4\left(\frac1{4t}\right)^2\int2td\,d\mu=\frac d{2t}.\]

Lean first proves joint integrability of the density-weighted displacement energy as \texttt{hprod}, using \texttt{integrable\_\allowbreak{}prod\_\allowbreak{}iff} and the constant component moment. The right side of pointwise Jensen is consequently integrable. Continuity of the mixture score supplies measurability of the left integrand, so \texttt{Integrable.\allowbreak{}mono'} proves \texttt{h\allowbreak{}W}. Only after establishing both integrability statements does the proof apply \texttt{integral\_\allowbreak{}mono}, exchange the integrals, and evaluate the Gaussian moment. The bound is independent of the initial distribution's location, density, or state moments.

\subsection{\texorpdfstring{Propagate finite gradient energy without a state second moment}{Propagate finite gradient energy without a state second moment}}


Write

\[\nabla U(Y_t)=\bigl(\nabla U(Y_t)-\nabla U(X)\bigr)+\nabla U(X).\]

The first term is in \(L^2\) because Lipschitzness bounds it by \(L\|Y_t-X\|\), and the actual frozen displacement has finite second moment whenever \(\nabla U(X)\in L^2\). The second term is in \(L^2\) by assumption. Lean uses \texttt{frozen\_\allowbreak{}gradient\_\allowbreak{}error\_\allowbreak{}energy} and the \(L^2\) characterization for the discrepancy, adds the pulled-back initial gradient, and transfers the sum to the output measure with \texttt{mem\allowbreak{}Lp\_\allowbreak{}map\_\allowbreak{}measure\_\allowbreak{}iff}. The initial position cancels from the displacement, which is why no \(\mathbb E\|X\|^2\) premise appears.

\subsection{\texorpdfstring{Transfer the smoothing Fisher estimate to the actual update law}{Transfer the smoothing Fisher estimate to the actual update law}}


The generic mixture Fisher theorem applies with \(c_t(x)=x-t\nabla U(x)\). Rewriting its integral through \(\rho_t=p_t\,dy\) gives

\[\nabla\log p_t\in L^2(\rho_t),\qquad\int\|\nabla\log p_t\|^2\,d\rho_t\le\frac d{2t}.\]

The local identity \texttt{heq} matches the actual density with the generic kernel integral. \texttt{hf} is the generic integrability-and-bound result. The density-law identity and \texttt{integrable\_\allowbreak{}with\allowbreak{}Density\_\allowbreak{}iff\_\allowbreak{}integrable\_\allowbreak{}smul'} convert its integrability to the actual law; the integral-with-density formula converts the numerical bound. Pointwise finiteness of the density ensures that real and extended density representations agree.

\subsection{\texorpdfstring{Retain everywhere positivity for the real density}{Retain everywhere positivity for the real density}}


The extended mixture density is positive and finite at every point for \(t>0\). Its real conversion is consequently positive:

\[p_t(y)>0.\]

The proof is exactly \texttt{ENNReal.\allowbreak{}to\allowbreak{}Real\_\allowbreak{}pos}, with the two previously proved guards \texttt{interpolation\allowbreak{}Density\_\allowbreak{}pos} and \texttt{interpolation\allowbreak{}Density\_\allowbreak{}ne\_\allowbreak{}top}. This real positivity witness is used whenever Lean differentiates \(\log p_t\) or rewrites the logarithm of a product.

\subsection{\texorpdfstring{Split the relative score into potential and Lebesgue scores}{Split the relative score into potential and Lebesgue scores}}


Since \(r_t=Ze^Up_t\), positivity and finite normalization give

\[\log r_t=\log Z+U+\log p_t,\qquad\nabla\log r_t=\nabla U+\nabla\log p_t.\]

Lean first proves that the real partition function is positive and finite. The function equality \texttt{heq} follows from \texttt{Real.\allowbreak{}log\_\allowbreak{}mul} and \texttt{Real.\allowbreak{}log\_\allowbreak{}exp}. Spatial differentiability of \(p_t\) then justifies the logarithmic derivative; differentiability of \(U\) handles the other summand. The proof differentiates the displayed sum and uses linearity of the Riesz inverse. The normalization constant disappears from the gradient exactly because it is spatially constant.

\subsection{\texorpdfstring{Construct finite relative-\allowbreak{}score energy from two proved L2 fields}{Construct finite relative-score energy from two proved L2 fields}}


The potential gradient remains in \(L^2(\rho_t)\), while the Gaussian smoothing bound gives \(\nabla\log p_t\in L^2(\rho_t)\). Their sum is therefore in \(L^2\), and the score decomposition identifies it with the actual relative score:

\[\nabla\log r_t\in L^2(\rho_t).\]

Lean obtains continuity of the Lebesgue score from \(C^1\) regularity of its positive density, converts the finite Fisher integral to \texttt{Mem\allowbreak{}Lp}, and adds it to \texttt{mem\allowbreak{}Lp\_\allowbreak{}gradient\_\allowbreak{}interpolation\allowbreak{}Law}. The final function extensionality uses \texttt{gradient\_\allowbreak{}log\_\allowbreak{}interpolation\allowbreak{}Relative\allowbreak{}Density} pointwise. This is the finite-energy fact used for the bounded logarithmic entropy tests.

\subsection{\texorpdfstring{Prove finite relative Fisher information before taking real values}{Prove finite relative Fisher information before taking real values}}


Square integrability of the actual relative score implies

\[I(\rho_t\mid\nu)=\int\|\nabla\log r_t\|^2\,d\rho_t<\infty.\]

Lean starts from the \(L^2\) score theorem and obtains the integrable squared norm. It unfolds \texttt{relative\allowbreak{}Fisher}, replaces the tilted measure by the actual interpolation law, and rewrites the nonnegative integral as \texttt{ENNReal.\allowbreak{}of\allowbreak{}Real} of the ordinary energy integral. Such an extended value is finite. This explicit proof prevents later uses of \texttt{to\allowbreak{}Real} from discarding an infinite Fisher information.

\subsection{\texorpdfstring{Derive finite actual KL from LSI and the constructed finite Fisher energy}{Derive finite actual KL from LSI and the constructed finite Fisher energy}}


At positive time, the constructed density is an admissible LSI test, so \(2\alpha\,\mathrm{KL}(\rho_t\mid\nu)\le I(\rho_t\mid\nu)\). With \(\alpha>0\) and finite Fisher information, the KL divergence must be finite.



Lean uses the target's LSI probability instance, \texttt{interpolation\allowbreak{}Relative\allowbreak{}Fisher\_\allowbreak{}ne\_\allowbreak{}top}, and \texttt{interpolation\allowbreak{}Law\_\allowbreak{}log\allowbreak{}Sobolev}. If KL were infinity, its product with the positive nonzero factor \texttt{ENNReal.\allowbreak{}of\allowbreak{}Real (2 * \ensuremath{\alpha})} would be infinity. The inequality would then force Fisher information to be infinity, contradicting the proved finiteness. No entropy continuity or initial density smoothness is assumed to obtain this positive-time finiteness.

\section{\texorpdfstring{Time derivatives and spacetime integrability of the Gaussian mixture}{Time derivatives and spacetime integrability of the Gaussian mixture}}
\label{proof-module-UnadjustedLangevin-LSI-GaussianTimeDerivative-lean}
{\small\color{gray}Source module: \texttt{Unadjusted\allowbreak{}Langevin/\allowbreak{}LSI/\allowbreak{}Gaussian\allowbreak{}Time\allowbreak{}Derivative.\allowbreak{}lean}.\par}


Let \(g_x=\nabla U(x)\), \(c_t(x)=x-tg_x\), \(d=\dim E\), and \(K_t(x,y)=(4\pi t)^{-d/2}\exp(-\|y-c_t(x)\|^2/(4t))\). This module differentiates the actual normalized proposal, constructs \(\dot p_t(y)=\int\partial_tK_t(x,y)\,\mu(dx)\), and proves its local spacetime integrability. The only initial moment used is \(G_0=\int\|g_x\|^2d\mu<\infty\); this is derived from finite initial KL elsewhere.

\subsection{\texorpdfstring{gaussian\allowbreak{}Proposal\allowbreak{}Normalizer eq rpow}{gaussianProposalNormalizer eq rpow}}


The proposal normalizer is defined as the reciprocal of an actual Gaussian integral, not as a free scalar parameter. Gaussian Fourier integration evaluates \(\int e^{-a\|z\|^2}dz=(\pi/a)^{d/2}\) for \(a>0\). Substitute \(a=1/(4t)\), simplify \(\pi/a=4\pi t\), and invert the real power to obtain \(N_t=(4\pi t)^{-d/2}\). Positivity of \(t\) justifies the Gaussian formula and the real-power inversion.

\subsection{\texorpdfstring{has\allowbreak{}Deriv\allowbreak{}At gaussian\allowbreak{}Proposal\allowbreak{}Normalizer}{hasDerivAt gaussianProposalNormalizer}}


On a neighborhood of \(t>0\), all times remain positive, so the actual normalizer agrees locally with \((4\pi s)^{-d/2}\). The real-power chain rule gives derivative \((-d/2)(4\pi t)^{-d/2-1}(4\pi)\). Use \((4\pi t)^{p-1}=(4\pi t)^p/(4\pi t)\) to rewrite this as \(-dN_t/(2t)\). The local equality \texttt{heq} transfers the derivative back to the normalizer defined by the Gaussian integral.

\subsection{\texorpdfstring{frozen\allowbreak{}Gaussian\allowbreak{}Kernel}{frozenGaussianKernel}}


The real component kernel is \(K_t(x,y)=N_t\exp(-\|y-x+tg_x\|^2/(4t))\). The gradient \(g_x\) is evaluated at the initial state \(x\), and is constant when differentiating in \(t\) or \(y\). Its integral over \(y\) is one at positive time by the previously proved Gaussian proposal normalization.

\subsection{\texorpdfstring{frozen\allowbreak{}Gaussian\allowbreak{}Kernel\allowbreak{}Dot}{frozenGaussianKernelDot}}


Define the explicit candidate derivative \(\dot K_t=K_t[-d/(2t)+\|y-x\|^2/(4t^2)-\|g_x\|^2/4]\). The dimension term comes from the normalizer; the remaining terms come from differentiating the exponent. In the expanded exponent, the mixed inner-product term is independent of time, which explains its absence from this expression. The following theorem verifies that this formula is the derivative.

\subsection{\texorpdfstring{frozen exponent eq}{frozen exponent eq}}


Set \(z=y-x\). The squared-norm identity gives \(\|z+sg_x\|^2=\|z\|^2+2s\langle z,g_x\rangle+s^2\|g_x\|^2\). For \(s\ne0\), division by \(-4s\) yields \(-\|z+sg_x\|^2/(4s)=-\|z\|^2/(4s)-\langle z,g_x\rangle/2-s\|g_x\|^2/4\). The vector identity and scalar denominator cancellation are proved separately in the Lean calculation.

\subsection{\texorpdfstring{has\allowbreak{}Deriv\allowbreak{}At frozen\allowbreak{}Gaussian\allowbreak{}Kernel}{hasDerivAt frozenGaussianKernel}}


Write the expanded exponent as \(-A/s-B-sC\), where \(A=\|y-x\|^2/4\), \(B=\langle y-x,g_x\rangle/2\), and \(C=\|g_x\|^2/4\). Its derivative at positive \(t\) is \(A/t^2-C\). The expansion agrees with the original exponent on a neighborhood avoiding zero, so this is also the derivative of the original expression.



Apply the exponential chain rule and multiply by the actual normalizer. Its derivative contributes \(-dK_t/(2t)\); the exponential derivative contributes \(K_t(A/t^2-C)\). Collecting terms gives exactly \texttt{frozen\allowbreak{}Gaussian\allowbreak{}Kernel\allowbreak{}Dot}. No derivative of \(g_x\), and therefore no Hessian, is taken.

\subsection{\texorpdfstring{norm frozen\allowbreak{}Gaussian\allowbreak{}Kernel\allowbreak{}Dot le}{norm frozenGaussianKernelDot le}}


Put \(r=y-x+tg_x\) and \(e=e^{-\|r\|^2/(4t)}\). The triangle inequality and \((a+b)^2\le2a^2+2b^2\) give \(\|y-x\|^2\le2\|r\|^2+2t^2\|g_x\|^2\). The Gaussian envelope gives \(\|r\|^2e\le4t\), while \(e\le1\). Consequently \(\|y-x\|^2e\le8t+2t^2\|g_x\|^2\).



Bound the absolute value of each derivative-coefficient term by their nonnegative sum, multiply by \(e\), and insert these bounds. The resulting convenient envelope is \(|\dot K_t(x,y)|\le N_t[(d+4)/t+\|g_x\|^2]\). It is uniform in \(y\), and is integrable in the initial variable whenever \(G_0<\infty\). Its purpose is time differentiation under the initial-state integral, rather than output-space integration.

\subsection{\texorpdfstring{gaussian\allowbreak{}Proposal\allowbreak{}Normalizer measurable}{gaussianProposalNormalizer measurable}}


The map \((t,z)\mapsto\exp(-(1/(4t))\|z\|^2)\) is jointly measurable. Parameter integration makes its real Gaussian integral measurable, even at parameters where the total integral takes its default value. Taking its reciprocal preserves measurability and gives the actual proposal normalizer. The all-time measurable extension is used only to form measurable parameter integrals; positive-time analysis uses the genuine finite normalization.

\subsection{\texorpdfstring{frozen\allowbreak{}Gaussian\allowbreak{}Kernel\allowbreak{}Dot measurable}{frozenGaussianKernelDot measurable}}


Lipschitzness makes \(x\mapsto g_x\) continuous and measurable. Combine this with the measurable normalizer and the explicit algebraic formula for \(\dot K_t\). Products, norms, scalar inverses, and exponentials preserve measurability, giving joint measurability in \((t,y,x)\). This is the input needed to integrate out the initial variable while retaining measurable time-space dependence.

\subsection{\texorpdfstring{interpolation\allowbreak{}Lebesgue\allowbreak{}Density\allowbreak{}Time\allowbreak{}Derivative}{interpolationLebesgueDensityTimeDerivative}}


Define \(\dot p_t(y)=\int\dot K_t(x,y)\,d\mu(x)\). At this point it is an ordinary parameter integral. The subsequent differentiation theorem proves it is the actual derivative of \(p_t(y)\), and the energy estimates prove its integrability, so no analytic assertion is built into this definition.

\subsection{\texorpdfstring{interpolation\allowbreak{}Lebesgue\allowbreak{}Density\allowbreak{}Time\allowbreak{}Derivative measurable}{interpolationLebesgueDensityTimeDerivative measurable}}


The component derivative is jointly measurable in \((t,y,x)\). Its strongly measurable parameter integral over the probability measure \(\mu\) is therefore measurable in \((t,y)\). This gives a single measurable version of \(\dot p\) for spacetime Fubini and the entropy chain rule.

\subsection{\texorpdfstring{has\allowbreak{}Deriv\allowbreak{}At interpolation\allowbreak{}Lebesgue\allowbreak{}Density time}{hasDerivAt interpolationLebesgueDensity time}}


Fix \(t>0\). Continuity of \(N_t\) supplies a common neighborhood \(S\) on which \(s>t/2\) and \(N_s<N_t+1\). The pointwise derivative envelope then gives \(|\dot K_s(x,y)|\le(N_t+1)[2(d+4)/t+\|g_x\|^2]\) for every \(s\in S\). This bound is integrable in \(x\) by the initial gradient-energy hypothesis, and the component kernel at \(t\) is integrable because it is bounded by \(N_t\).



The parameter-integral derivative theorem now supplies both integrability of \(\dot K_t(\cdot,y)\) and the derivative of \(\int K_s(x,y)d\mu(x)\). On a positive-time neighborhood this integral equals the real density of the actual interpolation law. The local equality \texttt{heq} transfers the derivative to that density, proving \(\partial_t p_t(y)=\dot p_t(y)\) for every output state. No invertibility of the Euler map or step-size restriction is needed.

\subsection{\texorpdfstring{frozen\allowbreak{}Gaussian\allowbreak{}Kernel eq real\allowbreak{}Gaussian\allowbreak{}Kernel}{frozenGaussianKernel eq realGaussianKernel}}


The spatial Fisher calculation uses precision notation \(q_a(c,y)=Z_a^{-1}e^{-a\|y-c\|^2}\). Substitute \(a=1/(4t)\) and \(c=x-tg_x\). The normalizers agree by definition, and \(-a\|y-c\|^2=-\|y-c\|^2/(4t)\). This pointwise identity makes the time and spatial calculations use precisely the same normalized Gaussian kernel.

\subsection{\texorpdfstring{norm frozen\allowbreak{}Gaussian\allowbreak{}Kernel\allowbreak{}Dot le weighted}{norm frozenGaussianKernelDot le weighted}}


Retain the Gaussian factor instead of replacing it by its supremum. With \(r=y-x+tg_x\), the bound \(\|y-x\|^2\le2\|r\|^2+2t^2\|g_x\|^2\) implies that the absolute derivative coefficient is at most \(d/(2t)+\|r\|^2/(2t^2)+\|g_x\|^2\). Multiplication by \(K_t\ge0\) yields \(|\dot K_t|\le K_t[d/(2t)+\|r\|^2/(2t^2)+\|g_x\|^2]\). Unlike the initial-variable envelope, this bound remains integrable in the output variable.

\subsection{\texorpdfstring{frozen\allowbreak{}Gaussian\allowbreak{}Energy}{frozenGaussianEnergy}}


Use the common nonnegative envelope \(\mathcal E_t(x,y)=K_t(x,y)(1+\|g_x\|^2+\|y-c_t(x)\|^2)\). Its three terms correspond to Gaussian mass, initial gradient energy, and Gaussian displacement energy. This envelope can dominate the derivative with a time-dependent scalar coefficient and can also be reused in spatial integration by parts.

\subsection{\texorpdfstring{frozen\allowbreak{}Gaussian\allowbreak{}Energy nonneg}{frozenGaussianEnergy nonneg}}


At positive time the actual normalizer and the Gaussian exponential are positive, so \(K_t\ge0\). Every term in \(1+\|g_x\|^2+\|y-c_t(x)\|^2\) is nonnegative. Their product is nonnegative; this identifies its norm with itself in the Fubini integrability proof.

\subsection{\texorpdfstring{frozen\allowbreak{}Gaussian\allowbreak{}Energy integral}{frozenGaussianEnergy integral}}


For fixed \(x\), Gaussian normalization gives \(\int K_t(x,y)dy=1\), and the centered Gaussian covariance gives \(\int K_t(x,y)\|y-c_t(x)\|^2dy=2td\). Both integrands are genuinely integrable. The initial gradient factor is constant in \(y\), so its integral is \(\|g_x\|^2\). Expand the common energy into these three integrable terms to obtain \(\int\mathcal E_t(x,y)dy=1+\|g_x\|^2+2td\).

\subsection{\texorpdfstring{integrable frozen\allowbreak{}Gaussian\allowbreak{}Energy}{integrable frozenGaussianEnergy}}


The common energy is jointly measurable. For every initial \(x\), its output section is integrable by the component energy theorem. Since the energy is nonnegative, the integral of its norm is exactly \(1+\|g_x\|^2+2td\), an integrable function of \(x\) under the probability law \(\mu\). The product integrability criterion therefore proves \(\mathcal E_t\in L^1(\mu\otimes dy)\).

\subsection{\texorpdfstring{integral frozen\allowbreak{}Gaussian\allowbreak{}Energy}{integral frozenGaussianEnergy}}


Joint integrability allows Fubini, giving \(\int\mathcal E_t\,d(\mu\otimes dy)=\int(1+\|g_x\|^2+2td)d\mu(x)\). Probability normalization and the definition of \(G_0\) reduce this to \(1+G_0+2td\). The displayed equality supplies a finite quantitative bound, rather than relying on an unspecified integrability witness.

\subsection{\texorpdfstring{gaussian\allowbreak{}Time\allowbreak{}Coefficient}{gaussianTimeCoefficient}}


Set \(C_t=d/(2t)+1/(2t^2)+1\). It is positive for \(t>0\), and is bounded above on any interval \([a,b]\) with \(a>0\) by \(C_a\). The coefficient is deliberately coarse: only local spacetime integrability, rather than a sharp density-derivative estimate, is needed for the entropy chain rule.

\subsection{\texorpdfstring{norm frozen\allowbreak{}Gaussian\allowbreak{}Kernel\allowbreak{}Dot le energy}{norm frozenGaussianKernelDot le energy}}


In the weighted derivative bound, denote the three nonnegative coefficients by \(A=d/(2t)\), \(B=1/(2t^2)\), and \(1\). Each is at most \(C_t=A+B+1\). Therefore \(A+B\|r\|^2+\|g_x\|^2\le C_t(1+\|g_x\|^2+\|r\|^2)\). Multiply by \(K_t\ge0\) to prove \(|\dot K_t|\le C_t\mathcal E_t\).

\subsection{\texorpdfstring{integrable frozen\allowbreak{}Gaussian\allowbreak{}Kernel\allowbreak{}Dot}{integrable frozenGaussianKernelDot}}


The component time derivative is jointly measurable, and its norm is bounded by \(C_t\mathcal E_t\). The common energy is integrable on \(\mu\otimes dy\), and \(C_t\) is a finite scalar at positive time. Domination therefore proves joint integrability of \(\dot K_t\), which licenses the changes of integration order used in the time weak equation.

\subsection{\texorpdfstring{integrable interpolation\allowbreak{}Lebesgue\allowbreak{}Density\allowbreak{}Time\allowbreak{}Derivative}{integrable interpolationLebesgueDensityTimeDerivative}}


Fubini applied to the jointly integrable component derivative first gives \(\dot p_t\in L^1(dy)\). The norm inequality for Bochner integration gives \(|\dot p_t(y)|\le\int|\dot K_t(x,y)|d\mu(x)\). Integrate in \(y\), exchange the two integrals, and dominate by \(C_t\mathcal E_t\). The exact envelope integral yields \(\int|\dot p_t(y)|dy\le C_t(1+G_0+2td)\).

\subsection{\texorpdfstring{integrable interpolation\allowbreak{}Lebesgue\allowbreak{}Density\allowbreak{}Time\allowbreak{}Derivative spacetime}{integrable interpolationLebesgueDensityTimeDerivative spacetime}}


Joint measurability of \(\dot p\) permits the product integrability criterion on \([a,b]\times E\), where \(a>0\). Every time section is in \(L^1(dy)\). For times in the interval, \(C_t\le C_a\) and \(1+G_0+2td\le1+G_0+2bd\), giving the common finite scalar bound \(\int|\dot p_t|dy\le C_a(1+G_0+2bd)\). The time interval has finite measure, so the section norms are integrable. Hence \(\dot p\in L^1(dt\,dy)\) locally away from zero. If \(a>b\), the restricted time measure is zero and the same argument remains valid.

\section{\texorpdfstring{Jointly measurable density and relative score}{Jointly measurable density and relative score}}
\label{proof-module-UnadjustedLangevin-LSI-GaussianMeasurability-lean}
{\small\color{gray}Source module: \texttt{Unadjusted\allowbreak{}Langevin/\allowbreak{}LSI/\allowbreak{}Gaussian\allowbreak{}Measurability.\allowbreak{}lean}.\par}


Spacetime integration requires joint measurability in time and state, rather than measurability separately at each fixed time. This module constructs a single measurable extension of the actual positive-time relative score and proves that it agrees with the gradient used in Fisher information.

\subsection{\texorpdfstring{Integrate a jointly measurable Gaussian kernel}{Integrate a jointly measurable Gaussian kernel}}


The density is a nonnegative parameter integral \(p_t(y)=\int K_t(x,y)d\mu(x)\). The normalizer, affine center, norm, and exponential are measurable in \((t,y,x)\). The parameterized nonnegative integral is therefore measurable. Passing from the extended integral to its real value preserves measurability. Multiplication by \(Z e^{U(y)}\) gives the joint measurability of \(r_t(y)\).



The definitions cover all real times using Lean’s total real operations; only positive times are used for density and score identities. This avoids defining a different measurable representative for each interval.

\subsection{\texorpdfstring{Measurability of the spatial derivative}{Measurability of the spatial derivative}}


The derivative kernel is a continuous linear functional proportional to the inner product with the Gaussian residual:

\[D_yK_t(x,y)[v]=-\frac{K_t(x,y)}{2t}\langle y-x+t\nabla U(x),v\rangle.\]

The inner-product-to-dual map is continuous, so this is jointly strongly measurable as a continuous-linear-map-valued function. Its Bochner integral in \(x\), denoted \(A_t(y)\), remains strongly measurable. At positive times the earlier differentiation theorem identifies \(A_t=D p_t\).

\subsection{\texorpdfstring{One relative-\allowbreak{}score field for all spacetime integrals}{One relative-score field for all spacetime integrals}}


Define

\[\mathcal S(t,y)=\nabla U(y)+p_t(y)^{-1}\mathcal R^{-1}(A_t(y)),\]

where \(\mathcal R^{-1}\) is the inverse Riesz identification of the continuous dual with the Hilbert space. Continuous scalar multiplication, inversion as a measurable real function, and the previous measurability results make \(\mathcal S\) jointly strongly measurable.



For \(t>0\), \(p_t>0\), so the log chain rule and \(r_t=Z e^U p_t\) give

\[\mathcal S(t,y)=\nabla U(y)+\nabla\log p_t(y)=\nabla\log r_t(y).\]

Thus the measurable field is the actual relative score wherever it is used. Its auxiliary extension at nonpositive time introduces no analytic premise into the proof.

\section{\texorpdfstring{Data processing and a finite-\allowbreak{}entropy change of reference}{Data processing and a finite-entropy change of reference}}
\label{proof-module-UnadjustedLangevin-LSI-EntropyEndpoint-lean}
{\small\color{gray}Source module: \texttt{Unadjusted\allowbreak{}Langevin/\allowbreak{}LSI/\allowbreak{}Entropy\allowbreak{}Endpoint.\allowbreak{}lean}.\par}


A Gaussian ULA interpolation applies the same transition to the input law and the Gibbs law. Data processing controls their mutual entropy. A one-sided upper bound on the updated Gibbs density then changes the reference back to the original target, proving finiteness as well as an upper bound. Propagated gradient moments make the logarithmic density correction integrable, without any initial position-moment assumption.

\subsection{\texorpdfstring{kl\allowbreak{}Div prod same right}{klDiv prod same right}}


Adjoin the same independent probability innovation \(\eta\) to two finite measures. The log likelihood of \(\mu\otimes\eta\) relative to \(\nu\otimes\eta\) depends only on the first coordinate, so KL is unchanged. The formal proof uses Mathlib's conditional-product KL identity with the constant Markov kernel \(x\mapsto\eta\); its conditional products are exactly the ordinary product measures. This equality supplies the independent-noise step in data processing.

\subsection{\texorpdfstring{llr eq add of absolutely\allowbreak{}Continuous}{llr eq add of absolutelyContinuous}}


For \(\mu\ll\rho\ll\nu\), the Radon-Nikodym chain rule gives \(d\mu/d\nu=(d\mu/d\rho)(d\rho/d\nu)\) almost everywhere under the integrated law \(\mu\). Each factor is positive and finite there: positivity is transported along absolute continuity, and finiteness uses the probability measures' finite Radon-Nikodym derivatives.



Convert the product of extended densities to real numbers, then apply the logarithm product rule. This proves \(\log(d\mu/d\nu)=\log(d\mu/d\rho)+\log(d\rho/d\nu)\), as an almost-everywhere identity under \(\mu\). The positivity and finiteness checks prevent using logarithm algebra at a zero or infinite density.

\subsection{\texorpdfstring{integrable exp neg llr}{integrable exp neg llr}}


For \(\mu\ll\nu\), the reciprocal likelihood \(e^{-\log(d\mu/d\nu)}\) agrees \(\mu\)-almost everywhere with the real reverse Radon-Nikodym derivative \((d\nu/d\mu)^{\mathrm{real}}\). This derivative is integrable under \(\mu\), even when \(\nu\) has a singular component relative to \(\mu\). Thus the negative exponential of the log likelihood is always integrable; finite KL has not been assumed at this stage.

\subsection{\texorpdfstring{integrable llr of integrable upper bound}{integrable llr of integrable upper bound}}


Suppose the log likelihood \(\ell=\log(d\mu/d\nu)\) has an integrable upper bound \(b\). Where \(\ell\ge0\), \(|\ell|\le b\le|b|\). Where \(\ell<0\), the elementary exponential bound gives \(|\ell|=-\ell\le e^{-\ell}\). Hence globally \(|\ell|\le|b|+e^{-\ell}\). The first term is integrable by hypothesis, and the second by reciprocal-likelihood integrability. Measurable domination proves \(\ell\in L^1(\mu)\), controlling its negative tail without a lower density bound.

\subsection{\texorpdfstring{kl\allowbreak{}Div finite and le of reference density bound}{klDiv finite and le of reference density bound}}


Assume finite \(H(\mu\mid\rho)\), \(\rho\ll\nu\), and \(d\rho/d\nu\le e^b\), with \(b\in L^1(\mu)\). Finite KL gives \(\mu\ll\rho\), hence \(\mu\ll\nu\). The density bound transfers to \(\mu\)-almost every point, where positivity of the intermediate density permits logarithms. Thus \(\log(d\rho/d\nu)\le b\).



The log-likelihood chain rule gives \(\log(d\mu/d\nu)\le\log(d\mu/d\rho)+b\). Its right side is integrable. The preceding one-sided integrability theorem therefore proves the new log likelihood integrable and hence \(H(\mu\mid\nu)<\infty\). Only after proving finiteness do we integrate the inequality, yielding \(H(\mu\mid\nu)\le H(\mu\mid\rho)+\mathbb E_\mu b\). Finite entropy against the new reference is a conclusion, rather than an additional premise.

\subsection{\texorpdfstring{interpolation\allowbreak{}Law kl\allowbreak{}Div le}{interpolationLaw klDiv le}}


The interpolation is the measurable image of the initial law times standard Gaussian noise under \((x,z)\mapsto x-t\nabla U(x)+\sqrt{2t}z\). Apply KL data processing to this same map for both initial laws. The independent-noise product identity removes the common Gaussian factor, giving \(H(\rho_t^\mu\mid\rho_t^\nu)\le H(\mu\mid\nu)\). This argument compares two evolved laws and does not assert that either law is invariant under the discrete step.

\subsection{\texorpdfstring{interpolation\allowbreak{}Law gradient moment}{interpolationLaw gradient moment}}


Write \(\nabla U(Y_t)=[\nabla U(Y_t)-\nabla U(X)]+\nabla U(X)\). The initial gradient is \(L^2\), and the frozen-error energy theorem makes the bracketed discrepancy \(L^2\) on the independent initial/\allowbreak{}noise space. Their sum is \(L^2\); the map-measure equivalence transfers this to the actual interpolation law.



Pointwise, the squared norm of the sum is at most twice each squared norm. Integrate and insert \(\mathbb E\|\nabla U(Y_t)-\nabla U(X)\|^2\le L^2(t^2G_0+2td)\). The initial-coordinate marginal gives the remaining \(G_0\), so \(G_t\le2(1+L^2t^2)G_0+4tdL^2\). This proves both integrability and the quantitative estimate without a position moment.

\subsection{\texorpdfstring{interpolation\allowbreak{}Law entropy bound of density bound}{interpolationLaw entropy bound of density bound}}


Let \(\rho_t\) be the evolved input and \(\sigma_t\) the evolved Gibbs target. Data processing gives \(H(\rho_t\mid\sigma_t)\le H_0<\infty\). Suppose the actual reference law satisfies \(d\sigma_t/d\nu\le\exp(b+c\|\nabla U\|^2)\), with its absolute continuity retained explicitly.



Finite initial KL first supplies finite \(G_0\) through the direct entropy/\allowbreak{}score estimate. The propagated gradient-moment theorem then proves \(G_t<\infty\). Consequently the logarithmic reference correction \(b+c\|\nabla U\|^2\) is integrable under \(\rho_t\). Apply the finite-entropy change-of-reference theorem to obtain \(H(\rho_t\mid\nu)<\infty\) and \(H(\rho_t\mid\nu)\le H_0+b+cG_t\). The pointwise density comparison itself is established in the Gibbs interpolation module.

\section{\texorpdfstring{A Gaussian density comparison at the initial endpoint}{A Gaussian density comparison at the initial endpoint}}
\label{proof-module-UnadjustedLangevin-LSI-GibbsInterpolationBound-lean}
{\small\color{gray}Source module: \texttt{Unadjusted\allowbreak{}Langevin/\allowbreak{}LSI/\allowbreak{}Gibbs\allowbreak{}Interpolation\allowbreak{}Bound.\allowbreak{}lean}.\par}


The entropy limit at zero is obtained by comparing the updated Gibbs density with the original Gibbs density. This calculation uses the differentiable potential and its Lipschitz gradient directly, with no initial density smoothness or change-of-variables theorem.

\subsection{\texorpdfstring{The Gaussian normalizer ratio}{The Gaussian normalizer ratio}}


For \(t>0\), integration of the radial Gaussian gives \(N_t=(4\pi t)^{-d/2}\). Therefore for positive \(s,t\),

\[N_t/N_s=(s/t)^{d/2}.\]

The proof obtains the individual integrals from the finite-dimensional Gaussian integral, uses positivity to invert them, and applies the multiplication and division identities for real powers.

\subsection{\texorpdfstring{Complete the square using only first derivatives}{Complete the square using only first derivatives}}


Write \(g_x=\nabla U(x)\), \(g_y=\nabla U(y)\), and choose positive \(s,t\) with \(s(1-4tL)=t\). The smooth upper model at \(y\) bounds \(U(y)-U(x)\) by \(\langle g_y,y-x\rangle+(L/2)\|y-x\|^2\). Lipschitzness gives \(\|g_x-g_y\|\le L\|x-y\|\). Expanding squares and bounding their cross term yields

\[U(y)-U(x)-\frac{\|y-x+tg_x\|^2}{4t}\le-\frac{\|x-y+sg_y\|^2}{4s}+\frac s4\|g_y\|^2.\]

The relation between \(s\) and \(t\) absorbs the Lipschitz error in the negative quadratic. Exponentiation preserves order. Multiplying by the positive Gaussian normalizer rewrites this as an upper bound on the Gibbs-weighted transition kernel.

\subsection{\texorpdfstring{Integrate the comparison against the Gibbs start}{Integrate the comparison against the Gibbs start}}


The relative density of the updated Gibbs law is

\[r_{\nu,t}(y)=\int N_t\exp\!\left(U(y)-U(x)-\frac{\|y-x+t\nabla U(x)\|^2}{4t}\right)dx.\]

Insert the pointwise comparison. The remaining \(x\)-dependent expression is a normalized Gaussian with time \(s\) centered at \(y-s\nabla U(y)\), so its integral is one. The factors independent of \(x\) give

\[r_{\nu,t}(y)\le(s/t)^{d/2}\exp\!\left(\frac s4\|\nabla U(y)\|^2\right).\]

Normalization of the original Gibbs probability cancels its partition function. Taking \(s=t/(1-4tL)\) later supplies an explicit reference-density bound whose logarithmic cost tends to zero.

\section{\texorpdfstring{The initial entropy endpoint from an explicit upper envelope}{The initial entropy endpoint from an explicit upper envelope}}
\label{proof-module-UnadjustedLangevin-LSI-InitialEntropy-lean}
{\small\color{gray}Source module: \texttt{Unadjusted\allowbreak{}Langevin/\allowbreak{}LSI/\allowbreak{}Initial\allowbreak{}Entropy.\allowbreak{}lean}.\par}


The scalar entropy comparison needs an upper boundary at time zero, even if the initial density is irregular. This module constructs a continuous explicit function \(B(t)\) with \(B(0)=H_0\) and actual interpolation entropy \(H_t\le B(t)\) for sufficiently small positive \(t\). The proof uses the updated Gibbs density bound and data processing; no initial Fisher information, state second moment, or assumed entropy continuity is required.

\subsection{\texorpdfstring{interpolation\allowbreak{}Entropy\allowbreak{}Majorant}{interpolationEntropyMajorant}}


Put \(\delta(t)=1-4tL\) and \(G_0=\mathbb E_\mu\|\nabla U\|^2\). Define the explicit envelope \(B(t)=H_0-\frac d2\log\delta(t)+\frac{t}{4\delta(t)}[2(1+L^2t^2)G_0+4tdL^2]\). The logarithmic term is the normalizer correction from the updated Gibbs density; the final term uses the propagated gradient moment. The definition is a scalar function at all times, while its entropy interpretation below requires \(t>0\) and \(\delta(t)>0\).

\subsection{\texorpdfstring{interpolation\allowbreak{}Entropy\allowbreak{}Majorant zero}{interpolationEntropyMajorant zero}}


At \(t=0\), the logarithm has argument \(\delta(0)=1\), hence is zero. The coefficient \(t/(4\delta(t))\) is zero as well. Both correction terms vanish exactly, proving \(B(0)=H_0\). This is a direct algebraic evaluation of the constructed envelope, not an assumed limit of the actual entropy path.

\subsection{\texorpdfstring{interpolation\allowbreak{}Entropy\allowbreak{}Majorant continuous\allowbreak{}At zero}{interpolationEntropyMajorant continuousAt zero}}


The denominator \(\delta(t)\) has value one at zero, so it stays nonzero locally. The logarithm is continuous at that same positive value. All other terms in \(B\) are polynomial or constant scalar operations. Their continuity proves \(B(t)\to B(0)=H_0\) as \(t\to0\). The initial gradient integral is an already defined scalar; its genuine finiteness in the entropy application is established by the finite-KL moment theorem.

\subsection{\texorpdfstring{interpolation\allowbreak{}Law entropy le majorant}{interpolationLaw entropy le majorant}}


For positive \(t\) with \(\delta=1-4tL>0\), choose the comparison Gaussian time \(s=t/\delta\). The updated Gibbs law has an explicit relative density bounded by \((s/t)^{d/2}\exp(s\|\nabla U(y)\|^2/4)\). Since \(s/t=\delta^{-1}\), this is \(\exp[-(d/2)\log\delta+(s/4)\|\nabla U(y)\|^2]\).



The actual relative density identifies the Radon-Nikodym derivative and proves absolute continuity of the updated Gibbs reference law. The change-of-reference entropy estimate then gives finite \(H_t\) and \(H_t\le H_0-(d/2)\log\delta+(s/4)G_t\). The coefficient \(s/4\) is nonnegative, so inserting \(G_t\le2(1+L^2t^2)G_0+4tdL^2\) proves \(H_t\le B(t)\).

\subsection{\texorpdfstring{interpolation\allowbreak{}Law entropy eventually le majorant}{interpolationLaw entropy eventually le majorant}}


For all sufficiently small positive \(t\), specifically \(t<1/(8L)\), one has \(1-4tL>0\). The previous positive-time entropy bound therefore applies eventually in the right neighborhood of zero, giving \(H_t\le B(t)\). Combined with continuity of \(B\) and \(B(0)=H_0\), this supplies the upper initial boundary used in scalar comparison. It does not insert an assumption that \(H_t\) is differentiable or continuous at zero.

\section{\texorpdfstring{Integrating entropy dissipation with almost-\allowbreak{}everywhere derivatives}{Integrating entropy dissipation with almost-everywhere derivatives}}
\label{proof-module-UnadjustedLangevin-LSI-IntegratedEntropy-lean}
{\small\color{gray}Source module: \texttt{Unadjusted\allowbreak{}Langevin/\allowbreak{}LSI/\allowbreak{}Integrated\allowbreak{}Entropy.\allowbreak{}lean}.\par}


The actual entropy chain rule yields local absolute continuity and a derivative inequality almost everywhere. These scalar lemmas use precisely that regularity, together with a continuous one-sided upper envelope at time zero. They do not require pointwise derivatives at every time.

\subsection{\texorpdfstring{Integrate the affine differential inequality}{Integrate the affine differential inequality}}


Suppose \(H\) is absolutely continuous on \([s,t]\) and \(H^{\prime}\le-aH+b\) almost everywhere, with \(a>0\). Set \(F(u)=e^{au}(H(u)-b/a)\). Products of absolutely continuous functions on a compact interval are absolutely continuous. The product rule holds almost everywhere and gives \(F^{\prime}(u)=e^{au}(a(H(u)-b/a)+H^{\prime}(u))\le0\).



The fundamental theorem of calculus for absolutely continuous functions gives \(F(t)-F(s)=\int_s^tF^{\prime}\le0\). Multiply by \(e^{-at}>0\), cancel the exponentials, and rearrange:

\[H(t)\le e^{-a(t-s)}H(s)+\frac ba(1-e^{-a(t-s)}).\]

The proof removes no exceptional set by assuming a derivative there; both differentiability and the inequality are used only almost everywhere.

\subsection{\texorpdfstring{Send the positive initial time to zero}{Send the positive initial time to zero}}


Suppose \(H\) is absolutely continuous on every \([s,T]\), \(0<s<T\), and \(H(s)\le M(s)\) for all sufficiently small positive \(s\). If \(M\) is continuous at zero, the preceding estimate gives

\[H(T)\le e^{-a(T-s)}M(s)+\frac ba(1-e^{-a(T-s)}).\]

The right side tends to \(e^{-aT}M(0)+(b/a)(1-e^{-aT})\). Order is preserved under limits, giving the desired initial-endpoint bound. Single endpoints have zero Lebesgue measure, so an inequality on the open interval suffices. Continuity or differentiability of \(H\) at zero is not needed.

\subsection{\texorpdfstring{Absorb the ULA forcing with its exact constants}{Absorb the ULA forcing with its exact constants}}


Assume \(2\alpha H\le I\), \(H^{\prime}\le-3I/4+t^2L^2G_0+2tdL^2\), and \(G_0\le4L^2M(0)/\alpha+2dL\). Let \(0<h\le\alpha/(4L^2)\), \(0<\alpha\le L\), and \(M(0)\ge0\). Then \(hL\le1/4\), \(\alpha h\le1/4\),

\[h^2L^2(4L^2/\alpha)\le\alpha/4,\qquad h^2L^2(2dL)\le hdL^2/2.\]

The derivative is therefore at most \(-3\alpha H/2+B\), where \(B=\alpha M(0)/4+(5/2)hdL^2\). The affine comparison and \(1-e^{-ah}\le ah\) imply \(H(h)\le e^{-3\alpha h/2}M(0)+Bh\). Finally \(e^{-x}\ge e^{-3x/2}+x/4\) for \(0\le x\le1/4\) absorbs the entropy part of \(Bh\). Bounding \(5/2\) by 6 gives

\[H(h)\le e^{-\alpha h}M(0)+6h^2dL^2.\]

The final actual-law theorem proves every analytic premise before invoking this scalar result.

\section{\texorpdfstring{An integrated chain rule from spacetime integrability}{An integrated chain rule from spacetime integrability}}
\label{proof-module-UnadjustedLangevin-LSI-EntropyChainRule-lean}
{\small\color{gray}Source module: \texttt{Unadjusted\allowbreak{}Langevin/\allowbreak{}LSI/\allowbreak{}Entropy\allowbreak{}Chain\allowbreak{}Rule.\allowbreak{}lean}.\par}


To differentiate entropy, bounded approximations are first integrated in time. A pointwise time derivative together with joint spacetime \(L^1\) suffices for Fubini and the fundamental theorem of calculus; no common spatial majorant for every time derivative is assumed. The resulting spatial integral is absolutely continuous and has the integrated derivative almost everywhere. A dominated-limit theorem then passes truncated integral identities to the actual entropy identity.

\subsection{\texorpdfstring{integral eq initial add integral of spacetime integrable}{integral eq initial add integral of spacetime integrable}}


Assume \(f(a,\cdot)\in L^1(\nu)\), \(\dot f\in L^1((a,b]\times\nu)\), and \(\partial_t f(t,x)=\dot f(t,x)\) for every \(t\in[a,b]\) and almost every \(x\). Fubini gives \(\dot f(\cdot,x)\in L^1(a,b)\) for almost every \(x\). Apply the scalar fundamental theorem to these sections: \(f(b,x)=f(a,x)+\int_a^b\dot f(t,x)dt\).



The time-integrated derivative is integrable in \(x\) by spacetime integrability. The section identity therefore proves terminal integrability of \(f(b,\cdot)\), rather than requiring it separately. Integrate the identity and exchange time and space to obtain \(\int f(b)d\nu=\int f(a)d\nu+\int_a^b\int\dot f(t,x)d\nu(x)dt\).

\subsection{\texorpdfstring{absolutely\allowbreak{}Continuous integral of spacetime integrable}{absolutelyContinuous integral of spacetime integrable}}


Apply the integrated fundamental theorem on every subinterval \([a,t]\subseteq[a,b]\). Restricting the spacetime measure preserves integrability, giving both integrability of \(f(t,\cdot)\) and \(F(t)=F(a)+\int_a^t q(u)du\), where \(F(t)=\int f(t,x)d\nu(x)\) and \(q(t)=\int\dot f(t,x)d\nu(x)\).



Fubini gives \(q\in L^1(a,b)\). An indefinite integral of an \(L^1\) function is absolutely continuous; adding the constant \(F(a)\) preserves this property. Equality on the interval transfers absolute continuity to \(F\). Thus regularity of the spatial integral is a conclusion of spacetime integrability and pointwise differentiation.

\subsection{\texorpdfstring{integral has\allowbreak{}Deriv\allowbreak{}At ae of spacetime integrable}{integral hasDerivAt ae of spacetime integrable}}


The subinterval argument gives \(F(t)=F(a)+\int_a^t q(u)du\) for every \(t\in[a,b]\), with \(q\in L^1(a,b)\). Lebesgue differentiation supplies derivative \(q(t)\) for the indefinite integral at almost every interior time. For each such \(t\), an open neighborhood stays inside \((a,b)\), where the displayed identity holds. The local equality transfers the derivative to \(F\), yielding \(F'(t)=\int\dot f(t,x)d\nu(x)\) almost everywhere. No pointwise derivative assertion at every time, or at the endpoints, is inferred.

\subsection{\texorpdfstring{integral identity of dominated limits}{integral identity of dominated limits}}


Suppose truncated identities have the form \(B_n-A_n=\int q_n d\nu\). Assume the two endpoint sequences converge to their intended values, \(q_n\to q\) almost everywhere, and \(|q_n|\le g\) with \(g\in L^1(\nu)\). Dominated convergence gives convergence of the right-hand integrals, while the endpoint limits give convergence of the left-hand differences. Their equality for every \(n\) and uniqueness of limits prove \(B-A=\int qd\nu\).



In the entropy application, the endpoints are truncated entropy integrals and \(q_n\) are bounded-test fluxes on a fixed time interval. Spatial energy estimates supply the pointwise flux limit and a time-integrable dominating bound. This lemma passes the integrated identities to the limit without assuming in advance that the full entropy has already been differentiated.

\subsection{\texorpdfstring{regularity of integral identity}{regularity of integral identity}}


Assume only \(q\in L^1(a,b)\) and the exact identity \(H(t)=H(a)+\int_a^t q(u)du\) throughout \([a,b]\). The primitive of \(q\) is absolutely continuous; adding \(H(a)\) and using the identity proves that \(H\) is absolutely continuous.



Lebesgue differentiation gives derivative \(q(t)\) for the primitive at almost every interior time. The identity holds in an interior neighborhood of each such time, so the derivative transfers to \(H\). This gives precisely the final entropy comparison inputs: local absolute continuity and the actual flux as derivative almost everywhere.

\section{\texorpdfstring{Bounded logarithmic tests and their entropy primitives}{Bounded logarithmic tests and their entropy primitives}}
\label{proof-module-UnadjustedLangevin-LSI-BoundedLog-lean}
{\small\color{gray}Source module: \texttt{Unadjusted\allowbreak{}Langevin/\allowbreak{}LSI/\allowbreak{}Bounded\allowbreak{}Log.\allowbreak{}lean}.\par}


The unbounded logarithm in entropy production is replaced by a smooth bounded function. The replacement has derivative between zero and one, converges to the identity, and admits an entropy primitive dominated by the actual KL integrand. These properties separately justify the spatial test, the time chain rule, and the limit in the entropy identity.

\subsection{\texorpdfstring{Choose the bounded test and its weight}{Choose the bounded test and its weight}}


For a natural number \(m\), write \(c=m+1>0\). Define

\[\chi_m(z)=c\arctan(z/c),\qquad w_m(z)=\frac1{1+(z/c)^2},\qquad\Phi_m(r)=\int_1^r\chi_m(\log u)du.\]

The positive offset in \(c\) prevents division by zero even for \(m=0\). The interval integral is oriented: when \(r<1\), its sign reverses. Lean defines these expressions for every real argument; the entropy claims below explicitly require \(r>0\).

\subsection{\texorpdfstring{Differentiate and bound the test}{Differentiate and bound the test}}


The arctangent chain rule gives \(\chi_m(0)=0\) and \(\chi_m^{\prime}(z)=w_m(z)\). The denominator is at least one, so \(0\le w_m\le1\). The derivative bound implies that \(\chi_m\) is 1-Lipschitz, hence \(|\chi_m(z)|\le|z|\). Independently, \(|\arctan u|<\pi/2\) gives \(|\chi_m(z)|\le c\pi/2\). This last bound controls the time derivative integrand without any log-integrability requirement. Since arctangent has the sign of its argument, \(\chi_m\) does too.



\texttt{Cont\allowbreak{}Diff \ensuremath{\mathbb R} \ensuremath{\top}} records smoothness of the scalar test. \texttt{Has\allowbreak{}Deriv\allowbreak{}At} gives its exact derivative; the proof multiplies the derivative of arctangent by \(1/c\) and cancels the outer factor \(c\).

\subsection{\texorpdfstring{Recover the logarithm and its primitive}{Recover the logarithm and its primitive}}


As \(m\to\infty\), \(z/(m+1)\to0\). The derivative of arctangent at zero is one; its difference quotient therefore gives \(\chi_m(z)\to z\). At \(z=0\) the claim is exact at every index. Continuity of inversion at the positive denominator one gives \(w_m(z)\to1\).



The function \(u\mapsto\chi_m(\log u)\) is measurable and bounded, hence integrable on every finite interval, including intervals containing zero. When \(r>0\), continuity in a neighborhood of \(r\) and the fundamental theorem of calculus imply \(\Phi_m^{\prime}(r)=\chi_m(\log r)\).

\subsection{\texorpdfstring{Dominate the entropy endpoints}{Dominate the entropy endpoints}}


We have \(\Phi_m(1)=0\). For \(r\ge1\), both \(\log u\) and its test are nonnegative on \([1,r]\), giving \(\Phi_m(r)\ge0\). For \(0<r\le1\), the test is nonpositive on \([r,1]\), and reversing the interval again gives nonnegativity.



On the positive side, \(0\le\chi_m(\log u)\le\log u\). On the negative side, \(\log u\le\chi_m(\log u)\le0\); reversing the interval preserves the corresponding upper bound for the primitive. In both cases,

\[0\le\Phi_m(r)\le\int_1^r\log u\,du=r\log r-r+1.\]

Pointwise convergence of the test and domination by \(|\log u|\), integrable on the finite interval between 1 and \(r\), prove \(\Phi_m(r)\to r\log r-r+1\). These are exactly the hypotheses for dominated convergence against the finite-KL relative density later.

\subsection{\texorpdfstring{Identify the spatial gradient of the test}{Identify the spatial gradient of the test}}


For a differentiable scalar field \(f\), the Fréchet chain rule gives \(D(\chi_m\circ f)=w_m(f)Df\). Applying the Hilbert-space identification of linear derivatives with gradients yields

\[\nabla(\chi_m\circ f)(x)=w_m(f(x))\nabla f(x),\qquad\|\nabla(\chi_m\circ f)(x)\|\le\|\nabla f(x)\|.\]

The inequality uses \(0\le w_m\le1\). With \(f=\log r_t\), it supplies square-integrability of the bounded log test gradient from the actual Fisher bound.

\section{\texorpdfstring{Derive the actual Gaussian weak equation for bounded finite-\allowbreak{}energy tests}{Derive the actual Gaussian weak equation for bounded finite-energy tests}}
\label{proof-module-UnadjustedLangevin-LSI-GaussianWeakEquation-lean}
{\small\color{gray}Source module: \texttt{Unadjusted\allowbreak{}Langevin/\allowbreak{}LSI/\allowbreak{}Gaussian\allowbreak{}Weak\allowbreak{}Equation.\allowbreak{}lean}.\par}


Write \(g(x)=\nabla U(x)\), \(c_t(x)=x-tg(x)\), and \(K_t(x,y)=N_t e^{-\|y-c_t(x)\|^2/(4t)}\). The spatial flux \(F_{t,x}(y)=K_t(x,y)(g(x)-(y-c_t(x))/(2t))\) has divergence equal to the actual Gaussian kernel time derivative. Global integration by parts therefore gives \(\int\dot p_t\psi=-\int F_{t,x}\cdot\nabla\psi\,d\mu\,dy\) for every bounded \(C^1\) test with \(\nabla\psi\in L^2(\rho_t)\). This module proves all integrability and boundary-product obligations from finite initial gradient energy and Gaussian second moments. It differentiates no Hessian of \(U\) and assumes no weak equation.

\subsection{\texorpdfstring{Choose the spatial flux with the correct Fokker–Planck sign}{Choose the spatial flux with the correct Fokker–Planck sign}}


The component density evolves under frozen drift \(-g(x)\). Its time derivative is \(\Delta_yK_t+\operatorname{div}_y(K_tg(x))\), so define

\[F_{t,x}(y)=\nabla_yK_t(x,y)+K_t(x,y)g(x)=K_t(x,y)\left(g(x)-\frac{y-c_t(x)}{2t}\right).\]

\texttt{frozen\allowbreak{}Gaussian\allowbreak{}Flux} is the last expression. Its divergence equals the density's time derivative, while the weak equation consequently carries a minus sign after integration by parts. The frozen gradient \(g(x)\) is constant in the spatial variable \(y\), which avoids differentiating the possibly merely Lipschitz gradient field.

\subsection{\texorpdfstring{Differentiate each coordinate of the flux explicitly}{Differentiate each coordinate of the flux explicitly}}


Let \(r=y-c_t(x)\). For a direction \(v\), differentiating the kernel and the affine residual gives

\[D_v\langle F_{t,x}(y),v\rangle=K_t(x,y)\left(\frac{\langle r,v\rangle^2}{4t^2}-\frac{\langle r,v\rangle\langle g(x),v\rangle}{2t}-\frac{\|v\|^2}{2t}\right).\]

The Lean proof constructs the derivative of the affine vector factor as \texttt{hlin}, pairs it with the constant vector \(v\) as \texttt{hi}, and multiplies it by the already certified derivative of the real Gaussian kernel. The function equality \texttt{heq} moves the scalar density outside the inner product. Evaluation of the resulting linear map at \(v\), the inner-product rules, and scalar algebra give exactly the three displayed terms. This identity is valid without assuming the potential has a second derivative.

\subsection{\texorpdfstring{Sum the directional derivatives to obtain a coordinate-\allowbreak{}free divergence}{Sum the directional derivatives to obtain a coordinate-free divergence}}


Sum over an orthonormal basis \((b_i)_{i=1}^d\). Parseval gives \(\sum_i\langle r,b_i\rangle^2=\|r\|^2\), completeness gives \(\sum_i\langle r,b_i\rangle\langle g,b_i\rangle=\langle r,g\rangle\), and \(\sum_i\|b_i\|^2=d\). Thus

\[\operatorname{div}F_{t,x}=K_t\left(-\frac d{2t}+\frac{\|r\|^2}{4t^2}-\frac{\langle r,g\rangle}{2t}\right)=K_t\left(-\frac d{2t}+\frac{\|y-x\|^2}{4t^2}-\frac{\|g\|^2}{4}\right).\]

Lean records the two basis identities as \texttt{hsq} and \texttt{hcross}. The residual equality \texttt{hres} is \(r=y-x+tg\); expanding its squared norm cancels the mixed term. The explicit nonzero-time assumption justifies the denominator algebra. This is the same coefficient as the independently derived kernel time derivative.

\subsection{\texorpdfstring{Transport the actual independent pair to its Gaussian Lebesgue density}{Transport the actual independent pair to its Gaussian Lebesgue density}}


The pair \((X,Y_t)\), constructed from independent initial state and standard noise, satisfies

\[\mathcal L(X,Y_t)(dx,dy)=K_t(x,y)\,\mu(dx)\,dy.\]

The theorem \texttt{ULA.\allowbreak{}gaussian\_\allowbreak{}proposal\_\allowbreak{}pair\_\allowbreak{}density} already proves this transport for a measurable center map. Lean applies it to \(x\mapsto x-t\nabla U(x)\), whose measurability follows from the Lipschitz gradient. It then checks that the real kernel here is exactly the previously normalized proposal density. This identity is essential: it lets the proof establish integrability under a simple independent probability law and transfer the result to the Lebesgue pair integrals needed for integration by parts.

\subsection{\texorpdfstring{Prove flux pairings integrable by transporting two L2 factors}{Prove flux pairings integrable by transporting two L2 factors}}


For a continuous scalar field \(\varphi\in L^2(\rho_t)\), the coordinate flux pairing is integrable:

\[\int K_t(x,y)\left|\left\langle g(x)-\frac{y-c_t(x)}{2t},v\right\rangle\varphi(y)\right|\,d\mu(x)\,dy<\infty.\]

Under the independent pair \((X,Z)\), the vector factor becomes \(g(X)-(\sqrt{2t}/(2t))Z\). Both terms are in \(L^2\), by the initial gradient hypothesis and the Gaussian second moment. Pulling back \(\varphi\) through the actual update preserves its \(L^2\) property. Lean's \texttt{hflux} and \texttt{htest} are these two facts; \texttt{Mem\allowbreak{}Lp.\allowbreak{}mul} puts their scalar product in \(L^1\). The pair-law density identity and the density-integrability equivalence then yield the displayed weighted Lebesgue integrability. No unproved decay of the test at infinity is used.

\subsection{\texorpdfstring{Admit the gradient pairing of a C1 finite-\allowbreak{}energy test}{Admit the gradient pairing of a C1 finite-energy test}}


Choose the scalar field \(\varphi(y)=D\psi(y)[v]=\langle\nabla\psi(y),v\rangle\). For \(\psi\in C^1\), this field is continuous. If \(\nabla\psi\in L^2(\rho_t)\), its coordinate pairing is also in \(L^2\).



Lean applies the preceding general flux-pairing theorem. Continuity follows from continuity of the Fréchet derivative of a \(C^1\) function, evaluated at \(v\). The equality \texttt{inner\_\allowbreak{}gradient\_\allowbreak{}left} identifies the coordinate derivative with the inner product, and \texttt{Mem\allowbreak{}Lp.\allowbreak{}inner\_\allowbreak{}const} supplies its square integrability. This establishes the \(F_i\partial_i\psi\) integrability premise of global integration by parts.

\subsection{\texorpdfstring{Control the boundary product from boundedness of the test}{Control the boundary product from boundedness of the test}}


A bounded continuous test satisfies \(|\psi|\le M\), hence belongs to \(L^2(\rho_t)\) because \(\rho_t\) is a probability measure. The general flux-pairing result therefore gives integrability of \(\langle F_{t,x}(y),v\rangle\psi(y)\).



Lean installs the actual interpolation probability instance, uses \texttt{Mem\allowbreak{}Lp.\allowbreak{}of\_\allowbreak{}bound}, and applies \texttt{frozen\allowbreak{}Gaussian\allowbreak{}Flux\_\allowbreak{}mul\_\allowbreak{}integrable}. This is precisely the product \(F_i\psi\) required by the integration-by-parts theorem on the whole Euclidean space. Boundedness and finite-energy flux replace a separate compact-support assumption.

\subsection{\texorpdfstring{Match the divergence to the independently differentiated Gaussian kernel}{Match the divergence to the independently differentiated Gaussian kernel}}


The explicit time derivative from \texttt{Gaussian\allowbreak{}Time\allowbreak{}Derivative.\allowbreak{}lean} is

\[\dot K_t(x,y)=K_t(x,y)\left(-\frac d{2t}+\frac{\|y-x\|^2}{4t^2}-\frac{\|\nabla U(x)\|^2}{4}\right).\]

The divergence calculation gives exactly this expression. Lean rewrites the sum of coordinate derivatives with \texttt{sum\_\allowbreak{}fderiv\_\allowbreak{}frozen\allowbreak{}Gaussian\allowbreak{}Flux}, unfolds the independently defined \texttt{frozen\allowbreak{}Gaussian\allowbreak{}Kernel\allowbreak{}Dot}, and matches the real kernel definitions. Thus the time/\allowbreak{}spatial relation used in the PDE is an equality between actual derivatives, not an asserted Fokker–Planck premise.

\subsection{\texorpdfstring{Dominate each unit-\allowbreak{}coordinate flux derivative by a quadratic Gaussian energy}{Dominate each unit-coordinate flux derivative by a quadratic Gaussian energy}}


For \(\|v\|=1\), Cauchy–Schwarz bounds \(|\langle r,v\rangle|\le\|r\|\) and \(|\langle g,v\rangle|\le\|g\|\). The inequality \(\|r\|\|g\|\le\|r\|^2+\|g\|^2\) then gives

\[\left|D_v\langle F_{t,x},v\rangle\right|\le\left(\frac1{4t^2}+\frac1{2t}\right)K_t(x,y)(1+\|g(x)\|^2+\|r\|^2).\]

Lean records the coordinate, square, and mixed-product bounds as \texttt{hi}, \texttt{hg}, \texttt{hsq}, and \texttt{hcross}. The local \texttt{hbound} applies the triangle inequality to the three derivative terms and dominates their coefficients by the displayed envelope. Multiplication by the nonnegative Gaussian density preserves the bound. Only quadratic energies occur, so no fourth moment of the initial gradient is introduced.

\subsection{\texorpdfstring{Integrate the coordinate derivative envelope under the actual pair measure}{Integrate the coordinate derivative envelope under the actual pair measure}}


The common envelope \(K_t(1+\|g(X)\|^2+\|Y-c_t(X)\|^2)\) is jointly integrable: its integral is \(1+G_0+2td\), where \(G_0=\int\|\nabla U\|^2d\mu\). The preceding pointwise bound therefore proves joint integrability of every unit-coordinate flux derivative.



Lean obtains this envelope from \texttt{integrable\_\allowbreak{}frozen\allowbreak{}Gaussian\allowbreak{}Energy} and scales it by the finite time coefficient. To supply measurability, \texttt{heq} rewrites the derivative as its explicit polynomial-Gaussian formula; continuity of the gradient then lets \texttt{fun\_\allowbreak{}prop} prove the required measurability. \texttt{Integrable.\allowbreak{}mono'} finishes using the coordinate energy bound. These are the actual derivative integrals needed before applying Fubini or integration by parts.

\subsection{\texorpdfstring{Apply global integration by parts after proving all three product integrabilities}{Apply global integration by parts after proving all three product integrabilities}}


For one unit coordinate \(v\), global integration by parts gives

\[\int D_v\langle F_{t,x}(y),v\rangle\,\psi(y)\,d\mu(x)\,dy=-\int\langle F_{t,x}(y),v\rangle D_v\psi(y)\,d\mu(x)\,dy.\]

The proof first constructs all three required integrable products: \texttt{h\allowbreak{}Df} is the derivative of the flux times the bounded test, \texttt{h\allowbreak{}Fg} is the flux times the finite-energy test derivative, and \texttt{h\allowbreak{}Ff} is the boundary product. Fubini gives the corresponding section integrability for almost every initial state. At each such state, the Gaussian flux coordinate is differentiable and the test is differentiable because it is \(C^1\). Lean applies \texttt{integral\_\allowbreak{}mul\_\allowbreak{}fderiv\_\allowbreak{}eq\_\allowbreak{}neg\_\allowbreak{}fderiv\_\allowbreak{}mul\_\allowbreak{}of\_\allowbreak{}integrable} on the whole space, then integrates the section identity back over the initial law. The theorem's boundary condition is integrability of the actual product, not an unproved pointwise decay statement.

\subsection{\texorpdfstring{Derive the actual weak PDE for every bounded C1 finite-\allowbreak{}energy test}{Derive the actual weak PDE for every bounded C1 finite-energy test}}


Sum the coordinate identities over an orthonormal basis. The divergence identity turns the sum on the left into \(\dot K_t\), while completeness of the basis turns the sum on the right into \(\langle F_{t,x},\nabla\psi\rangle\). Fubini then yields

\[\int\dot p_t(y)\psi(y)\,dy=-\int\langle F_{t,x}(y),\nabla\psi(y)\rangle\,d\mu(x)\,dy.\]

Lean calls the coordinate integrands \texttt{A} and \texttt{B}. Their integrability is \texttt{h\allowbreak{}A} and \texttt{h\allowbreak{}B}; their integral equality is \texttt{h\allowbreak{}Eq}. The identities \texttt{hsum\allowbreak{}A} and \texttt{hsum\allowbreak{}B} perform the two basis sums pointwise. Kernel time-derivative integrability times boundedness of the test gives \texttt{h\allowbreak{}Dot}, which justifies exchanging its initial-state and proposed-state integrals. Finally \texttt{integral\_\allowbreak{}finset\allowbreak{}Sum} commutes the finite sums and integrals. Every premise has been derived from the explicit Gaussian law, initial gradient energy, and the ordinary bounded finite-energy test conditions.

\section{\texorpdfstring{Relative Fisher energy and the sharp frozen-\allowbreak{}gradient error}{Relative Fisher energy and the sharp frozen-gradient error}}
\label{proof-module-UnadjustedLangevin-LSI-GaussianEnergyIntegrability-lean}
{\small\color{gray}Source module: \texttt{Unadjusted\allowbreak{}Langevin/\allowbreak{}LSI/\allowbreak{}Gaussian\allowbreak{}Energy\allowbreak{}Integrability.\allowbreak{}lean}.\par}


Let \(\rho_t\) be the actual frozen Gaussian interpolation, \(r_t=d\rho_t/d\nu\), \(s_t=\nabla\log r_t\), and \(G_0=\mathbb E_\mu\|\nabla U\|^2\). This module proves \(I_t\le4(1+L^2t^2)G_0+8tdL^2+d/t\), local spacetime integrability of the weighted relative-score energy, and the exact pair-law error bound \(\mathbb E\|\nabla U(Y_t)-\nabla U(X)\|^2\le L^2(t^2G_0+2td)\). These estimates provide domination for the unbounded entropy test and its truncation limits.

\subsection{\texorpdfstring{norm add sq le two}{norm add sq le two}}


The triangle inequality gives \(\|x+y\|\le\|x\|+\|y\|\). Square both nonnegative sides and use \(2\|x\|\|y\|\le\|x\|^2+\|y\|^2\), which follows from \((\|x\|-\|y\|)^2\ge0\). This proves \(\|x+y\|^2\le2\|x\|^2+2\|y\|^2\), the pointwise estimate used to split relative Fisher energy.

\subsection{\texorpdfstring{interpolation\allowbreak{}Relative\allowbreak{}Score\allowbreak{}Field energy le}{interpolationRelativeScoreField energy le}}


The jointly measurable score field agrees at positive time with the actual relative score. Its decomposition is \(s_t=\nabla U+\nabla\log p_t\). The propagated gradient is square integrable, and Gaussian smoothing gives \(\mathbb E_{\rho_t}\|\nabla\log p_t\|^2\le d/(2t)\). These facts also prove integrability of the relative-score square.



Apply the squared sum bound pointwise and integrate: \(I_t\le2G_t+2\mathbb E_{\rho_t}\|\nabla\log p_t\|^2\le2G_t+d/t\). The proof keeps the integrability witnesses for all terms before using integral monotonicity.

\subsection{\texorpdfstring{interpolation\allowbreak{}Relative\allowbreak{}Score\allowbreak{}Field energy bound}{interpolationRelativeScoreField energy bound}}


The interpolation gradient-moment estimate is \(G_t\le2(1+L^2t^2)G_0+4tdL^2\). Substitute it into \(I_t\le2G_t+d/t\) and distribute the factor two. This gives \(I_t\le4(1+L^2t^2)G_0+8tdL^2+d/t\). Every quantity on the right is determined by the initial gradient energy and the original potential; no entropy-production assumption is used.

\subsection{\texorpdfstring{interpolation\allowbreak{}Relative\allowbreak{}Score\allowbreak{}Field weighted energy}{interpolationRelativeScoreField weighted energy}}


The explicitly constructed relative density satisfies \(r_t\nu=\rho_t\). The with-density integrability equivalence therefore transports integrability of \(\|s_t\|^2\) under \(\rho_t\) to integrability of \(r_t\|s_t\|^2\) under the fixed target \(\nu\). The corresponding integral identity is \(\int r_t\|s_t\|^2d\nu=\int\|s_t\|^2d\rho_t\). Strict positivity and measurability of the constructed density justify the ordinary weighted-integral formulas.

\subsection{\texorpdfstring{integrable interpolation\allowbreak{}Relative\allowbreak{}Score\allowbreak{}Field energy spacetime}{integrable interpolationRelativeScoreField energy spacetime}}


The relative density and the chosen relative-score field are jointly measurable in time and space. Their weighted energy is nonnegative. On a positive interval \([a,b]\), the section at every time is integrable by the density-change theorem.



Since \(0<a\le t\le b\), the quantitative estimate is bounded uniformly by \(4(1+L^2b^2)G_0+8bdL^2+d/a\). The integral of the section norm is exactly its nonnegative energy, so this finite scalar bound makes it integrable in time. The product criterion yields \(r_t\|s_t\|^2\in L^1(dt\,d\nu)\) on the interval. This is the spacetime dominating function required to pass from bounded entropy tests to the full relative Fisher term.

\subsection{\texorpdfstring{frozen\allowbreak{}Gradient\allowbreak{}Error\allowbreak{}Field}{frozenGradientErrorField}}


For the fixed independent pair \((X,Z)\sim\mu\otimes\mathcal N(0,I)\), define \(e_t(X,Z)=\nabla U(X-t\nabla U(X)+\sqrt{2t}Z)-\nabla U(X)\). The underlying probability space is independent of time. This avoids assuming a jointly measurable choice of a state-space conditional expectation when integrating the error energy over time.

\subsection{\texorpdfstring{integrable frozen\allowbreak{}Gradient\allowbreak{}Error\allowbreak{}Field energy spacetime}{integrable frozenGradientErrorField energy spacetime}}


The proposal and Lipschitz gradient make \((t,X,Z)\mapsto\|e_t(X,Z)\|^2\) jointly measurable. For \(t\ge0\), the exact Gaussian calculation already gives integrability and \(\mathbb E\|e_t\|^2\le L^2(t^2G_0+2td)\). On a compact nonnegative time interval this is at most \(L^2(b^2G_0+2bd)\). The time measure is finite, so Fubini's integrability criterion proves the squared error is integrable on the time/\allowbreak{}initial-state/\allowbreak{}noise product.

\subsection{\texorpdfstring{frozen\allowbreak{}Gaussian\allowbreak{}Pair gradient error energy}{frozenGaussianPair gradient error energy}}


Map the independent initial/\allowbreak{}noise pair by \(T_t(x,z)=(x,x-t\nabla U(x)+\sqrt{2t}z)\). The squared error as a function of the output pair is \(F(x,y)=\|\nabla U(y)-\nabla U(x)\|^2\). Its composition with \(T_t\) is exactly the frozen-gradient discrepancy from the Gaussian energy calculation. Measurability of \(T_t\) and continuity of \(F\) let the map-measure integrability and integral theorems transport both the integrability witness and the bound \(L^2(t^2G_0+2td)\) to the actual initial/\allowbreak{}output pair law.

\subsection{\texorpdfstring{frozen\allowbreak{}Gaussian\allowbreak{}Kernel gradient error energy}{frozenGaussianKernel gradient error energy}}


The actual pair law is identified with \(K_t(x,y)\,\mu(dx)dy\) by the previously proved Gaussian proposal pair-density theorem. The proof specializes its measurable center to \(x-t\nabla U(x)\); the precision and step-size versions of the kernel agree pointwise.



Rewrite the pair-law energy theorem using this exact density identity. The with-density integrability equivalence gives \(K_t\|\nabla U(y)-\nabla U(x)\|^2\in L^1(\mu\otimes dy)\), and its integral formula gives the sharp bound \(\int K_t\|\nabla U(y)-\nabla U(x)\|^2d\mu\,dy\le L^2(t^2G_0+2td)\). This transport preserves the exact constants needed in the final entropy recurrence; a cruder pointwise Gaussian envelope is not substituted for the error energy.

\section{\texorpdfstring{The complete entropy chain rule for the actual Gaussian interpolation}{The complete entropy chain rule for the actual Gaussian interpolation}}
\label{proof-module-UnadjustedLangevin-LSI-GaussianEntropyIdentity-lean}
{\small\color{gray}Source module: \texttt{Unadjusted\allowbreak{}Langevin/\allowbreak{}LSI/\allowbreak{}Gaussian\allowbreak{}Entropy\allowbreak{}Identity.\allowbreak{}lean}.\par}


This module proves the entropy identity and positive-time absolute continuity for the actual frozen-drift Gaussian interpolation. The proof starts with bounded smooth log tests, applies the global Gaussian weak equation and a spacetime \(L^1\) time chain rule, then removes both spatial and time truncations by dominated convergence. All regularity and integrability needed in these steps are constructed from the original Gaussian update and the initial gradient energy.



Write \(\pi=Z^{-1}e^{-U}dx\), \(\rho_t=\operatorname{Law}(X-t\nabla U(X)+\sqrt{2t}Z_0)\), \(r_t=d\rho_t/d\pi\), and \(S_t=\nabla\log r_t\). Here \(Z_0\) is an independent standard Gaussian. The eventual conclusion is \(H(\rho_b\mid\pi)-H(\rho_a\mid\pi)=\int_a^b q(t)dt\), with the explicit raw Gaussian flux \(q\), for every \(0<a\le b\). This identity proves local absolute continuity and the a.e. derivative; they are not assumed.

\subsection{\texorpdfstring{Removing bounded log weights by dominated convergence}{Removing bounded log weights by dominated convergence}}


Let \(w_m(z)=[1+(z/(m+1))^2]^{-1}\). For a measurable scalar field \(z\) and an integrable field \(f\),

\[\int w_m(z(x))f(x)\,d\nu(x)\longrightarrow\int f(x)\,d\nu(x).\]

The pointwise limit is \(w_m(z(x))\to1\), proved in \texttt{bounded\allowbreak{}Log\allowbreak{}Weight\_\allowbreak{}tendsto}. Each weighted integrand is measurable, because the weight is obtained by division, squaring and inversion of measurable real functions. Furthermore \(0\le w_m\le1\), so \(|w_m(z)f|\le|f|\). The Lean proof invokes dominated convergence with the single integrable majorant \texttt{hf.\allowbreak{}norm}. This lemma later removes the truncation from the actual Gaussian flux, whose two factors are square integrable.

\subsection{\texorpdfstring{Why the square-\allowbreak{}integrable flux pairing is integrable}{Why the square-integrable flux pairing is integrable}}


If \(v,s\in L^2(\nu;E)\), then

\[|\langle v,s\rangle|\le\|v\|\|s\|,\qquad \|v\|\|s\|\in L^1(\nu).\]

The second assertion follows from Hölder's inequality with exponents \(2,2,1\). In Lean, \texttt{hf.\allowbreak{}norm.\allowbreak{}mul hg.\allowbreak{}norm} supplies membership of the norm product in \(L^1\); \texttt{mem\allowbreak{}Lp\_\allowbreak{}one\_\allowbreak{}iff\_\allowbreak{}integrable} converts it to ordinary integrability without assuming that the ambient measure is finite. The measurable inner product is then dominated by this product using \texttt{norm\_\allowbreak{}inner\_\allowbreak{}le\_\allowbreak{}norm}. This establishes actual integrability rather than relying on the default value of an undefined Bochner integral.

\subsection{\texorpdfstring{The real Gibbs density and its normalization}{The real Gibbs density and its normalization}}


Write \(Z=\int e^{-U(x)}\,dx\), and let \(\pi\) be the normalized Gibbs law. Define

\[J(y)=Z^{-1}e^{-U(y)}.\]

The existing probability-mass argument proves \(0<Z<\infty\); differentiability makes \(U\) measurable and gives strict positivity of the partition function. Consequently \(J(y)>0\) at every point. The measure identity is

\[\pi(dy)=J(y)\,dy.\]

The formal representation of \(\pi\) is initially an extended scalar multiple of \texttt{volume.\allowbreak{}with\allowbreak{}Density}. \texttt{with\allowbreak{}Density\_\allowbreak{}smul} absorbs that scalar into the density. The conversion \texttt{ENNReal.\allowbreak{}of\allowbreak{}Real\_\allowbreak{}to\allowbreak{}Real} is justified by finite nonzero normalization, so the real formula agrees exactly with the original measure. These declarations introduce no strong-convexity assumption.

\subsection{\texorpdfstring{The time derivative of the weighted entropy primitive}{The time derivative of the weighted entropy primitive}}


Let \(p_t\) be the actual Lebesgue density of the frozen Gaussian update, and let \(r_t=p_t/J\) be its relative density. For the bounded entropy primitive \(\Phi_m\), one has \(\Phi_m'(r)=\chi_m(\log r)\) for \(r>0\). At every positive time, the actual kernel differentiation theorem gives \(\partial_t p_t=\dot p_t\), hence

\[\partial_t r_t=Ze^{U}\dot p_t,\qquad \partial_t[J\Phi_m(r_t)]=\chi_m(\log r_t)\dot p_t.\]

The Lean proof applies the time chain rule first to the constant multiple defining \(r_t\), then to \texttt{bounded\allowbreak{}Entropy\_\allowbreak{}has\allowbreak{}Deriv\allowbreak{}At}, then to multiplication by \(J\). The remaining coefficient cancels because \(J(y)Ze^{U(y)}=1\), using the inverse of the strictly positive real partition function and \(e^{-U}e^U=1\). This cancellation is what permits the spacetime estimate to use the Lebesgue \(L^1\) bound on \(\dot p_t\).

\subsection{\texorpdfstring{Identifying finite entropy through the chosen density}{Identifying finite entropy through the chosen density}}


Suppose \(\mu=r\nu\) is a probability measure, \(r\ge0\) is measurable, and \(H(\mu\mid\nu)<\infty\). Then

\[\mathrm{kl}(r)=r\log r+1-r\in L^1(\nu),\qquad \int\mathrm{kl}(r)\,d\nu=H(\mu\mid\nu).\]

Mathlib expresses KL using the Radon–Nikodym derivative. The equality \texttt{Measure.\allowbreak{}rn\allowbreak{}Deriv\_\allowbreak{}with\allowbreak{}Density} identifies that derivative almost everywhere with \texttt{ENNReal.\allowbreak{}of\allowbreak{}Real \ensuremath{\circ} r}. Nonnegativity of \(r\) makes its real part exactly \(r\). Finite KL yields integrability of the log-likelihood ratio, which the KL library converts to integrability of \texttt{kl\allowbreak{}Fun} of the derivative. An almost-everywhere change of representative gives both displayed conclusions for the chosen density itself.

\subsection{\texorpdfstring{Integrability of every actual bounded entropy}{Integrability of every actual bounded entropy}}


At \(t>0\), the actual interpolation has a positive \(C^1\) relative density \(r_t\). Its Fisher information is finite from the initial score moment and Gaussian smoothing; LSI then gives finite actual KL. The bounded primitive satisfies

\[0\le\Phi_m(r_t(y))\le\mathrm{kl}(r_t(y)).\]

The preceding density lemma proves that the right-hand side is integrable. The left-hand side is measurable: continuity of \(r_t\) and the derivative of \(\Phi_m\) at positive arguments give continuity of their composition. The Lean \texttt{mono'} argument therefore derives integrability for each truncation from a verified KL majorant. Initial density differentiability, initial Fisher information and a second moment of the state are not required.

\subsection{\texorpdfstring{Convergence of bounded entropy to actual KL}{Convergence of bounded entropy to actual KL}}


For every fixed positive time,

\[\int\Phi_m(r_t)\,d\pi\longrightarrow H(\rho_t\mid\pi).\]

Pointwise convergence \(\Phi_m(r)\to\mathrm{kl}(r)\) holds for \(r>0\). The domination \(0\le\Phi_m(r_t)\le\mathrm{kl}(r_t)\), together with integrability proved above, gives dominated convergence under the fixed Gibbs measure. Before taking the limit, Lean identifies the integral of \texttt{kl\allowbreak{}Fun r\_\allowbreak{}t} with the actual KL through \texttt{tilted\allowbreak{}Measure\_\allowbreak{}interpolation\allowbreak{}Relative\allowbreak{}Density}. Thus the limiting quantity is the KL of the actual Gaussian law, rather than an independently specified entropy path.

\subsection{\texorpdfstring{The bounded entropy and its time flux}{The bounded entropy and its time flux}}


The two functions used in the time integration are

\[H_m(t)=\int\Phi_m(r_t(y))\,d\pi(y),\qquad q_m(t)=\int\chi_m(\log r_t(y))\dot p_t(y)\,dy.\]

\texttt{interpolation\allowbreak{}Bounded\allowbreak{}Entropy} uses the fixed Gibbs reference measure; \texttt{interpolation\allowbreak{}Bounded\allowbreak{}Entropy\allowbreak{}Flux} uses Lebesgue measure. Their different measures match the exact cancellation in the weighted chain rule. The definitions are explicit functions of the actual mixture density and its actual kernel time derivative.

\subsection{\texorpdfstring{Admissibility of the actual bounded log test}{Admissibility of the actual bounded log test}}


The spatial weak test is \(\psi_m(y)=\chi_m(\log r_t(y))\). Positivity and \(C^1\) regularity of \(r_t\), combined with smoothness of \(\chi_m\), show that \(\psi_m\in C^1\). Its uniform bound is

\[|\psi_m(y)|\le(m+1)\pi/2.\]

The spatial chain rule gives

\[\nabla\psi_m(y)=w_m(\log r_t(y))\nabla\log r_t(y),\qquad \|\nabla\psi_m(y)\|\le\|\nabla\log r_t(y)\|.\]

The relative score is already in \(L^2(\rho_t)\). Its norm domination therefore proves square integrability of the test gradient via \texttt{Mem\allowbreak{}Lp.\allowbreak{}of\_\allowbreak{}le}. Lean explicitly constructs measurability of that gradient from its continuous Fréchet derivative. These are precisely the three hypotheses of the global Gaussian weak equation: a bounded \(C^1\) test and an \(L^2\) gradient.

\subsection{\texorpdfstring{The bounded entropy flux is the actual spatial flux}{The bounded entropy flux is the actual spatial flux}}


Let \(K_t(x,y)\) be the Gaussian transition density, \(g=\nabla U\), and \(S_t=\nabla\log r_t\). The proved Gaussian weak equation, applied to the admissible test above, gives

\[q_m(t)=-\iint K_t(x,y)w_m(\log r_t(y))\left\langle g(x)-\frac{y-x+tg(x)}{2t},S_t(y)\right\rangle\,\mu(dx)\,dy.\]

The underlying flux is \(K_t[g(x)-(y-x+tg(x))/(2t)]\). Its divergence equals the actual kernel time derivative; global integration by parts is justified by the previously proved finite-gradient-energy estimates. The code substitutes the test-gradient chain rule and the equality between the explicit measurable score field and the actual relative score. The remaining rearrangement is bilinearity of the inner product and scalar multiplication.

\subsection{\texorpdfstring{The bounded time chain rule without a spatial time-\allowbreak{}majorant}{The bounded time chain rule without a spatial time-majorant}}


On any \(0<a\le b\), the weighted primitive has derivative \(\chi_m(\log r_t)\dot p_t\). The kernel theorem proves \(\dot p\in L^1([a,b]\times E)\), while \(\chi_m\) is bounded by \((m+1)\pi/2\). Thus the derivative of the weighted primitive is jointly integrable.



Integrability at time \(a\) follows from the bounded-entropy lemma after rewriting Gibbs measure as \(J(y)dy\). The generic spacetime FTC/\allowbreak{}Fubini theorem then yields

\[H_m(b)-H_m(a)=\int_a^b q_m(t)\,dt.\]

Applying the same identity on every subinterval proves absolute continuity of \(H_m\); Lebesgue differentiation gives \(H_m'(t)=q_m(t)\) almost everywhere. In code, \texttt{hdot} is obtained by multiplying the actual spacetime \(L^1\) derivative by a bounded measurable function, \texttt{hderiv} supplies the pointwise chain rule, and \texttt{hweight} changes the integrated weighted primitive back to \(H_m\). No common bound on all time derivatives by a single spatial function is assumed.

\subsection{\texorpdfstring{Expressing the flux on one independent probability space}{Expressing the flux on one independent probability space}}


Take \(X\sim\mu\) and \(Z\sim N(0,I)\) independently, and set \(Y_t=X-tg(X)+\sqrt{2t}Z\). Define

\[V_t(X,Z)=g(X)-\frac{\sqrt{2t}}{2t}Z,\qquad \widehat S_t(X,Z)=S_t(Y_t),\qquad q(t)=-\mathbb E\langle V_t,\widehat S_t\rangle.\]

The three Lean definitions give exactly these fields and the raw entropy flux. A fixed product probability space is used at every time, which makes joint measurability and time dominated convergence explicit. The field \(V_t\) is the frozen Gaussian spatial flux divided by its transition density after substituting the actual update.

\subsection{\texorpdfstring{Square integrability of the two actual Gaussian factors}{Square integrability of the two actual Gaussian factors}}


The initial gradient belongs to \(L^2(\mu)\), and the standard Gaussian noise belongs to \(L^2\). Pulling both to the independent product space and taking their scalar combination proves \(V_t\in L^2(\mu\otimes\gamma)\). The actual relative score belongs to \(L^2(\rho_t)\). Since \(\rho_t\) is the pushforward of \(\mu\otimes\gamma\) by \(Y_t\), the measurable-map characterization of \(L^2\) gives

\[\widehat S_t=S_t\circ Y_t\in L^2(\mu\otimes\gamma).\]

\texttt{frozen\allowbreak{}Gaussian\allowbreak{}Velocity\allowbreak{}Field\_\allowbreak{}mem\allowbreak{}Lp} performs the product-space linear combination. \texttt{interpolation\allowbreak{}Score\allowbreak{}Along\allowbreak{}Update\_\allowbreak{}mem\allowbreak{}Lp} first identifies the explicit score with the actual gradient of \(\log r_t\), then applies \texttt{mem\allowbreak{}Lp\_\allowbreak{}map\_\allowbreak{}measure\_\allowbreak{}iff}. Consequently their inner product is integrable.

\subsection{\texorpdfstring{The exact frozen velocity energy}{The exact frozen velocity energy}}


Independence and the zero Gaussian mean cancel the drift-noise cross term:

\[\mathbb E\|V_t\|^2=\mathbb E\|g(X)\|^2+\frac{d}{2t},\qquad t>0.\]

The proof specializes the already verified affine-Gaussian moment formula with coefficients \(1\) and \(-\sqrt{2t}/(2t)\). The squared coefficient is \((2t)/(4t^2)=1/(2t)\), using \((\sqrt{2t})^2=2t\). This exact finite energy gives a compact-time bound that is uniform away from zero. No moment of \(X\) itself is used.

\subsection{\texorpdfstring{Transporting the bounded flux to independent noise}{Transporting the bounded flux to independent noise}}


The law of \((X,Y_t)\) has density \(K_t(x,y)\) with respect to \(\mu(dx)dy\). Therefore the joint flux identity becomes

\[q_m(t)=-\mathbb E\!\left[w_m(\log r_t(Y_t))\langle V_t,\widehat S_t\rangle\right].\]

The code introduces the measurable map \texttt{T} sending \((x,z)\) to \((x,Y_t)\), the kernel \texttt{K}, and the weighted pairing \texttt{F}. \texttt{frozen\allowbreak{}Gaussian\allowbreak{}Joint\allowbreak{}Law\_\allowbreak{}eq\_\allowbreak{}density} identifies the pushforward law. The integral under its real density becomes \(\iint K F\), and \texttt{integral\_\allowbreak{}map\_\allowbreak{}of\_\allowbreak{}strongly\allowbreak{}Measurable} changes that integral back to the independent product space. Finally \(Y_t-(X-tg(X))=\sqrt{2t}Z\) reduces the residual to the declared velocity field. The change of variables concerns the actual Gaussian law and does not assume an invariant Euler step.

\subsection{\texorpdfstring{Removing the truncation at fixed positive time}{Removing the truncation at fixed positive time}}


The raw pairing \(\langle V_t,\widehat S_t\rangle\) is integrable by the two \(L^2\) estimates. The scalar map \((x,z)\mapsto\log r_t(Y_t)\) is measurable. Since \(0\le w_m\le1\) and \(w_m\to1\), dominated convergence on the independent probability space yields

\[q_m(t)\longrightarrow q(t),\qquad t>0.\]

The dominating function is \(|\langle V_t,\widehat S_t\rangle|\). Crucially this limit is performed after spatial integration by parts; it does not ask for integrability of \(|\log r_t|\,|\dot p_t|\), which is not supplied by the original assumptions.

\subsection{\texorpdfstring{A fixed-\allowbreak{}reference formula for the relative Fisher energy}{A fixed-reference formula for the relative Fisher energy}}


Define

\[I(t)=\int r_t(y)\|S_t(y)\|^2\,d\pi(y).\]

The actual relative-density identity gives \(I(t)=\int\|S_t\|^2\,d\rho_t\). Applying the actual Gaussian pushforward once more yields

\[\mathbb E\|\widehat S_t\|^2=I(t).\]

The first definition fixes the reference measure for time measurability and spacetime integrability. The equality uses \texttt{interpolation\allowbreak{}Relative\allowbreak{}Score\allowbreak{}Field\_\allowbreak{}weighted\_\allowbreak{}energy} and the measurable-map Bochner integral formula. The proof supplies strong measurability of the score norm square explicitly, fixing time in the joint measurable score field.

\subsection{\texorpdfstring{A uniform bound for every truncated flux}{A uniform bound for every truncated flux}}


For any vectors \(v,s\) and any \(0\le w\le1\),

\[|w\langle v,s\rangle|\le\|v\|\|s\|\le\tfrac12(\|v\|^2+\|s\|^2).\]

The second inequality is \((\|v\|-\|s\|)^2\ge0\). Integrate on the independent Gaussian space, use the norm inequality for Bochner integrals, and substitute the two exact energies:

\[|q_m(t)|\le\frac12\left(G_0+\frac d{2t}+I(t)\right),\qquad G_0=\int\|g\|^2\,d\mu.\]

In Lean, \texttt{hi} proves integrability of the bounded weighted pairing, \texttt{hmajor} proves integrability of the quadratic majorant, and \texttt{hpoint} is the pointwise Cauchy–Schwarz/\allowbreak{}Young estimate. Thus the bound is justified for every truncation without assuming convergence or integrability of the eventual entropy derivative.

\subsection{\texorpdfstring{Joint measurability of both time fluxes}{Joint measurability of both time fluxes}}


The bounded flux is a parameter integral of \(\chi_m(\log r_t(y))\dot p_t(y)\). Both factors are jointly measurable in \((t,y)\), so parameter integration makes \(t\mapsto q_m(t)\) measurable. The raw flux uses the fixed independent space: the map \((t,x,z)\mapsto(t,Y_t)\) is measurable, the explicit relative-score field is jointly strongly measurable, and the velocity field is a measurable algebraic combination of \(g(X)\), \(t^{-1}\), \(\sqrt{2t}\), and \(Z\). Their inner product and its parameter integral are measurable. The definitions at nonpositive times serve only to make these globally specified measurable functions; all analytic conclusions are restricted to positive times.

\subsection{\texorpdfstring{The integrable majorant for the time limit}{The integrable majorant for the time limit}}


Fix \(0<a\le b\). The actual Gaussian energy theorem proves \(I\in L^1([a,b])\). Set

\[D(t)=\tfrac12\left(G_0+\frac d{2a}+I(t)\right).\]

The interval has finite Lebesgue measure, so the constant term is integrable; hence \(D\in L^1([a,b])\). Since \(t\ge a>0\), \(d/(2t)\le d/(2a)\). The uniform flux estimate gives \(|q_m(t)|\le D(t)\) for every \(m\). The fixed-time limit \(q_m(t)\to q(t)\), followed by continuity of the absolute value, also gives \(|q(t)|\le D(t)\). The formal statement supplies the single majorant together with all three verified integrability/\allowbreak{}domination conclusions, ready for dominated convergence in time.

\subsection{\texorpdfstring{The full actual entropy identity}{The full actual entropy identity}}


Let \(H(t)=H(\rho_t\mid\pi)\), represented by the actual finite KL. For each truncation,

\[H_m(b)-H_m(a)=\int_a^b q_m(t)\,dt.\]

At both endpoints, the spatial dominated-convergence theorem gives \(H_m\to H\). In time, the common integrable majorant \(D\) and the pointwise convergence \(q_m\to q\) give \(\int_a^b q_m\to\int_a^b q\). Uniqueness of limits therefore yields

\[H(b)-H(a)=\int_a^b q(t)\,dt.\]

\texttt{interpolation\allowbreak{}Entropy\_\allowbreak{}integral\_\allowbreak{}identity} applies the general \texttt{integral\_\allowbreak{}identity\_\allowbreak{}of\_\allowbreak{}dominated\_\allowbreak{}limits} under Lebesgue measure restricted to \((a,b]\). Its endpoint convergence hypotheses are the already established actual-KL limits, and all time domination comes from the previous theorem. No entropy-path regularity is supplied as a hypothesis.

\subsection{\texorpdfstring{Absolute continuity and the a.\allowbreak{}e.\allowbreak{} derivative of actual entropy}{Absolute continuity and the a.e. derivative of actual entropy}}


The raw flux is measurable and satisfies \(|q|\le D\), so it is integrable on every positive-time compact interval. Apply the full entropy identity with terminal time \(t\in[a,b]\):

\[H(t)=H(a)+\int_a^t q(u)\,du.\]

An indefinite integral of an \(L^1\) function is absolutely continuous, and the Lebesgue differentiation theorem identifies its derivative almost everywhere. Thus

\[H\in AC([a,b]),\qquad H'(t)=q(t)\quad\text{for almost every }t\in(a,b).\]

The Lean proof derives interval integrability from the actual raw-flux measurability and the verified majorant, then invokes \texttt{regularity\_\allowbreak{}of\_\allowbreak{}integral\_\allowbreak{}identity}. This completes the analytic time-regularity bridge from the original potential, LSI and finite initial score-energy assumptions. A later score-pairing identity rewrites the raw Gaussian flux as negative Fisher energy plus the frozen gradient discrepancy, yielding the quantitative entropy-production estimate.

\section{\texorpdfstring{Averaging the Gaussian component score}{Averaging the Gaussian component score}}
\label{proof-module-UnadjustedLangevin-LSI-GaussianScorePairing-lean}
{\small\color{gray}Source module: \texttt{Unadjusted\allowbreak{}Langevin/\allowbreak{}LSI/\allowbreak{}Gaussian\allowbreak{}Score\allowbreak{}Pairing.\allowbreak{}lean}.\par}


The Gaussian component score \(S_{t,x}(y)=-(y-x+t\nabla U(x))/(2t)\) averages to the actual mixture score. This module establishes the vector identity \(p_t\nabla\log p_t=\int K_tS_{t,x}d\mu\), proves all joint integrability needed for pairing it with a continuous \(L^2(\rho_t)\) vector field, and transfers the resulting identity to the fixed independent initial/\allowbreak{}noise space used by entropy production.

\subsection{\texorpdfstring{frozen\allowbreak{}Gaussian\allowbreak{}Score}{frozenGaussianScore}}


For a fixed initial state \(x\), the Gaussian proposal has center \(c_t(x)=x-t\nabla U(x)\). Its score in the proposed state is \(S_{t,x}(y)=\nabla_y\log K_t(x,y)=-(y-c_t(x))/(2t)\). This definition is an explicit affine vector field. It is different from the mixture score, which averages this field over the posterior initial-state distribution.

\subsection{\texorpdfstring{frozen\allowbreak{}Gaussian\allowbreak{}Score mem\allowbreak{}Lp}{frozenGaussianScore memLp}}


Represent the actual initial/\allowbreak{}output pair as \((X,Y_t)\), with \(Y_t=c_t(X)+\sqrt{2t}Z\) and independent standard Gaussian \(Z\). Substituting into the component score cancels the center: \(S_{t,X}(Y_t)=-(\sqrt{2t}/(2t))Z\). The Gaussian identity random variable is square integrable; scalar multiplication and pullback to the product preserve this property. The map-measure theorem transfers it to the actual pair law. No initial position or gradient moment is needed for this component score.

\subsection{\texorpdfstring{frozen\allowbreak{}Gaussian\allowbreak{}Joint mem\allowbreak{}Lp state field}{frozenGaussianJoint memLp state field}}


If \(V\in L^2(\rho_t)\), then \(V(Y_t)\in L^2(\mu\otimes\mathcal N(0,I))\), because \(\rho_t\) is the actual pushforward by the proposal. Map the independent pair to \((X,Y_t)\); a second application of the map-measure equivalence gives \((x,y)\mapsto V(y)\in L^2\) under the actual pair law. Continuity of \(V\) supplies the measurability required in both transfers.

\subsection{\texorpdfstring{integrable frozen\allowbreak{}Gaussian\allowbreak{}Score pairing}{integrable frozenGaussianScore pairing}}


The component score and the state field \(V(y)\) are both square integrable under the actual pair law. Holder's inequality makes the product of their norms integrable, and \(|\langle S,V\rangle|\le\|S\|\|V\|\) proves integrability of their inner product. The exact pair-density identity then rewrites this law as \(K_t(x,y)\mu(dx)dy\). The with-density equivalence gives joint integrability of \(K_t\langle S_{t,x},V(y)\rangle\) under \(\mu\otimes dy\), which is the hypothesis needed for Fubini in the averaging calculation.

\subsection{\texorpdfstring{interpolation\allowbreak{}Lebesgue\allowbreak{}Density weighted score}{interpolationLebesgueDensity weighted score}}


The already constructed spatial derivative of the Gaussian mixture identifies its score as the posterior mean of the component score. The posterior density in \(x\) is \(K_t(x,y)/p_t(y)\), where \(p_t(y)=\int K_t(x,y)d\mu(x)>0\). Substitute this explicit with-density formula into the posterior integral and multiply by \(p_t(y)\). The common denominator cancels pointwise, giving \(p_t(y)\nabla\log p_t(y)=\int K_t(x,y)S_{t,x}(y)d\mu(x)\). The proof uses the actual Gaussian normalizers in both precision and step-size notation, rather than assuming a formal Bayes identity.

\subsection{\texorpdfstring{integrable frozen\allowbreak{}Gaussian\allowbreak{}Score weighted}{integrable frozenGaussianScore weighted}}


The component score is \(L^2\) under the actual pair probability law, hence is also \(L^1\). Identify that law with its explicit density \(K_t\) under \(\mu\otimes dy\). The vector-valued with-density integrability theorem gives \((x,y)\mapsto K_t(x,y)S_{t,x}(y)\in L^1(\mu\otimes dy)\). In particular, Fubini supplies integrability of the vector section in \(x\) for almost every output \(y\); this is needed before moving an inner product through its integral.

\subsection{\texorpdfstring{integral frozen\allowbreak{}Gaussian\allowbreak{}Score pairing}{integral frozenGaussianScore pairing}}


For almost every \(y\), vector-section integrability permits the linear functional \(w\mapsto\langle w,V(y)\rangle\) to commute with integration. The pointwise weighted-score identity then gives \(\int K_t(x,y)\langle S_{t,x}(y),V(y)\rangle d\mu(x)=p_t(y)\langle\nabla\log p_t(y),V(y)\rangle\).



The scalar pairing is jointly integrable by the preceding Holder argument, so Fubini exchanges the initial-state and output integrals. Integrate the displayed identity in \(y\). Finally \(p_t(y)dy=\rho_t(dy)\), with its explicit measurable finite density, changes the weighted Lebesgue integral into the actual-law integral. Thus \(\int K_t\langle S_{t,x},V\rangle d\mu\,dy=\int\langle\nabla\log p_t,V\rangle d\rho_t\).

\subsection{\texorpdfstring{integral frozen\allowbreak{}Gaussian\allowbreak{}Score pairing independent}{integral frozenGaussianScore pairing independent}}


Apply the measurable pair map \(T_t(x,z)=(x,Y_t(x,z))\) to the continuous pairing \(\langle S_{t,x}(y),V(y)\rangle\). The map-integral theorem expresses its actual pair-law integral as an integral over the fixed independent initial/\allowbreak{}noise space. The exact pair-density theorem expresses the same integral as the Gaussian-weighted Lebesgue pairing just evaluated.



Substitute \(Y_t=c_t(x)+\sqrt{2t}z\) in the component score. The residual becomes \(\sqrt{2t}z\), so the result is \(\mathbb E\langle-(\sqrt{2t}/(2t))Z,V(Y_t)\rangle=\mathbb E_{\rho_t}\langle\nabla\log p_t,V\rangle\). With \(V=\nabla\log r_t\), this identity turns the Gaussian-noise part of the raw entropy flux into the Lebesgue-score pairing needed for the full Fisher term.

\section{\texorpdfstring{Prove the actual Fisher dissipation and sharp frozen-\allowbreak{}gradient forcing}{Prove the actual Fisher dissipation and sharp frozen-gradient forcing}}
\label{proof-module-UnadjustedLangevin-LSI-GaussianEntropyProduction-lean}
{\small\color{gray}Source module: \texttt{Unadjusted\allowbreak{}Langevin/\allowbreak{}LSI/\allowbreak{}Gaussian\allowbreak{}Entropy\allowbreak{}Production.\allowbreak{}lean}.\par}


The actual Gaussian entropy chain rule gives an almost-everywhere derivative represented by a flux on the independent initial/\allowbreak{}noise probability space. Averaging the genuine Gaussian component score identifies this flux with \(-I_t+\mathbb E\langle\nabla U(Y_t)-\nabla U(X),\nabla\log r_t(Y_t)\rangle\). Young's inequality absorbs a quarter of \(I_t\); the already proved exact Gaussian displacement estimate bounds the remaining error. The result is the unconditional analytic production inequality \(H'_t\le-\tfrac34I_t+t^2L^2G_0+2tdL^2\) almost everywhere at positive times, with no interpolation regularity or PDE data hypotheses.

\subsection{\texorpdfstring{Justify inner-\allowbreak{}product integrability from two L2 fields}{Justify inner-product integrability from two L2 fields}}


For vector fields \(F,G\in L^2(P)\), Hölder's inequality puts \(\|F\|\|G\|\) in \(L^1(P)\). Since \(|\langle F,G\rangle|\le\|F\|\|G\|\), their inner product is integrable.



This private Lean helper applies \texttt{Mem\allowbreak{}Lp.\allowbreak{}mul} to the two norm fields with resulting exponent one, converts it to integrability, and uses \texttt{Integrable.\allowbreak{}mono'} with \texttt{norm\_\allowbreak{}inner\_\allowbreak{}le\_\allowbreak{}norm}. Measurability follows from the two input fields. The helper is used for the initial gradient, current gradient, Gaussian score, and relative-score pairings before splitting or recombining their integrals.

\subsection{\texorpdfstring{Identify relative-\allowbreak{}score energy on the independent noise space with actual Fisher information}{Identify relative-score energy on the independent noise space with actual Fisher information}}


Let \(S_t(X,Z)=\nabla\log r_t(Y_t)\) and \(P=\mu\otimes\mathcal N(0,I)\). The exact update-law identity gives

\[\mathbb E_P\|S_t\|^2=\int\|\nabla\log r_t\|^2\,d\rho_t=I_t.\]

Lean first obtains integrability of the actual state-space relative score squared. \texttt{integral\_\allowbreak{}map} transports this integral through the measurable Gaussian update. The pointwise equality \texttt{hpoint} identifies the jointly measurable score field used along the update with the actual gradient of the log relative density. Finally the definition of \texttt{relative\allowbreak{}Fisher}, the actual tilted-law equality, and the ordinary/\allowbreak{}nonnegative integral conversion identify the finite real Fisher value. No energy equality is supplied as a hypothesis.

\subsection{\texorpdfstring{Match the measurable state-\allowbreak{}space score field to the same finite Fisher value}{Match the measurable state-space score field to the same finite Fisher value}}


The jointly measurable field used in time-space estimates agrees with the actual relative score at each positive time. Consequently

\[\int\|\mathrm{ScoreField}_t(y)\|^2\,d\rho_t(y)=I_t.\]

Lean converts the previously proved actual-score \(L^2\) result to \texttt{Mem\allowbreak{}Lp} for this explicit field using \texttt{interpolation\allowbreak{}Relative\allowbreak{}Score\allowbreak{}Field\_\allowbreak{}eq}. It then maps the squared norm to the independent update space and invokes the preceding Fisher-energy equality. This bridge allows the LSI inequality, the time-space domination bounds, and the final production estimate to use interchangeable but formally identified score representations.

\subsection{\texorpdfstring{Integrate Young absorption against the actual frozen-\allowbreak{}gradient discrepancy}{Integrate Young absorption against the actual frozen-gradient discrepancy}}


Write \(D_t=\nabla U(Y_t)-\nabla U(X)\). The pointwise inequality

\[\langle D_t,S_t\rangle\le\frac14\|S_t\|^2+\|D_t\|^2\]

is integrable because both squared norms have already been proved integrable. Their expectations satisfy \(\mathbb E\|S_t\|^2=I_t\) and

\[\mathbb E\|D_t\|^2\le L^2\bigl(t^2G_0+2td\bigr),\qquad G_0=\mathbb E\|\nabla U(X)\|^2.\]

Lean constructs discrepancy measurability as \texttt{h\allowbreak{}Dm}, score measurability as \texttt{h\allowbreak{}Sm}, and the two energy integrabilities as \texttt{h\allowbreak{}Ei} and \texttt{h\allowbreak{}Si}. The weighted-inner-product helper with weight one proves cross-term integrability before applying the integrated Young inequality. \texttt{h\allowbreak{}F} replaces joint score energy by finite actual Fisher information; \texttt{h\allowbreak{}Eb} is the previously proved sharp Gaussian discrepancy bound. Combining them gives exactly \(\mathbb E\langle D_t,S_t\rangle\le I_t/4+t^2L^2G_0+2tdL^2\).

\subsection{\texorpdfstring{Convert the actual Gaussian entropy flux into Fisher dissipation and gradient error}{Convert the actual Gaussian entropy flux into Fisher dissipation and gradient error}}


The raw entropy flux is \(-\mathbb E\langle g(X)-\sqrt{2t}Z/(2t),S_t\rangle\). The averaged Gaussian score identity gives

\[\mathbb E\left\langle-\frac{\sqrt{2t}}{2t}Z,S_t\right\rangle=\int\langle\nabla\log p_t,\nabla\log r_t\rangle\,d\rho_t.\]

Since \(\nabla\log r_t=\nabla U+\nabla\log p_t\), the right side equals \(I_t-\mathbb E\langle\nabla U(Y_t),S_t\rangle\). Substituting and collecting the two gradient terms yields

\[\mathrm{RawFlux}_t=-I_t+\mathbb E\langle\nabla U(Y_t)-\nabla U(X),S_t\rangle.\]

The Lean proof names the pulled-back initial gradient \texttt{GX}, current gradient \texttt{GY}, Gaussian component score \texttt{N}, and relative score \texttt{S}. All four are in \(L^2(P)\), so \texttt{h\allowbreak{}X}, \texttt{h\allowbreak{}G}, \texttt{h\allowbreak{}Noise}, and \texttt{h\allowbreak{}Out} establish the integrability required for every sum and difference. The map identities \texttt{h\allowbreak{}Map\allowbreak{}G} and \texttt{h\allowbreak{}Map\allowbreak{}S} identify output-space integrals with the actual independent-noise integrals. \texttt{h\allowbreak{}Leb} is the pointwise score decomposition, \texttt{h\allowbreak{}Avg'} is its integrated Gaussian-score consequence, and \texttt{h\allowbreak{}Raw}/\allowbreak{}\texttt{h\allowbreak{}Err} perform the two final algebraic expansions. This connects the analytic entropy flux to actual Fisher information without assuming a conditional drift equation.

\subsection{\texorpdfstring{Retain three quarters of actual Fisher dissipation}{Retain three quarters of actual Fisher dissipation}}


Insert the integrated Young estimate into the exact entropy-flux decomposition:

\[\mathrm{RawFlux}_t=-I_t+\mathbb E\langle D_t,S_t\rangle\le-\frac34I_t+t^2L^2G_0+2tdL^2.\]

The Lean proof rewrites the flux with \texttt{interpolation\allowbreak{}Raw\allowbreak{}Entropy\allowbreak{}Flux\_\allowbreak{}eq\_\allowbreak{}fisher\_\allowbreak{}error}, obtains the already justified numerical cross-term bound, and uses linear arithmetic. The target's probability normalization and finite initial gradient energy suffice for this pointwise flux estimate. LSI is needed later to relate Fisher information to entropy, rather than to justify Young absorption.

\subsection{\texorpdfstring{Obtain the genuine positive-\allowbreak{}time entropy derivative inequality on one common full-\allowbreak{}measure set}{Obtain the genuine positive-time entropy derivative inequality on one common full-measure set}}


The actual KL entropy is absolutely continuous on every compact interval inside \((0,\infty)\), and its derivative equals the actual raw flux almost everywhere there. To use this fact on an entire finite range \((0,T)\), choose the countable intervals

\[\left(\frac1{n+1},T+1\right),\qquad n=0,1,2,\ldots.\]

Intersect their full-measure derivative sets. Every \(t\in(0,T)\) belongs to at least one such interval because \(1/(n+1)\to0\). At that time the actual derivative equals the raw flux, so the pointwise flux bound gives

\[H'_t\le-\frac34 I_t+t^2L^2G_0+2tdL^2\quad\text{for almost every }t\in(0,T).\]

Lean's \texttt{hderiv} applies the fully proved local entropy chain rule for each interval; \texttt{ae\_\allowbreak{}all\_\allowbreak{}iff} makes their countable intersection. \texttt{exists\_\allowbreak{}nat\_\allowbreak{}one\_\allowbreak{}div\_\allowbreak{}lt} chooses an interval containing the current positive time. \texttt{hd.\allowbreak{}deriv} identifies the ordinary derivative, and \texttt{interpolation\allowbreak{}Raw\allowbreak{}Entropy\allowbreak{}Flux\_\allowbreak{}le} supplies the numerical bound. This argument avoids asserting a derivative at time zero, where the initial law need not have a smooth positive density.

\section{\texorpdfstring{The full Gaussian ULA entropy theorem}{The full Gaussian ULA entropy theorem}}
\label{proof-module-UnadjustedLangevin-LSI-UnconditionalTheorem-lean}
{\small\color{gray}Source module: \texttt{Unadjusted\allowbreak{}Langevin/\allowbreak{}LSI/\allowbreak{}Unconditional\allowbreak{}Theorem.\allowbreak{}lean}.\par}


For a differentiable potential with an \(L\)-Lipschitz gradient, its normalized Gibbs target satisfying LSI with constant \(\alpha>0\), and a finite-KL initial probability law, this module proves the actual ULA bound

\[\mathrm{KL}(\rho_k\mid\nu)\le e^{-\alpha hk}\mathrm{KL}(\rho_0\mid\nu)+8hdL^2/\alpha,\quad0<h\le\alpha/(4L^2).\]

The final conclusion includes finiteness of the actual divergence. All score moments, density calculus, weak equations, entropy endpoints, and numerical compatibility estimates are supplied by the preceding modules. The theorem applies also in dimension zero.

\subsection{\texorpdfstring{Handle the singleton state space and a zero smoothness parameter}{Handle the singleton state space and a zero smoothness parameter}}


On a singleton state space every probability measure is equal: any measurable set either contains the unique point and is the whole space, or is empty. Measure extensionality and probability normalization therefore identify any two laws. In particular their KL divergence is zero.



Independently of dimension, a positive admissible step forces \(L>0\). If \(L=0\), Lean’s total-field convention gives \(\alpha/(4L^2)=0\), contradicting \(0<h\le0\). This elementary guard permits the score-exponential estimates to divide by \(L\). It is derived from the step condition, not assumed in the final theorem.

\subsection{\texorpdfstring{Prove finiteness at every actual ULA iterate}{Prove finiteness at every actual ULA iterate}}


Start from the supplied finite initial KL and induct on \(k\). If \(\rho_k\) has finite KL, the direct entropy/\allowbreak{}score estimate gives \(\nabla U\in L^2(\rho_k)\). For the actual positive-time Gaussian interpolation, the Lebesgue score energy is at most \(d/(2h)\), and the propagated gradient moment is finite. Thus relative Fisher information is finite. LSI with \(\alpha>0\) then gives finite KL of the output.



The endpoint identity identifies that output with \(\rho_{k+1}\). This induction uses no step restriction beyond positivity; it is a separate guard that prevents the extended value infinity from disappearing under \texttt{to\allowbreak{}Real}.

\subsection{\texorpdfstring{Derive all scalar step restrictions}{Derive all scalar step restrictions}}


In positive dimension, the bounded sinusoidal tests in \texttt{Compatibility} give \(\alpha\le L\). Since \(\alpha>0\), also \(L>0\). Multiplying the original restriction by \(4L^2>0\) gives \(4hL^2\le\alpha\le L\), hence \(hL\le1/4\). Multiplying \(\alpha\le L\) by \(h>0\) also gives \(\alpha h\le1/4\).



The final theorem splits off the singleton case before invoking this positive-dimensional argument. No separate compatibility assumption is required.

\subsection{\texorpdfstring{Apply LSI to the actual positive-\allowbreak{}time density}{Apply LSI to the actual positive-time density}}


Finite initial KL supplies the initial gradient moment. The actual Gaussian mixture is positive C\ensuremath{{}^{1}} relative to Gibbs, has finite KL, and has finite Fisher information at every \(t>0\). Apply the original extended-valued LSI to this density. Since both sides are finite, order is preserved on passing to their real values:

\[2\alpha\,\mathrm{KL}(\rho_t\mid\nu)\le I_t.\]

The integral form identifies \(I_t\) with \(\int\|\mathcal S(t,y)\|^2d\rho_t(y)\), using the proved equality between the measurable score field and \(\nabla\log r_t\), and the equality \(r_t\nu=\rho_t\). The second version keeps the extended-Fisher notation for the entropy-production assembly.

\subsection{\texorpdfstring{Supply the actual analytic inputs to scalar comparison}{Supply the actual analytic inputs to scalar comparison}}


Fix any finite-KL probability law \(\mu\). The score-exponential argument gives the integrability of \(\|\nabla U\|^2\) and

\[G_0\le4L^2H_\nu(\mu)/\alpha+2dL.\]

The bounded-log entropy identity proves local absolute continuity of the actual entropy \(H(t)\) on every \([s,h]\) with \(s>0\). The actual Gaussian weak equation and Young estimate give \(H^{\prime}\le-3I/4+t^2L^2G_0+2tdL^2\) almost everywhere. The preceding LSI estimate supplies \(2\alpha H\le I\).



The explicit initial majorant \(M\) is continuous at zero, satisfies \(M(0)=H_\nu(\mu)\), and bounds \(H(t)\) for sufficiently small positive \(t\). These are exactly the premises of \texttt{one\_\allowbreak{}step\_\allowbreak{}entropy\_\allowbreak{}of\_\allowbreak{}ae\_\allowbreak{}interpolation}, giving

\[H_\nu(\mu P_h)\le e^{-\alpha h}H_\nu(\mu)+6h^2dL^2.\]

The singleton case instead uses equality of probability laws and nonnegativity of the right side. No density path, weak equation, entropy derivative, or moment estimate is a theorem argument.

\subsection{\texorpdfstring{Derive the actual discrete recurrence}{Derive the actual discrete recurrence}}


The law \(\rho_k\) is a probability measure by the proved Gaussian Markov construction. Its KL is finite by the induction above. Apply the single-transition estimate with \(\mu=\rho_k\). The interpolation endpoint is precisely \(\rho_{k+1}\), yielding

\[H_{k+1}\le e^{-\alpha h}H_k+6h^2dL^2.\]

\texttt{ula\allowbreak{}Entropy} is the real value of the actual KL of \texttt{ula\allowbreak{}Laws}. The recurrence is derived here; it is not supplied as a premise.

\subsection{\texorpdfstring{Sum the recurrence and retain the finite-\allowbreak{}KL guard}{Sum the recurrence and retain the finite-KL guard}}


Return the already-proved output finiteness as the first component of the conclusion. In positive dimension, the step estimates give \(\alpha h\le1/4\). The scalar geometric-sum theorem applied to the actual recurrence gives

\[H_k\le e^{-\alpha hk}H_0+\frac{6h^2dL^2}{1-e^{-\alpha h}}\le e^{-\alpha hk}H_0+\frac{8hdL^2}{\alpha}.\]

The final inequality uses \(1-e^{-x}\ge3x/4\) for \(0\le x\le1/4\). In dimension zero, equality of probability measures makes the left side zero and the result follows by nonnegativity.



The complete binder list contains only differentiability, global gradient Lipschitzness, the original Gibbs LSI, positivity of \(\alpha\) and the step, the paper’s step restriction, and finite initial KL. The target probability normalization is included in the LSI definition. Lean’s kernel checks the entire chain, and \texttt{Check.\allowbreak{}lean} prints its axiom dependencies.

\section{\texorpdfstring{Unconditional IPLA entropy from the actual lifted LSI}{Unconditional IPLA entropy from the actual lifted LSI}}
\label{proof-module-IPLA-LSI-UnconditionalTheorem-lean}
{\small\color{gray}Source module: \texttt{IPLA/\allowbreak{}LSI/\allowbreak{}Unconditional\allowbreak{}Theorem.\allowbreak{}lean}.\par}


The actual IPLA state laws are ULA laws for the scaled potential \(V_N\), with step \(\gamma/N\), smoothness \(NL\), and dimension \(D=d_\theta+Nd_x\). Applying the shared unconditional ULA theorem proves finite joint KL and its convergence estimate from smoothness, the actual lifted-target LSI, the positive step restriction, and finite initial KL. The module also derives normalized, linear-LSI, and parameter-projection consequences.

\subsection{\texorpdfstring{Finite joint entropy and the full IPLA bound}{Finite joint entropy and the full IPLA bound}}


Let \(N>0\), \(D=d_\theta+Nd_x\), and \(V_N(z)=\sum_iU(\theta,\sqrt N z_i)\). Assume differentiable \(U\) with \(L\)-Lipschitz gradient, a genuine LSI with constant \(\beta>0\) for the actual lifted Gibbs target \(\Pi_N\), a positive step \(\gamma\le\beta/(4NL^2)\), and finite initial joint entropy \(H_0\). Then every actual particle law has finite KL and

\[H_k\le e^{-\beta\gamma k/N}H_0+\frac{8\gamma ND L^2}{\beta}.\]

The lifted potential is differentiable and has \(NL\)-Lipschitz gradient. The particle law is exactly its ULA law with step \(h=\gamma/N\). The step equivalence gives \(h\le\beta/(4(NL)^2)\), so the shared unconditional ULA theorem applies. Substituting \(h=\gamma/N\), dimension \(D\), and smoothness \(NL\) transforms its rate into \(e^{-\beta\gamma k/N}\) and its bias into \(8\gamma ND L^2/\beta\). The proof derives positive \(L\) from the positive admissible step before dividing.



The shared theorem constructs all Gaussian interpolation and entropy calculus. Neither compatibility nor regular interpolation data is a premise here. Finiteness is retained as the first component of the conclusion, before real-valued KL is used.

\subsection{\texorpdfstring{Normalize the full joint entropy}{Normalize the full joint entropy}}


Start with the finite full-state entropy theorem. Divide its real inequality by the positive number \(N\), preserving the inequality direction:

\[\frac{H_k}{N}\le e^{-\beta\gamma k/N}\frac{H_0}{N}+\frac{8\gamma D L^2}{\beta}.\]

The helper \texttt{normalize\_\allowbreak{}entropy\_\allowbreak{}bound} divides both the initial contribution and the bias, canceling one particle factor in the latter. The same finite-KL witness is retained for the unnormalized divergence. This quantity is the entropy of the full state divided by particle number; parameter projection is handled separately.

\subsection{\texorpdfstring{An actual lifted LSI of strength N alpha}{An actual lifted LSI of strength N alpha}}


Suppose the actual lifted target satisfies LSI with constant \(\beta=N\alpha\), where \(\alpha>0\). The positive particle count gives \(\beta>0\), and the step equivalence simplifies \(\gamma\le\beta/(4NL^2)\) to \(\gamma\le\alpha/(4L^2)\). Apply the unconditional normalized theorem at this lifted constant.



The rate simplifies to \(e^{-\alpha\gamma k}\). Since \(D/N=d_x+d_\theta/N\), its normalized bias becomes \((8\gamma L^2/\alpha)(d_x+d_\theta/N)\). Thus

\[\frac{H_k}{N}\le e^{-\alpha\gamma k}\frac{H_0}{N}+\frac{8\gamma L^2}{\alpha}\left(d_x+\frac{d_\theta}{N}\right).\]

The LSI premise is for \(\Pi_N\) itself. The proof performs these exact substitutions and retains finite output KL; it does not infer the lifted LSI by tensorizing a parameter–latent pair model.

\subsection{\texorpdfstring{Project the actual joint law to the parameter}{Project the actual joint law to the parameter}}


Apply the parameter projection to the actual particle law and to the same lifted Gibbs target. KL data processing gives \(\operatorname{KL}(\mu_k^\theta\Vert\Pi_N^\theta)\le H_k\). The unconditional joint theorem first supplies \(H_k<\infty\), which also proves finiteness of the projected KL before real conversion.



Transitivity with the full joint estimate yields

\[\operatorname{KL}(\mu_k^\theta\Vert\Pi_N^\theta)\le e^{-\beta\gamma k/N}H_0+\frac{8\gamma ND L^2}{\beta}.\]

The initial entropy on the right is the full joint \(H_0\). The bound is unnormalized, so the factor \(N\) in the bias is retained. A per-particle joint estimate must be multiplied by \(N\) before it bounds this unnormalized parameter divergence.


\end{document}
